\documentclass[11pt]{article}
\pdfoutput=1

\usepackage[margin=1in]{geometry}
\usepackage{amsmath,amssymb,amsthm,amsfonts,bm,physics,mathrsfs,upgreek,bbm}
\usepackage{graphicx}
\usepackage{float}
\usepackage{enumitem}
\usepackage{array}
\usepackage{booktabs}
\usepackage{multirow}
\usepackage{afterpage}
\usepackage{listings}
\usepackage{subcaption}
\usepackage{algorithm}
\usepackage{algorithmic}
\usepackage{changepage}

\newcommand{\INPUT}{\item[\textbf{Input:}]}

\usepackage[hang,flushmargin]{footmisc}
\usepackage{setspace}
\usepackage{xcolor}
\usepackage[hypertexnames=false]{hyperref}
\usepackage{xr-hyper}
\usepackage[capitalize,noabbrev]{cleveref}
\usepackage[round]{natbib}
\setlength{\bibsep}{0pt}
\definecolor{Bleu}{RGB}{0,0,204}
\hypersetup{
  colorlinks,
  citecolor=Bleu,
  linkcolor=Bleu,
  urlcolor=Bleu
}
\makeatletter
\newcommand*{\addFileDependency}[1]{%
  \typeout{(#1)}%
  \@addtofilelist{#1}%
  \IfFileExists{#1}{}{\typeout{No file #1.}}%
}
\makeatother

\newcommand*{\myexternaldocument}[1]{%
  \externaldocument{#1}%
  \addFileDependency{#1.tex}%
  \addFileDependency{#1.aux}%
}

\myexternaldocument{main_jasa}

\DeclareMathOperator*{\argmin}{argmin}
\DeclareMathOperator*{\argmax}{argmax}
\DeclareMathOperator*{\esssup}{ess\,sup}
\DeclareMathOperator*{\essinf}{ess\,inf}
\renewcommand{\P}{\mathbb{P}}

\newcommand{\blue}[1]{{\color{blue}#1}}
\newcommand{\addTK}{\widetilde T_{\mathcal K}}
\newcommand{\adk}{\widetilde{\partial}_k}
\newcommand\independent{\protect\mathpalette{\protect\independenT}{\perp}}
\def\independenT#1#2{\mathrel{\rlap{$#1#2$}\mkern2mu{#1#2}}}

\DeclareFontFamily{U}{jkpmia}{}
\DeclareFontShape{U}{jkpmia}{m}{it}{<->s*jkpmia}{}
\DeclareFontShape{U}{jkpmia}{bx}{it}{<->s*jkpbmia}{}
\DeclareMathAlphabet{\mathfrak}{U}{jkpmia}{m}{it}
\SetMathAlphabet{\mathfrak}{bold}{U}{jkpmia}{bx}{it}

\theoremstyle{plain}
\newtheorem{theorem}{Theorem}
\newtheorem{lemma}[theorem]{Lemma}
\newtheorem{corollary}{Corollary}
\newtheorem{definition}{Definition}
\newtheorem{assumption}{Assumption}
\newtheorem{condition}{Condition}
\newtheorem{example}{Example}

\def\spacingset#1{\renewcommand{\baselinestretch}{#1}\small\normalsize}
\usepackage[compact]{titlesec}

\newcommand{\JASATitleText}{Efficient Inference for Inverse Reinforcement Learning and Dynamic Discrete Choice Models}
\title{\bf Supplementary Appendix for \\ \JASATitleText}
\author{
  Lars van der Laan \\
  University of Washington \\
  \and
  Aur\'elien Bibaut \\
  Netflix Research \\
  \and
  Nathan Kallus \\
  Cornell Tech, Cornell University
}
\date{}

\begin{document}
\maketitle
\spacingset{1.15}
\appendix
 

This appendix is organized around the main reader tasks left after the core
results. Section~\ref{app::additionalapplications} collects the
reward-normalization extensions and the assumptions used by
Theorem~\ref{theorem::EIFnorm}. Section~\ref{app:estimator:details} gives
self-contained detailed statements for the example estimators, so each
example's estimator, additional assumptions, and asymptotic conclusion can be
read together. Section~\ref{app::normsim} reports additional simulation
results. The remaining sections contain proofs, technical lemmas for the soft
Bellman equation, and a brief expanded related-work discussion.


 
 
 

\section{Extended applications under reward normalization}
\label{app::additionalapplications}

This section collects the appendix-only extensions of the normalized-reward
framework and the detailed assumptions referenced in
Theorem~\ref{theorem::EIFnorm}.

\subsection{Detailed assumptions for normalized-reward functionals}
\label{appendix::conditionsEIFnorm}

We state here the assumptions used by Theorem~\ref{theorem::EIFnorm} and by
the normalized-reward extensions below.

\begin{enumerate}[label=\textbf{(C\arabic**)}, ref=C\arabic**, series=cond2, start=1]
\item \label{cond::boundedfunc::norm} For all $(r,k) \in \mathcal{H} \otimes \mathcal{K}$ in an $L^2$-neighborhood of $(r_{0,\nu,\gamma,g}, k_0)$, there exist Gâteaux derivatives
$\partial_r \tilde{m}(\cdot, r, k) : T_{\mathcal{H}} \to L^2(\lambda)$
and $\partial_k \tilde{m}(\cdot, r, k) : T_{\mathcal{K}}(k) \to L^2(\lambda)$
such that, for all $h \in T_{\mathcal{H}}$ and $\eta \in T_{\mathcal{K}}(k)$,
\[
   \frac{d}{dt}\, E_0\!\left[\tilde{m}\big(A,S,\, r + t h,\,(1+t\eta)k\big)\right]\Big|_{t=0}
   = E_0\!\left[\partial_r \tilde{m}(A,S,r,k)(h) + \partial_k \tilde{m}(A,S,r,k)(\eta)\right].
\]
\item (\textit{Continuity of derivatives}) \label{cond::boundedfunc2::norm} The linear functionals $h \mapsto E_0[\partial_r \tilde{m}(A,S,r_{0,\nu,\gamma,g},k_0)(h)]$ and
    $\beta \mapsto E_0[\partial_k \tilde{m}(A,S,r_{0,\nu,\gamma,g},k_0)(\beta)]$ are continuous,
    with respective domains $(T_{\mathcal{H}}, \|\cdot\|_{\mathcal{H}, 0})$ and
    $(T_{\mathcal{K}}, \|\cdot\|_{\mathcal{K}, 0})$.
\item (\textit{Lipschitz continuity}) \label{cond::lipschitz::norm} The map $(r,k) \mapsto E_0[\tilde{m}(A,S,r,k)]$ is locally Lipschitz continuous at $(r_{0,\nu,\gamma,g},k_0)$ with respect to $(\mathcal{H}, \|\cdot\|_{\mathcal{H}, 0}) \times (\mathcal{K}, \|\cdot\|_{\mathcal{K}})$.
\end{enumerate}

\subsection{Softmax values under reward normalization}

Write \(\bar r_0 := r_{0,\nu,\gamma,g}\). Let
\(\pi_{\bar r_0,k_0}^{\star,\gamma'}\),
\(q_{\bar r_0,k_0}^{\star,\gamma'}\),
\(V_{\bar r_0,k_0}^{\star,\gamma'}\),
\(v_{\bar r_0,k_0}^{\star,\gamma'}\),
\(\Phi_{\bar r_0,k_0}^{\star,\gamma'}\),
\(\rho_{\bar r_0,k_0}^{\star,\gamma'}\),
\(\widetilde \rho_{\bar r_0,k_0}^{\star,\gamma'}\),
\(d_{\bar r_0,k_0}^{\star,\gamma'}\), and
\(\widetilde d_{\bar r_0,k_0}^{\star,\gamma'}\) denote the analogues of the
objects in Theorem~\ref{theorem::EIFsoftmax}, but with reward \(\bar r_0\) and
target discount \(\gamma'\).

\begin{theorem}[EIF for softmax value under reward normalization]
\label{theorem::EIFsoftmaxnorm}
Suppose the analogue of Condition~\ref{cond::stationary3} holds for
\(\pi_{\bar r_0,k_0}^{\star,\gamma'}\), and suppose the soft Bellman fixed point
is locally Lipschitz in the sense of Condition~\ref{cond::lipschitzvalue}, with
\(r_0\) replaced by \(\bar r_0\). Assume also that
Conditions~\ref{cond::boundedfunc::norm}--\ref{cond::lipschitz::norm} hold for
\(\tilde m(a,s,r,k)=V_{r,k}^{\star,\gamma'}(s)\). Then the normalized softmax
value parameter from Example~\ref{example::norm} is pathwise differentiable with
efficient influence function
\begin{align*}
\chi_0(s,a,s')
={}& \alpha_0(a,s)
+ \alpha_0(a,s)
   \bigl\{r_0(a,s)-g(s)+\gamma V_{r_0-g,k_0}^{\nu,\gamma}(s')
   -q_{r_0-g,k_0}^{\nu,\gamma}(a,s)\bigr\} \\
&+ \frac{\gamma'}{\tau^\star}
   \widetilde d_{\bar r_0,k_0}^{\star,\gamma'}(a,s)
   \bigl\{\Phi_{\bar r_0,k_0}^{\star,\gamma'}(s')
   -v_{\bar r_0,k_0}^{\star,\gamma'}(a,s)\bigr\} \\
&+ d_{\bar r_0,k_0}^{\star,\gamma'}(a,s)
   \bigl\{\bar r_0(a,s)+\gamma' V_{\bar r_0,k_0}^{\star,\gamma'}(s')
   -q_{\bar r_0,k_0}^{\star,\gamma'}(a,s)\bigr\}
+ V_{\bar r_0,k_0}^{\star,\gamma'}(s) - \Psi(P_0),
\end{align*}
where
\[
\alpha_0(a,s)
:=
\left\{\tau^{\star -1}
\widetilde \rho_{\bar r_0,k_0}^{\star,\gamma'}(s)
+ \rho_{\bar r_0,k_0}^{\star,\gamma'}(s)\right\}
\left\{
\frac{\pi_{\bar r_0,k_0}^{\star,\gamma'}(a\mid s)-\nu(a\mid s)}
{\pi_0(a\mid s)}
\right\}.
\]
\end{theorem}

\subsection{Structural coefficients under reward normalization}

The representation in Lemma \ref{lemma::linearfuncidentnorm} underlies the normalization argument from
Section~\ref{sec::eifnorm}. It shows that linear averages of the normalized
reward depend only on the anchored centered log-behavior-policy reward.

 
 

Under the classical normalization of \citet{rust1987optimal},
\(\nu(a\mid s)=\mathbbm{1}\{a=a^\star\}\), this gives
\[
E_0[w(A,S) r_0^\dagger(A,S)]
= E_0\!\big[w(A,S)\{r_0(A,S)-r_0(a^\star,S)\}\big].
\]
Thus linear functionals of the normalized reward, and the structural quantities
derived from them, are identified directly from \(r_0\), without imposing
parametric structure on \(r_0^\dagger\).

\begin{theorem}[EIF for structural coefficients under reward normalization]
\label{theorem::EIFcoefnorm}
Assume \(x:\mathcal A \times \mathcal S \to \mathbb R^p\) is bounded and
\(G_0 := E_0[x(A,S)x(A,S)^\top]\) is nonsingular. Let
\(\bar r_0 := r_{0,\nu,\gamma,g}\) and
\[
\alpha_0^x(a,s)
:=
x(a,s) - \frac{\nu(a\mid s)}{\pi_0(a\mid s)}
\sum_{\tilde a \in \mathcal A}\pi_0(\tilde a\mid s)x(\tilde a,s).
\]
Then the EIF for \(b_0 := E_0[x(A,S)\bar r_0(A,S)]\) is
\[
\begin{aligned}
\chi_{b,0}(s,a,s')
:={}& \alpha_0^x(a,s)
+ \alpha_0^x(a,s)\bigl\{r_0(a,s)-g(s)
+\gamma V_{r_0-g,k_0}^{\nu,\gamma}(s')
-q_{r_0-g,k_0}^{\nu,\gamma}(a,s)\bigr\} \\
&+ x(a,s)\bar r_0(a,s)-b_0,
\end{aligned}
\]
and the EIF for \(\theta_0 := G_0^{-1}b_0\) is
\[
\chi_{\theta,0}(s,a,s')
=
G_0^{-1}\!\left[
\chi_{b,0}(s,a,s')
- \{x(a,s)x(a,s)^\top - G_0\}\theta_0
\right].
\]
\end{theorem}

 


\section{Detailed statements for example estimators}
\label{app:estimator:details}
\label{app::examplesestimators}

This section groups the example-specific estimator displays, the additional
assumptions deferred from Section~\ref{sec::examplesestimators}, and the
resulting asymptotic conclusions. The main-text theorem labels remain
unchanged; we restate the first three results only in detailed, unnumbered
form for readability.

\subsection{Example~\ref{example::1a}: value of a fixed policy}

Let \(\pi_n:=\exp\{r_n\}\), let \(q_n^{\pi,\gamma}\) estimate
\(q_{r_0,k_0}^{\pi,\gamma}\), and define
\(V_n^{\pi,\gamma}(s):=\sum_{\tilde a\in\mathcal A}\pi(\tilde a\mid s)
q_n^{\pi,\gamma}(\tilde a,s)\). Let \(\rho_n^{\pi,\gamma}\) estimate the state
occupancy ratio \(\rho_0^{\pi,\gamma}\). The estimator of
\(\psi_0=E_0[V_{r_0,k_0}^{\pi,\gamma}(S)]\) is
\begin{align*}
\psi_n
&:= \frac{1}{n}\sum_{i=1}^n V_n^{\pi,\gamma}(S_i)
  + \frac{1}{n}\sum_{i=1}^n
  \rho_n^{\pi,\gamma}(S_i)
  \frac{\pi(A_i\mid S_i)}{\pi_n(A_i\mid S_i)}
  \left\{
  r_n(A_i,S_i)+\gamma V_n^{\pi,\gamma}(S_i')
  -q_n^{\pi,\gamma}(A_i,S_i)
  \right\} \\
&\quad
  + \frac{1}{n}\sum_{i=1}^n
  \rho_n^{\pi,\gamma}(S_i)
  \left\{
  1-\frac{\pi(A_i\mid S_i)}{\pi_n(A_i\mid S_i)}
  \right\}.
\end{align*}

The additional conditions for Theorem~\ref{theorem::ALex1} are stated below.
\begin{enumerate}[label=\textbf{(E\arabic*)}, ref=E\arabic*, resume = efficiency1a]
    \item \label{assump:boundedness::policy}(\textit{Boundedness})
    There exists $M \in (0,\infty)$ such that
    $\|\rho_n^{\pi,\gamma}\|_{L^\infty(\mu)}
      + \|q_n^{\pi,\gamma}\|_{L^\infty(\lambda)}
      + \|\pi^{-1}\|_{L^\infty(\lambda)} < M$
    for all $n$ almost surely.

    \item \label{assump:splitting::policy}(\textit{Sample splitting})
    The estimators $r_n$, $\rho_n^{\pi,\gamma}$, $q_n^{\pi,\gamma}$ are constructed using data independent of $\mathcal{D}_n$.

    \item \label{assump:consistency::policy}(\textit{Consistency})
   $ \|\rho_n^{\pi,\gamma} - \rho_0^{\pi,\gamma}\|_{L^2(\rho_0)}
      + \|r_n - r_0\|_{\mathcal{H},r_0} +
    \bigl\|
        q_n^{\pi,\gamma} - q_0^{\pi,\gamma}
    \bigr\|_{\mathcal{H},r_0}
    = o_p(1)$.
\end{enumerate}

\medskip\noindent\textbf{Detailed version of Theorem~\ref{theorem::ALex1}.}
Assume \ref{cond::stationary}, \ref{cond::rates}\ref{cond::reward},
\ref{assump:rates::policy}, and
\ref{assump:boundedness::policy}--\ref{assump:consistency::policy}. Then
Theorem~\ref{thm:asymptotic_linearity} holds for the estimator above with
influence function \(\chi_0\) given in Theorem~\ref{theorem::EIFvalue}.

\subsection{Example~\ref{example::1b}: value of a counterfactual softmax policy}

Let \(r_n\) estimate \(r_0\), let \(v_n^\star\) estimate the soft Bellman
fixed point \(v_{r_0,k_0}^\star\), and define
\[
\pi_n^\star(a\mid s)
\propto
1\{a\in\mathcal A^\star\}
\exp\!\left\{
\tfrac{1}{\tau^\star}\{r_n(a,s)+\gamma v_n^\star(a,s)\}
\right\}.
\]
Let \(q_n^\star\) estimate \(q_{r_0,k_0}^\star\), set
\[
V_n^\star(s)
:=
\sum_{\tilde a\in\mathcal A^\star}
\pi_n^\star(\tilde a\mid s)q_n^\star(\tilde a,s),
\]
and let \(\widetilde\rho_n^\star\) and \(\rho_n^\star\) estimate
\(\widetilde\rho_{r_0,k_0}^\star\) and \(\rho_{r_0,k_0}^\star\). The resulting
estimator of \(\psi_0=E_0\{V_{r_0,k_0}^\star(S)\}\) is
\begin{align*}
\psi_n
&:=
\frac{1}{n}\sum_{i=1}^n V_n^\star(S_i)
+
\frac{1}{n}\sum_{i=1}^n
\rho_n^\star(S_i)
\frac{\pi_n^\star(A_i\mid S_i)}{\pi_n(A_i\mid S_i)}
\{r_n(A_i,S_i)+\gamma V_n^\star(S_i')-q_n^\star(A_i,S_i)\}
\\
&\quad+
\frac{\gamma}{n\tau^\star}\sum_{i=1}^n
\widetilde\rho_n^\star(S_i)
\frac{\pi_n^\star(A_i\mid S_i)}{\pi_n(A_i\mid S_i)}
\left[
\tau^\star
\log\!\left\{
\sum_{a\in\mathcal A^\star}
\exp\!\left(
\tfrac{1}{\tau^\star}\{r_n(a,S_i')+\gamma v_n^\star(a,S_i')\}
\right)
\right\}
-
v_n^\star(A_i,S_i)
\right]
\\
&\quad+
\frac{1}{n}\sum_{i=1}^n
\left\{
\tfrac{1}{\tau^\star}\widetilde\rho_n^\star(S_i)
+
\rho_n^\star(S_i)
\right\}
\left\{
\frac{\pi_n^\star(A_i\mid S_i)}{\pi_n(A_i\mid S_i)}-1
\right\}.
\end{align*}

The additional conditions for Theorem~\ref{theorem::ALex1b} are stated below.
\begin{enumerate}[label=\textbf{(E\arabic**)}, ref=E\arabic**, resume = efficiency1b]
    \item \label{assump:boundedness::softstar}(\textit{Boundedness})
    There exists $M \in (0,\infty)$ such that, for all $n$ almost surely,
    \[
      \|\rho_n^\star\|_{L^\infty(\mu)}
      + \|\widetilde{\rho}_n^\star\|_{L^\infty(\mu)}
      + \|q_n^\star\|_{L^\infty(\lambda)}
      + \|v_n^\star\|_{L^\infty(\lambda)}
      + \|r_n\|_{L^\infty(\lambda)}
      + \|\pi^{-1}\|_{L^\infty(\lambda)}
      < M.
    \]

    \item \label{assump:splitting::softstar}(\textit{Sample splitting})
    The nuisance estimators $\rho_n^\star$, $\widetilde{\rho}_n^\star$, $q_n^\star$, $v_n^\star$, and $r_n$
    are constructed using data independent of $\mathcal{D}_n$ (or via cross-fitting).

    \item \label{assump:consistency::softstar}(\textit{Consistency})
    \[
      \|\rho_n^\star - \rho_{r_0,k_0}^\star\|_{L^2(\rho_0)}
      + \|\widetilde{\rho}_n^\star - \widetilde{\rho}_{r_0,k_0}^\star\|_{L^2(\rho_0)}
      + \|q_n^\star - q_{r_0,k_0}^\star\|_{\mathcal{H},r_0}
      + \|v_n^\star - v_{r_0,k_0}^\star\|_{\mathcal{H},r_0}
      + \|r_n - r_0\|_{\mathcal{H},r_0}
      = o_p(1).
    \]
\end{enumerate}

\medskip\noindent\textbf{Detailed version of Theorem~\ref{theorem::ALex1b}.}
Assume \ref{cond::stationary3}, \ref{cond::rates}\ref{cond::reward},
\ref{assump:rates::softstar}, and
\ref{assump:boundedness::softstar}--\ref{assump:consistency::softstar}. Then
Theorem~\ref{thm:asymptotic_linearity} holds for the estimator above with
influence function \(\chi_0\) given in Theorem~\ref{theorem::EIFsoftmax}.

\subsection{Example~\ref{example::norm}: fixed-policy value under reward normalization}

In addition to the example-specific conditions below, the normalized-reward EIF
used in Theorem~\ref{theorem::EIFvaluenorm} relies on
Conditions~\ref{cond::boundedfunc::norm}--\ref{cond::lipschitz::norm} from
Appendix~\ref{appendix::conditionsEIFnorm}.

Let \(\pi_n:=\exp\{r_n\}\), let \(q_n^{\nu,\gamma}\) estimate
\(q_{r_0-g,k_0}^{\nu,\gamma}\), and define
\(V_n^{\nu,\gamma}(s):=\sum_{\tilde a\in\mathcal A}\nu(\tilde a\mid s)\,
q_n^{\nu,\gamma}(\tilde a,s)\) and
\(r_{n,\nu,\gamma,g}(a,s):=g(s)+q_n^{\nu,\gamma}(a,s)-V_n^{\nu,\gamma}(s)\).
Let \(q_{n,\nu}^{\pi,\gamma'}\) estimate
\(q_{r_0,k_0,\nu}^{\pi,\gamma'}\), set
\(V_{n,\nu}^{\pi,\gamma'}(s):=\sum_{\tilde a\in\mathcal A}\pi(\tilde a\mid s)\,
q_{n,\nu}^{\pi,\gamma'}(\tilde a,s)\), and let \(\rho_n^{\pi,\gamma'}\)
estimate \(\rho_0^{\pi,\gamma'}\). The corresponding estimator of
\(\psi_0=E_0[V_{r_0,k_0,\nu}^{\pi,\gamma'}(S)]\) is
\begin{align*}
\psi_n
&:= \frac{1}{n}\sum_{i=1}^n V_{n,\nu}^{\pi,\gamma'}(S_i) \\
&\quad + \frac{1}{n}\sum_{i=1}^n \rho_n^{\pi,\gamma'}(S_i)\,
        \frac{\pi(A_i \mid S_i)}{\pi_n(A_i \mid S_i)}
        \bigl\{r_{n,\nu,\gamma,g}(A_i,S_i)
        + \gamma' V_{n,\nu}^{\pi,\gamma'}(S_i')
        - q_{n,\nu}^{\pi,\gamma'}(A_i,S_i)\bigr\} \\
&\quad + \frac{1}{n}\sum_{i=1}^n \rho_n^{\pi,\gamma'}(S_i)
        \left(
            \frac{\pi(A_i \mid S_i)}{\pi_n(A_i \mid S_i)}
            -
            \frac{\nu(A_i \mid S_i)}{\pi_n(A_i \mid S_i)}
        \right)
        \bigl\{r_n(A_i,S_i) - g(S_i)
        + \gamma V_n^{\nu,\gamma}(S_i')
        - q_n^{\nu,\gamma}(A_i,S_i)\bigr\} \\
&\quad + \frac{1}{n}\sum_{i=1}^n \rho_n^{\pi,\gamma'}(S_i)
        \left(
            \frac{\pi(A_i \mid S_i)}{\pi_n(A_i \mid S_i)}
            -
            \frac{\nu(A_i \mid S_i)}{\pi_n(A_i \mid S_i)}
        \right).
\end{align*}

The additional conditions for Theorem~\ref{theorem::ALex2} are stated below.
\begin{enumerate}[label=\textbf{(E\arabic***)}, ref=E\arabic***, resume = efficiency2]
    \item \label{assump:boundedness::policynorm}(\textit{Boundedness})
    There exists $M \in (0,\infty)$ such that, for all $n$ almost surely,
    \[
      \|\rho_n^{\pi,\gamma'}\|_{L^\infty(\mu)}
      + \|q_{n,\nu}^{\pi,\gamma'}\|_{L^\infty(\lambda)}
      + \|q_n^{\nu,\gamma}\|_{L^\infty(\lambda)}
      + \|r_n\|_{L^\infty(\lambda)}
      + \|\pi^{-1}\|_{L^\infty(\lambda)}
      + \|\nu^{-1}\|_{L^\infty(\lambda)}
      < M.
    \]

    \item \label{assump:splitting::policynorm}(\textit{Sample splitting})
    The nuisance estimators
    $r_n$, $\rho_n^{\pi,\gamma'}$, $q_{n,\nu}^{\pi,\gamma'}$, and $q_n^{\nu,\gamma}$
    are constructed using data independent of $\mathcal{D}_n$.

    \item \label{assump:consistency::policynorm}(\textit{Consistency})
   $ \|\rho_n^{\pi,\gamma'} - \rho_0^{\pi,\gamma'}\|_{L^2(\rho_0)}
      + \|r_n - r_0\|_{\mathcal{H},r_0} +
    \bigl\|
       q_n^{\nu,\gamma} - q_0^{\nu,\gamma}
    \bigr\|_{\mathcal{H},r_0} +   \bigl\|
      q_{n,\nu}^{\pi,\gamma'} - q_{0,\nu}^{\pi,\gamma'}
    \bigr\|_{\mathcal{H},r_0}
    = o_p(1)$.
\end{enumerate}

\medskip\noindent\textbf{Detailed version of Theorem~\ref{theorem::ALex2}.}
Assume \ref{cond::stationary}, \ref{cond::stationary2},
\ref{cond::rates}\ref{cond::reward},
\ref{assump:rates::policynorm}, and
\ref{assump:boundedness::policynorm}--\ref{assump:consistency::policynorm}.
Then Theorem~\ref{thm:asymptotic_linearity} holds for the estimator above with
influence function \(\chi_0\) given in Theorem~\ref{theorem::EIFvaluenorm}.

\subsection{Example~\ref{example::normcoef}: structural coefficients under reward normalization}

To estimate \(\theta_0\) in Example~\ref{example::normcoef}, define
\(\bar r_n(a,s):=g(s)+q_n^{\nu,\gamma}(a,s)-V_n^{\nu,\gamma}(s)\),
\(\pi_n:=\exp\{r_n\}\),
\(V_n^{\nu,\gamma}(s):=\sum_{\tilde a}\nu(\tilde a\mid s)
q_n^{\nu,\gamma}(\tilde a,s)\), and
\[
\alpha_n^x(a,s)
:=
x(a,s) - \frac{\nu(a\mid s)}{\pi_n(a\mid s)}
\sum_{\tilde a \in \mathcal A}\pi_n(\tilde a\mid s)x(\tilde a,s).
\]
Set
\[
\begin{aligned}
b_n
:=
\frac{1}{n}\sum_{i=1}^n \Big[
&x(A_i,S_i)\bar r_n(A_i,S_i)
+ \alpha_n^x(A_i,S_i) \\
&+ \alpha_n^x(A_i,S_i)\{r_n(A_i,S_i)-g(S_i)
+\gamma V_n^{\nu,\gamma}(S_i')
-q_n^{\nu,\gamma}(A_i,S_i)\}
\Big],
\end{aligned}
\]
and \(G_n:=n^{-1}\sum_{i=1}^n x(A_i,S_i)x(A_i,S_i)^\top\),
\(\theta_n:=G_n^{-1}b_n\). The example-specific rate condition from the main
text is
\begin{enumerate}[label=\textbf{(E\arabic****)}, ref=E\arabic****, series=efficiency3]
\item \label{assump:rates::coefnorm} (\textit{Nuisance rates})
\[
\|r_n-r_0\|_{\mathcal H,r_0}\,
\bigl\|\mathcal T_{k_0,\nu,\gamma}(q_n^{\nu,\gamma}-q_0^{\nu,\gamma})\bigr\|_{\mathcal H,r_0}
= o_p(n^{-1/2}).
\]
\end{enumerate}

The additional conditions for Theorem~\ref{theorem::ALex3} are stated below.
\begin{enumerate}[label=\textbf{(E\arabic****)}, ref=E\arabic****, resume = efficiency3]
    \item \label{assump:boundedness::coefnorm}(\textit{Boundedness})
    There exists $M \in (0,\infty)$ such that, for all $n$ almost surely,
    \[
      \|q_n^{\nu,\gamma}\|_{L^\infty(\lambda)}
      + \|r_n\|_{L^\infty(\lambda)}
      < M.
    \]

    \item \label{assump:splitting::coefnorm}(\textit{Sample splitting})
    The nuisance estimators $r_n$ and $q_n^{\nu,\gamma}$ are constructed using
    data independent of $\mathcal{D}_n$.

    \item \label{assump:consistency::coefnorm}(\textit{Consistency})
    \[
      \|r_n - r_0\|_{\mathcal{H},r_0}
      + \|q_n^{\nu,\gamma} - q_0^{\nu,\gamma}\|_{\mathcal{H},r_0}
      = o_p(1).
    \]
\end{enumerate}

\begin{theorem}[Asymptotic linearity for Example~\ref{example::normcoef}]
\label{theorem::ALex3}
Assume \ref{cond::rates}\ref{cond::reward} and
\ref{assump:rates::coefnorm}--\ref{assump:consistency::coefnorm}. Then
\[
\theta_n - \theta_0
= \frac{1}{n}\sum_{i=1}^n \chi_{\theta,0}(S_i,A_i,S_i') + o_p(n^{-1/2}),
\qquad
\sqrt n(\theta_n-\theta_0)\rightsquigarrow N(0,\Sigma_0),
\]
where \(\Sigma_0 := \operatorname{Var}_{P_0}\{\chi_{\theta,0}(S,A,S')\}\).
\end{theorem}

\section{Additional simulation results for Example~\ref{example::norm}}
\label{app::normsim}

We report three versions of the estimator for Example~\ref{example::norm}: a
\emph{practical} version that estimates the behaviour policy and solves the two
Bellman equations using the true transition law, a \emph{quasi-oracle} version
that uses the true \(Q\) nuisances and true occupancy ratio, and a
\emph{full-oracle} benchmark that also uses the true behaviour policy. For the
oracle rows, we report only the D-IRL estimator, since a plug-in estimator
with true nuisances is not an informative comparison.
Table~\ref{tab:simulation-ex2} reports the main \(n \in \{5000,10000\}\)
results.

\begin{table}[htb]
\centering
\caption{Example~\ref{example::norm}: practical and oracle benchmarks.
Coverage refers to the nominal $95\%$ Wald interval for D-IRL.}
\label{tab:simulation-ex2}
\begingroup
\small
\setlength{\tabcolsep}{4pt}
\begin{tabular}{llrrrrrrr}
\toprule
Method & $n$ & Plug-in bias & Plug-in SD & Plug-in RMSE & D-IRL bias & D-IRL SD & D-IRL RMSE & Coverage \\
\midrule
Practical & 5000 & -0.032 & 0.480 & 0.480 & 0.100 & 0.443 & 0.454 & 0.826 \\
Practical & 10000 & -0.017 & 0.372 & 0.372 & 0.084 & 0.300 & 0.311 & 0.770 \\
Quasi-oracle & 5000 & --- & --- & --- & -0.064 & 0.316 & 0.322 & 0.862 \\
Quasi-oracle & 10000 & --- & --- & --- & -0.035 & 0.165 & 0.169 & 0.930 \\
Full oracle & 5000 & --- & --- & --- & -0.001 & 0.215 & 0.214 & 0.950 \\
Full oracle & 10000 & --- & --- & --- & 0.007 & 0.143 & 0.143 & 0.974 \\
\bottomrule
\end{tabular}
\endgroup
\end{table}

The normalization experiment is the most challenging setting in our study. In
the practical implementation, D-IRL improves on the plug-in estimator in RMSE
at both \(n=5{,}000\) and \(n=10{,}000\), but interval performance remains
weaker than in Examples~\ref{example::1a} and~\ref{example::1b}. By contrast,
the quasi-oracle and full-oracle benchmarks show that inference improves
markedly once the hard nuisance components are supplied accurately: coverage
rises to \(0.930\)--\(0.974\) at \(n=10{,}000\), and the full-oracle estimator
already reaches \(0.950\) at \(n=5{,}000\). The small-sample stress tests at
\(n \in \{1000,2500\}\) show the same pattern with much larger variance.




 
\section{Expanded related work}
\label{app::relatedwork}

 

\noindent \textbf{Dynamic discrete choice models.} Existing econometric work on
DDC and dynamic games is predominantly parametric, focusing on identification
and efficient estimation of finite-dimensional structural reward (utility)
parameters. Classical approaches include the nested fixed-point method of
\citet{rust1987optimal}, the two-step conditional choice probability approach
of \citet{hotz1993conditional}, and simulation-based estimators
\citep{hotz1994simulation,aguirregabiria2010dynamic}. These methods do not
allow for fully nonparametric reward models or target the reward-dependent
value functionals—such as policy values and value differences—central to IRL.
Our work extends this line of research to settings with nonparametric reward
representations and modern machine-learning-based nuisance estimation. Unlike
classical DDC estimators, our framework avoids repeated dynamic programming or
Monte Carlo simulation inside the estimator, instead reducing inference to
standard supervised-learning and offline reinforcement learning components. To
our knowledge, there is no existing derivation of efficient influence functions
for such functionals in softmax IRL or in DDC models with nonparametric
rewards, nor any debiased, machine-learning-based inference framework for this
setting.

While identification and semiparametric efficiency have been studied extensively
in DDC models, this work concerns largely parametric utility specifications. For
dynamic discrete games, \citet{bajari2015dynamic} derive identification and
efficiency results when per-period rewards are parametrized by a
finite-dimensional vector, treating equilibrium objects and transition
probabilities nonparametrically. Foundational identification results include
\citet{hotz1993conditional}, \citet{hotz1994simulation}, and
\citet{magnac2002identifying}. \citet{blundell2007non} develop nonparametric
specification tests for dynamic discrete choice models, and
\citet{blevins2014nonparametric} study nonparametric identification of dynamic
decision processes with discrete and continuous choices. Other work addresses
latent heterogeneity or shock distributions: \citet{kasahara2009nonparametric}
analyze identification with unobserved heterogeneity,
\citet{buchinsky2010semiparametric} propose a diagnostic for zero semiparametric
information in fixed-effects DDC models, and \citet{buchholz2021semiparametric}
consider models with unspecified shock distributions. Additional semiparametric
results include \citet{arcidiacono2011conditional} and \citet{norets2014semiparametric}.

A complementary line of research studies identification of counterfactuals under
minimal structure. \citet{kalouptsidi2021identification} show that certain
counterfactuals may be point- or set-identified even when utilities are not.
Other identification results include \citet{abbring2020identifying} on
discount-factor identification and \citet{kalouptsidi2017non} on
nonidentification of counterfactuals in dynamic games. These papers consider
broader DDC settings but do not yield the exact identification of value
differences or normalized rewards that arises under the softmax (Gumbel-shock)
framework we impose. Following \citet{rust1987optimal}, we adopt the widely used Gumbel
specification. Extending our analysis to more general shock distributions is an
interesting direction for future work.


 




 

\section{Proofs for Section \ref{sec:ident}}

 
 


\begin{proof}[Proof of Lemma \ref{theorem::trivialsol}]
The key observation is that when \(r_0=\log\pi_0\), the log-partition term is
zero, so \(v_{r_0,k_0}\equiv 0\) solves the Bellman constraint. We then compare
the induced softmax representation with that of \(r_0^\dagger\) to identify the
shaping potential. By definition, we have
\begin{align*}
\sum_{a \in \mathcal{A}} \exp\{ r_0(a,s) \}
&= \sum_{a \in \mathcal{A}} \exp\{ \log \pi_0(a \mid s) \}  \\
&= \sum_{a \in \mathcal{A}} \pi_0(a \mid s) \\
&= 1.
\end{align*}
Hence, $\log \sum_{a \in \mathcal{A}} \exp\{ r_0(a,s) \} = 0,$
and
\begin{align*}
0 = \int \log \sum_{a' \in \mathcal{A}} \exp\{ r_0(a',s') + \gamma \cdot 0 \} \,
k_0(s' \mid a, s) \,\mu(ds').
\end{align*}
It follows that \(0\) solves the soft Bellman equation with respect to the
reward \(r_0\). Moreover, the induced softmax policy satisfies
\[
\frac{\exp\{r_0(a,s)\}}{\sum_{\tilde a \in \mathcal{A}} \exp\{r_0(\tilde a,s)\}}
= \exp\{r_0(a,s)\}
= \pi_0(a \mid s),
\]
so it coincides with the true behaviour policy and therefore maximizes the
population log-likelihood. Thus, \(r_0\) is also a solution to the soft
Bellman-constrained log-likelihood in \eqref{eqn::softbellmanloglik}.


Furthermore, we have
\[
\exp\{r_0(a,s)\} = \pi_0(a \mid s) 
= \frac{\exp\{r_0^\dagger(a,s) + \gamma v_0^\dagger(a,s)\}}{\sum_{\tilde a \in \mathcal{A}} \exp\{r_0^\dagger(\tilde a,s) + \gamma v_0^\dagger(\tilde a,s)\}}.
\]
Taking logarithms of both sides yields
\[
r_0(a,s) = r_0^\dagger(a,s) + \gamma v_0^\dagger(a,s) 
- \log \sum_{\tilde a \in \mathcal{A}} \exp\{r_0^\dagger(\tilde a,s) + \gamma v_0^\dagger(\tilde a,s)\}.
\]
The soft value function is defined as
\[
V_{0}^{\dagger}(s) = \log \Bigg( \sum_{a' \in \mathcal{A}} 
\exp\big\{ r_0^\dagger(a',s) + \gamma v_0^\dagger(a', s)\big\} \Bigg).
\]   
By the soft Bellman equation,
\[
v_0^\dagger(a,s)  = \int  V_{0}^{\dagger}(s') \, k_0(s' \mid a , s) \, \mu(ds').
\]
Substituting this expression, we obtain
\[
r_0(a,s) = r_0^\dagger(a,s) + \gamma \int V_{0}^{\dagger}(s') \, k_0(s' \mid a , s) \, \mu(ds') - V_{0}^{\dagger}(s).
\]


 

    
\end{proof}




\begin{proof}[Proof of Theorem \ref{theorem::identWithNormalization}]
For any reward \(r\), let \(v_{r,k_0}\) solve the soft Bellman equation and set
\(q_r(a,s):=r(a,s)+\gamma v_{r,k_0}(a,s)\) and
\(\Lambda_r(s):=\log\sum_{a\in\mathcal A}\exp\{q_r(a,s)\}\). Then the criterion in
\eqref{eqn::softbellmanloglik} equals
\[
E_0\{\Lambda_r(S)-q_r(A,S)\}
=
E_0\{-\log \pi_r(A\mid S)\},
\qquad
\pi_r(a\mid s):=\exp\{q_r(a,s)-\Lambda_r(s)\}.
\]
Since \(A\mid S\sim\pi_0(\cdot\mid S)\),
\[
E_0\{-\log \pi_r(A\mid S)\}
=
E_0\!\left[\mathrm{KL}\{\pi_0(\cdot\mid S),\pi_r(\cdot\mid S)\}\right]
-
E_0\!\left[\sum_a \pi_0(a\mid S)\log \pi_0(a\mid S)\right].
\]
Thus \(r\) minimizes \eqref{eqn::softbellmanloglik} if and only if
\(\pi_r(\cdot\mid s)=\pi_0(\cdot\mid s)\) for \(P_0\)-a.e. \(s\).

Let \(\bar r_0:=r_0-g\) and
\(q:=q_{r_0-g,k_0}^{\nu,\gamma}
=\mathcal T_{k_0,\nu,\gamma}^{-1}(\bar r_0)\), so
\(q=\bar r_0+\gamma\mathcal P_{k_0}\nu q\). Define
\(c^\star:=-\nu q\), \(r^\star:=g+(I-\nu)q\), and
\(v^\star:=\mathcal P_{k_0}c^\star=-\mathcal P_{k_0}\nu q\). Since
\(q-\gamma\mathcal P_{k_0}\nu q=\bar r_0\),
\[
r^\star+\gamma v^\star
=
g+(I-\nu)q-\gamma\mathcal P_{k_0}\nu q
=
r_0+c^\star .
\]
Therefore
\[
\Lambda_{r^\star}(s)
=
\log\sum_{a\in\mathcal A}\exp\{r_0(a,s)+c^\star(s)\}
=
c^\star(s),
\]
because \(\sum_a e^{r_0(a,s)}=\sum_a\pi_0(a\mid s)=1\). Hence
\(v^\star=\mathcal P_{k_0}c^\star=\mathcal P_{k_0}\Lambda_{r^\star}\), so
\(v^\star\) solves the soft Bellman equation for \(r^\star\). Also
\[
\pi_{r^\star}(a\mid s)
=
\exp\{r_0(a,s)+c^\star(s)-\Lambda_{r^\star}(s)\}
=
\pi_0(a\mid s),
\]
so \(r^\star\) minimizes \eqref{eqn::softbellmanloglik}. Finally,
\(\nu r^\star=\nu g+\nu(I-\nu)q=g\), and therefore \(r^\star\) is normalized.

For uniqueness, let \(\tilde r\in\mathcal H\) be any normalized minimizer, and
let \(\tilde v\) solve its soft Bellman equation. Since \(\tilde r\) induces
\(\pi_0\), we have
\[
\tilde r(a,s)+\gamma\tilde v(a,s)-\tilde\Lambda(s)=r_0(a,s),
\qquad
\tilde\Lambda(s):=\log\sum_{a\in\mathcal A}
\exp\{\tilde r(a,s)+\gamma\tilde v(a,s)\}.
\]
Thus \(\tilde r=r_0+\tilde\Lambda-\gamma\tilde v\). Since
\(\tilde v=\mathcal P_{k_0}\tilde\Lambda\), the normalization
\(\nu\tilde r=g\) gives
\[
\tilde\Lambda
=
-\nu(r_0-g)+\gamma\nu\mathcal P_{k_0}\tilde\Lambda .
\]
The right-hand side is a \(\gamma\)-contraction in sup norm. Thus this
fixed-point equation has a unique bounded solution. Applying \(\nu\) to
\(q=\bar r_0+\gamma\mathcal P_{k_0}\nu q\) shows that
\(c^\star=-\nu q\) satisfies the same equation,
\(c^\star=-\nu(r_0-g)+\gamma\nu\mathcal P_{k_0}c^\star\). Hence
\(\tilde\Lambda=c^\star\), \(\tilde v=\mathcal P_{k_0}\tilde\Lambda
=-\mathcal P_{k_0}\nu q\), and
\[
\tilde r
=
r_0+\tilde\Lambda-\gamma\tilde v
=
r_0-\nu q+\gamma\mathcal P_{k_0}\nu q
=
g+q-\nu q
=
r^\star .
\]
Thus the normalized solution is unique.

Writing \(q=q_{r_0-g,k_0}^{\nu,\gamma}\) statewise gives
\[
r^\star(a,s)
=
g(s)+q_{r_0-g,k_0}^{\nu,\gamma}(a,s)
-\sum_{\tilde a\in\mathcal A}\nu(\tilde a\mid s)
q_{r_0-g,k_0}^{\nu,\gamma}(\tilde a,s).
\]
Moreover, if \(\gamma>0\), then
\[
v^\star(a,s)
=
-\mathcal P_{k_0}\nu q(a,s)
=
\gamma^{-1}\{r_0(a,s)-g(s)-q_{r_0-g,k_0}^{\nu,\gamma}(a,s)\}.
\]
Hence \(r^\star=r_{0,\nu,\gamma,g}\) and
\(v^\star=v_{0,\nu,\gamma,g}\), as claimed.
\end{proof}


\begin{proof}[Proof of Theorem \ref{theorem::identvalue}]
 By Lemma~\ref{theorem::trivialsol},
\[
r_0(a,s) \;=\; r_0^\dagger(a,s) + \gamma \int V_0^\dagger(s') \, k_0(s' \mid a,s)\,\mu(ds') - V_0^\dagger(s).
\]
By the definition of the $Q$-function under policy $\pi$, we have
\begin{align*}
q^{\pi,\gamma}_{r_0,k_0}(a,s) 
&= r_0(a,s) + \gamma \int \sum_{a' \in \mathcal{A}} \pi(a' \mid s') \, q^{\pi,\gamma}_{r_0,k_0}(a',s') \, k_0(s' \mid a,s) \,\mu(ds').
\end{align*}
Substituting the expression for $r_0(a,s)$ gives
\begin{align*}
q^{\pi,\gamma}_{r_0,k_0}(a,s) 
&= r_0^\dagger(a,s) + \gamma \int V_0^\dagger(s') \, k_0(s' \mid a,s)\,\mu(ds') - V_0^\dagger(s) \\
&\quad + \gamma \int \sum_{a' \in \mathcal{A}} \pi(a' \mid s') \, q^{\pi,\gamma}_{r_0,k_0}(a',s') \, k_0(s' \mid a,s)\,\mu(ds').
\end{align*}
Rearranging terms,
\[
q^{\pi,\gamma}_{r_0,k_0}(a,s) + V_0^\dagger(s)
= r_0^\dagger(a,s) + \gamma \int \sum_{a' \in \mathcal{A}} \pi(a' \mid s') \big(q^{\pi,\gamma}_{r_0,k_0}(a',s') + V_0^\dagger(s')\big)\,k_0(s' \mid a,s)\,\mu(ds').
\]
Thus, \(q^{\pi,\gamma}_{r_0,k_0}(a,s) + V_0^\dagger(s)\) satisfies the Bellman equation with reward \(r_0^\dagger\).  
By uniqueness of the solution, it follows that
\[
q^{\pi,\gamma}_{r_0,k_0}(a,s) + V_0^\dagger(s) \;=\; q^{\pi,\gamma}_{r_0^\dagger,k_0}(a,s),
\]
or equivalently,
\[
q^{\pi,\gamma}_{r_0,k_0}(a,s) \;=\; q^{\pi,\gamma}_{r_0^\dagger,k_0}(a,s) - V_0^\dagger(s).
\]

Taking expectations with respect to the initial state distribution $\rho_0$ and the policy $\pi$, we obtain
\[
V_{r_0,k_0}(\pi; \gamma) \;=\; V_{r_0^\dagger,k_0}(\pi; \gamma) - E_0[V_0^\dagger(S)].
\]

Finally, for any two policies $\pi_1$ and $\pi_2$,
\begin{align*}
V_{r_0,k_0}(\pi_1; \gamma) - V_{r_0,k_0}(\pi_2; \gamma) 
&= \big(V_{r_0^\dagger,k_0}(\pi_1; \gamma) - E_0[V_0^\dagger(S)]\big) 
 - \big(V_{r_0^\dagger,k_0}(\pi_2; \gamma) - E_0[V_0^\dagger(S)]\big) \\
&= V_{r_0^\dagger,k_0}(\pi_1; \gamma) - V_{r_0^\dagger,k_0}(\pi_2; \gamma).
\end{align*}
This establishes the claim.
\end{proof}
 




\section{Proofs for Section \ref{sec::EIFsub}}

 \label{appendix::EIFsub}

We first derive general differentiability facts for the parameter map and then
assemble them into the EIF and von Mises expansion stated in
Section~\ref{sec::EIFsub}.


%In this case, we may consider the nonparametrically defined log-likelihood projection parameter $\Psi: \mathcal{M} \to \mathbb{R}$,
%\[
%\Psi(P) := E_P[m(A,S,r_P,k_P)],
%\]
%where $k_P$ is the state transition kernel under $P$ and $r_P$ is the solution to the multinomial logit problem
%\begin{equation*} 
%\begin{aligned}
%r_P 
%&:= \argmin_{r \in \mathcal{H}} \; -E_P[r(A,S)] \\
%&\quad \text{subject to} \quad 
%\sum_{a \in \mathcal{A}} \exp\{r(a,s)\} = 1,
%\quad \mu\text{-a.e.\ } s \in \mathcal{S}.
%\end{aligned}
%\end{equation*}
%Note that $r_P$ is defined exactly as in \eqref{sec::statobj}, except that $\mathcal{H}$ may be a strict subset of $L^\infty(\lambda)$.

The following lemma shows that $r_P$ is a solution to a moment-matching problem. In the main text, we set $\mathcal{H} := L^\infty(\lambda)$. For generality, the lemma below allows \(\mathcal{H} \subseteq L^\infty(\lambda)\) to be a linear subspace that includes the state-only functions, i.e., $L^\infty(\mu)$. For example, $\mathcal{H}$ may correspond to an additive or partially linear model.



 
We will use that $r_P$ is the solution to the multinomial logit 
problem
\begin{equation}\label{eqn::Mestimand}
\begin{aligned}
r_P 
&:= \argmin_{r \in \mathcal{H}} \; -E_P[r(A,S)] \\
&\quad \text{subject to} \quad 
\sum_{a \in \mathcal{A}} \exp\{r(a,s)\} = 1,
\qquad \mu\text{-a.e.\ } s \in \mathcal{S}.
\end{aligned}
\end{equation}






\begin{lemma}[Softmax moment matching and identification]\label{thm:softmax-moment-matching}
Let $\mathcal H \subset L^\infty(\lambda)$ be a linear class that contains all state-only functions, i.e., $L^\infty(\mu)$. Let $r_P$ denote the solution to \eqref{eqn::Mestimand} for this choice of $\mathcal{H}$. Then
\[
E_P[h(A,S)]
= E_{P}\!\Big[\sum_{a\in\mathcal A} \exp\{ r_P(a,S) \}\,h(a,S)\Big],
\quad \forall\, h\in\mathcal H .
\]
\end{lemma}

\begin{proof}[Proof of Lemma \ref{thm:softmax-moment-matching}]
The equality-constrained optimization problem defining $r_P$ in \eqref{eqn::Mestimand} is equivalent to the convex inequality-constrained problem
\begin{align}
    r_P 
    &:= \argmin_{r \in \mathcal{H}} \; - E_P\!\left[r(A,S)\right] \nonumber \\
    &\quad \text{subject to} \quad 
    \sum_{a \in \mathcal{A}} \exp\!\left\{ r(a,s) \right\} \;\leq\; 1,
    \quad \mu\text{-a.e. } s \in \mathcal{S}. \label{eq:mm-main}
\end{align}
The solution satisfies the boundary constraint
$$ \sum_{a \in \mathcal{A}} \exp\!\left\{ r_P(a,s) \right\} \;=\; 1,
    \quad \mu\text{-a.e. } s \in \mathcal{S}. $$

 
Since $\mathcal H$ contains constants, Slater’s condition holds (take $r\equiv -M$ with $M$ large), so strong duality and KKT apply. Let $\lambda\in L^1_+(\mu)$ be the Lagrange multiplier for the pointwise constraint, with Lagrangian
\[
\mathcal L(r,\lambda)= - E_P[r(A,S)] + \int \lambda(s)\!\left(\sum_{a} e^{r(a,s)} - 1\right)\mu(ds).
\]
Stationarity at $r_P$ in any direction $h\in\mathcal H$ gives
\[
-\,E_P[h(A,S)]+\int \lambda(s)\!\left(\sum_{a} e^{r_P(a,s)}h(a,s)\right)\mu(ds)=0.
\]
In particular, since $\mathcal{H}$ is assumed to contain all state-only functions. We have, for state-only directions $h(a,s)\equiv c(s)$ with $c\in L^\infty(\mu)$,
\[
\int\!\Big\{-\rho_P(s)+\lambda(s)\sum_a e^{r_P(a,s)}\Big\}c(s)\,\mu(ds)=0,
\]
hence
\[
\rho_P(s)=\lambda(s)\sum_{a\in\mathcal A} e^{r_P(a,s)}\quad \mu\text{-a.e.}
\]
Complementary slackness is pointwise: $\lambda(s)\big(\sum_a e^{r_P(a,s)}-1\big)=0$ $\mu$-a.e.  
Since $\rho_P>0$ $\mu$-a.e., the previous display implies $\lambda(s)>0$ $\mu$-a.e., so the constraint binds $\mu$-a.e., i.e., $\sum_a e^{r_P(a,s)}=1$ and consequently $\lambda(s)=\rho_P(s)$ $\mu$-a.e.

Returning to the stationarity identity for general $h\in\mathcal H$,
\[
E_P[h(A,S)]
=\int \lambda(s)\sum_a e^{r_P(a,s)}h(a,s)\,\mu(ds)
=\int \rho_P(s)\sum_a e^{r_P(a,s)}h(a,s)\,\mu(ds)
=E_P\!\Big[\sum_a e^{r_P(a,S)}h(a,S)\Big],
\]
which is the claim.
\end{proof}


 


 
For each $r \in \mathcal{H}$, define the integral operator 
$\mathcal{L}_{r}: L^\infty(\lambda) \to L^2(\lambda)$ by  
\[
\mathcal{L}_{r}(h)(a,s) := \left\{\sum_{\tilde a \in \mathcal A} \exp\{ r(\tilde a,s) \}\,h(\tilde a,s)\right\} - h(a,s).
\]
The first-order optimality conditions for $r_P$, given in \eqref{eq:mm-main}, can then be written as
\[
E_P\!\Big[\mathcal{L}_{r_P}(h)(A,S)\Big] = 0,
\quad \forall\, h\in\mathcal H.
\]
The Gâteaux derivative of $\mathcal{L}_{r}(h)$ in the direction $h' \in L^\infty(\lambda)$ is
\[
\dot{\mathcal{L}}_r(h,h')(a,s) 
:= \frac{d}{dt}\, \mathcal{L}_{r+th'}(h)(a,s)\Big|_{t=0}
=  \sum_{\tilde a \in \mathcal A} \exp\{ r(\tilde a,s) \}\,h'(\tilde a,s)\,h(\tilde a,s).
\]
%$Observe that the map $(h, h') \mapsto E_P[\dot{\mathcal{L}}_r(h,h')(A,S)]$ is a bilinear, positive definite functional and, hence, defines an inner product on $\mathcal{H}$.


\begin{lemma}[First-order expansion of $\mathcal{L}_r$]
\label{lemma::quadraticloss}
Fix $r,h,h' \in L^\infty(\lambda)$, and set $r' := r+h'$.  
Suppose $r$, $r'$, and $h$ are essentially bounded by $M$. Then
\[
\Big| E_P\!\Big[ \big\{ \mathcal L_{r'}(h) - \mathcal L_{r}(h) - \dot{\mathcal L}_{r}(h,h') \big\}(A,S) \Big] \Big|
\;\lesssim\; E_P[\dot{\mathcal L}_{r}(h',h')(A,S)].
\]
where the implicit constant depends only on $M$. 
\end{lemma}
\begin{proof}[Proof of Lemma \ref{lemma::quadraticloss}]
By definition,
\[
\mathcal{L}_{r}(h)(a,s)
= \sum_{\tilde a} e^{\,r(\tilde a,s)} h(\tilde a,s) - h(a,s),
\qquad
\dot{\mathcal L}_{r}(h,h')(a,s)
= \sum_{\tilde a} e^{\,r(\tilde a,s)} h'(\tilde a,s)\,h(\tilde a,s).
\]
Hence,
\begin{align*}
\mathcal L_{r'}(h)(a,s) - \mathcal L_{r}(h)(a,s) - \dot{\mathcal L}_{r}(h,h')(a,s)
&= \sum_{\tilde a} \Big( e^{\,r(\tilde a,s)+h'(\tilde a,s)} - e^{\,r(\tilde a,s)} - e^{\,r(\tilde a,s)}h'(\tilde a,s) \Big)\,h(\tilde a,s) \\
&= \sum_{\tilde a} e^{\,r(\tilde a,s)}\big( e^{\,h'(\tilde a,s)} - 1 - h'(\tilde a,s) \big)\,h(\tilde a,s).
\end{align*}
Using the inequality $|e^{u} - 1 - u| \le \tfrac{1}{2} e^{|u|} u^2$ for all $u \in \mathbb{R}$, we obtain
\[
\big| \mathcal L_{r'}(h)(a,s) - \mathcal L_{r}(h)(a,s) - \dot{\mathcal L}_{r}(h,h')(a,s) \big|
\;\le\; \tfrac{1}{2} \sum_{\tilde a} e^{\,r(\tilde a,s)} e^{\,|h'(\tilde a,s)|}\,|h(\tilde a,s)|\,|h'(\tilde a,s)|^2.
\]
If $r$, $r'$, and $h$ are essentially bounded by $M$, then
\begin{align*}
\big| \mathcal L_{r'}(h)(a,s) - \mathcal L_{r}(h)(a,s) - \dot{\mathcal L}_{r}(h,h')(a,s) \big|
&\le \sum_{\tilde a} e^{\,r(\tilde a,s)} e^{2M}\,M\,|h'(\tilde a,s)|^2 \\
&\lesssim \dot{\mathcal L}_{r}(h',h')(a,s),
\end{align*}
where the implicit constant depends only on $M$. Taking expectation under $P$, we conclude that
\[
\Big| E_P\!\Big[ \big\{ \mathcal L_{r'}(h) - \mathcal L_{r}(h) - \dot{\mathcal L}_{r}(h,h') \big\}(A,S) \Big] \Big|
\;\lesssim\; E_P[\dot{\mathcal L}_{r}(h',h')(A,S)].
\]
 
\end{proof}

\begin{lemma}
\label{lemma::frechetimplieslip}
Assume 
\((r,k) \mapsto E_0[m(A,S,r,k)]\) is Fréchet differentiable at $(r_0,k_0)$ in the additive kernel argument; that is,
\[
E_0\!\left[m(A,S,r_0+h,k_0+g) - m(A,S,r_0,k_0)
- \partial_r m(A,S,r_0,k_0)(h) - \adk m(A,S,r_0,k_0)(g)\right]
= o\!\left(\|h\|_{\mathcal{H}, 0} + \|g\|_{\mathcal{K}}\right)
\]
as \((h,g) \to (0,0)\) in the product norm 
\(\|h\|_{\mathcal{H}, 0} + \|g\|_{\mathcal{K}}\).
Then Conditions \ref{cond::boundedfunc}-\ref{cond::lipschitz} hold.
\end{lemma}
\begin{proof}[Proof of Lemma \ref{lemma::frechetimplieslip}]
Gâteaux differentiability in Condition~\ref{cond::boundedfunc} follows by
evaluating the Fréchet expansion along multiplicative kernel submodels, that is,
with additive kernel direction \(g = k_0\beta\) for \(\beta \in T_{\mathcal K}\).
Since multiplication by \(k_0\) is continuous from
\((T_{\mathcal K},\|\cdot\|_{\mathcal K,0})\) to \((\addTK,\|\cdot\|_{\mathcal K})\),
the boundedness of the derivative functionals in Condition~\ref{cond::boundedfunc2}
also follows from Fréchet differentiability. Let \(F(r,k) := E_0[m(A,S,r,k)]\).
By the triangle inequality and boundedness of the derivatives,
\begin{align*}
\big|F(r_0+h,k_0+g) - F(r_0,k_0)\big|
&= \left|E_0\!\left[\partial_r m(A,S,r_0,k_0)(h)
      + \adk m(A,S,r_0,k_0)(g)\right]\right|
      + o\!\left(\|h\|_{\mathcal{H},0} + \|g\|_{\mathcal{K}}\right) \\
&= O\!\left(\|h\|_{\mathcal{H},0} + \|g\|_{\mathcal{K}}\right)
   + o\!\left(\|h\|_{\mathcal{H},0} + \|g\|_{\mathcal{K}}\right).
\end{align*}

Thus, by the definition of the \(o(\cdot)\) term, for any $\varepsilon>0$
there exists $\delta>0$ such that whenever
\(\|h\|_{\mathcal{H},0}+\|g\|_{\mathcal{K}}<\delta\),
\[
|F(r_0+h,k_0+g)-F(r_0,k_0)|
\le (C_r + C_k + \varepsilon)\,
(\|h\|_{\mathcal{H},0}+\|g\|_{\mathcal{K}}),
\]
which establishes local Lipschitz continuity at $(r_0,k_0)$.
\end{proof}




 \begin{proof}[Proof of Theorem \ref{theorem::EIF}]

By Lemma \ref{thm:softmax-moment-matching}, the solution $r_P$ to
\eqref{eqn::Mestimand} satisfies the first-order optimality conditions:
\[
E_P[h(A,S)]
= E_{P}\!\Big[\sum_{a\in\mathcal A} \exp\{ r_P(a,S) \}\,h(a,S)\Big],
\quad \forall\,  h \in \mathcal{H}.
\]
Equivalently,
\[
E_P\!\Big[\mathcal{L}_{r_P}(h)(A,S)\Big] = 0,
\quad \forall\, h\in\mathcal H.
\]

For $\delta \in \mathbb{R}$, let $(P_t : t \in (-\delta, \delta))$ be a regular (quadratic mean differentiable) submodel through $P_0$, satisfying $P_t = P_0$ at $t=0$, with score function $\varphi \in L^2_0(P_0)$. Then, for all $t \in (-\delta, \delta)$,
\[
E_{P_t}\!\big[\mathcal{L}_{r_{P_t}}(h)(A,S)\big] = 0,
\quad \forall\, h \in \mathcal{H}.
\]
Differentiating both sides at $t=0$ gives
\[
\frac{d}{dt} E_{P_t}\!\big[\mathcal{L}_{r_{P_t}}(h)(A,S)\big]\Big|_{t=0} = 0,
\quad \forall\, h \in \mathcal{H}.
\]
Computing the derivative and applying the chain rule yields
\begin{equation}
\label{eqn::eifproof_chainrule}
\frac{d}{dt} E_{P_t}\!\big[\mathcal{L}_{r_0}(h)(A,S)\big]\Big|_{t=0}
+ \frac{d}{dt} E_{P_0}\!\big[\mathcal{L}_{r_{P_t}}(h)(A,S)\big]\Big|_{t=0}
= 0,
\quad \forall\, h \in \mathcal{H}.
\end{equation}
By the smoothness of the map $r \mapsto \mathcal{L}_r$
(Lemma~\ref{lemma::quadraticloss}) and the implicit function theorem, the path
$t \mapsto r_{P_t}$ is smooth with derivative $\dot{r}_{P_0}(\varphi) \in
\mathcal{H}$. Hence,
\[
\frac{d}{dt} E_{P_0}\!\big[\mathcal{L}_{r_{P_t}}(h)(A,S)\big]\Big|_{t=0}
= E_{P_0}\!\big[\dot{\mathcal{L}}_{r_0}(h, \dot{r}_{P_0}(\varphi))(A,S)\big],
\]
where, by definition, 
\[
E_{P_0}\!\big[\dot{\mathcal{L}}_{r_0}(h, \dot{r}_{P_0}(\varphi))(A,S)\big]
= \langle h, \dot{r}_{P_0}(\varphi) \rangle_{\mathcal{H}, 0}.
\]
Moreover,
\[
\frac{d}{dt} E_{P_t}\!\big[\mathcal{L}_{r_0}(h)(A,S)\big]\Big|_{t=0}
= E_{0}\!\big[\varphi(S',A,S)\,\mathcal{L}_{r_0}(h)(A,S)\big].
\]
Thus, rearranging terms,
\begin{equation}
    \langle h, \dot{r}_{P_0}(\varphi) \rangle_{\mathcal{H}, 0}
= - E_{0}\!\big[\varphi(S',A,S)\,\mathcal{L}_{r_0}(h)(A,S)\big],
\quad \forall\, h \in \mathcal{H}. \label{eqn::hessaninner}
\end{equation}
 
 



With the above identity in hand, we now turn to establishing pathwise differentiability of the functional. Adding and subtracting, we obtain
\begin{align*}
    E_0[m(A,S,r_{P_t},k_{P_t})] 
    &= E_0[m(A,S,r_0 + t \dot{r}_{P_0}(\varphi),\, (1+t\dot{k}_{P_0}(\varphi))k_0)] \\
    &\quad + \Big\{E_0[m(A,S,r_{P_t},k_{P_t}) - m(A,S,r_0 + t \dot{r}_{P_0}(\varphi),\, (1+t\dot{k}_{P_0}(\varphi))k_0)]\Big\}.
\end{align*} 
By \ref{cond::lipschitz}, it follows that
\begin{align*}
   &\left|E_0[m(A,S,r_{P_t},k_{P_t}) - m(A,S,r_0 + t \dot{r}_{P_0}(\varphi),\, (1+t\dot{k}_{P_0}(\varphi))k_0)]\right|  \\
   &\quad \lesssim \|r_{P_t} - r_0 - t \dot{r}_{P_0}(\varphi)\|_{\mathcal{H},0} 
   + \left\|\frac{k_{P_t}-k_0}{k_0} - t \dot{k}_{P_0}(\varphi)\right\|_{L^2(P_0)} \\
   &\quad = o(t),
\end{align*}
where we used that $\|r_{P_t} - r_0 - t \dot{r}_{P_0}(\varphi)\|_{\mathcal{H},0} = o(t)$ and $\left\|\frac{k_{P_t}-k_0}{k_0} - t \dot{k}_{P_0}(\varphi)\right\|_{L^2(P_0)} = o(t)$ by differentiability of $t \mapsto r_{P_t}$ and $t \mapsto k_{P_t}$. We conclude that
\[
E_0[m(A,S,r_{P_t},k_{P_t})] 
= E_0[m(A,S,r_0 + t \dot{r}_{P_0}(\varphi),\, (1+t\dot{k}_{P_0}(\varphi))k_0)] + o(t),
\]
and hence,
\[
\frac{d}{dt} E_0[m(A,S,r_{P_t},k_{P_t})]\Big|_{t=0} 
= \frac{d}{dt} E_0[m(A,S,r_0 + t \dot{r}_{P_0}(\varphi),\, (1+t\dot{k}_{P_0}(\varphi))k_0)]\Big|_{t=0},
\]
where this derivative exists by \ref{cond::boundedfunc}. Furthermore, by \ref{cond::boundedfunc}, 
\begin{align*}
    \frac{d}{dt} E_0[m(A,S,r_{P_t},k_{P_t})]\Big|_{t=0} =E_0\left[\partial_r m(A,S,r_0,k_0)(\dot{r}_{P_0}(\varphi))\right] + E_0\left[\partial_k m(A,S,r_0,k_0)(\dot{k}_{P_0}(\varphi))\right].
\end{align*}


By boundedness of the partial derivative functionals in
\ref{cond::boundedfunc}, the Riesz representers $\alpha_0$ and $\beta_0$
defined above Theorem~\ref{theorem::EIF} exist. Therefore,
\[
E_0\!\left[\partial_r m(A,S,r_0,k_0)(\dot{r}_{P_0}(\varphi))\right] 
= \langle \alpha_0, \dot{r}_{P_0}(\varphi)\rangle_{\mathcal{H}, 0},
\qquad
E_0\!\left[\partial_k m(A,S,r_0,k_0)(\dot{k}_{P_0}(\varphi))\right] 
= \langle \beta_0, \dot{k}_{P_0}(\varphi)\rangle_{\mathcal{K}, 0}.
\]
Note that $\dot{k}_{P_0}(\varphi)$ is the score along the submodel $k_{P_t}$ through the conditional density $k_0$. Hence, it coincides with the orthogonal projection of the score $\varphi$ onto the tangent space associated with the likelihood component $k_0$, which consists of square-integrable functions of $(S',A,S)$ with mean zero under $k_0(\cdot \mid A,S)\mu(ds')$ conditional on $(A,S)$. Since $\beta_0$ lies in this space, because $\beta_0 \in T_{\mathcal{K}}$, we have
\begin{align*}
     \langle \beta_0, \dot{k}_{P_0}(\varphi) \rangle_{\mathcal{K}, 0} 
     &=  \langle \beta_0, \dot{k}_{P_0}(\varphi) \rangle_{L^2(P_0)} \\
     &=  \langle \beta_0, \varphi \rangle_{L^2(P_0)}.
\end{align*}
Furthermore, applying \eqref{eqn::hessaninner} with $h = \dot{r}_{P_0}(\varphi)$, we obtain
\[
E_0\!\left[\partial_r m(A,S,r_0,k_0)(\dot{r}_{P_0}(\varphi))\right] 
= - E_{0}\!\big[\varphi(S',A,S)\,\mathcal{L}_{r_0}(\alpha_0)(A,S)\big]
= \langle - \mathcal{L}_{r_0}(\alpha_0), \varphi \rangle_{L^2(P_0)}.
\]
Hence,
\[
\frac{d}{dt} E_0[m(A,S,r_{P_t},k_{P_t})]\Big|_{t=0} 
= \Big\langle - \mathcal{L}_{r_0}(\alpha_0) + \beta_0, \varphi \Big\rangle_{L^2(P_0)}.
\]
Moreover, by the chain rule,
\begin{align*}
    \frac{d}{dt} \Psi(P_t)\Big|_{t=0} 
    &= \frac{d}{dt} E_{P_t}[m(A,S,r_{P_t},k_{P_t})]\Big|_{t=0} \\
    &= E_0[\varphi(S',A,S)\,m(A,S,r_0,k_0)] 
    + \frac{d}{dt} E_0[m(A,S,r_{P_t},k_{P_t})]\Big|_{t=0} \\
    &= \Big\langle m(\cdot,\cdot,r_0,k_0) - \Psi(P_0), \varphi \Big\rangle_{L^2(P_0)} 
    + \Big\langle - \mathcal{L}_{r_0}(\alpha_0) + \beta_0, \varphi \Big\rangle_{L^2(P_0)} \\
    &= \Big\langle m(\cdot,\cdot,r_0,k_0) - \Psi(P_0) - \mathcal{L}_{r_0}(\alpha_0) + \beta_0, \varphi \Big\rangle_{L^2(P_0)}.
\end{align*}
The function 
\[
m(\cdot,\cdot,r_0,k_0) - \Psi(P_0) - \mathcal{L}_{r_0}(\alpha_0) + \beta_0
\]
has mean zero and therefore belongs to the nonparametric tangent space. Hence, it is the EIF under the nonparametric model. By definition,
\[
- \mathcal{L}_{r_0}(\alpha_0)
= \alpha_0 - \sum_{\tilde a \in \mathcal{A}} \exp\{r_0(\tilde a,\cdot)\}\,\alpha_0(\tilde a,\cdot).
\]
It follows that
\[
\chi_0 = m(\cdot,\cdot,r_0,k_0) - \Psi(P_0) - \mathcal{L}_{r_0}(\alpha_0) + \beta_0
\]
is the EIF.



 
 
\end{proof}

 


\begin{proof}[Proof of Theorem \ref{theorem::vonmises}]
By \ref{cond::boundedfunc} and the definition of the linearization remainder $\operatorname{Rem}_0(r_0,k_0;\,r,k)$, we have
\begin{align*}
    E_0[m(A,S,r,k)] - E_0[m(A,S,r_0,k_0)] = E_0[\partial_r m(A,S,r_0,k_0)(r- r_0)] + E_0\!\left[\partial_k m(A,S,r_0,k_0)\!\left(\tfrac{k - k_0}{k_0}\right)\right] + \operatorname{Rem}_0(r_0,k_0;\,r,k).
\end{align*}
By the definition of the Riesz representers $\alpha_0$ and $\beta_0$, which exist by \ref{cond::boundedfunc}, we have 
$$ E_0[\partial_r m(A,S,r_0,k_0)(r- r_0)] = \langle \alpha_0, r - r_0 \rangle_{\mathcal{H}, 0}$$
   $$ E_0\!\left[\partial_k m(A,S,r_0,k_0)\!\left(\tfrac{k - k_0}{k_0}\right)\right] = \left\langle \beta_0, \tfrac{k - k_0}{k_0} \right\rangle_{\mathcal{K}, 0}.$$
Moreover, using this fact, by Lemma \ref{lemma::vonmiseskernel}, which we state and prove below,
$$E_0\!\left[\partial_k m(A,S,r_0,k_0)\!\left(\tfrac{k - k_0}{k_0}\right)\right] 
=  - P_0 \{\beta\} 
\;+\; O\left(\|\beta_0
- \beta\|_{L^2(P_0)} \;
   \Big\|\tfrac{k-k_0}{k_0}\Big\|_{L^2(P_0)}\right).$$
 Hence,
   \begin{align*}
    E_0[m(A,S,r,k)] - E_0[m(A,S,r_0,k_0)]
    &=  \langle \alpha_0, r - r_0 \rangle_{\mathcal{H}, 0}  - P_0 \{\beta\} \\
    &\quad +  O\left(\|\beta_0
    - \beta\|_{L^2(P_0)} \;
   \Big\|\tfrac{k-k_0}{k_0}\Big\|_{L^2(P_0)}\right) \\
   & \quad + \operatorname{Rem}_0(r_0,k_0;\,r,k).
\end{align*}
We now turn to  the term $ \langle \alpha_0, r - r_0 \rangle_{\mathcal{H}, 0}$.
Adding and subtracting,
   \begin{align*}
   \langle \alpha_0, r - r_0 \rangle_{\mathcal{H}, 0} &=  \langle \alpha, r - r_0 \rangle_{\mathcal{H}, 0} +  \langle \alpha_0 - \alpha, r - r_0 \rangle_{\mathcal{H}, 0}.
\end{align*}
By definition, we have 
\begin{align*}
     \langle \alpha, r - r_0 \rangle_{\mathcal{H}, 0} = E_0[\dot{\mathcal{L}}_{r_0}(\alpha, r - r_0)].
\end{align*}
Thus, applying the second-order expansion in Lemma \ref{lemma::quadraticloss},
\begin{align*}
 E_0[\dot{\mathcal{L}}_{r_0}(\alpha, r - r_0)] &= E_0[\mathcal L_{r}(\alpha)] - E_0[\mathcal L_{r_0}(\alpha)]  + O\left(E_0[\dot{\mathcal L}_{r_0}(r- r_0,r - r_0)(A,S)]\right)\\
 &= E_0[\mathcal L_{r}(\alpha)] - E_0[\mathcal L_{r_0}(\alpha)]  + O(\|r - r_0\|_{\mathcal{H}, 0}^2),
\end{align*}
where the implicit constant depends only on $M$. By Lemma \ref{thm:softmax-moment-matching}, the optimality conditions for $r_0$ imply that $E_0[\mathcal L_{r_0}(\alpha)] = 0$, since $\alpha \in T_{\mathcal{H}}(r_0)$. Thus,
\begin{align*}
 E_0[\dot{\mathcal{L}}_{r_0}(\alpha, r - r_0)] &= E_0[\mathcal L_{r}(\alpha)]  + O(\|r - r_0\|_{\mathcal{H}, 0}^2) \\
 &= P_0 \mathcal L_{r}(\alpha)  + O(\|r - r_0\|_{\mathcal{H}, 0}^2)
\end{align*}
Hence, 
\begin{align*}
   \langle \alpha_0, r - r_0 \rangle_{\mathcal{H}, 0} &=  P_0 \mathcal L_{r}(\alpha) +  \langle \alpha_0 - \alpha, r - r_0 \rangle_{\mathcal{H}, 0}  + O(\|r - r_0\|_{\mathcal{H}, 0}^2).
\end{align*}


Putting everything together, we obtain
\begin{align*}
    E_0[m(A,S,r,k)] - E_0[m(A,S,r_0,k_0)] 
    &= P_0\!\left\{\mathcal{L}_{r}(\alpha) - \beta\right\}  
    + \langle \alpha_0 - \alpha,\, r - r_0 \rangle_{\mathcal{H}, 0}  \\
    &\quad + O\!\left(\|\beta_0
    - \beta\|_{L^2(P_0)} 
    \cdot \Big\|\tfrac{k-k_0}{k_0}\Big\|_{L^2(P_0)}\right) 
    + \operatorname{Rem}_0(r_0,k_0;\,r,k).
\end{align*}
Rearranging terms and using the definition 
$\varphi_{r,k}(\cdot;\alpha,\beta) = -\mathcal{L}_{r}(\alpha) + \beta$, we obtain
\begin{align*}
&E_0[m(A,S,r,k) ]- E_0[m(A,S,r_0,k_0)] 
\;+\; \int \varphi_{r,k}(s',a,s;\alpha,\beta)\, P_0(ds,da,ds') \\
 & \qquad   = -\langle r - r_0,\, \alpha - \alpha_0 \rangle_{\mathcal{H}, 0}  + O(\|r - r_0\|_{\mathcal{H},0}^2) 
   + O\!\Big(\|\beta_0
    - \beta\|_{L^2(P_0)} 
    \cdot \Big\|\tfrac{k-k_0}{k_0}\Big\|_{L^2(P_0)}\Big) + \operatorname{Rem}_0(r_0,k_0;\,r,k).
\end{align*}
where the implicit constants depend only on $M$.
This completes the proof.



\end{proof}




\begin{lemma}[von Mises expansion for kernel partial derivative]
\label{lemma::vonmiseskernel}
For all $k \in \mathcal{K}$ and $\beta \in T_{\mathcal K}(k)$, it holds that
\[
E_0\!\left[\partial_k m(A,S,r_0,k_0)\!\left(\tfrac{k - k_0}{k_0}\right)\right] 
=  - E_0[\beta(S',A,S)] 
\;+\; O\!\left(\|\beta_0
- \beta\|_{L^2(P_0)} \;
   \Big\|\tfrac{k-k_0}{k_0}\Big\|_{L^2(P_0)}\right).
\]
\end{lemma}

\begin{proof}
By the first-order optimality conditions defining the Riesz representer 
$\beta_0$ in \eqref{eqn::rieszkernel}, for all $k \in \mathcal{K}$,
\[
E_0\!\left[\partial_k m(A,S,r_0,k_0)\!\left(\tfrac{k-k_0}{k_0}\right)\right] 
= \left\langle \tfrac{k - k_0}{k_0}, \beta_0 \right\rangle_{L^2(P_0)}.
\]
Hence,
\begin{align*}
E_0\!\left[\partial_k m(A,S,r_0,k_0)\!\left(\tfrac{k - k_0}{k_0}\right)\right] 
&= \left\langle \tfrac{k - k_0}{k_0}, \beta_0 \right\rangle_{L^2(P_0)} \\[4pt]
&= E_0\!\left[ \int \beta_0(s',A,S)\,
\big(k(s'\!\mid\!A,S) - k_0(s'\!\mid\!A,S)\big)\,\mu(ds') \right] \\[4pt]
&= E_0\!\left[ \int \beta_0(s',A,S)\,k(s'\!\mid\!A,S)\,\mu(ds') \right],
\end{align*}
where we used that 
\(
\int \beta_0(s',A,S) k_0(s' \mid A, S)\,\mu(ds') = 0
\)
$P_0$-almost surely (since $\beta_0 \in T_{\mathcal K}(k_0)$).

Adding and subtracting $\beta$,
\begin{align*}
E_0\!\left[\partial_k m(A,S,r_0,k_0)\!\left(\tfrac{k - k_0}{k_0}\right)\right] 
&= 
E_0\!\left[
    \int 
        \Bigg\{
            \beta_0
            - 
            \beta
        \Bigg\}(s',A,S)\,
        k(s' \mid A,S)\,
    \mu(ds')
\right]  \\
&\quad
+\;
E_0\!\left[
    \int 
        \beta(s',A,S)\,
        k(s' \mid A,S)\,
    \mu(ds')
\right] 
\\[6pt]
&=
E_0\!\left[
    \int 
        \Bigg\{
            \beta_0
            - 
            \beta
        \Bigg\}(s',A,S)\,
        k(s' \mid A,S)\,
    \mu(ds')
\right],
\end{align*}
where we used that
\[
\int \beta(s',A,S)\,k(s' \mid A,S)\,\mu(ds') = 0
\]
$P_0$-almost surely, since $\beta \in T_{\mathcal K}(k)$.

Decomposing $k = (k-k_0)+k_0$, we obtain
\begin{align*}
E_0\!\left[\partial_k m(A,S,r_0,k_0)\!\left(\tfrac{k - k_0}{k_0}\right)\right] 
&= 
E_0\!\left[
    \int 
        \Bigg\{
            \beta_0
            - 
            \beta
        \Bigg\}(s',A,S)\,
        (k - k_0)(s' \mid A,S)\,
    \mu(ds')
\right] \\
&\quad
-\;
E_0\!\left[
    \beta(S',A,S)
\right].
\end{align*}
Here we used that
\[
\int \beta_0(s',A,S)\,k_0(s'\!\mid\!A,S)\,\mu(ds') = 0
\]
$P_0$-almost surely, since $\beta_0 \in T_{\mathcal K}(k_0)$, and
\[
E_0\!\left[\int \beta(s',A,S)\,k_0(s'\!\mid\!A,S)\,\mu(ds')\right]
= E_0[\beta(S',A,S)].
\]

Rearranging,
\[
E_0\!\left[\partial_k m(A,S,r_0,k_0)\!\left(\tfrac{k - k_0}{k_0}\right)\right] 
=  - E_0[\beta(S',A,S)]
\;+\; E_0\!\left[ \int \Bigg\{\beta_0
- \beta\Bigg\}(s',A,S)\,(k-k_0)(s'\!\mid\!A,S)\,\mu(ds') \right].
\]
Finally, the remainder is bounded by Cauchy–Schwarz:
\begin{align*}
&\; E_0\!\left[ \int \Bigg\{\beta_0
- \beta\Bigg\}(s',A,S)\,(k-k_0)(s'\!\mid\!A,S)\,\mu(ds') \right] \\[4pt]
&= \Big\langle \beta_0
- \beta,\; \tfrac{k-k_0}{k_0} \Big\rangle_{L^2(P_0)}
\;\le\; \|\beta_0
- \beta\|_{L^2(P_0)} \;
   \Big\|\tfrac{k-k_0}{k_0}\Big\|_{L^2(P_0)}.
\end{align*}
This completes the proof.
\end{proof}



\subsection{Proofs for Section \ref{sec::eifnorm}}




\begin{proof}[Proof of Lemma \ref{lemma::linearfuncidentnorm}]
Without loss of generality, assume that $w$ is nonnegative. Otherwise, we can decompose it into its positive and negative parts and apply the argument to each separately. Furthermore, we may assume without loss of generality that $w$ has unit expectation, $E_0[w(A,S)] = 1$, since the result of the theorem is invariant to rescaling. Denote the $w$-weighted expectation of $f(A,S)$ by $E_{w}[f(A,S)] := E_0[w(A,S)\,f(A,S)]$.

By Theorem~\ref{theorem::identWithNormalization}, \(r_{0}^\dagger = g + (I - \nu)\,\mathcal{T}_{k_0,\nu,\gamma}^{-1}(r_{0}-g)\). Hence,
\begin{align*}
E_0\!\left[w(A,S)\,r_{0}^\dagger(A,S)\right] 
&= E_{w}\!\left[r_{0}^\dagger(A,S)\right]\\
&= E_w[g(S)] + E_{w}\!\left[(I - \nu)\,\mathcal{T}_{k_0,\nu,\gamma}^{-1}(r_{0}-g)(A,S)\right].
\end{align*}
Applying the final identity in Lemma~\ref{lemma::adjointident} with $\pi := \nu$ and $P(da,ds) = w(a,s)\,\rho_0(s)\,\pi_0(a \mid s)\,\mu(ds)$ (so that $E_P = E_{w}$), we obtain
\[
E_{w}\!\left[(I - \nu)\,\mathcal{T}_{k_0,\nu,\gamma}^{-1}(r_{0}-g)(A,S)\right]
= E_{w}\!\Big[r_0(A,S)-g(S) - \sum_{a \in \mathcal{A}} \nu(a \mid S)\,\{r_0(a,S)-g(S)\}\Big].
\]
Therefore,
\begin{align*}
E_0\!\left[w(A,S)\,r_{0}^\dagger(A,S)\right] 
= E_0\!\Big[w(A,S)\left\{g(S) + r_0(A,S) - \sum_{a \in \mathcal{A}} \nu(a \mid S)\,r_0(a,S)\right\}\Big].
\end{align*}
\end{proof}



\begin{proof}[Proof of Theorem \ref{theorem::EIFnorm}]
We apply Theorem \ref{theorem::EIF} to the map
\[
m(a,s,r,k) := \tilde m(a,s,r_{r,k,\nu,\gamma,g},k).
\]
Conditions \ref{cond::boundedfunc}--\ref{cond::lipschitz} hold because
\ref{cond::boundedfunc} is satisfied with $m := \tilde{m}$, while
\ref{cond::boundedfunc2::norm} and \ref{cond::lipschitz::norm} also hold.
Moreover, the map
\[
(r,k) \mapsto r_{r,k,\nu,\gamma,g} = g + (I - \nu)\,\mathcal{T}_{k,\nu,\gamma}^{-1}(r-g)
\]  
is Fréchet differentiable and Lipschitz continuous by
Lemma \ref{lemma::bellmanfrechet} and the proof of
Lemma \ref{lemma::normalizedvaluedifferentiable}. Hence, the EIF is given by
$\chi_0$ in Theorem \ref{theorem::EIF}. We now turn to deriving $\alpha_0$ and
$\beta_0$.


 

We have
\begin{align}
E_0[m(A, S, r, k)] 
= E_0[\tilde{m}(A,S, r_{r,k,\nu,\gamma,g}, k)],
\label{eq:step1}
\end{align}
which rewrites $m$ in terms of $\tilde{m}$ and the transformed reward.  

We begin with computing the reward partial derivative. By the chain rule,
\begin{align*}
E_0[\partial_r m(A, S, r, k)[\alpha]] 
&= E_0\!\left[\partial_r \tilde{m}(A,S, g + (I - \nu)\mathcal{T}_{k, \nu,\gamma}^{-1}(r-g), k)
  \bigl[(I - \nu)\mathcal{T}_{k, \nu,\gamma}^{-1}(\alpha)\bigr]\right] \\
= E_0\!\left[\partial_r \tilde{m}(A,S, r_{r, k, \nu, \gamma,g}, k)
  \bigl[(I - \nu)\mathcal{T}_{k, \nu,\gamma}^{-1}(\alpha)\bigr]\right].
\end{align*}

Evaluating at $r = r_0$ and $k = k_0$, we obtain
\begin{align}
E_0[\partial_r m(A, S, r_0, k_0)[\alpha]] 
&= E_0\!\left[\partial_r \tilde{m}(A,S, r_{0,\nu,\gamma,g}, k_0)
  \bigl[(I - \nu)\mathcal{T}_{k_0, \nu,\gamma}^{-1}(\alpha)\bigr]\right] \notag \\
= \langle \tilde{\alpha}_0, (I - \nu)\mathcal{T}_{k_0, \nu,\gamma}^{-1}(\alpha) \rangle_{\mathcal{H}, 0} ,
\label{eq:step2}
\end{align}
where we used the Riesz representation: for any $\tilde \alpha \in L^2(\rho_0 \otimes \pi_0)$,
\[
E_0\!\left[\partial_r \tilde{m}(A,S, r_{0,\nu,\gamma,g}, k_0)[\tilde \alpha]\right] 
= \langle \tilde{\alpha}_0, \tilde \alpha \rangle_{\mathcal{H}, 0},
\]
and the fact that $\langle \cdot, \cdot \rangle_{\mathcal{H}, 0} = \langle \cdot, \cdot \rangle_{L^2(\rho_0 \otimes \pi_0)}$ when $\mathcal{H} = L^\infty(\lambda)$.  

By the proof of Lemma \ref{lemma::linearfuncidentnorm}, the inner product in \eqref{eq:step2} simplifies further:
\begin{align}
\langle \tilde{\alpha}_0,\,(I - \nu)\mathcal{T}_{k_0,\nu,\gamma}^{-1}(\alpha) \rangle_{L^2(\rho_0 \otimes \pi_0)} 
&= E_0\!\left[\tilde{\alpha}_0(A,S)\,(I - \nu)\mathcal{T}_{k_0,\nu,\gamma}^{-1}(\alpha)(A,S)\right] \notag \\
&= E_0\!\left[\tilde{\alpha}_0(A,S)\,(I - \nu)(\alpha)(A,S)\right] \label{eq:step2a} \\
&= E_0\!\left[\tilde{\alpha}_0(A,S)\,\alpha(A,S)\right] 
   - E_0\!\left[\tilde{\alpha}_0(A,S)\sum_{a \in \mathcal{A}}\nu(a \mid S)\,\alpha(a,S)\right].
\label{eq:step3}
\end{align}
Moreover,
\begin{align*}
E_0\!\left[\tilde{\alpha}_0(A,S)\sum_{a \in \mathcal{A}} \nu(a \mid S)\,\alpha(a,S)\right]
&= E_0\!\left[\Big\{\sum_{\tilde a \in \mathcal{A}}\pi_0(\tilde a \mid S)\,\tilde{\alpha}_0(\tilde a,S)\Big\}
                 \Big\{\sum_{a \in \mathcal{A}} \nu(a \mid S)\,\alpha(a,S)\Big\}\right] \\
&= E_0\!\left[\Big\{\sum_{\tilde a \in \mathcal{A}}\pi_0(\tilde a \mid S)\,\tilde{\alpha}_0(\tilde a,S)\Big\}
              \,\frac{\nu(A \mid S)}{\pi_0(A \mid S)}\,\alpha(A,S)\right].
\end{align*}
Hence,
\begin{align}
\langle \tilde{\alpha}_0,\,(I - \nu)\mathcal{T}_{k_0,\nu,\gamma}^{-1}(\alpha) \rangle_{L^2(\rho_0 \otimes \pi_0)}
&= E_0\!\left[\tilde{\alpha}_0(A,S)\,\alpha(A,S)\right]
  - E_0\!\left[\Big\{\sum_{\tilde a \in \mathcal{A}}\pi_0(\tilde a \mid S)\,\tilde{\alpha}_0(\tilde a,S)\Big\}
                 \frac{\nu(A \mid S)}{\pi_0(A \mid S)}\,\alpha(A,S)\right] \notag \\
&= E_0\!\big[\alpha_0(A,S)\,\alpha(A,S)\big],
\label{eq:alpha0}
\end{align}
where
\[
\alpha_0(a,s)
:= \tilde{\alpha}_0(a,s)\;-\;\frac{\nu(a \mid s)}{\pi_0(a \mid s)}
   \sum_{\tilde a \in \mathcal{A}} \pi_0(\tilde a \mid s)\,\tilde{\alpha}_0(\tilde a,s).
\]

Combining \eqref{eq:step1}–\eqref{eq:step3}, we conclude that
\begin{equation}
E_0[\partial_r m(A, S, r_0, k_0)[\alpha]] = E_0[\alpha_0(A,S)\,\alpha(A,S)].
\label{eq:final}
\end{equation}



We now turn to the kernel partial derivative. By the chain rule,
\begin{align}
E_0[\partial_k m(A,S,r,k)[\beta]]
&= E_0\!\left[\partial_r \tilde{m}\!\left(A,S,\,r_{r,k,\nu,\gamma,g},\,k\right)
      \bigl[\,(I-\nu)\,\adk\!\left\{\mathcal{T}^{-1}_{k,\nu,\gamma}\right\}(r-g)[k\beta]\bigr]\right] \notag \\
&\quad + E_0\!\left[\partial_k \tilde{m}\!\left(A,S,\,r_{r,k,\nu,\gamma,g},\,k\right)[\beta]\right].
\label{eq:chainrule-kernel}
\end{align}
By the resolvent derivative identity
\[
\adk\!\left(\mathcal{T}^{-1}_{k,\nu,\gamma}\right)[k\beta]
= \gamma\,\mathcal{T}^{-1}_{k,\nu,\gamma}\,(\mathcal{P}_{k\beta}\nu)\,\mathcal{T}^{-1}_{k,\nu,\gamma},
\]
we obtain
\begin{align}
E_0[\partial_k m(A,S,r,k)[\beta]]
&= E_0\!\left[\partial_r \tilde{m}\!\left(A,S,\,r_{r,k,\nu,\gamma,g},\,k\right)
      \bigl[\,(I-\nu)\,\gamma\,\mathcal{T}^{-1}_{k,\nu,\gamma}(\mathcal{P}_{k\beta}\nu)\,\mathcal{T}^{-1}_{k,\nu,\gamma}(r-g)\bigr]\right] \notag \\
&\quad + E_0\!\left[\partial_k \tilde{m}\!\left(A,S,\,r_{r,k,\nu,\gamma,g},\,k\right)[\beta]\right],
\label{eq:split-kernel}
\end{align}
where \(r_{r,k,\nu,\gamma,g} := g + (I-\nu)\mathcal{T}^{-1}_{k,\nu,\gamma}(r-g)\).
Hence, evaluating at $r = r_0$ and $k = k_0$,
\begin{align}
E_0[\partial_k m(A,S,r_0,k_0)[\beta]]
&= E_0\!\left[\partial_r \tilde{m}\!\left(A,S,\,r_{0,\nu,\gamma,g},\,k_0\right)
      \bigl[\,(I-\nu)\,\gamma\,\mathcal{T}^{-1}_{k_0,\nu,\gamma}(\mathcal{P}_{k_0\beta}\nu)\,\mathcal{T}^{-1}_{k_0,\nu,\gamma}(r_0-g)\bigr]\right] \notag \\
&\quad + E_0\!\left[\partial_k \tilde{m}\!\left(A,S,\,r_{0,\nu,\gamma,g},\,k_0\right)[\beta]\right].
\label{eq:eval-kernel}
\end{align}

By the Riesz representation property of $\tilde{\beta}_{0}$, the second term on the right-hand side of \eqref{eq:eval-kernel} satisfies:
\begin{equation}
\label{eq:eval-kernel:secondterm}
    E_0\!\left[\partial_k \tilde{m}\!\left(A,S,\,r_{0,\nu,\gamma,g},\,k_0\right)[\beta]\right] = \langle \tilde{\beta}_{0}, \beta \rangle_{\mathcal{K},0}.
\end{equation}
    
 
 

By the Riesz representation property of $\tilde{\alpha}_0$, the first term on the right-hand side of \eqref{eq:eval-kernel} satisfies:
\begin{align}
    E_0\!\left[\partial_r \tilde{m}\!\left(A,S,\,r_{0,\nu,\gamma,g},\,k_0\right)
      \bigl[\,(I-\nu)\,\gamma\,\mathcal{T}^{-1}_{k_0,\nu,\gamma}(\mathcal{P}_{k_0\beta}\nu)\,\mathcal{T}^{-1}_{k_0,\nu,\gamma}(r_0-g)\bigr]\right]
    = \langle \tilde{\alpha}_0, \, (I-\nu)\,\gamma\,\mathcal{T}^{-1}_{k_0,\nu,\gamma}(\mathcal{P}_{k_0\beta}\nu)\,\mathcal{T}^{-1}_{k_0,\nu,\gamma}(r_0-g)\rangle_{L^2(\rho_0 \otimes \pi_0)}.
\label{eq:riesz-step}
\end{align}
By the proof of Lemma \ref{lemma::linearfuncidentnorm}, the inner product in \eqref{eq:riesz-step} simplifies as
\begin{align}
\big\langle \tilde{\alpha}_0,\,(I-\nu)\,\gamma\,\mathcal{T}^{-1}_{k_0,\nu,\gamma}(\mathcal{P}_{k_0\beta}\nu)\,
     \mathcal{T}^{-1}_{k_0,\nu,\gamma}(r_0-g)\big\rangle_{L^2(\rho_0 \otimes \pi_0)} 
&= \gamma\,\big\langle \alpha_0,\,(\mathcal{P}_{k_0\beta}\nu)\,q_{r_0-g, k_0}^{\nu, \gamma} \big\rangle_{L^2(\rho_0 \otimes \pi_0)} ,
\label{eq:alpha0-step}
\end{align}
where we used that $q_{r_0-g, k_0}^{\nu, \gamma} =  \mathcal{T}^{-1}_{k_0,\nu,\gamma}(r_0-g)$.

Writing out the action of $\mathcal{P}_{k_0\beta}\nu$, equation~\eqref{eq:alpha0-step} becomes
\begin{align}
\gamma\,\big\langle \alpha_0,\;(\mathcal{P}_{k_0\beta}\nu)\,
     q_{r_0-g, k_0}^{\nu, \gamma}\big\rangle_{L^2(\rho_0 \otimes \pi_0)}
&= \gamma\,E_0\!\Bigg[ 
     \alpha_0(A,S)\,
     \Big\{V_{r_0-g, k_0}^{\nu, \gamma}(S') \notag\\
&\qquad\qquad\qquad\quad
     - (\mathcal{P}_{k_0}V_{r_0-g, k_0}^{\nu, \gamma})(A,S)\Big\}
     \,\beta(S' \mid A,S) \Bigg].
\label{eq:bellman-simplify}
\end{align}
By Bellman’s equation, the bracket in \eqref{eq:bellman-simplify} equals
\begin{align}
    r_0(A,S) - g(S) +  \gamma V_{r_0-g, k_0}^{\nu, \gamma}(S') -  q_{r_0-g, k_0}^{\nu, \gamma}(A,S).
\label{eq:bellman-id}
\end{align}
Hence, 
\begin{align}
   \gamma\,\big\langle \alpha_0,\;(\mathcal{P}_{k_0\beta}\nu)\,
     q_{r_0-g, k_0}^{\nu, \gamma}\big\rangle_{L^2(\rho_0 \otimes \pi_0)} 
   &= E_0\!\left[ 
     \beta_{0,1}(S' \mid A,S)\,
     \beta(S' \mid A,S) \right],
\label{eq:final-kernel}
\end{align}
where
\[
  \beta_{0,1}(s' \mid a,s)
  := \alpha_0(a,s)\,
     \Big\{r_0(a,s) - g(s) + \gamma V_{r_0-g, k_0}^{\nu, \gamma}(s')
        - q_{r_0-g, k_0}^{\nu, \gamma}(a,s)\Big\}.
\]

Combining \eqref{eq:chainrule-kernel}–\eqref{eq:final-kernel}, we conclude that
\[
E_0[\partial_k m(A,S,r_0,k_0)[\beta]] =  \langle \beta_0, \beta \rangle_{\mathcal{K}, 0},
\]
where
\[
\beta_0 = \tilde \beta_0 + \beta_{0,1}.
\]




\end{proof}

\section{Proofs for Section \ref{sec::est}}




\begin{proof}[Proof of Theorem \ref{thm:asymptotic_linearity}]

 
Write $\psi_n = P_n\varphi_n$, where we define the estimated uncentered influence function by
\[
\varphi_n(s',a,s)
:= m(a,s,r_n,k_n)
   + \varphi_{r_n,k_n}(s',a,s;\alpha_n,\beta_n),
\]
with $\varphi_{r,k}$ as in Section~\ref{sec::EIFsub}.  
Let
\[
\varphi_0(s',a,s)
:= m(a,s,r_0,k_0)
   + \varphi_{r_0,k_0}(s',a,s;\alpha_0,\beta_0)
\]
so that 
$P_0\varphi_0 = \psi_0$ and $P_0\chi_0 = 0$, where $\chi_0$ is the EIF in Theorem~\ref{theorem::EIF}.  
Then
\begin{align*}
    \psi_n - \psi_0
    &= P_n\varphi_n - P_0\varphi_0 \\
    &= P_n(\varphi_0 - P_0\varphi_0) + P_n(\varphi_n - \varphi_0) \\
    &= P_n\chi_0 + (P_n - P_0)(\varphi_n - \varphi_0) + \{P_0\varphi_n - \psi_0\}.
\end{align*}
By Conditions~\ref{cond::bounded}–\ref{cond::consistency},  
$\|\varphi_n - \varphi_0\|_{L^2(P_0)} = o_P(1)$.  
By sample splitting (Condition~\ref{cond::samplesplit}) and Markov’s inequality,
\[
(P_n - P_0)(\varphi_n - \varphi_0)
    = O_P\!\left(n^{-1/2}\,\|\varphi_n - \varphi_0\|_{L^2(P_0)}\right)
    = o_P(n^{-1/2}).
\]
Thus
\[
\psi_n - \psi_0
= P_n\chi_0 + o_P(n^{-1/2}) + \{P_0\varphi_n - \psi_0\}.
\]
It remains to control the remainder term $P_0\varphi_n - \psi_0$.  
By definition,
\[
P_0\varphi_n - \psi_0
= E_0[m(A,S,r_n,k_n)] - E_0[m(A,S,r_0,k_0)]
  + E_0[\varphi_{r_n,k_n}(\cdot;\alpha_n,\beta_n)].
\]
Applying Theorem~\ref{theorem::vonmises} with $(r,k,\alpha,\beta)=(r_n,k_n,\alpha_n,\beta_n)$ yields
\begin{align*}
\{P_0\varphi_n - \psi_0\}
&= - \langle r_n - r_0,\;\alpha_n - \alpha_0\rangle_{\mathcal{H},r_0}
   + O\!\big(\|r_n - r_0\|_{\mathcal{H}}^2\big) \\
&\quad 
   + O\!\Big(
      \Big\|\beta_0 - \beta_n\Big\|_{L^2(P_0)}
      \Big\|\tfrac{k_n - k_0}{k_0}\Big\|_{L^2(P_0)}
     \Big)
   + \mathrm{Rem}_0(r_0,k_0;\,r_n,k_n).
\end{align*}
Each term on the right-hand side is $o_P(n^{-1/2})$ by the rate conditions  
\ref{cond::reward}–\ref{cond::remainder}.  
Therefore,
\[
\{P_0\varphi_n - \psi_0\} = o_P(n^{-1/2}),
\]
and we conclude
\[
\psi_n - \psi_0
= P_n\chi_0 + o_P(n^{-1/2}).
\]
Finally, by the central limit theorem,
\[
\sqrt{n}(\psi_n - \psi_0)
= \frac{1}{\sqrt{n}}\sum_{i=1}^n \chi_0(S_i,A_i,S_i') + o_P(1)
\rightsquigarrow N\!\left(0,\;\text{Var}_{P_0}\!\big(\chi_0(S,A,S')\big)\right),
\]
establishing asymptotic linearity and efficiency.  
\end{proof}


 

\section{Proofs for Section \ref{sec::applications}}
 
\label{appendix::applications}
\subsection{Technical lemmas}

We begin with operator and occupancy-ratio identities that are reused across the
example-specific EIF derivations and asymptotic expansions.



 
 
 
\begin{lemma}[Invertibility of the Bellman operator under stationarity overlap]
\label{lemma::invertibility}
Assume $\gamma < 1$ and \ref{cond::stationary}. Then $\mathcal{T}_{k_0,\pi}:L^2(\lambda) \to L^2(\lambda)$ admits a bounded inverse, both as an operator on $L^2(\lambda)$ and on $L^2(\pi_0 \otimes \rho_0)$.
\end{lemma}
\begin{proof}
Recall that $p_{k_0}^{\pi,\gamma}$ is the stationary state $\mu$-density for the counterfactual MDP with policy $\pi$ and kernel $k_0$. We claim that $\mathcal{T}_{k_0,\pi}$ is a $\gamma$-contraction in $L^2(\pi \otimes p_{k_0}^{\pi,\gamma})$. Let $\mathcal{P}_{k_0,\pi}$ denote the Markov (transition) operator under $k_0$ and $\pi$, acting on state-only functions by
\[
\mathcal{P}_{k_0,\pi} f(s) \;:=\; \int f(s')\,k_0(\mathrm ds' \mid s,a)\,\pi(\mathrm da\mid s),
\]
so that $\mathcal{T}_{k_0,\pi} := I - \gamma\,\mathcal{P}_{k_0,\pi}$.  

For any $f \in L^2(p_{k_0}^{\pi,\gamma})$, Jensen’s inequality gives
\[
\|\mathcal{P}_{k_0,\pi} f\|_{L^2(p_{k_0}^{\pi,\gamma})}^2
= E_{S \sim p_{k_0}^{\pi,\gamma}}\!\left[\big(E[f(S') \mid S]\big)^2\right]
\le E_{S \sim p_{k_0}^{\pi,\gamma}}\!\left[E\!\left(f(S')^2 \mid S\right)\right]
= E_{S' \sim p_{k_0}^{\pi,\gamma}}\!\left[f(S')^2\right]
= \|f\|_{L^2(p_{k_0}^{\pi,\gamma})}^2,
\]
where all expectations are taken under the stationary distribution $p_{k_0}^{\pi,\gamma}$, and the third equality follows from stationarity of the chain. Hence $\|\mathcal{P}_{k_0,\pi}\|_{L^2(p_{k_0}^{\pi,\gamma}) \to L^2(p_{k_0}^{\pi,\gamma})} \le 1$, and therefore
\[
\|\gamma\,\mathcal{P}_{k_0,\pi}\|_{L^2(p_{k_0}^{\pi,\gamma}) \to L^2(p_{k_0}^{\pi,\gamma})} \le \gamma < 1.
\]
Thus $\mathcal{T}_{k_0,\pi}$ is a perturbation of the identity by a $\gamma$-contraction with $\gamma < 1$. By Banach’s inversion theorem, $\mathcal{T}_{k_0,\pi}$ is therefore invertible on $L^2(\pi \otimes p_{k_0}^{\pi,\gamma})$.

By \ref{cond::stationary}, the spaces $L^2(\pi \otimes p_{k_0}^{\pi,\gamma})$ and $L^2(\pi_0 \otimes \rho_0)$ are equivalent, with norms that are uniformly comparable. Hence $\mathcal{T}_{k_0,\pi}$ is also invertible on $L^2(\pi_0 \otimes \rho_0)$. Similarly, by the assumptions in Section \ref{sec::notation}, both $\pi_0$ and $\rho_0$ are uniformly bounded above and below. It follows that $L^2(\pi_0 \otimes \rho_0)$ and $L^2(\lambda)$ are equivalent in norm and as spaces. Therefore $\mathcal{T}_{k_0,\pi}$ is also invertible on $L^2(\lambda)$.

\end{proof}




In the following lemma, let $P$ be a distribution over $(A,S)$, and factor it as $P(\{a\},ds) = \pi_P(a \mid s)\,\rho_P(s)\,\mu(ds),$
where $\pi_P$ is the induced policy and $\rho_P$ is the marginal density of $S$. Let $\gamma \in [0,1)$ and let $k \in \mathcal{K}$ be a transition kernel.  We define the (unnormalized) discounted state–action occupancy density ratio as
\begin{equation}
    d_{P,k}^{\pi}(a,s) := \rho_{P,k}^\pi(s)\,\frac{\pi(a \mid s)}{\pi_P(a \mid s)},
    \qquad 
    \rho_{P,k}^\pi(s) :=  \sum_{t=0}^\infty \gamma^t \,\frac{d\mathbb{P}_{k,\pi}}{d\rho_P}(S_t = s),
\end{equation}
where $\rho_{P,k}^\pi$ denotes the discounted state occupancy distribution under policy $\pi$ and dynamics $k$.





\begin{lemma}[Adjoint identities]
\label{lemma::adjointident}
Let $\gamma \in [0,1)$. Then, for any $f \in L^2(\lambda)$,
\[
E_P\!\big[(\Pi \mathcal{T}_{k,\pi,\gamma}^{-1} f)(S)\big] 
= E_P\!\big[d_{P,k}^{\pi}(A,S)\,f(A,S)\big],
\]
\[
E_P\!\big[(\mathcal{T}_{k,\pi,\gamma}^{-1} f)(A,S)\big] 
= E_P\!\Big[\Big\{1 - \tfrac{\pi(A \mid S)}{\pi_P(A \mid S)} 
+ d_{P,k}^{\pi}(A,S)\Big\} f(A,S)\Big],
\]
and consequently,
\begin{align*}
   E_P\!\big[(I - \Pi)\,\mathcal{T}_{k,\pi,\gamma}^{-1} f(A,S)\big] 
   &= E_P\!\Big[\Big\{1 - \tfrac{\pi(A \mid S)}{\pi_P(A \mid S)}\Big\} f(A,S)\Big] \\
   &= E_P\!\Big[f(A,S) - \sum_{a \in \mathcal{A}} \pi(a \mid S)\,f(a,S)\Big].
\end{align*}
\end{lemma}

\begin{proof}
Consider the inner product $\langle g,h\rangle_{P} := E_P[gh]$. Then
\[
E_P[\Pi \mathcal{T}_{k,\pi,\gamma}^{-1}f]
= \langle 1, \Pi \mathcal{T}_{k,\pi,\gamma}^{-1}f\rangle_{P}
= \langle \Pi^*1, \mathcal{T}_{k,\pi,\gamma}^{-1}f\rangle_{P}
= \big\langle (\mathcal{T}_{k,\pi,\gamma}^{-1})^*(\Pi^*1),\,f\big\rangle_{P}.
\]
Under $E_P$ we have $\Pi^*1(a,s) = \pi(a\mid s)/\pi_P(a\mid s)$. Expanding $(\mathcal{T}_{k,\pi,\gamma}^{-1})^*$ as a resolvent series gives
\[
(\mathcal{T}_{k,\pi,\gamma}^{-1})^*(\Pi^*1)(a,s)
= \sum_{t\ge0} \gamma^t \,\frac{p_{k,\pi}(S_t=s)}{\rho_P(s)}\,
   \frac{\pi(a\mid s)}{\pi_P(a\mid s)}
= d_{P,k}^{\pi}(a,s),
\]
which proves the first identity.

For the second identity, note that
\[
E_P[\mathcal{T}_{k,\pi,\gamma}^{-1}f]
= \sum_{t\ge0} \gamma^t\,\mathbb{E}[f(A_t,S_t)],
\]
where the expectation is taken with $S_0\sim\rho_P$, $A_0\sim\pi_P(\cdot\mid S_0)$, the dynamics follow $k$, and for $t\ge1$ the actions follow $\pi$. Split the sum into the $t=0$ term and $\sum_{t\ge1}$. The $t=0$ contribution is $E_P[f(A,S)]$.

For $t\ge1$, a change of measure back to $E_P$ yields
\[
\sum_{t\ge1} \gamma^t \, E_P\!\left[
  \frac{p_{k,\pi}(S_t=S)}{\rho_P(S)} \,
  \frac{\pi(A\mid S)}{\pi_P(A\mid S)} \, f(A,S)
\right]
= E_P\!\Big[ \big(\rho_{P,k}^{\pi}(S)-1\big)
   \frac{\pi(A\mid S)}{\pi_P(A\mid S)}\,f(A,S)\Big].
\]
Adding the $t=0$ term, we obtain
\[
E_P[\mathcal{T}_{k,\pi,\gamma}^{-1}f]
= E_P\!\Big[\Big\{1 - \frac{\pi(A\mid S)}{\pi_P(A\mid S)}
   + \rho_{P,k}^{\pi}(S)\frac{\pi(A\mid S)}{\pi_P(A\mid S)}\Big\}
   f(A,S)\Big].
\]
Since $d_{P,k}^{\pi}(A,S) = \rho_{P,k}^{\pi}(S)\,\frac{\pi(A\mid S)}{\pi_P(A\mid S)}$, this gives the desired expression.  
Subtracting the first identity from the second establishes the final identity.
\end{proof}

 

\begin{lemma}[Smoothness of Bellman operator]
\label{lemma::bellmanfrechet}
Suppose $\mathcal{T}_{k,\pi,\gamma}:L^2(\lambda)\to L^2(\lambda)$ is invertible with bounded inverse $\mathcal{T}_{k,\pi,\gamma}^{-1}$. Then the map $(r,k)\mapsto \mathcal{T}_{k,\pi,\gamma}^{-1} r$ is Fréchet differentiable in $L^2$, with partial derivatives, for $g \in \addTK$ and $h\in L^2(\lambda)$,
\[
\adk\bigl(\mathcal{T}_{k,\pi,\gamma}^{-1}\bigr)[g]
= \gamma \,\mathcal{T}_{k,\pi,\gamma}^{-1}\bigl(\mathcal{P}_{g}\Pi \bigr)\mathcal{T}_{k,\pi,\gamma}^{-1},
\qquad
\partial_r\bigl(\mathcal{T}_{k,\pi,\gamma}^{-1} r\bigr)[h] = \mathcal{T}_{k,\pi,\gamma}^{-1} h.
\]
Moreover,
\[
\|\mathcal T_{k+g,\pi,\gamma}^{-1}(r+h)-\mathcal T_{k,\pi,\gamma}^{-1}(r)-\adk\mathcal T_{k,\pi,\gamma}^{-1}[g](r)- \partial_r (\mathcal T_{k,\pi,\gamma}^{-1} r)[h]\|_{L^2(\lambda)}
\;\lesssim\; \|g\|_{\mathcal K}^2\,\|r\|_{L^2(\lambda)} \;+\; \|g\|_{\mathcal K}\,\|h\|_{L^2(\lambda)}.
\]
\end{lemma}
\begin{proof}
We have
\[
\adk\bigl(\mathcal{T}_{k,\pi,\gamma}^{-1}\bigr)[g]
= -\,\mathcal{T}_{k,\pi,\gamma}^{-1}\bigl(\adk \mathcal{T}_{k,\pi,\gamma}[g]\bigr)\mathcal{T}_{k,\pi,\gamma}^{-1},
\]
where $\adk \mathcal{T}_{k,\pi,\gamma}[g] = - \gamma \mathcal{P}_{g}\Pi $.
Adding and subtracting terms gives
\begin{align*}
&\mathcal{T}_{k+g,\pi,\gamma}^{-1} - \mathcal{T}_{k,\pi,\gamma}^{-1} 
- \bigl(\adk \mathcal{T}_{k,\pi,\gamma}^{-1}\bigr)[g] \\
&\quad= \mathcal{T}_{k,\pi,\gamma}^{-1}\,\bigl(\mathcal{T}_{k+g,\pi,\gamma} - \mathcal{T}_{k,\pi,\gamma}\bigr)\,
\mathcal{T}_{k,\pi,\gamma}^{-1}\,\bigl(\mathcal{T}_{k+g,\pi,\gamma} - \mathcal{T}_{k,\pi,\gamma}\bigr)\,
\mathcal{T}_{k+g,\pi,\gamma}^{-1} \\
&\qquad\;+\;
\mathcal{T}_{k,\pi,\gamma}^{-1}\Bigl(\adk \mathcal{T}_{k,\pi,\gamma}[g] - 
\bigl(\mathcal{T}_{k+g,\pi,\gamma} - \mathcal{T}_{k,\pi,\gamma}\bigr)\Bigr)\mathcal{T}_{k,\pi,\gamma}^{-1}.
\end{align*}
\[
\mathcal{T}_{k+g,\pi,\gamma} - \mathcal{T}_{k,\pi,\gamma} = - \gamma \mathcal{P}_{g}\Pi .
\]
\[
\adk \mathcal{T}_{k,\pi,\gamma}[g] - 
\bigl(\mathcal{T}_{k+g,\pi,\gamma} - \mathcal{T}_{k,\pi,\gamma}\bigr) = 0,
\]
since \(k \mapsto \mathcal{T}_{k,\pi,\gamma} = I - \gamma \mathcal{P}_k \Pi\) is affine in \(k\).  
Thus, the final term vanishes, and we obtain
\begin{align*}
&\mathcal{T}_{k+g,\pi,\gamma}^{-1} - \mathcal{T}_{k,\pi,\gamma}^{-1} 
- \bigl(\adk \mathcal{T}_{k,\pi,\gamma}^{-1}\bigr)[g] \\
&\quad=  \gamma^2 \mathcal{T}_{k,\pi,\gamma}^{-1}\,\bigl( \mathcal{P}_{g}\Pi \bigr)\,
\mathcal{T}_{k,\pi,\gamma}^{-1}\,\bigl( \mathcal{P}_{g}\Pi \bigr)\,
\mathcal{T}_{k+g,\pi,\gamma}^{-1}.
\end{align*}
Taking operator norms and using submultiplicativity,
 \[
\bigl\|\mathcal{T}_{k+g,\pi,\gamma}^{-1} - \mathcal{T}_{k,\pi,\gamma}^{-1} 
- (\adk \mathcal{T}_{k,\pi,\gamma}^{-1})[g]\bigr\|_{L^2(\lambda)\to L^2(\lambda)}
\;\le\;  \,\|\mathcal{T}_{k,\pi,\gamma}^{-1}\|_{L^2(\lambda)\to L^2(\lambda)}^2 \,\|\mathcal{T}_{k+g,\pi,\gamma}^{-1}\|_{L^2(\lambda)\to L^2(\lambda)}\,
\|\mathcal{P}_{g}\Pi \|_{L^2(\lambda)\to L^2(\lambda)}^{2}.
\]
We claim that $\|\mathcal{P}_{g}\Pi \|_{L^2(\lambda) \rightarrow L^2(\lambda)} \lesssim \|g\|_{\mathcal{K}}$. Assuming this bound, we find that
\[
\bigl\|\mathcal{T}_{k+g,\pi,\gamma}^{-1} - \mathcal{T}_{k,\pi,\gamma}^{-1} 
- (\adk \mathcal{T}_{k,\pi,\gamma}^{-1})[g]\bigr\|_{L^2(\lambda)\to L^2(\lambda)}
\;\lesssim\; \,  \,\|\mathcal{T}_{k,\pi,\gamma}^{-1}\|_{L^2(\lambda)\to L^2(\lambda)}^{2}\,\|\mathcal{T}_{k+g,\pi,\gamma}^{-1}\|_{L^2(\lambda)\to L^2(\lambda)}\;
\|g\|_{\mathcal K}^{2}.
\]
Finally, $\|\mathcal{T}_{k,\pi,\gamma}^{-1}\|_{L^2(\lambda)\to L^2(\lambda)} < \infty$ by definition, and by continuity of 
$k \mapsto \mathcal{T}_{k,\pi,\gamma}$ in operator norm we have that 
$\|\mathcal{T}_{k+g,\pi,\gamma}^{-1}\|_{L^2(\lambda)\to L^2(\lambda)} < \infty$ 
for all $g$ in an $L^2(\lambda)$ ball around zero. Thus,
\[
\bigl\|\mathcal{T}_{k+g,\pi,\gamma}^{-1} - \mathcal{T}_{k,\pi,\gamma}^{-1} 
- (\adk \mathcal{T}_{k,\pi,\gamma}^{-1})[g]\bigr\|_{L^2(\lambda)\to L^2(\lambda)}
= O\!\left(\|g\|_{\mathcal K}^{2}\right),
\]
as claimed.



Finally, note that since $r \mapsto \mathcal T_{k,\pi,\gamma}^{-1} r$ is linear, we have $\partial_r \big(\mathcal T_{k,\pi,\gamma}^{-1} r\big)[h]=\mathcal T_{k,\pi,\gamma}^{-1}h$. Adding and subtracting,
\[
\mathcal T_{k+g,\pi,\gamma}^{-1}(r+h)-\mathcal T_{k,\pi,\gamma}^{-1}(r)-\adk\mathcal T_{k,\pi,\gamma}^{-1}[g](r)-\partial_r \big(\mathcal T_{k,\pi,\gamma}^{-1} r\big)[h]
=\underbrace{\big(\mathcal T_{k+g,\pi,\gamma}^{-1}-\mathcal T_{k,\pi,\gamma}^{-1}-\adk\mathcal T_{k,\pi,\gamma}^{-1}[g]\big)r}_{\text{2nd-order in }g}
\;+\;\underbrace{\big(\mathcal T_{k+g,\pi,\gamma}^{-1}-\mathcal T_{k,\pi,\gamma}^{-1}\big)h}_{\text{1st-order in }g}.
\]
Hence, by the triangle inequality and operator–norm submultiplicativity,
\begin{align*}
\|\mathcal T_{k+g,\pi,\gamma}^{-1}(r+h)-\mathcal T_{k,\pi,\gamma}^{-1}(r)-\adk\mathcal T_{k,\pi,\gamma}^{-1}[g](r)-\mathcal T_{k,\pi,\gamma}^{-1}h\|_{L^2(\lambda)}
&\le \|\mathcal T_{k+g,\pi,\gamma}^{-1}-\mathcal T_{k,\pi,\gamma}^{-1}-\adk\mathcal T_{k,\pi,\gamma}^{-1}[g]\|_{L^2(\lambda)\to L^2(\lambda)}\,\|r\|_{L^2(\lambda)} \\
&\quad + \|\mathcal T_{k+g,\pi,\gamma}^{-1}-\mathcal T_{k,\pi,\gamma}^{-1}\|_{L^2(\lambda)\to L^2(\lambda)}\,\|h\|_{L^2(\lambda)}.
\end{align*}
From above, we have
\[
\|\mathcal T_{k+g,\pi,\gamma}^{-1}-\mathcal T_{k,\pi,\gamma}^{-1}-\adk\mathcal T_{k,\pi,\gamma}^{-1}[g]\|_{L^2(\lambda)\to L^2(\lambda)} \lesssim \|g\|_{\mathcal K}^2,
\]
and, by the triangle inequality,
\[
\|\mathcal T_{k+g,\pi,\gamma}^{-1}-\mathcal T_{k,\pi,\gamma}^{-1}\|_{L^2(\lambda)\to L^2(\lambda)} \lesssim \|\adk\mathcal T_{k,\pi,\gamma}^{-1}[g]\|_{L^2(\lambda)\to L^2(\lambda)} + \|g\|_{\mathcal K}^2 \lesssim \|g\|_{\mathcal K} + \|g\|_{\mathcal K}^2.
\]
Hence,
\[
\|\mathcal T_{k+g,\pi,\gamma}^{-1}(r+h)-\mathcal T_{k,\pi,\gamma}^{-1}(r)-\adk\mathcal T_{k,\pi,\gamma}^{-1}[g](r)- \partial_r (\mathcal T_{k,\pi,\gamma}^{-1} r)[h]\|_{L^2(\lambda)}
\;\lesssim\; \|g\|_{\mathcal K}^2\,\|r\|_{L^2(\lambda)} \;+\; \|g\|_{\mathcal K}\,\|h\|_{L^2(\lambda)}.
\]




It remains to prove the claim that $\|\mathcal{P}_{g}\Pi \|_{L^2(\lambda) \rightarrow L^2(\lambda)} \lesssim \|g\|_{\mathcal{K}}$. For $(a,s)\in\mathcal{A}\times\mathcal{S}$, write
\[
(\mathcal{P}_{g}\Pi f)(a,s)
= \int_{\mathcal S} g(s' \mid a,s)\,h(s')\,\mu(ds'),
\qquad
h(s') := \sum_{a'\in\mathcal A} \pi(a' \mid s')\, f(a',s').
\]
By Cauchy--Schwarz in $L^2(\mu)$,
\[
|(\mathcal{P}_{g}\Pi f)(a,s)|
\;\le\; \Big(\int g(s'\mid a,s)^2\,\mu(ds')\Big)^{1/2}
       \Big(\int h(s')^2\,\mu(ds')\Big)^{1/2}.
\]
Next, by Cauchy--Schwarz over $a'\in\mathcal A$,
\[
h(s')^2 \;\le\; \Big(\sum_{a'} \pi(a'\mid s')^2\Big)\Big(\sum_{a'} f(a',s')^2\Big).
\]
Integrating in $s'$ and letting $c_\pi := \sup_{s'}\sum_{a'} \pi(a'\mid s')^2\le 1$, we get
\[
\int h(s')^2\,\mu(ds') \;\le\; c_\pi \int \sum_{a'} f(a',s')^2\,\mu(ds')
= c_\pi\,\|f\|_{L^2(\#_{\mathcal A}\otimes \mu)}^2.
\]
Therefore,
\[
|(\mathcal{P}_{g}\Pi f)(a,s)|^2
\;\le\; c_\pi\Big(\int g(s'\mid a,s)^2\,\mu(ds')\Big)\,\|f\|_{L^2(\#_{\mathcal A}\otimes \mu)}^2.
\]
Integrating over $(a,s)$ with respect to $\lambda(da,ds)$ yields
\[
\|\mathcal{P}_{g}\Pi f\|_{L^2(\lambda)}^2
\;\le\; c_\pi \Big(\int\!\!\int g(s'\mid a,s)^2\,\mu(ds')\,\lambda(da,ds)\Big)
          \|f\|_{L^2(\#_{\mathcal A}\otimes \mu)}^2
= c_\pi\,\|g\|_{L^2(\lambda\otimes \mu)}^2\,\|f\|_{L^2(\#_{\mathcal A}\otimes \mu)}^2.
\]
Taking square roots and using that $\|f\|_{L^2(\#_{\mathcal A}\otimes \mu)}^2 \lesssim |\mathcal{A}| \cdot \|f\|_{L^2(\lambda)}^2$ gives
\[
\|\mathcal{P}_{g}\Pi f\|_{L^2(\lambda)}
\lesssim \|g\|_{L^2(\lambda\otimes \mu)}\,\|f\|_{L^2(\lambda)},
\]
where $\lesssim$ only depends on  $|\mathcal{A}|$. Thus, $\|\mathcal{P}_{g}\Pi \|_{L^2(\lambda) \rightarrow L^2(\lambda)} \lesssim \|g\|_{\mathcal{K}}$.



\end{proof}


 

\subsubsection{Lemmas for Example~\ref{example::norm}}
\label{lemmas::norm}

For $(r,k) \in L^\infty(\lambda) \otimes \mathcal{K}$, define the two–resolvent value operator:
\[
\mathcal{V}_k^g(r) := \Pi\,\mathcal{T}_{k,\pi,\gamma'}^{-1}\!\left\{g + (I - \nu)\,\mathcal{T}_{k,\nu,\gamma}^{-1}(r-g)\right\}.
\]
Note that $\mathcal{V}_k^g(r)$ coincides with the value function of policy $\pi$ under the 
$\nu$-normalized reward \(r_{r,k,\nu,\gamma,g} = g + (I - \nu)\,\mathcal{T}_{k,\nu,\gamma}^{-1}(r-g)\).

\begin{lemma}[Differentiability of the two–resolvent value operator]
 \label{lemma::normalizedvaluedifferentiable}
Let $\pi,\nu$ be policies. Suppose $\mathcal{T}_{k,\pi,\gamma'}:L^2(\lambda)\to L^2(\lambda)$ and 
$\mathcal{T}_{k,\nu,\gamma}:L^2(\lambda)\to L^2(\lambda)$ are invertible with bounded inverses. 
Then the map $(r,k) \mapsto \mathcal{V}_k^g(r)$ defined by
\[
\mathcal{V}_k^g(r) := \Pi\,\mathcal{T}_{k,\pi,\gamma'}^{-1}\!\left\{g + (I-\nu)\,\mathcal{T}_{k,\nu,\gamma}^{-1}(r-g)\right\}
\]
is Fréchet differentiable in $L^2$, with derivatives
\begin{align*}
\partial_r \mathcal{V}_k^g(r)[h] 
&= \Pi\,\mathcal{T}_{k,\pi,\gamma'}^{-1}(I-\nu)\,\mathcal{T}_{k,\nu,\gamma}^{-1}(h), \\[6pt]
\adk \mathcal{V}_k^g(r)[\eta]
&= \gamma'\,\Pi\,\mathcal{T}_{k,\pi,\gamma'}^{-1}(\mathcal{P}_\eta \Pi)\,
      \mathcal{T}_{k,\pi,\gamma'}^{-1}\!\left\{g + (I-\nu)\,\mathcal{T}_{k,\nu,\gamma}^{-1}(r-g)\right\} \\
&\quad+\; \gamma\,\Pi\,\mathcal{T}_{k,\pi,\gamma'}^{-1}(I-\nu)\,
      \mathcal{T}_{k,\nu,\gamma}^{-1}(\mathcal{P}_\eta \nu)\,
      \mathcal{T}_{k,\nu,\gamma}^{-1}(r-g).
\end{align*}
Moreover, the second-order remainder satisfies
\[
\|\mathcal V_{k+\eta}^g(r+h)-\mathcal V_{k}^g(r)-\adk\mathcal V_{k}^g(r)[\eta]- \partial_r \mathcal V_{k}^g(r)[h]\|_{L^2(\lambda)}
\;\lesssim\; \|\eta\|_{\mathcal K}^2\,\|r\|_{L^2(\lambda)} \;+\; \|\eta\|_{\mathcal K}\,\|h\|_{L^2(\lambda)},
\]
where the implicit constant depends on operator norms of $\Pi$, $(I-\nu)$, and the inverses 
$\mathcal{T}_{k,\pi,\gamma'}^{-1}$ and $\mathcal{T}_{k,\nu,\gamma}^{-1}$.
\end{lemma}
\begin{proof}
Since \(g\) is fixed, it follows immediately that, for $h \in L^2(\lambda)$,
\begin{align*}
    \partial_r \mathcal{V}_k^g(r)[h]
    \;=\; \Pi\,\mathcal{T}_{k,\pi,\gamma'}^{-1}(I-\nu)\,\mathcal{T}_{k,\nu,\gamma}^{-1}(h).
\end{align*}

We now turn to the additive kernel derivative. By Lemma~\ref{lemma::bellmanfrechet}, we have, for $\eta \in \addTK$,
\[
\adk\bigl(\mathcal{T}_{k,\nu,\gamma}^{-1}\bigr)[\eta]
= \gamma \,\mathcal{T}_{k,\nu,\gamma}^{-1}\bigl(\mathcal{P}_{\eta}\nu \bigr)\mathcal{T}_{k,\nu,\gamma}^{-1},
\qquad
\adk \mathcal{P}_{k,\nu}[\eta] \;=\; \mathcal{P}_{\eta}\nu,
\]
and similarly,
\[
\adk\bigl(\mathcal{T}_{k,\pi,\gamma'}^{-1}\bigr)[\eta]
\;=\; \gamma' \,\mathcal{T}_{k,\pi,\gamma'}^{-1}\bigl(\mathcal{P}_{\eta}\Pi \bigr)\mathcal{T}_{k,\pi,\gamma'}^{-1},
\qquad
\adk (\mathcal{P}_{k}\Pi)[\eta] \;=\; \mathcal{P}_{\eta}\Pi.
\]
By the chain rule and the above identities ,
\begin{align*}
\adk \mathcal{V}_k^g(r)[\eta]
&= \Pi\Bigl(\adk \mathcal{T}_{k,\pi,\gamma'}^{-1}[\eta]\Bigr)
      \left\{g + (I-\nu)\,\mathcal{T}_{k,\nu,\gamma}^{-1}(r-g)\right\}
  \;+\; \Pi\,\mathcal{T}_{k,\pi,\gamma'}^{-1}
      (I-\nu)\,\Bigl(\adk \mathcal{T}_{k,\nu,\gamma}^{-1}[\eta]\Bigr)(r-g) \\
&= \Pi\Bigl[
      \gamma'\,\mathcal{T}_{k,\pi,\gamma'}^{-1}\bigl(\mathcal{P}_\eta \Pi\bigr)\,
      \mathcal{T}_{k,\pi,\gamma'}^{-1}\!\left\{g + (I-\nu)\,\mathcal{T}_{k,\nu,\gamma}^{-1}(r-g)\right\}
      \;+\;
      \mathcal{T}_{k,\pi,\gamma'}^{-1}(I-\nu)\,
      \gamma\,\mathcal{T}_{k,\nu,\gamma}^{-1}\bigl(\mathcal{P}_\eta \nu\bigr)\,
      \mathcal{T}_{k,\nu,\gamma}^{-1}(r-g)
    \Bigr].
\end{align*}
Equivalently, using \(r_{r,k,\nu,\gamma,g}= g + (I-\nu)\,\mathcal{T}_{k,\nu,\gamma}^{-1}(r-g)\) and \(q_{r-g,k}^{\nu,\gamma} =  \mathcal{T}_{k,\nu,\gamma}^{-1}(r-g)\),
\begin{align*}
\adk \mathcal{V}_k^g(r)[\eta]
&= \gamma'\,\Pi\,\mathcal{T}_{k,\pi,\gamma'}^{-1}(\mathcal{P}_\eta \Pi)\,
      \mathcal{T}_{k,\pi,\gamma'}^{-1}\bigl(r_{r,k,\nu,\gamma,g}\bigr) \\
&\quad+\; \gamma\,\Pi\,\mathcal{T}_{k,\pi,\gamma'}^{-1}(I-\nu)\,
      \mathcal{T}_{k,\nu,\gamma}^{-1}(\mathcal{P}_\eta \nu)\,
      q_{r-g,k}^{\nu,\gamma}.
\end{align*}


 
We now turn to the second–order remainder. Adding and subtracting, we have
\begin{align*}
\mathcal{V}_{k+\eta}^g(r)-\mathcal{V}_k^g(r)
&= \Pi\Big\{\mathcal{T}_{k+\eta,\pi,\gamma'}^{-1}(I-\nu)\,\mathcal{T}_{k+\eta,\nu,\gamma}^{-1}
      - \mathcal{T}_{k,\pi,\gamma'}^{-1}(I-\nu)\,\mathcal{T}_{k,\nu,\gamma}^{-1}\Big\}(r-g) \\
&= \Pi\Big\{\big(\mathcal{T}_{k+\eta,\pi,\gamma'}^{-1}-\mathcal{T}_{k,\pi,\gamma'}^{-1}\big)(I-\nu)\,\mathcal{T}_{k,\nu,\gamma}^{-1}
      \;+\; \mathcal{T}_{k+\eta,\pi,\gamma'}^{-1}(I-\nu)\big(\mathcal{T}_{k+\eta,\nu,\gamma}^{-1}-\mathcal{T}_{k,\nu,\gamma}^{-1}\big)\Big\}(r-g) .
\end{align*}
Subtracting the linear term
\[
\adk \mathcal{V}_k^g(r)[\eta]
= \Pi\Big\{\big(\adk \mathcal{T}_{k,\pi,\gamma'}^{-1}[\eta]\big)(I-\nu)\,\mathcal{T}_{k,\nu,\gamma}^{-1}
\;+\; \mathcal{T}_{k,\pi,\gamma'}^{-1}(I-\nu)\big(\adk \mathcal{T}_{k,\nu,\gamma}^{-1}[\eta]\big)\Big\}(r-g),
\]
and adding and subtracting $\Pi\,\mathcal{T}_{k,\pi,\gamma'}^{-1}(I-\nu)\big(\mathcal{T}_{k+\eta,\nu,\gamma}^{-1}-\mathcal{T}_{k,\nu,\gamma}^{-1}\big)(r-g)$, we obtain
\begin{align*}
&\mathcal{V}_{k+\eta}^g(r) - \mathcal{V}_k^g(r) - \adk \mathcal{V}_k^g(r)[\eta] \\
&= \Pi\Big\{\big(\mathcal{T}_{k+\eta,\pi,\gamma'}^{-1}-\mathcal{T}_{k,\pi,\gamma'}^{-1}
      - \adk \mathcal{T}_{k,\pi,\gamma'}^{-1}[\eta]\big)(I-\nu)\,\mathcal{T}_{k,\nu,\gamma}^{-1}\Big\}(r-g) \\
&\quad+\; \Pi\Big\{\mathcal{T}_{k,\pi,\gamma'}^{-1}(I-\nu)\big(\mathcal{T}_{k+\eta,\nu,\gamma}^{-1}-\mathcal{T}_{k,\nu,\gamma}^{-1}
      - \adk \mathcal{T}_{k,\nu,\gamma}^{-1}[\eta]\big)\Big\}(r-g) \\
&\quad+\; \Pi\Big\{\big(\mathcal{T}_{k+\eta,\pi,\gamma'}^{-1}-\mathcal{T}_{k,\pi,\gamma'}^{-1}\big)
      (I-\nu)\,\big(\mathcal{T}_{k+\eta,\nu,\gamma}^{-1}-\mathcal{T}_{k,\nu,\gamma}^{-1}\big)\Big\}(r-g) .
\end{align*}



\noindent \textbf{Term 1.}  
By Lemma~\ref{lemma::bellmanfrechet},
\[
\|\mathcal T_{k+\eta,\pi, \gamma'}^{-1}-\mathcal T_{k,\pi, \gamma'}^{-1}-(\adk\mathcal T_{k,\pi,\gamma'}^{-1})[\eta]\|_{L^2(\lambda)\to L^2(\lambda)}
= O\!\big(\|\eta\|_{\mathcal K}^2\big).
\]
Hence,
\begin{align*}
    &\Big\|\Pi\Big\{\big(\mathcal{T}_{k+\eta,\pi,\gamma'}^{-1}-\mathcal{T}_{k,\pi,\gamma'}^{-1}
      - (\adk \mathcal{T}_{k,\pi,\gamma'}^{-1})[\eta]\big)(I-\nu)\,\mathcal{T}_{k,\nu,\gamma}^{-1}\Big\}(r-g) \Big\|_{L^2(\lambda)} \\
    &\qquad \le \|\mathcal{T}_{k+\eta,\pi,\gamma'}^{-1}-\mathcal{T}_{k,\pi,\gamma'}^{-1}
      - (\adk \mathcal{T}_{k,\pi,\gamma'}^{-1})[\eta]\|_{L^2(\lambda)\to L^2(\lambda)} \,
      \|(I-\nu)\,\mathcal{T}_{k,\nu,\gamma}^{-1}(r-g)\|_{L^2(\lambda)} \\
    &\qquad = O\!\big(\|\eta\|_{\mathcal K}^2 \, \|r\|_{L^2(\lambda)}\big),
\end{align*}
where the big-oh notation absorbs the constant 
\(\|\mathcal{T}_{k,\nu,\gamma}^{-1}\|_{L^2(\lambda)\to L^2(\lambda)}\), and we use that 
\(\|(I-\nu)\|_{L^2(\lambda)\to L^2(\lambda)} \leq 1 + \|\nu\|_{L^2(\lambda)\to L^2(\lambda)} < \infty\), 
since \(\lambda\) induces a uniform policy on \(\mathcal{A}\).



\noindent \textbf{Term 2.}  
By Lemma~\ref{lemma::bellmanfrechet},
\[
\|\mathcal T_{k+\eta,\nu, \gamma}^{-1}-\mathcal T_{k,\nu, \gamma}^{-1}-(\adk\mathcal T_{k,\nu,\gamma}^{-1})[\eta]\|_{L^2(\lambda)\to L^2(\lambda)}
= O\!\big(\|\eta\|_{\mathcal K}^2\big).
\]
Hence,
\begin{align*}
    &\Big\|\Pi\Big\{\mathcal{T}_{k,\pi,\gamma'}^{-1}(I-\nu)\big(\mathcal{T}_{k+\eta,\nu,\gamma}^{-1}-\mathcal{T}_{k,\nu,\gamma}^{-1}
      - (\adk \mathcal{T}_{k,\nu,\gamma}^{-1})[\eta]\big)\Big\}(r-g) \Big\|_{L^2(\lambda)} \\
    &\qquad \le \|\mathcal{T}_{k+\eta,\nu,\gamma}^{-1}-\mathcal{T}_{k,\nu,\gamma}^{-1}
      - (\adk \mathcal{T}_{k,\nu,\gamma}^{-1})[\eta]\|_{L^2(\lambda)\to L^2(\lambda)} \,
      \|\mathcal{T}_{k,\pi,\gamma'}^{-1}(I-\nu)(r-g)\|_{L^2(\lambda)} \\
    &\qquad = O\!\big(\|\eta\|_{\mathcal K}^2 \, \|r\|_{L^2(\lambda)}\big),
\end{align*}
where the big-oh notation absorbs the constant 
\(\|\mathcal{T}_{k,\pi,\gamma'}^{-1}\|_{L^2(\lambda)\to L^2(\lambda)} \).



\noindent \textbf{Term 3.} By Lemma~\ref{lemma::bellmanfrechet},
\[
\|\mathcal T_{k+\eta,\pi, \gamma'}^{-1}-\mathcal T_{k,\pi, \gamma'}^{-1}-(\adk\mathcal T_{k,\pi,\gamma'}^{-1})[\eta]\|_{L^2(\lambda)\to L^2(\lambda)}
= O\!\big(\|\eta\|_{\mathcal K}^2\big).
\]
Moreover, by the triangle inequality,
\[
\|\mathcal T_{k+\eta,\pi,\gamma'}^{-1}-\mathcal T_{k,\pi,\gamma'}^{-1}\|_{L^2(\lambda)\to L^2(\lambda)}
\lesssim \|\eta\|_{\mathcal K},
\]
since \(\|(\adk\mathcal T_{k,\pi,\gamma'}^{-1})[\eta]\|_{L^2(\lambda)\to L^2(\lambda)} \le C\,\|\eta\|_{\mathcal K}\) for some constant \(C\). By a symmetric argument, we also obtain
\[
\|\mathcal T_{k+\eta,\nu,\gamma}^{-1}-\mathcal T_{k,\nu,\gamma}^{-1}\|_{L^2(\lambda)\to L^2(\lambda)}
\lesssim \|\eta\|_{\mathcal K}.
\]
Hence, taking operator norms, 
\begin{align*}
&\Big\|\Pi\Big\{\mathcal T_{k,\pi,\gamma'}^{-1}(I-\nu)\big(\mathcal T_{k+\eta,\nu,\gamma}^{-1}-\mathcal T_{k,\nu,\gamma}^{-1}-(\adk \mathcal T_{k,\nu,\gamma}^{-1})[\eta]\big) \\
&\hspace{2.5cm}+ \big(\mathcal T_{k+\eta,\pi,\gamma'}^{-1}-\mathcal T_{k,\pi,\gamma'}^{-1}\big)(I-\nu)\big(\mathcal T_{k+\eta,\nu,\gamma}^{-1}-\mathcal T_{k,\nu,\gamma}^{-1}\big)\Big\}(r-g)\Big\|_{L^2(\lambda)} \\
&\qquad \lesssim \|\mathcal T_{k+\eta,\pi,\gamma'}^{-1}-\mathcal T_{k,\pi,\gamma'}^{-1}\|_{L^2(\lambda)\to L^2(\lambda)} \,
\|\mathcal T_{k+\eta,\nu,\gamma}^{-1}-\mathcal T_{k,\nu,\gamma}^{-1}\|_{L^2(\lambda)\to L^2(\lambda)} \,\|r\|_{L^2(\lambda)} \\
&\qquad \lesssim \|\eta\|_{\mathcal K}^2 \,\|r\|_{L^2(\lambda)}.
\end{align*}
Finally, note that 
\begin{align*}
   & \mathcal{V}_{k+\eta}^g(r+ h) - \mathcal{V}_k^g(r) - \adk \mathcal{V}_k^g(r)[\eta] - \partial_r \mathcal{V}_k^g(r)[h] \\
   &= \Big\{\mathcal{V}_{k+\eta}^g(r) - \mathcal{V}_k^g(r) - \adk \mathcal{V}_k^g(r)[\eta]\Big\} 
      + \Big\{\mathcal{V}_{k+\eta}^g(r+h) - \mathcal{V}_{k+\eta}^g(r) - \partial_r \mathcal{V}_k^g(r)[h]\Big\}.
\end{align*}
We claim that
\[
\|\mathcal{V}_{k+\eta}^g(r+h) - \mathcal{V}_{k+\eta}^g(r) - \partial_r \mathcal{V}_k^g(r)[h]\|_{L^2(\lambda)} \;\lesssim\; \|\eta\|_{\mathcal K}\,\|h\|_{L^2(\lambda)},
\]
in which case
\begin{align*}
   & \|\mathcal{V}_{k+\eta}^g(r+ h) - \mathcal{V}_k^g(r) - \adk \mathcal{V}_k^g(r)[\eta] - \partial_r \mathcal{V}_k^g(r)[h]\|_{L^2(\lambda)} \\
   &\leq \|\mathcal{V}_{k+\eta}^g(r) - \mathcal{V}_k^g(r) - \adk \mathcal{V}_k^g(r)[\eta]\|_{L^2(\lambda)}
      + \|\mathcal{V}_{k+\eta}^g(r+h) - \mathcal{V}_{k+\eta}^g(r) - \partial_r \mathcal{V}_k^g(r)[h]\|_{L^2(\lambda)}\\
   &= O(\|\eta\|_{\mathcal K}^2\,\|r\|_{L^2(\lambda)}) +  O(\|\eta\|_{\mathcal K}\,\|h\|_{L^2(\lambda)}),
\end{align*}
and the result follows.

\medskip
\noindent \emph{Proof of the claim.}  
To bound the difference \(\mathcal{V}_{k+\eta}^g(r+h) - \mathcal{V}_{k+\eta}^g(r) - \partial_r \mathcal{V}_k^g(r)[h]\), recall that
\[
\partial_r \mathcal{V}_k^g(r)[h] = \Pi\,\mathcal{T}_{k,\pi,\gamma'}^{-1}(I-\nu)\,\mathcal{T}_{k,\nu,\gamma}^{-1}(h).
\]
Hence,
\begin{align*}
\mathcal{V}_{k+\eta}^g(r+h) - \mathcal{V}_{k+\eta}^g(r) - \partial_r \mathcal{V}_k^g(r)[h]
&= \Pi\Big(\mathcal{T}_{k+\eta,\pi,\gamma'}^{-1}-\mathcal{T}_{k,\pi,\gamma'}^{-1}\Big)(I-\nu)\,\mathcal{T}_{k+\eta,\nu,\gamma}^{-1}(h) \\
&\quad+\; \Pi\,\mathcal{T}_{k,\pi,\gamma'}^{-1}(I-\nu)\Big(\mathcal{T}_{k+\eta,\nu,\gamma}^{-1}-\mathcal{T}_{k,\nu,\gamma}^{-1}\Big)(h).
\end{align*}
Taking $L^2(\lambda)$ norms gives
\begin{align*}
\big\|\mathcal{V}_{k+\eta}^g(r+h) - \mathcal{V}_{k+\eta}^g(r)
&\qquad - \partial_r \mathcal{V}_k^g(r)[h]\big\|_{L^2(\lambda)} \\
&\le \|\mathcal{T}_{k+\eta,\pi,\gamma'}^{-1}-\mathcal{T}_{k,\pi,\gamma'}^{-1}\|_{L^2(\lambda)\to L^2(\lambda)}
   \,\|(I-\nu)\,\mathcal{T}_{k+\eta,\nu,\gamma}^{-1}(h)\|_{L^2(\lambda)} \\
&\quad+\; \|\mathcal{T}_{k,\pi,\gamma'}^{-1}\|_{L^2(\lambda)\to L^2(\lambda)} \,
   \|(I-\nu)\|_{L^2(\lambda)\to L^2(\lambda)} \,
   \|\mathcal{T}_{k+\eta,\nu,\gamma}^{-1}-\mathcal{T}_{k,\nu,\gamma}^{-1}\|_{L^2(\lambda)\to L^2(\lambda)} \,
   \|h\|_{L^2(\lambda)}.
\end{align*}

By the Lipschitz continuity of the resolvent maps
$k \mapsto \mathcal{T}_{k,\pi,\gamma'}^{-1}$ and
$k \mapsto \mathcal{T}_{k,\nu,\gamma}^{-1}$, we have
\[
\begin{aligned}
\|\mathcal{T}_{k+\eta,\pi,\gamma'}^{-1}-\mathcal{T}_{k,\pi,\gamma'}^{-1}\|_{L^2(\lambda)\to L^2(\lambda)}
&\lesssim \|\eta\|_{\mathcal K}, \\
\|\mathcal{T}_{k+\eta,\nu,\gamma}^{-1}-\mathcal{T}_{k,\nu,\gamma}^{-1}\|_{L^2(\lambda)\to L^2(\lambda)}
&\lesssim \|\eta\|_{\mathcal K}.
\end{aligned}
\]
Substituting these bounds and absorbing the operator norms of $(I-\nu)$, $\Pi$, and 
$\mathcal{T}_{k,\pi,\gamma'}^{-1}$ into the implied constant, we conclude that
\[
\begin{aligned}
\|\mathcal{V}_{k+\eta}^g(r+h) - \mathcal{V}_{k+\eta}^g(r)
- \partial_r \mathcal{V}_k^g(r)[h]\|_{L^2(\lambda)}
\lesssim \|\eta\|_{\mathcal K}\,\|h\|_{L^2(\lambda)}.
\end{aligned}
\]



\end{proof}

 



 
 \subsubsection{Lemmas for Example~\ref{example::1b}}

 

In the following lemma, we recall the softmax policy $\pi_{r,k}^\star$ and the solution $v_{r,k}^\star$ of the soft Bellman equation from Example~\ref{example::1b}. For each $(r,k)$, define the operator $(\Pi_{r,k}^\star f)(s) := \sum_{\tilde a \in \mathcal{A}^\star} \pi_{r,k}^\star(\tilde a \mid s)\,f(\tilde a,s)$, and the log-partition function $\Phi_{r,k}^\star(s) := \tau^\star \log\!\big(\sum_{\tilde a \in \mathcal{A}^\star} \exp((r(\tilde a,s) + \gamma\,v_{r,k}^\star(\tilde a,s))/\tau^\star)\big)$. Define the conditional covariance under $\pi_{r,k}^\star$ by
\[
\operatorname{Cov}_{\pi_{r,k}^\star}\!\big(f,g\big)(s)
:= \big(\Pi_{r,k}^\star[fg]\big)(s)
- \big(\Pi_{r,k}^\star f\big)(s)\,\big(\Pi_{r,k}^\star g\big)(s), 
\qquad f,g \in L^\infty(\lambda).
\]

 
 



\begin{lemma}[Derivatives for soft Bellman operators]
\label{lemma::derivSoftBellmanStar}
The map $(r,k) \mapsto v_{r,k}^\star$ is sup-norm Fréchet differentiable as a map from $L^\infty(\lambda) \otimes L^\infty(\mu \otimes \lambda)$ to $L^\infty(\lambda)$, with
\[
\partial_r v_{r,k}^\star[h] = (\mathcal{T}_{r,k}^\star)^{-1}\,\mathcal{P}_k\,\Pi_{r,k}^\star h,
\qquad
\adk v_{r,k}^\star[w] = (\mathcal{T}_{r,k}^\star)^{-1}\,\mathcal{P}_w\,\Phi_{r,k}^\star.
\]
Moreover, for any test function $f$,
\[
\partial_r \Pi_{r,k}^\star[h] f
= \tfrac{1}{\tau^\star}\,
\operatorname{Cov}_{\pi_{r,k}^\star}\!\big(f,\,(\mathcal{T}_{r,k}^\star)^{-1}h\big),
\qquad
\adk \Pi_{r,k}^\star[w] f
= \tfrac{1}{\tau^\star}\,
\operatorname{Cov}_{\pi_{r,k}^\star}\!\big(f,\,\gamma\,\adk v_{r,k}^\star[w]\big).
\]
Consequently,
\[
\begin{aligned}
\partial_r \mathcal{T}_{r,k}^\star[h] f
&= -\,\tfrac{\gamma}{\tau^\star}\,
\mathcal{P}_k\,
\operatorname{Cov}_{\pi_{r,k}^\star}\!\big(f,\,(\mathcal{T}_{r,k}^\star)^{-1}h\big), \\[6pt]
\adk \mathcal{T}_{r,k}^\star[w] f
&= -\,\gamma\,\mathcal{P}_w\,\Pi_{r,k}^\star f
- \tfrac{\gamma}{\tau^\star}\,
\mathcal{P}_k\,
\operatorname{Cov}_{\pi_{r,k}^\star}\!\big(f,\,\gamma\,\adk v_{r,k}^\star[w]\big),
\end{aligned}
\]
and
\[
\partial_r \!\left[ \Pi_{r,k}^\star (\mathcal{T}_{r,k}^\star)^{-1}\right][h]
= \Big(I + \gamma\,\Pi_{r,k}^\star (\mathcal{T}_{r,k}^\star)^{-1}\mathcal{P}_k\Big)\,
\big(\partial_r \Pi_{r,k}^\star[h]\big)\,(\mathcal{T}_{r,k}^\star)^{-1}.
\]
\[
\adk \!\left[ \Pi_{r,k}^\star (\mathcal{T}_{r,k}^\star)^{-1}\right][w]
= \big(I-\gamma\,\Pi_{r,k}^\star \mathcal{P}_k\big)^{-1}
\Big[\,\big(\adk \Pi_{r,k}^\star[w]\big) + \gamma\,\Pi_{r,k}^\star \mathcal{P}_w \Pi_{r,k}^\star\,\Big]
(\mathcal{T}_{r,k}^\star)^{-1}.
\]
\end{lemma}

\begin{proof}
This is the direct specialization of Appendix~\ref{appendix::technicalsoftmax} to the counterfactual system in Example~\ref{example::1b}: replace the ambient action set $\mathcal{A}$ by $\mathcal{A}^\star$ and the fixed temperature $\tau$ by $\tau^\star$ in Lemmas~\ref{lemma::IFT}, \ref{lemma::qderiv}, and \ref{lemma::operatorderivatives}. This yields the displayed formulas for $\partial_r v_{r,k}^\star$, $\adk v_{r,k}^\star$, $\partial_r \Pi_{r,k}^\star$, and $\adk \Pi_{r,k}^\star$.

Since $\mathcal{T}_{r,k}^\star = I - \gamma \mathcal{P}_k \Pi_{r,k}^\star$, applying the chain rule yields
\[
\begin{aligned}
\partial_r \mathcal{T}_{r,k}^\star[h] f
&= -\,\gamma\,\mathcal{P}_k\,\big(\partial_r \Pi_{r,k}^\star[h]\big)f
= -\,\tfrac{\gamma}{\tau^\star}\,\mathcal{P}_k\,\operatorname{Cov}_{\pi_{r,k}^\star}\!\big(f,\,(\mathcal{T}_{r,k}^\star)^{-1}h\big), \\[6pt]
\adk \mathcal{T}_{r,k}^\star[w] f
&= -\,\gamma\,\mathcal{P}_w\,\Pi_{r,k}^\star f
- \gamma\,\mathcal{P}_k\,\big(\adk \Pi_{r,k}^\star[w]\big)f \\[3pt]
&= -\,\gamma\,\mathcal{P}_w\,\Pi_{r,k}^\star f
- \tfrac{\gamma}{\tau^\star}\,\mathcal{P}_k\,\operatorname{Cov}_{\pi_{r,k}^\star}\!\big(f,\,\gamma\,\adk v_{r,k}^\star[w]\big).
\end{aligned}
\]
This establishes the derivative identities for $\mathcal{T}_{r,k}^\star$. 




Let $G(r,k):=\Pi_{r,k}^\star(\mathcal{T}_{r,k}^\star)^{-1}$. Using the product rule and the inverse-derivative identity
\[
\adk(\mathcal{T}_{r,k}^\star)^{-1}[w]
= -\,(\mathcal{T}_{r,k}^\star)^{-1}\,\big(\adk \mathcal{T}_{r,k}^\star[w]\big)\,(\mathcal{T}_{r,k}^\star)^{-1},
\]
we get
\[
\adk G[w]
= \big(\adk \Pi_{r,k}^\star[w]\big)(\mathcal{T}_{r,k}^\star)^{-1}
+ \Pi_{r,k}^\star\,\adk(\mathcal{T}_{r,k}^\star)^{-1}[w].
\]
Since $\mathcal{T}_{r,k}^\star = I-\gamma\,\mathcal{P}_k \Pi_{r,k}^\star$,
\[
\adk \mathcal{T}_{r,k}^\star[w]
= -\gamma\big(\mathcal{P}_w \Pi_{r,k}^\star + \mathcal{P}_k\,\adk \Pi_{r,k}^\star[w]\big).
\]
Substituting gives
\[
\adk G[w]
= \big(\adk \Pi_{r,k}^\star[w]\big)(\mathcal{T}_{r,k}^\star)^{-1}
+ \gamma\,\Pi_{r,k}^\star(\mathcal{T}_{r,k}^\star)^{-1}
\big(\mathcal{P}_w \Pi_{r,k}^\star + \mathcal{P}_k\,\adk \Pi_{r,k}^\star[w]\big)
(\mathcal{T}_{r,k}^\star)^{-1}.
\]
Regrouping,
\[
\adk G[w]
= \Big(I+\gamma\,\Pi_{r,k}^\star(\mathcal{T}_{r,k}^\star)^{-1}\mathcal{P}_k\Big)\,
\big(\adk \Pi_{r,k}^\star[w]\big)\,(\mathcal{T}_{r,k}^\star)^{-1}
\;+\;
\gamma\,\Pi_{r,k}^\star(\mathcal{T}_{r,k}^\star)^{-1}\mathcal{P}_w \Pi_{r,k}^\star(\mathcal{T}_{r,k}^\star)^{-1}.
\]
Finally, by the resolvent identity $(I-AB)^{-1}=I+A(I-BA)^{-1}B$ applied with $A=\gamma\,\Pi_{r,k}^\star$, $B=\mathcal{P}_k$,
\[
I+\gamma\,\Pi_{r,k}^\star(\mathcal{T}_{r,k}^\star)^{-1}\mathcal{P}_k
= (I-\gamma\,\Pi_{r,k}^\star \mathcal{P}_k)^{-1},
\quad
\Pi_{r,k}^\star(\mathcal{T}_{r,k}^\star)^{-1}
= (I-\gamma\,\Pi_{r,k}^\star \mathcal{P}_k)^{-1}\Pi_{r,k}^\star,
\]
which yields the stated formula.



   
\end{proof}


 
\begin{lemma}[Linearization remainder for Example 1b]
\label{lemma::linearizationExamplestar}
Assume \ref{cond::stationary3}. Then there exists an $L^2$-ball around $(r_0,k_0) \in \mathcal{H} \otimes \mathcal{K}$ such that $(\mathcal{T}_{r,k}^\star)^{-1}: L^2(\lambda) \to L^2(\lambda)$ exists. Let $m(a,s,r,k) = \Pi_{r,k}^\star\,(\mathcal{T}_{r,k}^\star)^{-1}(r)(s)$.  
Then
\[
\operatorname{Rem}_0(r_0,k_0; r,k) \;\lesssim\; \|k - k_0\|_{\mathcal{K}}^{2} 
+ \|r - r_0\|_{\mathcal{H}}^{2} 
+ \|v_{r_0,k}-v_{r_0,k_0}\|_{L^2(\lambda)}^{2},
\]
with the implicit constant depending on $\gamma$, $\|r_0\|_{L^\infty(\lambda)}$, and $\|(\mathcal{T}_{r_0,k_0}^\star)^{-1}\|_{L^2(\lambda)\to L^2(\lambda)}$.
\end{lemma}
\begin{proof}
For this proof, suppress the star and write $v_{r,k}$, $\Pi_{r,k}$, and $\mathcal{T}_{r,k}$ for the soft Bellman objects associated with the restricted action set $\mathcal{A}^\star$ and fixed temperature $\tau^\star$ from Example~\ref{example::1b}. Then $m(a,s,r,k) = \Pi_{r,k}\,\mathcal{T}_{r,k}^{-1}(r)(s)$.

By \ref{cond::stationary3} and Lemma \ref{lemma::invertibility} applied with $\pi = \pi_{r_0,k_0}^\star$, the inverse $\mathcal{T}_{r_0,k_0}^{-1}: L^2(\lambda) \rightarrow L^2(\lambda)$ exists. Moreover, the proofs of Lemmas~\ref{lemma::valuestarderivativereward}--\ref{lemma::jointderivative} apply verbatim to this restricted-action, temperature-$\tau^\star$ system, since Appendix~\ref{appendix::technicalsoftmax} was established for an arbitrary fixed temperature $\tau > 0$. In particular, there exists an $L^2$-ball around $(r_0,k_0) \in \mathcal{H} \otimes \mathcal{K}$ such that $\mathcal{T}_{r,k}^{-1}: L^2(\lambda) \to L^2(\lambda)$ exists. By the assumed uniform boundedness of the state density $\rho_0$, we have 
\begin{align*}
    &\left|E_0\!\left[\Pi_{r,k}\,\mathcal{T}_{r,k}^{-1}(r) 
- \Pi_{r_0,k_0}\,\mathcal{T}_{r_0,k_0}^{-1}(r_0) 
- \partial_r\!\big[\Pi_{r_0,k_0}\,\mathcal{T}_{r_0,k_0}^{-1}(r_0)\big][\,r-r_0\,]
- \adk\!\big[\Pi_{r_0,k_0}\,\mathcal{T}_{r_0,k_0}^{-1}(r_0)\big][\,k-k_0\,]\right]\right| \\
&\;\;\lesssim\; 
\big\|\Pi_{r,k}\,\mathcal{T}_{r,k}^{-1}(r) 
- \Pi_{r_0,k_0}\,\mathcal{T}_{r_0,k_0}^{-1}(r_0) 
- \partial_r\!\big[\Pi_{r_0,k_0}\,\mathcal{T}_{r_0,k_0}^{-1}(r_0)\big][\,r-r_0\,]
- \adk\!\big[\Pi_{r_0,k_0}\,\mathcal{T}_{r_0,k_0}^{-1}(r_0)\big][\,k-k_0\,]\big\|_{L^1(\lambda)}.
\end{align*}
By Lemma \ref{lemma::jointderivative}, we have
\begin{align*}
&\big\|\Pi_{r,k}\,\mathcal{T}_{r,k}^{-1}(r) 
- \Pi_{r_0,k_0}\,\mathcal{T}_{r_0,k_0}^{-1}(r_0) 
- \partial_r\!\big[\Pi_{r_0,k_0}\,\mathcal{T}_{r_0,k_0}^{-1}(r_0)\big][\,r-r_0\,]
- \adk\!\big[\Pi_{r_0,k_0}\,\mathcal{T}_{r_0,k_0}^{-1}(r_0)\big][\,k-k_0\,]\big\|_{L^1(\lambda)} \\
&\;\;=\; O\!\left(\|r-r_0\|_{L^2(\lambda)}^{2} 
+ \|v_{r,k}-v_{r_0,k_0}\|_{L^2(\lambda)}^{2} 
+ \|k-k_0\|_{\mathcal K}^{2}\right).
\end{align*}
Furthermore, by Lemma \ref{lemma::lipschitzvalue}, 
\[
\|v_{r,k}-v_{r_0,k_0}\|_{L^2(\lambda)}^{2} 
\;\lesssim\; \|v_{r,k}-v_{r_0,k}\|_{L^2(\lambda)}^{2} 
+ \|v_{r_0,k}-v_{r_0,k_0}\|_{L^2(\lambda)}^{2} 
\;\lesssim\; \|r-r_0\|_{L^2(\lambda)}^{2} 
+ \|v_{r_0,k}-v_{r_0,k_0}\|_{L^2(\lambda)}^{2}.
\]
Hence,
\begin{align*}
&\big\|\Pi_{r,k}\,\mathcal{T}_{r,k}^{-1}(r) 
- \Pi_{r_0,k_0}\,\mathcal{T}_{r_0,k_0}^{-1}(r_0) 
- \partial_r\!\big[\Pi_{r_0,k_0}\,\mathcal{T}_{r_0,k_0}^{-1}(r_0)\big][\,r-r_0\,]
- \adk\!\big[\Pi_{r_0,k_0}\,\mathcal{T}_{r_0,k_0}^{-1}(r_0)\big][\,k-k_0\,]\big\|_{L^1(\lambda)} \\
&\;\;=\; O\!\left(\|r-r_0\|_{L^2(\lambda)}^{2} 
+ \|k-k_0\|_{\mathcal K}^{2} 
+ \|v_{r_0,k}-v_{r_0,k_0}\|_{L^2(\lambda)}^{2}\right).
\end{align*}
We conclude that
\begin{align*}
&\Big|E_0\!\Big[\Pi_{r,k}\,\mathcal{T}_{r,k}^{-1}(r) 
- \Pi_{r_0,k_0}\,\mathcal{T}_{r_0,k_0}^{-1}(r_0) 
- \partial_r\!\big[\Pi_{r_0,k_0}\,\mathcal{T}_{r_0,k_0}^{-1}(r_0)\big][\,r-r_0\,] \\
&\qquad\qquad
- \adk\!\big[\Pi_{r_0,k_0}\,\mathcal{T}_{r_0,k_0}^{-1}(r_0)\big][\,k-k_0\,]\Big]\Big| \\
&\quad=\; O\!\left(\|r-r_0\|_{L^2(\lambda)}^{2} 
+ \|k-k_0\|_{\mathcal K}^{2} 
+ \|v_{r_0,k}-v_{r_0,k_0}\|_{L^2(\lambda)}^{2}\right).
\end{align*}
and, hence,
\[
\operatorname{Rem}_0(r_0,k_0; r,k)\;\lesssim\;\|k - k_0\|_{\mathcal{K}}^{2} 
+  \|r - r_0\|_{\mathcal{H}}^{2} 
+ \|v_{r_0,k}-v_{r_0,k_0}\|_{L^2(\lambda)}^{2}.
\]

\end{proof}


 






\subsection{Proofs of EIFs in examples}

 
\label{appendix::eifexproofs}

 \begin{proof}[Proof of Theorem \ref{theorem::EIFvalue}]

Write $\lambda_0 := \pi_0 \otimes \rho_0$ and $\mathcal{T}_{k,\pi} := \mathcal{T}_{k,\pi,\gamma} = I - \gamma \mathcal{P}_k \Pi$. By \ref{cond::stationary} and Lemma~\ref{lemma::invertibility}, the linear operator $\mathcal{T}_{k_0,\pi}:L^2(\lambda) \rightarrow L^2(\lambda)$ is bounded and invertible, with Neumann-series inverse
\[
\mathcal{T}_{k_0,\pi}^{-1} = \sum_{t=0}^\infty (\gamma \mathcal{P}_{k_0}\Pi)^t.
\]
Hence,
\[
m(a,s,r,k) 
= \sum_{\tilde a \in \mathcal{A}} \pi(\tilde a \mid s)\, (\mathcal{T}_{k,\pi})^{-1}(r)(\tilde a,s) 
= \Pi \bigl[(\mathcal{T}_{k,\pi})^{-1}(r)\bigr](s)
\]
is well defined for all $k \in \mathcal{K}$ and $r \in L^2(\lambda)$ in a ball around $(r_0, k_0)$.
 

By the chain rule, for each \(h \in T_{\mathcal{H}}\) and \(\beta \in T_{\mathcal{K}}\), we have
\begin{align*}
    \partial_r m(A,S,r_0,k_0)[h] 
    &= \Pi (\mathcal{T}_{k_0,\pi})^{-1}(h)(s), \\
    \partial_k m(A,S,r_0,k_0)[\beta] 
    &= \Pi  \bigl(\adk \mathcal{T}_{k,\pi}^{-1}[k_0\beta] \mid_{k = k_0} \bigr)(r_0)(s) \\
    &= \gamma\, \Pi \mathcal{T}_{k_0,\pi}^{-1}\,\mathcal{P}_{k_0\beta}\,\Pi\,\mathcal{T}_{k_0,\pi}^{-1}(r_0)(s),
\end{align*}
where we used the additive derivative formula with kernel direction \(g = k_0\beta\), so that \(\adk(\mathcal{T}_{k,\pi}^{-1})[k_0\beta] = -\,\mathcal{T}_{k,\pi}^{-1}\big(\adk \mathcal{T}_{k,\pi}[k_0\beta]\big)\mathcal{T}_{k,\pi}^{-1}\) and \(\adk \mathcal{T}_{k,\pi}[k_0\beta] = -\,\gamma\,\mathcal{P}_{k_0\beta}\Pi\). 

 


Moreover, since $\Pi$, $\mathcal{P}_{k_0\beta}$, and $\mathcal{T}_{k_0,\pi}^{-1}$ are all bounded linear operators on $L^2(\lambda)$, it follows that $h \mapsto E_0[ \partial_r m(A,S,r_0,k_0)[h]]$ is a bounded linear functional on $(\mathcal{H}, \| \cdot\|_{\mathcal{H}})$. In addition, $\esssup_{a,s} \left|\frac{\pi_0(a \mid s)}{\exp \{r_0(a, s)\}}\right| < \infty$ since $r_0 \in L^\infty(\lambda)$, and hence we have $\| \cdot\|_{\mathcal{H}} \lesssim  \| \cdot\|_{\mathcal{H}, 0}$, so it is also a bounded linear functional on $(T_{\mathcal{H}}, \| \cdot\|_{\mathcal{H}, 0})$. Furthermore, the map $\beta \mapsto \mathcal{P}_{k_0\beta}$ is bounded in operator norm from $(T_{\mathcal{K}}, \| \cdot\|_{\mathcal{K}, 0})$ to the space of bounded linear operators on $L^2(\lambda)$. Hence, $\beta \mapsto   E_0[ \partial_k m(A,S,r_0,k_0)[\beta] ]$ is a bounded linear functional, since, by properties of operator norms,
\begin{align*}
    \left|E_0[ \partial_k m(A,S,r_0,k_0)[\beta] ]\right|
    & \leq \| \Pi \mathcal{T}_{k_0,\pi}^{-1}\,\mathcal{P}_{k_0\beta}\,\Pi\,\mathcal{T}_{k_0,\pi}^{-1}(r_0)\|_{L^2(\rho_0 \otimes \pi_0)} \\
    & \lesssim \| \Pi \mathcal{T}_{k_0,\pi}^{-1}\,\mathcal{P}_{k_0\beta}\,\Pi\,\mathcal{T}_{k_0,\pi}^{-1}(r_0)\|_{L^2(\lambda)} \\
    & \lesssim \|\mathcal{P}_{k_0\beta}\|_{L^2(\lambda) \rightarrow L^2(\lambda)} \,\|\Pi\,\mathcal{T}_{k_0,\pi}^{-1}(r_0)\|_{L^2(\lambda)} \\
    & \lesssim  \|\beta\|_{\mathcal{K}, 0},
\end{align*}
where we used that the $\lambda$ and $\lambda_0$ measures are equivalent by assumption.
Hence, \ref{cond::boundedfunc} holds under the stated conditions. Moreover, \ref{cond::lipschitz} follows from Lemma \ref{lemma::bellmanfrechet} applied with the triangle inequality and boundedness of the derivative.

We now determine the precise forms of the Riesz representers $\alpha_0$ and $\beta_0$ appearing in Theorem \ref{theorem::EIF}. We begin with $\alpha_0$. For each $\alpha \in T_{\mathcal{H}}$, we have
\begin{align*}
     E_0[ \partial_r m(A,S,r_0,k_0)[\alpha]] 
     &= E_0\!\left[\Pi (\mathcal{T}_{k_0,\pi})^{-1}(\alpha)(S)\right] \\
     &= E_0[d_0^{\pi,\gamma}(A,S)\,\alpha(A,S)],
\end{align*}
where the second equality follows from Lemma~\ref{lemma::adjointident}, and $d_0^{\pi,\gamma}$ is the discounted state–action occupancy ratio defined previously.
It follows that
$$\alpha_0(a,s) = d_0^{\pi,\gamma}(a,s) = \rho_0^{\pi,\gamma}(s) \frac{\pi(a \mid s)}{\pi_0(a\mid s)}.$$

Next, we determine $\beta_0$. Note
\begin{align*}
    E_0 [\partial_k m(A,S,r_0,k_0)[\beta] ]  
    &= \gamma\, E_0\!\left[\Pi \mathcal{T}_{k_0,\pi}^{-1}\,\mathcal{P}_{k_0\beta}\,\Pi\,\mathcal{T}_{k_0,\pi}^{-1}(r_0)(S)\right] \\
    &=  \gamma \langle 1, \mathcal{T}_{k_0,\pi}^{-1}\,\mathcal{P}_{k_0\beta}\,\Pi\,\mathcal{T}_{k_0,\pi}^{-1}(r_0) \rangle_{\rho_0 \otimes \pi} \\
     &=  \gamma  \langle \tfrac{\pi}{\pi_0}, \mathcal{T}_{k_0,\pi}^{-1}\,\mathcal{P}_{k_0\beta}\,\Pi\,\mathcal{T}_{k_0,\pi}^{-1}(r_0) \rangle_{\lambda_0}\\
      &=  \gamma \langle (\mathcal{T}_{k_0,\pi}^{-1})^*(\tfrac{\pi}{\pi_0}),  \, \mathcal{P}_{k_0\beta}\,\Pi\,\mathcal{T}_{k_0,\pi}^{-1}(r_0) \rangle_{\lambda_0} \\
       &=  \gamma  \langle d_0^{\pi,\gamma},  \, \mathcal{P}_{k_0\beta}\,\Pi\,\mathcal{T}_{k_0,\pi}^{-1}(r_0) \rangle_{\lambda_0},
\end{align*}
where $d_0^{\pi,\gamma} =  (\mathcal{T}_{k_0,\pi}^{-1})^*(\tfrac{\pi}{\pi_0})$ by the adjoint Poisson equation. Recall that $V_{r_0, k_0}^{\pi,\gamma} = \Pi\,\mathcal{T}_{k_0,\pi}^{-1}(r_0)$ is the value function. Hence,
\begin{align*}
    E_0 \big[\partial_k m(A,S,r_0,k_0)[\beta] \big] 
    &= \gamma\, \langle d_0^{\pi,\gamma},\, \mathcal{P}_{k_0\beta}\,\Pi\,\mathcal{T}_{k_0,\pi}^{-1}(r_0) \rangle_{\lambda_0} \\
    &= \gamma\, \langle d_0^{\pi,\gamma},\, \mathcal{P}_{k_0\beta}\,V_{r_0, k_0}^{\pi,\gamma}  \rangle_{\lambda_0} \\
    &= E_0\!\left[\gamma\, d_0^{\pi,\gamma}(A,S)\, V_{r_0, k_0}^{\pi,\gamma}(S')\, \beta(S'\mid A,S)\right] \\
    &= E_0\!\left[\gamma\, d_0^{\pi,\gamma}(A,S)\,\bigl\{V_{r_0, k_0}^{\pi,\gamma}(S') - \mathcal{P}_{k_0}V_{r_0, k_0}^{\pi,\gamma}(A,S)\bigr\}\, \beta(S'\mid A,S)\right].
\end{align*}
 It follows that
$$\beta_0(s' \mid a, s) = \gamma d_0^{\pi,\gamma}(a, s) \{V_{r_0, k_0}^{\pi,\gamma} (s') - \mathcal{P}_{k_0} V_{r_0, k_0}^{\pi,\gamma}(a,s)\} .$$
By Bellman’s equation,
\[
(\mathcal{P}_{k_0} V_{r_0,k_0}^{\pi,\gamma})(a,s)
= \frac{q_{r_0,k_0}^{\pi,\gamma}(a,s) - r_0(a,s)}{\gamma}.
\]
Substituting into the integrand,
\[
\beta_0(s' \mid a, s) = d_0^{\pi,\gamma}(a,s)\Big\{r_0(a,s) + \gamma\,V_{r_0,k_0}^{\pi,\gamma}(s') - q_{r_0,k_0}^{\pi,\gamma}(a,s)  \Big\}.
\]


Putting everything together, we conclude that the EIF of $\Psi$ is
\begin{align*}
    \chi_0(s,a,s') 
    &= \alpha_0(a,s) 
      - \sum_{\tilde a \in \mathcal{A}} \pi_0(\tilde a \mid s)\,\alpha_0(\tilde a, s) 
      + \beta_0(s' \mid a, s) + V_{r_0, k_0}^{\pi,\gamma}(s) - \Psi(P_0) \\
      &= \alpha_0(a,s) 
      - \sum_{\tilde a \in \mathcal{A}} \pi_0(\tilde a \mid s)\,\alpha_0(\tilde a, s)  + d_0^{\pi,\gamma}(a,s)\Big\{r_0(a,s) + \gamma\,V_{r_0,k_0}^{\pi,\gamma}(s') - q_{r_0,k_0}^{\pi,\gamma}(a,s)  \Big\} \\
      & \quad + V_{r_0, k_0}^{\pi,\gamma}(s) - \Psi(P_0).
\end{align*}
In the nonparametric case, where $\mathcal{H} = L^\infty(\lambda)$, we have that $\alpha_0 = d_0^{\pi,\gamma}$ and, hence,
 \begin{align*}
    \chi_0(s,a,s') 
   &= d_0^{\pi,\gamma}(a,s) 
      - \sum_{\tilde a \in \mathcal{A}} \pi_0(\tilde a \mid s)\,d_0^{\pi,\gamma}(\tilde a, s)  \\
      &+ d_0^{\pi,\gamma}(a,s)\Big\{r_0(a,s) + \gamma\,V_{r_0,k_0}^{\pi,\gamma}(s') - q_{r_0,k_0}^{\pi,\gamma}(a,s)  \Big\}+ V_{r_0, k_0}^{\pi,\gamma}(s) - \Psi(P_0) .
\end{align*}




The bound for the linearization remainder follows from Lemma~\ref{lemma::bellmanfrechet}. In particular, by boundedness of $\rho_0$, we obtain
\begin{align*}
   |\operatorname{Rem}_0(r_0,k_0;\,r,k)| 
   &\leq \|\mathcal T_{k,\pi}^{-1}(r) - \mathcal T_{k_0,\pi}^{-1}(r_0) 
   - \adk \mathcal T_{k_0,\pi}^{-1}[\,k-k_0\,](r_0) 
   - \partial_r (\mathcal T_{k_0,\pi}^{-1} r_0)[\,r-r_0\,]\|_{L^1(\pi_0 \otimes \rho_0)} \\
   &\lesssim \|\mathcal T_{k,\pi}^{-1}(r) - \mathcal T_{k_0,\pi}^{-1}(r_0) 
   - \adk \mathcal T_{k_0,\pi}^{-1}[\,k-k_0\,](r_0) 
   - \partial_r (\mathcal T_{k_0,\pi}^{-1} r_0)[\,r-r_0\,]\|_{L^2(\lambda)} \\ 
   &\lesssim \|k-k_0\|_{\mathcal K}^2 
   + \|k-k_0\|_{\mathcal K}\,\|r-r_0\|_{L^2(\lambda)}.
\end{align*}



 \end{proof}

 


 

\begin{proof}[Proof of Theorem \ref{theorem::EIFvaluenorm}]
We define
\[
m(a,s,r,k) := \sum_{\tilde a \in \mathcal{A}} \pi(\tilde a \mid s)\,
   \mathcal{T}_{k,\pi,\gamma'}^{-1}(r_{r,k,\nu,\gamma,g})(\tilde a, s),
   \qquad r_{r,k,\nu,\gamma,g} = g + (I - \nu)\,\mathcal{T}_{k,\nu,\gamma}^{-1}(r-g).
\]
As shown in Appendix~\ref{lemmas::norm}, we can equivalently write
\[
m(a,s,r,k) = \mathcal{V}_k^g(r)(s),
\qquad 
\mathcal{V}_k^g(r) := \Pi\,\mathcal{T}_{k,\pi,\gamma'}^{-1}\!\left\{g + (I - \nu)\,\mathcal{T}_{k,\nu,\gamma}^{-1}(r-g)\right\}.
\]

By Conditions~\ref{cond::stationary} and~\ref{cond::stationary2}, together with 
Lemma~\ref{lemma::invertibility} applied with $\pi=\pi$ and $\pi=\nu$, the operators 
$\mathcal{T}_{k_0,\pi,\gamma'}:L^2(\lambda)\to L^2(\lambda)$ and 
$\mathcal{T}_{k_0,\nu,\gamma}:L^2(\lambda)\to L^2(\lambda)$ are bounded and invertible. 
Thus $\mathcal{V}_k^g(r)$ is well defined in an $L^2$-neighborhood of $(r_0,k_0)$. 
The Fréchet differentiability of $m$ is established formally in 
Lemma~\ref{lemma::normalizedvaluedifferentiable}. Specifically, applied with additive direction \(k\beta\), that lemma yields the score derivative
\begin{align*}
\partial_r \mathcal{V}_k^g(r)[\alpha]
&= \Pi\,\mathcal{T}_{k,\pi,\gamma'}^{-1}(I-\nu)\,\mathcal{T}_{k,\nu,\gamma}^{-1}(\alpha), \\[6pt]
\partial_k \mathcal{V}_k^g(r)[\beta]
&= \gamma'\,\Pi\,\mathcal{T}_{k,\pi,\gamma'}^{-1}(\mathcal{P}_{k\beta} \Pi)\,
      \mathcal{T}_{k,\pi,\gamma'}^{-1}\bigl(r_{r,k,\nu,\gamma,g}\bigr) \\
&\quad+\; \gamma\,\Pi\,\mathcal{T}_{k,\pi,\gamma'}^{-1}(I-\nu)\,
      \mathcal{T}_{k,\nu,\gamma}^{-1}(\mathcal{P}_{k\beta} \nu)\,
      \mathcal{T}_{k,\nu,\gamma}^{-1}(r-g).
\end{align*}
Hence, for each $\alpha \in T_{\mathcal{H}}$ and $\beta \in T_{\mathcal{K}}$,
\begin{align}
E_0\!\left[\partial_r m(A,S,r_0,k_0)[\alpha]\right]
&=\; E_0\!\left[\Pi\,\mathcal{T}_{k_0,\pi,\gamma'}^{-1}(I-\nu)\,\mathcal{T}_{k_0,\nu,\gamma}^{-1}(\alpha)(S)\right],
\label{eq:drm} \\[6pt]
E_0\!\left[\partial_k m(A,S,r_0,k_0)[\beta]\right]
\;=\; \gamma'\,E_0\!\left[\Pi\,\mathcal{T}_{k_0,\pi,\gamma'}^{-1}(\mathcal{P}_{k_0\beta} \Pi)\,
      \mathcal{T}_{k_0,\pi,\gamma'}^{-1}\bigl(r_{0,\nu,\gamma,g}\bigr)(S)\right] \nonumber \\
&\quad+\; \gamma\,E_0\!\left[\Pi\,\mathcal{T}_{k_0,\pi,\gamma'}^{-1}(I-\nu)\,
      \mathcal{T}_{k_0,\nu,\gamma}^{-1}(\mathcal{P}_{k_0\beta} \nu)\,
      \mathcal{T}_{k_0,\nu,\gamma}^{-1}(r_0-g)(S)\right].
\label{eq:dkm-split}
\end{align}
The linear functionals $\alpha \mapsto E_0\!\left[\partial_r m(A,S,r_0,k_0)[\alpha]\right]$ and 
$\beta \mapsto E_0\!\left[\partial_k m(A,S,r_0,k_0)[\beta]\right]$ are compositions of bounded operators and are therefore bounded. 
Hence, Condition~\ref{cond::boundedfunc} holds under the stated assumptions. 
Condition~\ref{cond::lipschitz} then follows from the Fréchet differentiability of the map in 
Lemma~\ref{lemma::normalizedvaluedifferentiable}, together with the triangle inequality and the 
boundedness of the derivatives.

We now turn to determining the specific form of the EIF. We begin by deriving the Riesz representer $\beta_0$ for the kernel partial derivative $E_0[\partial_k m(A,S,r_0,k_0)[\beta]]$.  
An argument identical to the proof of Theorem~\ref{theorem::EIFvalue} shows that the first term on the right-hand side of \eqref{eq:dkm-split} is
\begin{equation}
\label{eqnproof::beta1}
\begin{aligned}
\gamma'\,E_0\!\left[\Pi\,\mathcal{T}_{k_0,\pi,\gamma'}^{-1}(\mathcal{P}_{k_0\beta} \Pi)\,
      \mathcal{T}_{k_0,\pi,\gamma'}^{-1}\bigl(r_{0,\nu,\gamma,g}\bigr)(S)\right] &= \gamma'\,E_0\!\left[\Pi\,\mathcal{T}_{k_0,\pi,\gamma'}^{-1}(\mathcal{P}_{k_0\beta} \Pi)\,
      q_{r_0,k_0,\nu}^{\pi,\gamma'}(A,S)\right] \\
&= E_0\!\left[\beta_{0,1}(S',A,S)\,\beta(S', A, S)\right],
\end{aligned}
\end{equation}
where
\[
\beta_{0,1}(s',a,s) :=    d_{0,k_0}^{\pi,\gamma'}(a,s) \{r_{0,\nu,\gamma,g}(a,s) + \gamma' V_{r_0, k_0,\nu}^{\pi,\gamma'}(s') - q_{r_0, k_0,\nu}^{\pi,\gamma'}(a,s)\}.
\]

We now turn to the second term in \eqref{eq:dkm-split}:
\begin{equation}
\label{eq:second-term}
\gamma\,E_0\!\left[\Pi\,\mathcal{T}_{k_0,\pi,\gamma'}^{-1}(I-\nu)\,
      \mathcal{T}_{k_0,\nu,\gamma}^{-1}(\mathcal{P}_{k_0\beta} \nu)\,
      \mathcal{T}_{k_0,\nu,\gamma}^{-1}(r_0-g)(S)\right]
= \gamma\,E_0\!\left[\Pi\,\mathcal{T}_{k_0,\pi,\gamma'}^{-1}(I-\nu)\,
      \mathcal{T}_{k_0,\nu,\gamma}^{-1}(\mathcal{P}_{k_0\beta} V_{r_0-g,k_0}^{\nu,\gamma})(S)\right],
\end{equation}
where we used the identity \(V_{r_0-g,k_0}^{\nu,\gamma} = \nu\,\mathcal{T}_{k_0,\nu,\gamma}^{-1}(r_0-g)\).  
By Lemma~\ref{lemma::adjointident}, the right-hand side of \eqref{eq:second-term} satisfies
\begin{align*}
   \gamma\,E_0\!\left[\Pi\,\mathcal{T}_{k_0,\pi,\gamma'}^{-1}(I-\nu)\,
      \mathcal{T}_{k_0,\nu,\gamma}^{-1}(\mathcal{P}_{k_0\beta} V_{r_0-g,k_0}^{\nu,\gamma})(S)\right] 
    &=  \gamma\,E_0\!\left[d_{0,k_0}^{\pi,\gamma'}(A,S)\,(I - \nu)\,
      \mathcal{T}_{k_0,\nu,\gamma}^{-1}(\mathcal{P}_{k_0\beta} V_{r_0-g,k_0}^{\nu,\gamma})(A,S)\right].
\end{align*}
Moreover, by a change of measure and another application of Lemma~\ref{lemma::adjointident},
\begin{align}
E_0\!\left[d_{0,k_0}^{\pi,\gamma'}(A,S)\,(I - \nu)\,\mathcal{T}_{k_0,\nu,\gamma}^{-1}(\mathcal{P}_{k_0\beta} V_{r_0-g,k_0}^{\nu,\gamma})(A,S)\right]
&= \frac{1}{1-\gamma'}\,E_{d_{0,k_0}^{\pi,\gamma'}}\!\left[(I - \nu)\,\mathcal{T}_{k_0,\nu,\gamma}^{-1}(\mathcal{P}_{k_0\beta} V_{r_0-g,k_0}^{\nu,\gamma})(A,S)\right] \notag \\
&= \frac{1}{1-\gamma'}\,E_{d_{0,k_0}^{\pi,\gamma'}}\!\left[\left\{1 - \tfrac{\nu(A \mid S)}{\pi(A \mid S)}\right\}\,(\mathcal{P}_{k_0\beta} V_{r_0-g,k_0}^{\nu,\gamma})(A,S)\right] \notag \\
&= E_0\!\left[d_{0,k_0}^{\pi,\gamma'}(A,S)\left\{1 - \tfrac{\nu(A \mid S)}{\pi(A \mid S)}\right\}(\mathcal{P}_{k_0\beta} V_{r_0-g,k_0}^{\nu,\gamma})(A,S)\right],
\label{eq:change-measure}
\end{align}
where \(E_{d_{0,k_0}^{\pi,\gamma'}}\) denotes expectation under the normalized discounted state–action occupancy measure, so the \((1-\gamma')\) factors cancel.

Writing out the action of $\mathcal{P}_{k_0\beta}$, we have
\begin{align}
E_0\!\Big[ d_{0,k_0}^{\pi,\gamma'}(A,S)\Bigl\{1 - &\tfrac{\nu(A \mid S)}{\pi(A \mid S)} \Bigr\}
          (\mathcal{P}_{k_0\beta} V_{r_0-g,k_0}^{\nu,\gamma})(A,S) \Big]  \notag \\
&=   E_0\!\Big[ d_{0,k_0}^{\pi,\gamma'}(A,S)\Bigl\{1 - \tfrac{\nu(A \mid S)}{\pi(A \mid S)} \Bigr\}
           V_{r_0-g, k_0}^{\nu,\gamma}(S') \,\beta(S' \mid A, S) \Big] \notag \\
&=   E_0\!\Big[ d_{0,k_0}^{\pi,\gamma'}(A,S)\Bigl\{1 - \tfrac{\nu(A \mid S)}{\pi(A \mid S)} \Bigr\}
           \bigl\{V_{r_0-g, k_0}^{\nu,\gamma}(S') - \mathcal{P}_{k_0}V_{r_0-g, k_0}^{\nu,\gamma}(A,S)\bigr\}
           \beta(S' \mid A, S) \Big].
\label{eq:Pb-expanded}
\end{align} 
Hence, combining \eqref{eq:second-term}–\eqref{eq:Pb-expanded},
\begin{align*}
 \gamma E_0\!\Big[  \Pi \,\mathcal{T}_{k_0,\pi,\gamma'}^{-1}(I - \nu)\,
            \mathcal{T}_{k_0,\nu,\gamma}^{-1}(\mathcal{P}_{k_0\beta} \nu)\,
            \mathcal{T}_{k_0,\nu,\gamma}^{-1}(r_0-g) \Big]
&=  E_0\!\Big[ d_{0,k_0}^{\pi,\gamma'}(A,S)\Bigl\{1-\tfrac{\nu(A \mid S)}{\pi(A \mid S)}   \Bigr\} \\
&\qquad \times \gamma \,\bigl\{V_{r_0-g, k_0}^{\nu,\gamma}(S') - \mathcal{P}_{k_0}V_{r_0-g, k_0}^{\nu,\gamma}(A,S)\bigr\}
                 \beta(S' \mid A, S) \Big].
\end{align*}


As in the proof of Theorem \ref{theorem::EIFvalue}, Bellman’s equation shows that 
\begin{equation*} 
\gamma \{V_{r_0-g, k_0}^{\nu,\gamma}(s') - \mathcal{P}_{k_0}V_{r_0-g, k_0}^{\nu,\gamma}(a,s)\} 
= r_0(a,s) - g(s) + \gamma V_{r_0-g, k_0}^{\nu,\gamma}(s') - q_{r_0-g, k_0}^{\nu,\gamma}(a,s).
\end{equation*}
Thus,
\begin{equation}
  \gamma\, E_0\!\left[ \Pi \,\mathcal{T}_{k_0,\pi,\gamma'}^{-1}(I - \nu)\,
  \mathcal{T}_{k_0,\nu,\gamma}^{-1}(\mathcal{P}_{k_0\beta} \nu)\,
  \mathcal{T}_{k_0,\nu,\gamma}^{-1}(r_0-g) \right] 
=  E_0\!\left[\beta_{0,2}(S' \mid A, S)\,\beta(S' \mid A, S) \right], \label{eqnproof::beta2}
\end{equation}
where
\[
\beta_{0,2}(s' \mid a, s) :=  d_{0,k_0}^{\pi,\gamma'}(a,s)\left\{1-\tfrac{\nu(a \mid s)}{\pi(a \mid s)}\right\}
\{r_0(a,s) - g(s) + \gamma V_{r_0-g, k_0}^{\nu,\gamma}(s') - q_{r_0-g, k_0}^{\nu,\gamma}(a,s)\}.
\]


 


Combining \eqref{eqnproof::beta1} and \eqref{eqnproof::beta2}, we conclude that
\begin{align*}
E_0\!\big[\partial_k m(A,S,r_0,k_0)[\beta]\big]
= E_0\!\left[\beta_0(S' \mid A, S)\,\beta(S' \mid A, S)\right],
\end{align*}
where $\beta_0 := \beta_{0,1} + \beta_{0,2}$ and
\begin{equation}
\label{eqn::proofbetaEIF}
\begin{aligned}
\beta_0(s' \mid a, s) 
&= d_{0,k_0}^{\pi,\gamma'}(a,s)\,\big\{r_{0,\nu,\gamma,g}(a,s) + \gamma' V_{r_0, k_0,\nu}^{\pi,\gamma'}(s') - q_{r_0, k_0,\nu}^{\pi,\gamma'}(a,s)\big\} \\
&\quad + d_{0,k_0}^{\pi,\gamma'}(a,s)\left\{ 1-\tfrac{\nu(a \mid s)}{\pi(a \mid s)} \right\}
          \big\{r_0(a,s) - g(s) + \gamma V_{r_0-g, k_0}^{\nu,\gamma}(s') - q_{r_0-g, k_0}^{\nu,\gamma}(a,s)\big\}.
\end{aligned}
\end{equation}

Next, we turn to deriving the Riesz representer $\alpha_0$ for the reward partial derivative 
$E_0[\partial_r m(A,S,r_0,k_0)[\alpha]]$. 
Denote $\alpha_{k_0,\nu} := (I - \nu)\,\mathcal{T}_{k_0,\nu,\gamma}^{-1}(\alpha)$. 
Applying the identities in Lemma~\ref{lemma::adjointident} to \eqref{eq:drm}, we obtain
\begin{align}
E_0[\partial_r m(A,S,r_0,k_0)[\alpha]] 
&= E_0\!\left[\Pi\,\mathcal{T}_{k_0,\pi,\gamma'}^{-1}\!\bigl(\alpha_{k_0,\nu}\bigr)(S)\right]\notag \\ 
&= E_0\!\left[d_{0,k_0}^{\pi,\gamma'}(A,S)\,\alpha_{k_0,\nu}(A,S)\right]\notag \\
&= E_0\!\left[d_{0,k_0}^{\pi,\gamma'}(A,S)\,(I - \nu)\,\mathcal{T}_{k_0,\nu,\gamma}^{-1}(\alpha)(A,S)\right].
\end{align}
By the same argument used to derive \eqref{eq:change-measure}---namely, a change of measure together with Lemma~\ref{lemma::adjointident}---we further obtain
\begin{align*}
E_0\!\left[d_{0,k_0}^{\pi,\gamma'}(A,S)\,(I - \nu)\,\mathcal{T}_{k_0,\nu,\gamma}^{-1}(\alpha)(A,S)\right] 
&= E_0\!\left[d_{0,k_0}^{\pi,\gamma'}(A,S)\,(I - \nu)\,\alpha(A,S)\right] \\
&= E_0\!\left[d_{0,k_0}^{\pi,\gamma'}(A,S)\,\Bigl(1 - \tfrac{\nu(A \mid S)}{\pi(A \mid S)}\Bigr)\,\alpha(A,S)\right].
\end{align*}
Thus,
\begin{align}
\label{eqnproof::normrewardderiv}
  E_0[\partial_r m(A,S,r_0,k_0)[\alpha]] 
&= E_0\!\left[\alpha_0(A,S)\alpha(A,S)\right],
\end{align}
where we define
\begin{align*}
\alpha_0(A,S) 
&:= d_{0,k_0}^{\pi,\gamma'}(A,S)\,\Bigl(1 - \tfrac{\nu(A \mid S)}{\pi(A \mid S)}\Bigr) \\
&= \rho_0^{\pi,\gamma'}(S)\Bigl(\tfrac{\pi(A \mid S)}{\pi_0(A \mid S)} - \tfrac{\nu(A \mid S)}{\pi_0(A \mid S)}\Bigr).
\end{align*}
We have
\[
\sum_{\tilde a\in\mathcal A}\pi_0(\tilde a\mid S)\,\frac{\pi(\tilde a\mid S)}{\pi_0(\tilde a\mid S)}=1
\quad\text{and}\quad
\sum_{\tilde a\in\mathcal A}\pi_0(\tilde a\mid S)\,\frac{\nu(\tilde a\mid S)}{\pi_0(\tilde a\mid S)}=1.
\]
Hence,
\begin{equation}
\label{eqn::proofalphaEIF}
\begin{aligned}
\alpha_0(A,S) - \sum_{\tilde a \in \mathcal{A}} \pi_0(\tilde a \mid S)\,\alpha_0(\tilde a, S) 
&= \rho_0^{\pi,\gamma'}(S)\left[\frac{\pi(A\mid S)}{\pi_0(A\mid S)}-\frac{\nu(A\mid S)}{\pi_0(A\mid S)}
 - \Big(1-1\Big)\right] \\
&= \rho_0^{\pi,\gamma'}(S)\left(\frac{\pi(A\mid S)}{\pi_0(A\mid S)}-\frac{\nu(A\mid S)}{\pi_0(A\mid S)}\right) \\
&= \alpha_0(A,S).
\end{aligned}
\end{equation}




Thus, by substituting \eqref{eqn::proofalphaEIF} and \eqref{eqn::proofbetaEIF} into the expression for the EIF $\chi_0$ given in Theorem~\ref{theorem::EIF}, we obtain
\begin{equation}
\label{eqnproof:firstEIF}
\begin{aligned}
\chi_0(s,a,s')
&= \rho_0^{\pi,\gamma'}(s)\left(\frac{\pi(a\mid s)}{\pi_0(a\mid s)} - \frac{\nu(a\mid s)}{\pi_0(a\mid s)}\right) \\[4pt]
&\quad + d_{0,k_0}^{\pi,\gamma'}(a,s)\left(1 - \frac{\nu(a \mid s)}{\pi(a \mid s)}\right)
          \bigl\{r_0(a,s) - g(s) + \gamma V_{r_0-g, k_0}^{\nu,\gamma}(s') - q_{r_0-g, k_0}^{\nu,\gamma}(a,s)\bigr\} \\[4pt]
&\quad + d_{0,k_0}^{\pi,\gamma'}(a,s)\,\bigl\{r_{0,\nu,\gamma,g}(a,s) + \gamma' V_{r_0, k_0,\nu}^{\pi,\gamma'}(s') - q_{r_0, k_0,\nu}^{\pi,\gamma'}(a,s)\bigr\} \\[4pt]
&\quad + m(a,s, r_0, k_0) - \Psi(P_0).
\end{aligned}
\end{equation}
 The result then follows from Theorem \ref{theorem::identWithNormalization}, noting that
 $$r_{0,\nu,\gamma,g}(a,s) = g(s) + q_{r_0-g, k_0}^{\nu,\gamma}(a,s) - V_{r_0-g, k_0}^{\nu,\gamma}(s).$$



 
 



 



   
\end{proof}



 

 
\begin{proof}[Proof of Theorem \ref{theorem::EIFsoftmax}]
By Lemma~\ref{lemma::derivSoftBellmanStar}, the map
$(r,k) \mapsto E_0[m(A,S,r,k)]$ is Fréchet differentiable in the supremum norm
and, hence, Gâteaux differentiable over $r \in \mathcal{H}$ and
$k \in \mathcal{K}$. The derivatives are provided in
Lemma~\ref{lemma::derivSoftBellmanStar} and are stated later in this proof.
All Bellman operators and policies below are the starred counterfactual objects
from Example~\ref{example::1b}, and we keep the stars throughout. By
\ref{cond::stationary3} and Lemma~\ref{lemma::invertibility}, applied with
$\pi = \pi_{r_0,k_0}^\star$, the operator
$\mathcal{T}_{r_0,k_0}^\star: L^2(\lambda) \to L^2(\lambda)$ admits a bounded
inverse, both on $L^2(\lambda)$ and on $L^2(\pi_0 \otimes \rho_0)$. As a
consequence, the partial derivatives in the sup-norm space,
\[
h \mapsto E_0[\partial_r m(A,S,r,k)(h)] \quad \text{and} \quad 
\beta \mapsto E_0[\partial_k m(A,S,r,k)(\beta)],
\]
admit continuous extensions on the respective domains
$(T_{\mathcal{H}}, \|\cdot\|_{\mathcal{H}, 0})$ and
$(T_{\mathcal{K}}, \|\cdot\|_{\mathcal{K}, 0})$. Furthermore, by
Lemma~\ref{lemma::linearizationExamplestar}, for all $(r,k)$ in an $L^2$-ball
around $(r_0,k_0)$, the inverse $(\mathcal{T}_{r,k}^\star)^{-1}$ exists, and
the linearization remainder based on these derivatives satisfies
\[
\operatorname{Rem}_0(r_0,k_0; r,k) 
= O\!\left(\|k - k_0\|_{\mathcal{K}}^{2} 
+ \|r - r_0\|_{\mathcal{H}}^{2} 
+ \|v_{r_0,k}-v_{r_0,k_0}\|_{L^2(\lambda)}^{2}\right),
\]
with an implicit constant depending only on 
$\max\{\|r\|_{L^\infty(\lambda)}, \|r_0\|_{L^\infty(\lambda)}\}$ 
and the constants in our assumptions. 

Moreover, by \ref{cond::lipschitzvalue}, we obtain
\[
\operatorname{Rem}_0(r_0,k_0; r,k) 
= O\!\left(\|k - k_0\|_{\mathcal{K}}^{2} 
+ \|r - r_0\|_{\mathcal{H}}^{2}\right).
\]
Thus, the map $(r,k) \mapsto E_0[m(A,S,r,k)]$ is Fréchet differentiable, satisfying
\[
E_0\!\left[m(A,S,r+h,k+g) - m(A,S,r,k)
- \partial_r m(A,S,r,k)(h) - \adk m(A,S,r,k)(g)\right]
= o\!\left(\|h\|_{\mathcal{H}, 0} + \|g\|_{\mathcal{K}}\right)
\]
as $(h,g) \to (0,0)$ in the product norm 
$\|h\|_{\mathcal{H}, 0} + \|g\|_{\mathcal{K}}$. Fréchet differentiability implies \ref{cond::boundedfunc}, and by the triangle inequality, the boundedness of the derivatives also implies \ref{cond::lipschitz}. With these conditions verified, we now turn to determining the explicit form of the EIF.


 
 
 

\noindent \textbf{Identities.} We will use the following identities throughout. Arguing as in the proof of Lemma~\ref{lemma::adjointident}, for any state-only function $g\in L^2(\rho_0)$,
\begin{equation}
\label{eqn::identity1}
\begin{aligned}
E_0\!\left[\big(I-\gamma\,\Pi_{r_0,k_0}^\star \mathcal{P}_{k_0}\big)^{-1} g(S)\right]
&= E_0\!\left[\Big(I + \gamma\,\Pi_{r_0,k_0}^\star (\mathcal{T}_{r_0,k_0}^\star)^{-1}\mathcal{P}_{k_0}\Big) g(S)\right] \\
&= \sum_{t=0}^\infty \gamma^t\,\mathbb{E}_{k_0,\pi_{r_0,k_0}^\star}[g(S_t)]
= E_0\!\big[g(S)\,\rho_{r_0,k_0}^\star(S)\big].
\end{aligned}
\end{equation}
The first equality in \eqref{eqn::identity1} uses the resolvent identity,
the second uses the Neumann-series expansion of the inverse,
and the last applies the Radon--Nikodym derivative to express discounted
trajectory expectations as a reweighted expectation under the initial state distribution. 
We will also use that, for any state–action function \(h\),
\begin{equation}
    \label{eqn::identity2}
\big((\mathcal{T}_{r_0,k_0}^\star)^{-1} h\big)(A,S)
= \int \sum_{a' \in \mathcal{A}^\star} d_{r_0,k_0}^{\star}(a',s'\mid A,S)\,h(a',s')\,
\pi_0(a'\mid s')\,\rho_0(s')\,d\mu(s').
\end{equation}

\noindent \textbf{Determining the $\alpha_0$ component (partial derivative in $r$).}
\[
\partial_r\!\left[\Pi_{r,k}^\star (\mathcal T_{r,k}^\star)^{-1}(r)\right]\Big|_{(r_0,k_0)}[h]
= \left(\partial_r\!\left[\Pi_{r,k}^\star (\mathcal T_{r,k}^\star)^{-1}\right]\Big|_{(r_0,k_0)}[h]\right)\!r_0
\;+\; \Pi_{r_0,k_0}^\star (\mathcal T_{r_0,k_0}^\star)^{-1}(h).
\]
Arguing exactly as in the proof of Theorem \ref{theorem::EIFvalue}, it holds that
\begin{equation}
    E_0[\Pi_{r_0,k_0}^\star (\mathcal T_{r_0,k_0}^\star)^{-1}(h)(S)] = E_0[h(A,S) \alpha_{0,1}(A,S)],\label{eqnproof::alpha01}
\end{equation}
where
$$\alpha_{0,1}(a,s) = d_{r_0,k_0}^{\star}(a,s) = \rho_{r_0,k_0}^{\star}(s) \frac{\pi_{r_0, k_0}^\star(a \mid s)}{\pi_0(a\mid s)}.$$


By Lemma~\ref{lemma::derivSoftBellmanStar}, we have
\begin{align*}
\partial_r \!\left[ \Pi_{r,k}^\star (\mathcal{T}_{r,k}^\star)^{-1}\right][h](r_0)\Big|_{(r_0,k_0)}
&= \Big(I + \gamma\,\Pi_{r_0,k_0}^\star (\mathcal{T}_{r_0,k_0}^\star)^{-1}\mathcal{P}_{k_0}\Big)\,
\big(\partial_r \Pi_{r,k}^\star[h]\big)\Big|_{(r_0,k_0)}\,(\mathcal{T}_{r_0,k_0}^\star)^{-1}(r_0)\\
&= \Big(I + \gamma\,\Pi_{r_0,k_0}^\star (\mathcal{T}_{r_0,k_0}^\star)^{-1}\mathcal{P}_{k_0}\Big)\,
\frac{1}{\tau^\star}\,
\operatorname{Cov}_{\pi_{r_0,k_0}^\star}\!\Big((\mathcal{T}_{r_0,k_0}^\star)^{-1} r_0,\;(\mathcal{T}_{r_0,k_0}^\star)^{-1} h\Big).
\end{align*}
Taking expectations, we have
\begin{align*}
E_0 \!\left[
\partial_r \!\left[ \Pi_{r,k}^\star (\mathcal{T}_{r,k}^\star)^{-1}\right][h](r_0)
\Big|_{(r_0,k_0)}(S)
\right]
&= E_0\!\left[\Big(I + \gamma\,\Pi_{r_0,k_0}^\star (\mathcal{T}_{r_0,k_0}^\star)^{-1}\mathcal{P}_{k_0}\Big)\,
\frac{1}{\tau^\star}\,
\operatorname{Cov}_{\pi_{r_0,k_0}^\star}\!\Big((\mathcal{T}_{r_0,k_0}^\star)^{-1} r_0,\;(\mathcal{T}_{r_0,k_0}^\star)^{-1} h\Big)(S)\right].
\end{align*}
Using the identity in \eqref{eqn::identity1} with the state-only function
\[
g(s)\;=\;\frac{1}{\tau^\star}\,
\operatorname{Cov}_{\pi_{r_0,k_0}^\star}\!\Big((\mathcal{T}_{r_0,k_0}^\star)^{-1} r_0,\;(\mathcal{T}_{r_0,k_0}^\star)^{-1} h\Big)(s),
\]
we obtain
\begin{align*}
E_0 \!\left[\partial_r \!\left[ \Pi_{r,k}^\star (\mathcal{T}_{r,k}^\star)^{-1}\right][h](r_0)\Big|_{(r_0,k_0)}(S) \right]
&= \frac{1}{\tau^\star}\,
E_0\!\Bigg[\Big(I + \gamma\,\Pi_{r_0,k_0}^\star (\mathcal{T}_{r_0,k_0}^\star)^{-1}\mathcal{P}_{k_0}\Big) \\
&\qquad\qquad \times
\operatorname{Cov}_{\pi_{r_0,k_0}^\star}\!\Big((\mathcal{T}_{r_0,k_0}^\star)^{-1} r_0,\;(\mathcal{T}_{r_0,k_0}^\star)^{-1} h\Big)(S)\Bigg] \\
&= \frac{1}{\tau^\star}\,
E_0\!\left[\rho_{r_0,k_0}^\star(S)\,
\operatorname{Cov}_{\pi_{r_0,k_0}^\star}\!\Big((\mathcal{T}_{r_0,k_0}^\star)^{-1} r_0,\;(\mathcal{T}_{r_0,k_0}^\star)^{-1} h\Big)(S)\right].
\end{align*}
Equivalently,
\begin{align*}
& E_0 \!\left[\partial_r \!\left[ \Pi_{r,k}^\star (\mathcal{T}_{r,k}^\star)^{-1}\right][h](r_0)\Big|_{(r_0,k_0)}(S) \right]\\
&= \frac{1}{\tau^\star}\,
E_0\!\left[\rho_{r_0,k_0}^\star(S)\,
\operatorname{Cov}_{\pi_{r_0,k_0}^\star}\!\Big((\mathcal{T}_{r_0,k_0}^\star)^{-1} r_0,\;(\mathcal{T}_{r_0,k_0}^\star)^{-1} h\Big)(S)\right] \\
&= \frac{1}{\tau^\star}\,
E_0\!\Bigg[\rho_{r_0,k_0}^\star(S)\,
\Pi_{r_0,k_0}^\star\!\Big(\big(I-\Pi_{r_0,k_0}^\star\big)(\mathcal{T}_{r_0,k_0}^\star)^{-1} r_0 \\
&\qquad\qquad\qquad\qquad\qquad \times
\big(I-\Pi_{r_0,k_0}^\star\big)(\mathcal{T}_{r_0,k_0}^\star)^{-1} h\Big)(S)\Bigg] \\
&= \frac{1}{\tau^\star}\,
E_0\Big[d_{r_0,k_0}^\star(A,S)
\big((I-\Pi_{r_0,k_0}^\star)(\mathcal{T}_{r_0,k_0}^\star)^{-1} r_0\big)(A,S)\;
\big((I-\Pi_{r_0,k_0}^\star)(\mathcal{T}_{r_0,k_0}^\star)^{-1} h\big)(A,S)
\Big]\\
&= \frac{1}{\tau^\star}\,
E_0\Big[d_{r_0,k_0}^\star(A,S)
\big((I-\Pi_{r_0,k_0}^\star)(\mathcal{T}_{r_0,k_0}^\star)^{-1} r_0\big)(A,S)\;\big((\mathcal{T}_{r_0,k_0}^\star)^{-1} h\big)(A,S)
\Big],
\end{align*}
where the final equality uses that
\begin{align*}
& E_0\Big[d_{r_0,k_0}^\star(A,S)\,
\big((I-\Pi_{r_0,k_0}^\star)(\mathcal{T}_{r_0,k_0}^\star)^{-1} r_0\big)(A,S)\;
\big(\Pi_{r_0,k_0}^\star(\mathcal{T}_{r_0,k_0}^\star)^{-1} h\big)(S)
\Big] \\
&= (1-\gamma)\,E_{d_{r_0,k_0}^\star}\Big[
\big((I-\Pi_{r_0,k_0}^\star)(\mathcal{T}_{r_0,k_0}^\star)^{-1} r_0\big)(A,S)\;
\big(\Pi_{r_0,k_0}^\star(\mathcal{T}_{r_0,k_0}^\star)^{-1} h\big)(S)
\Big] \\
&= 0,
\end{align*}
where \(E_{d_{r_0,k_0}^\star}\) is the expectation under the normalized probability distribution
\(\tfrac{1}{1-\gamma}\,d_{r_0,k_0}^\star\), and the final equality follows from the law of total expectation, since
\[
E_{d_{r_0,k_0}^\star}\!\left[
\big((I-\Pi_{r_0,k_0}^\star)(\mathcal{T}_{r_0,k_0}^\star)^{-1} r_0\big)(A,S)\,\Big|\,S
\right]=0.
\]
This holds because \(d_{r_0,k_0}^\star\) is, up to normalization, a density ratio that
reweights \(E_0\) to the distribution induced by the policy \(\pi_{r_0,k_0}^\star\)
(with initial state distribution proportional to \(\rho_{r_0,k_0}^\star\)).



Using that \(q_{r_0,k_0}^\star - V_{r_0,k_0}^\star = (I-\Pi_{r_0,k_0}^\star)(\mathcal{T}_{r_0,k_0}^\star)^{-1} r_0\), we have
\begin{align*}
E_0 \!\left[\partial_r \!\left[ \Pi_{r,k}^\star (\mathcal{T}_{r,k}^\star)^{-1}\right][h](r_0)\Big|_{(r_0,k_0)}(S) \right]
&= \frac{1}{\tau^\star}\,
E_0\!\Big[d_{r_0,k_0}^\star(A,S)\,
\big(q_{r_0,k_0}^\star - V_{r_0,k_0}^\star\big)(A,S)\,
\big((\mathcal{T}_{r_0,k_0}^\star)^{-1} h\big)(A,S)\Big].
\end{align*}
Applying the kernel form \eqref{eqn::identity2} of \(\big(\mathcal{T}_{r_0,k_0}^\star\big)^{-1}h\), we obtain
\begin{align*}
E_0 \!\left[\partial_r \!\left[ \Pi_{r,k}^\star (\mathcal{T}_{r,k}^\star)^{-1}\right][h](r_0)\Big|_{(r_0,k_0)} \right]
&= \frac{1}{\tau^\star}\,
E_0\!\Bigg[d_{r_0,k_0}^\star(A,S)\,\big(q_{r_0,k_0}^\star - V_{r_0,k_0}^\star\big)(A,S) \\
&\qquad\qquad \times
\int \sum_{a'\in\mathcal{A}^\star} d_{r_0,k_0}^{\star}(a',s'\mid A,S)\,h(a',s')\,
\pi_0(a'\mid s')\,\rho_0(s')\,d\mu(s')\Bigg] \\
&= \frac{1}{\tau^\star}\,
\iint \sum_{a'\in\mathcal{A}^\star}
E_0\!\Big[d_{r_0,k_0}^\star(A,S)\,
\big(q_{r_0,k_0}^\star - V_{r_0,k_0}^\star\big)(A,S)\,
d_{r_0,k_0}^{\star}(a',s'\mid A,S)\Big] \\
&\qquad\qquad \times h(a',s')\,\pi_0(a'\mid s')\,\rho_0(s')\,d\mu(s'),
\end{align*}
where the second line uses Fubini’s theorem. Defining
\[
\widetilde{d}_{r_0,k_0}^\star(a',s')
\;:=\;E_0\!\Big[d_{r_0,k_0}^\star(A,S)\,\big(q_{r_0,k_0}^\star - V_{r_0,k_0}^\star\big)(A,S)\,
d_{r_0,k_0}^{\star}(a',s'\mid A,S)\Big],
\]
we obtain that
\[
E_0 \!\left[\partial_r \!\left[ \Pi_{r,k}^\star (\mathcal{T}_{r,k}^\star)^{-1}\right][h](r_0)\Big|_{(r_0,k_0)} \right]
= \frac{1}{\tau^\star}E_0[\widetilde{d}_{r_0,k_0}^\star(A,S)\,h(A,S)].
\]
Hence,
\begin{equation}
\label{eqnproof::alpha02}
    E_0 \!\left[\partial_r \!\left[ \Pi_{r,k}^\star (\mathcal{T}_{r,k}^\star)^{-1}\right][h](r_0)\Big|_{(r_0,k_0)} \right] = \frac{1}{\tau^\star}E_0[\widetilde{d}_{r_0,k_0}^\star(A,S)\,h(A,S)],
\end{equation}
 where
 \[
 \alpha_{0,2}(a,s)
 := \frac{1}{\tau^\star}\widetilde{d}_{r_0,k_0}^\star(a,s)
 = \frac{1}{\tau^\star}\widetilde{\rho}_{r_0,k_0}^\star(s)
 \frac{\pi_{r_0,k_0}^\star(a\mid s)}{\pi_0(a\mid s)}.
 \]
Combining \eqref{eqnproof::alpha01} and \eqref{eqnproof::alpha02}, we conclude that
\[
E_0\!\left[\partial_r\!\left[\Pi_{r,k}^\star (\mathcal T_{r,k}^\star)^{-1}(r)\right]\Big|_{(r_0,k_0)}[h]\right]
= \langle h, \alpha_0 \rangle_{\mathcal{H}, 0},
\]
where
\begin{equation}
    \label{eqn::alphastar}
    \alpha_0(a,s) := \left\{\frac{1}{\tau^\star}\,\widetilde{\rho}_{r_0,k_0}^\star(s)  + \rho_{r_0,k_0}^{\star}(s) \right\} \frac{\pi_{r_0, k_0}^\star(a \mid s)}{\pi_0(a\mid s)}.
\end{equation}
 



 





\noindent \textbf{Determining component $\beta_0$ for $k$ partial derivative.} By Lemma~\ref{lemma::derivSoftBellmanStar}, applied with additive kernel direction \(k_0\beta\), we have
\begin{align*}
\partial_k \!\left[ \Pi_{r,k}^\star (\mathcal{T}_{r,k}^\star)^{-1}\right][\beta]\Big|_{(r_0,k_0)}
= \big(I-\gamma\,\Pi_{r_0,k_0}^\star \mathcal{P}_{k_0}\big)^{-1}
\Big[\,\big(\adk \Pi_{r_0,k_0}^\star[k_0\beta]\big) + \gamma\,\Pi_{r_0,k_0}^\star \mathcal{P}_{k_0\beta} \Pi_{r_0,k_0}^\star\,\Big]
(\mathcal{T}_{r_0,k_0}^\star)^{-1}.
\end{align*}
Recall \eqref{eqn::identity1}, which says that
\begin{equation*}
    E_0 \!\left[\big(I-\gamma\,\Pi_{r_0,k_0}^\star \mathcal{P}_{k_0}\big)^{-1} g(S)\right]= E_0\!\big[g(S)\,\rho_{r_0,k_0}^\star(S)\big].
\end{equation*}
Thus,
\begin{equation}
\label{eq:dkPiT}
\begin{aligned}
E_0\!\left[\partial_k \!\left[\Pi_{r,k}^\star (\mathcal{T}_{r,k}^\star)^{-1}\right][\beta]\Big|_{(r_0,k_0)}\,r_0(S)\right]
&= E_0 \!\left[\big(I-\gamma\,\Pi_{r_0,k_0}^\star \mathcal{P}_{k_0}\big)^{-1}
\Big(\,\adk \Pi_{r_0,k_0}^\star[k_0\beta] + \gamma\,\Pi_{r_0,k_0}^\star \mathcal{P}_{k_0\beta} \Pi_{r_0,k_0}^\star\Big)
(\mathcal{T}_{r_0,k_0}^\star)^{-1} (r_0)(S)\right] \\
&= E_0 \!\left[\rho_{r_0,k_0}^\star(S)\,
\Big(\,\adk \Pi_{r_0,k_0}^\star[k_0\beta] + \gamma\,\Pi_{r_0,k_0}^\star \mathcal{P}_{k_0\beta} \Pi_{r_0,k_0}^\star\Big)
(\mathcal{T}_{r_0,k_0}^\star)^{-1} (r_0)(S)\right]\\
&= E_0 \!\left[\rho_{r_0,k_0}^\star(S)\,
\big(\adk \Pi_{r_0,k_0}^\star[k_0\beta]\big)
\big((\mathcal{T}_{r_0,k_0}^\star)^{-1} r_0\big)(S)\right] \\
&\quad+ E_0 \!\left[\rho_{r_0,k_0}^\star(S)\,
\big(\gamma\,\Pi_{r_0,k_0}^\star \mathcal{P}_{k_0\beta} \Pi_{r_0,k_0}^\star\big)
\big((\mathcal{T}_{r_0,k_0}^\star)^{-1} r_0\big)(S)\right].
\end{aligned}
\end{equation}

The first term on the right-hand side of \eqref{eq:dkPiT} can be written as
\begin{align*}
& E_0 \!\left[\rho_{r_0,k_0}^\star(S)\,
\big(\adk \Pi_{r_0,k_0}^\star[k_0\beta]\big)
\big((\mathcal{T}_{r_0,k_0}^\star)^{-1} r_0\big)(S)\right] \\
&\qquad= \tfrac{1}{\tau^\star}\,
E_0\!\left[\rho_{r_0,k_0}^\star(S)\,
\operatorname{Cov}_{\pi_{r_0,k_0}^\star}\!\Big((\mathcal{T}_{r_0,k_0}^\star)^{-1} r_0,\,
\gamma\,\adk v_{r_0,k_0}^\star[k_0\beta]\Big)(S)\right]\\
&\qquad= \tfrac{1}{\tau^\star}\,
E_0\!\left[\rho_{r_0,k_0}^\star(S)\,
\Big(\Pi_{r_0,k_0}^\star\!\big[((\mathcal{T}_{r_0,k_0}^\star)^{-1} r_0)\,
(\gamma\,\adk v_{r_0,k_0}^\star[k_0\beta])\big]
- \Pi_{r_0,k_0}^\star\!\big[(\mathcal{T}_{r_0,k_0}^\star)^{-1} r_0\big]\,
\Pi_{r_0,k_0}^\star\!\big[\gamma\,\adk v_{r_0,k_0}^\star[k_0\beta]\big]\Big)(S)\right]\\
&\qquad= \tfrac{\gamma}{\tau^\star}\,
E_0\!\left[d_{r_0,k_0}^{\star}(A,S)
\big((I-\Pi_{r_0,k_0}^\star)\big((\mathcal{T}_{r_0,k_0}^\star)^{-1} r_0\big)\big)(A,S)\;
\big(\adk v_{r_0,k_0}^\star[k_0\beta]\big)(A,S)
\right],
\end{align*}
where \(d_{r_0,k_0}^\star(a,s) = \rho_{r_0,k_0}^\star(s) \tfrac{\pi_{r_0,k_0}^\star(a \mid s)}{\pi_0(a\mid s)}\).

Recall that $\adk v_{r,k}^\star[k_0\beta] = (\mathcal{T}_{r,k}^\star)^{-1}\,\mathcal{P}_{k_0\beta}\,\Phi_{r,k}^\star$. Hence,
\begin{align*}
& E_0 \!\left[\rho_{r_0,k_0}^\star(S)\,
\big(\adk \Pi_{r_0,k_0}^\star[k_0\beta]\big)
\big((\mathcal{T}_{r_0,k_0}^\star)^{-1} r_0\big)(S)\right] \\
&\qquad= \tfrac{\gamma}{\tau^\star}\,
E_0\!\left[d_{r_0,k_0}^{\star}(A,S)\,
\big((I-\Pi_{r_0,k_0}^\star)(\mathcal{T}_{r_0,k_0}^\star)^{-1} r_0\big)(A,S)\;
\big((\mathcal{T}_{r_0,k_0}^\star)^{-1}\mathcal{P}_{k_0\beta}\,\Phi_{r_0,k_0}^\star\big)(A,S)
\right],
\end{align*}
where we recall that
$$\Phi_{r_0,k_0}^\star(s) = \tau^\star \log\!\big(\sum_{\tilde a \in \mathcal{A}^\star} \exp((r_0(\tilde a,s) + \gamma\,v_{r_0,k_0}^\star(\tilde a,s))/\tau^\star)\big).$$


We recall, by \eqref{eqn::identity2}, that for any state–action function \(h\),
\[
\big((\mathcal{T}_{r_0,k_0}^\star)^{-1} h\big)(A,S)
= \int \sum_{a' \in \mathcal{A}^\star} d_{r_0,k_0}^{\star}(a',s'\mid A,S)\,h(a',s')\,
\pi_0(a'\mid s')\,\rho_0(s')\,d\mu(s').
\]
Thus,
\begin{align*}
\big((\mathcal{T}_{r_0,k_0}^\star)^{-1}\mathcal{P}_{k_0\beta} \,\Phi_{r_0,k_0}^\star\big)(A,S)
&= \int \sum_{a' \in \mathcal{A}^\star} d_{r_0,k_0}^{\star}(a',s'\mid A,S)
\left\{\int \Phi_{r_0,k_0}^\star(\tilde s)\,k_0(\tilde s\mid a',s')\,\beta(\tilde s\mid a',s')\,d\mu(\tilde s)\right\}
\notag\\
&\qquad\qquad \times \pi_0(a'\mid s')\,\rho_0(s')\,d\mu(s')\\
&= \iint \sum_{a' \in \mathcal{A}^\star} d_{r_0,k_0}^{\star}(a',s'\mid A,S)\,
 \Phi_{r_0,k_0}^\star(\tilde s)\,k_0(\tilde s\mid a',s')\,\beta(\tilde s\mid a',s')\,
 \pi_0(a'\mid s')\,\rho_0(s')\,d\mu(s')\,d\mu(\tilde s) \\
&= \iint \sum_{a \in \mathcal{A}^\star} d_{r_0,k_0}^{\star}(a,s\mid A,S)\,
 \Phi_{r_0,k_0}^\star(s')\,k_0(s'\mid a,s)\,\beta(s'\mid a,s)\,
 \pi_0(a\mid s)\,\rho_0(s)\,d\mu(s)\,d\mu(s') \\
&= \iint \sum_{a \in \mathcal{A}^\star} d_{r_0,k_0}^{\star}(a,s\mid A,S)\,
 \Phi_{r_0,k_0}^\star(s') \notag\\
&\qquad\qquad \times
 \beta(s'\mid a,s)\,P_0(ds',a,ds),
\end{align*}
where the penultimate equality uses Fubini to swap the order of integration and renames the dummy variables \((a',s',\tilde s)\) as \((a,s,s')\).

It follows that
\[
\begin{aligned}
&\tfrac{\gamma}{\tau^\star}\,
E_0\!\left[d_{r_0,k_0}^{\star}(A,S)\,
\big((I-\Pi_{r_0,k_0}^\star)(\mathcal{T}_{r_0,k_0}^\star)^{-1} r_0\big)(A,S)\;
\big((\mathcal{T}_{r_0,k_0}^\star)^{-1}\mathcal{P}_{k_0\beta}\,\Phi_{r_0,k_0}^\star\big)(A,S)
\right] \\
&\qquad= \tfrac{\gamma}{\tau^\star}
\iint \sum_{a\in\mathcal{A}^\star} \widetilde{d}_{r_0,k_0}^\star(a,s)\,
\Phi_{r_0,k_0}^\star(s') \\
&\qquad\qquad\qquad \times
\beta(s'\mid a,s)\,P_0(da,ds,ds'),
\end{aligned}
\]
where
\[
\begin{aligned}
\widetilde{d}_{r_0,k_0}^\star(a,s)
&:= E_0\!\left[d_{r_0,k_0}^{\star}(A,S)\,
\big((I-\Pi_{r_0,k_0}^\star)(\mathcal{T}_{r_0,k_0}^\star)^{-1} r_0\big)(A,S)\,
d_{r_0,k_0}^{\star}(a,s\mid A,S)\right] \\
&= E_0\!\left[d_{r_0,k_0}^{\star}(A,S)\,
\{q_{r_0,k_0}^\star - V_{r_0,k_0}^\star\}(A,S)\right. \\
&\qquad\qquad \left.
\times d_{r_0,k_0}^{\star}(a,s\mid A,S)\right].
\end{aligned}
\]
Define 
\[
\begin{aligned}
\beta_{0,1}(s' \mid a, s)
:={}& \tfrac{\gamma}{\tau^\star}\,\widetilde{d}_{r_0,k_0}^\star(a,s) \\
&\times \left\{\Phi_{r_0,k_0}^\star(s')
- \int \Phi_{r_0,k_0}^\star(\tilde s)\,k_0(\tilde s \mid a, s)\,\mu(d\tilde s)\right\}.
\end{aligned}
\]
Then,
\[
\begin{aligned}
&\tfrac{\gamma}{\tau^\star}\,
E_0\!\left[d_{r_0,k_0}^{\star}(A,S)\,
\big((I-\Pi_{r_0,k_0}^\star)(\mathcal{T}_{r_0,k_0}^\star)^{-1} r_0\big)(A,S)\;
\big((\mathcal{T}_{r_0,k_0}^\star)^{-1}\mathcal{P}_{k_0\beta}\,\Phi_{r_0,k_0}^\star\big)(A,S)
\right] \\
&\qquad=  \langle \beta_{0,1}, \beta \rangle_{P_0}.
\end{aligned}
\]
Further, by the soft Bellman equation,
\[
V_{r_0,k_0}^{\star,\mathrm{soft}}(s)=\Phi_{r_0,k_0}^\star(s),
\]
and
\[
v_{r_0,k_0}^\star(a,s)
=\int V_{r_0,k_0}^\star(s')\,k_0(s'\mid a,s)\,d\mu(s')
=\int \Phi_{r_0,k_0}^\star(s')\,k_0(s'\mid a,s)\,d\mu(s').
\]
Hence,
\[
\begin{aligned}
\beta_{0,1}(s'\mid a,s)
&=\frac{\gamma}{\tau^\star}\,\widetilde{d}_{r_0,k_0}^\star(a,s)\left\{
\Phi_{r_0,k_0}^\star(s')
-\int \Phi_{r_0,k_0}^\star(\tilde s)\,k_0(\tilde s\mid a,s)\,d\mu(\tilde s)
\right\}\\
&=\frac{\gamma}{\tau^\star}\,\widetilde{d}_{r_0,k_0}^\star(a,s)\,
\Big(V_{r_0,k_0}^{\star,\mathrm{soft}}(s')-v_{r_0,k_0}^\star(a,s)\Big).
\end{aligned}
\]
 




Next, we consider the second term on the right-hand side of \eqref{eq:dkPiT}. Note that
\begin{align*}
E_0 \!\left[\rho_{r_0,k_0}^\star(S)\,
\big(\gamma\,\Pi_{r_0,k_0}^\star \mathcal{P}_{k_0\beta} \Pi_{r_0,k_0}^\star\big)
\big((\mathcal{T}_{r_0,k_0}^\star)^{-1} r_0\big)(S)\right]
&= \gamma\,E_0 \!\left[\rho_{r_0,k_0}^\star(S)\,
\big(\Pi_{r_0,k_0}^\star \mathcal{P}_{k_0\beta}\big)\,
\Big(\Pi_{r_0,k_0}^\star\big((\mathcal{T}_{r_0,k_0}^\star)^{-1} r_0\big)\Big)(S)\right] \\
&= \gamma\,E_0 \!\left[\rho_{r_0,k_0}^\star(S)\,
\big(\Pi_{r_0,k_0}^\star \mathcal{P}_{k_0\beta}\big)\,V_{r_0,k_0}^{\star}(S)\right]\\
&= \gamma\,E_0\!\left[\rho_{r_0,k_0}^\star(S)\,
\sum_{a\in\mathcal{A}^\star} \pi_{r_0,k_0}^\star(a\mid S)\,
\int V_{r_0,k_0}^{\star}(s')\,k_0(s'\mid a,S)\,\beta(s'\mid a,S) \mu(ds')\right]\\
&= \gamma\,E_0\!\left[d_{r_0,k_0}^\star(A,S) V_{r_0,k_0}^{\star}(S')\,\beta(S'\mid A,S)  \right]\\
&= \langle \beta_{0,2},\beta\rangle_{P_0},
\end{align*}
where
\[
\begin{aligned}
\beta_{0,2}(s'\mid a,s)
:={}& \gamma\,d_{r_0,k_0}^\star(a,s) \\
&\times \left[
V_{r_0,k_0}^{\star}(s')
- \int V_{r_0,k_0}^{\star}(\tilde s)\,k_0(\tilde s \mid a, s)\,\mu(d\tilde s)
\right].
\end{aligned}
\]
By Bellman’s equation, we have
\[
\gamma \left[ V_{r_0,k_0}^{\star}(S') - \int V_{r_0,k_0}^{\star}(s')\,k_0(s' \mid A,S)\,d\mu(s') \right]
= r_0(A,S) + \gamma\,V_{r_0,k_0}^{\star}(S') - q_{r_0,k_0}^\star(A,S).
\]
Thus,
\[
\beta_{0,2}(s' \mid a,s)
:= d_{r_0,k_0}^\star(a,s)\,\Big\{\,r_0(a,s) + \gamma\,V_{r_0,k_0}^{\star}(s') - q_{r_0,k_0}^\star(a,s)\,\Big\}.
\]

Combining the identities above, we obtain from \eqref{eq:dkPiT} that
\begin{equation}
\label{eq:first-term-inner-product}
E_0\!\left[\partial_k \!\left[\Pi_{r,k}^\star (\mathcal{T}_{r,k}^\star)^{-1}\right][\beta]\Big|_{(r_0,k_0)}\,r_0(S)\right]
= \Big\langle \beta_{0},\,\beta\Big\rangle_{P_0},
\end{equation}
where
\[
\beta_0(s'\mid a,s)
= \beta_{0,1}(s'\mid a,s)
\;+\;
\beta_{0,2}(s'\mid a,s).
\]


\noindent \textbf{Determining the nonparametric EIF.}
By Theorem~\ref{theorem::EIF}, the nonparametric EIF is given by
\begin{align*}
   \chi_0(s,a,s') = \alpha_0(a,s) - \sum_{a'\in\mathcal{A}} \pi_{0}(a'\mid s)\,\alpha_0(a',s)  + \beta_0(s'\mid a,s) + m(a,s,r_0,k_0) - \Psi(P_0).
\end{align*}
Moreover, using our expression for $\alpha_0$ in \eqref{eqn::alphastar}, we can write
\begin{align*}
 \alpha_0(a,s) - \sum_{a'\in\mathcal{A}} \pi_0(a'\mid s)\,\alpha_0(a',s) 
=  \left\{\frac{1}{\tau^\star}\,\widetilde{\rho}_{r_0,k_0}^\star(s)  + \rho_{r_0,k_0}^{\star}(s) \right\} \left\{\frac{\pi_{r_0,k_0}^\star(a\mid s)}{\pi_0(a\mid s)} - 1\right\}.
\end{align*}
where
\[
\widetilde{\rho}_{r_0,k_0}^\star(s)
:= E_0\!\big[d_{r_0,k_0}^\star(A,S)\,\big(q_{r_0,k_0}^\star - V_{r_0,k_0}^{\star}\big)(A,S)\,\rho_{r_0,k_0}^\star(s\mid A,S)\big].
\]
Thus, substituting our expression for $\beta_0$, we conclude that
\begin{align*}
   \chi_0(s,a,s') &= \left\{\frac{1}{\tau^\star}\,\widetilde{\rho}_{r_0,k_0}^\star(s)  + \rho_{r_0,k_0}^{\star}(s) \right\} \left\{\frac{\pi_{r_0,k_0}^\star(a\mid s)}{\pi_0(a\mid s)} - 1\right\}  \\
   &\quad + \frac{\gamma}{\tau^\star}\,\widetilde{d}_{r_0,k_0}^\star(a,s)\Big(V_{r_0,k_0}^{\star,\mathrm{soft}}(s')-v_{r_0,k_0}^\star(a,s)\Big) \\
   &\quad + d_{r_0,k_0}^\star(a,s)\,\Big\{\,r_0(a,s) + \gamma\,V_{r_0,k_0}^{\star}(s') - q_{r_0,k_0}^\star(a,s)\,\Big\}\\
   &\quad + m(a,s, r_0, k_0) - \Psi(P_0).
\end{align*}


 
 

\end{proof}




\begin{proof}[Proof of Theorem \ref{theorem::EIFsoftmaxnorm}]
Write \(\bar r_0 := r_{0,\nu,\gamma,g}\), and let
\[
\tilde m(a,s,r,k) := V_{r,k}^{\star,\gamma'}(s),
\qquad
m(a,s,r,k) := \tilde m(a,s,r_{r,k,\nu,\gamma,g},k).
\]
Under the stated assumptions, Theorem~\ref{theorem::EIFsoftmax} applies with
\(r_0\) replaced by \(\bar r_0\) and \(\gamma\) by \(\gamma'\). Hence the
Riesz representers of the partial derivatives of
\((r,k) \mapsto E_0[\tilde m(A,S,r,k)]\) at \((\bar r_0,k_0)\) are
\[
\tilde\alpha_0(a,s)
= \left\{\frac{1}{\tau^\star}\,\widetilde\rho_{\bar r_0,k_0}^{\star,\gamma'}(s)
+ \rho_{\bar r_0,k_0}^{\star,\gamma'}(s)\right\}
\frac{\pi_{\bar r_0,k_0}^{\star,\gamma'}(a\mid s)}{\pi_0(a\mid s)}
\]
and
\[
\begin{aligned}
\tilde\beta_0(s'\mid a,s)
&= \frac{\gamma'}{\tau^\star}\,
\widetilde d_{\bar r_0,k_0}^{\star,\gamma'}(a,s)\,
\bigl\{\Phi_{\bar r_0,k_0}^{\star,\gamma'}(s')
- v_{\bar r_0,k_0}^{\star,\gamma'}(a,s)\bigr\} \\
&\quad + d_{\bar r_0,k_0}^{\star,\gamma'}(a,s)\,
\bigl\{\bar r_0(a,s) + \gamma' V_{\bar r_0,k_0}^{\star,\gamma'}(s')
- q_{\bar r_0,k_0}^{\star,\gamma'}(a,s)\bigr\}.
\end{aligned}
\]

Applying Theorem~\ref{theorem::EIFnorm} to the map \(\tilde m\) gives
\[
\alpha_0(a,s)
= \tilde\alpha_0(a,s)
- \frac{\nu(a\mid s)}{\pi_0(a\mid s)}
\sum_{\tilde a \in \mathcal{A}} \pi_0(\tilde a\mid s)\,
\tilde\alpha_0(\tilde a,s),
\]
and
\[
\beta_0(s'\mid a,s)
= \tilde\beta_0(s'\mid a,s)
+ \alpha_0(a,s)\,
\bigl\{r_0(a,s) - g(s) + \gamma V_{r_0-g,k_0}^{\nu,\gamma}(s')
- q_{r_0-g,k_0}^{\nu,\gamma}(a,s)\bigr\}.
\]
Since
\[
\sum_{\tilde a \in \mathcal{A}} \pi_0(\tilde a\mid s)\,\tilde\alpha_0(\tilde a,s)
= \left\{\frac{1}{\tau^\star}\,\widetilde\rho_{\bar r_0,k_0}^{\star,\gamma'}(s)
+ \rho_{\bar r_0,k_0}^{\star,\gamma'}(s)\right\}
\sum_{\tilde a \in \mathcal{A}} \pi_{\bar r_0,k_0}^{\star,\gamma'}(\tilde a\mid s),
\]
we obtain
\[
\alpha_0(a,s)
= \left\{\frac{1}{\tau^\star}\,\widetilde\rho_{\bar r_0,k_0}^{\star,\gamma'}(s)
+ \rho_{\bar r_0,k_0}^{\star,\gamma'}(s)\right\}
\left\{\frac{\pi_{\bar r_0,k_0}^{\star,\gamma'}(a\mid s)}{\pi_0(a\mid s)}
- \frac{\nu(a\mid s)}{\pi_0(a\mid s)}\right\}.
\]
It follows immediately that
\[
\sum_{\tilde a \in \mathcal{A}} \pi_0(\tilde a\mid s)\,\alpha_0(\tilde a,s) = 0,
\]
because both \(\pi_{\bar r_0,k_0}^{\star,\gamma'}(\cdot\mid s)\) and
\(\nu(\cdot\mid s)\) sum to one.

Finally, Theorem~\ref{theorem::EIF} yields
\[
\chi_0(s,a,s')
= \alpha_0(a,s) + \beta_0(s'\mid a,s)
+ V_{\bar r_0,k_0}^{\star,\gamma'}(s) - \Psi(P_0),
\]
and substituting the expressions above gives the stated EIF.
\end{proof}



\subsection{Functional von Mises expansions for examples}
 

 

 
 

In the following lemmas, we define
\[
V_q^\pi(s) := \sum_{a \in \mathcal{A}} \pi(a \mid s)\, q(a,s), \qquad d_{\rho,r}^{\pi,\gamma}(a,s) := \rho(s)\,\frac{\pi(a \mid s)}{\pi_r(a \mid s) },
\qquad 
\pi_r(a \mid s) := \exp\{r(a,s)\}.
\]

\begin{lemma}[von Mises expansion for policy value]
\label{lemma::vonmisespolicyvalue}
For any $r, q \in L^\infty(\lambda)$ and $\rho \in L^\infty(\mu)$,
\begin{align*}
&\; E_0\!\Big[\,V_q^\pi(S) 
   + d_{\rho,r}^{\pi,\gamma}(A,S)\{r(A,S) + \gamma V_q^\pi(S') - q(A,S)\}  
   + \rho(S)\Bigl\{\tfrac{\pi}{\pi_r}(A \mid S) - 1\Bigr\}\,\Big] - E_0[ V_{r_0, k_0}^\pi(S) ] \\[6pt]
&= E_0\!\left[\{d_0^{\pi,\gamma}(A,S) - d_{\rho,r}^{\pi,\gamma}(A,S)\}
   \{q(A,S) - q_{r_0,k_0}^{\pi,\gamma}(A,S) - \gamma\,(V_q^\pi(S') - V_{r_0,k_0}^{\pi,\gamma}(S'))\}\right] \\
&\quad + O\!\left(\|r - r_0\|^2_{\mathcal{H}, 0}\right),
\end{align*}
where the big-oh notation depends on $\|r_0\|_{L^\infty(\lambda)}$, $\|r\|_{L^\infty(\lambda)}$, and $\|\rho\|_{L^\infty(\mu)}$. 
\end{lemma}
\begin{proof}
    In what follows, we drop explicit dependence on $\pi$ and $\gamma$ in our notation, writing
    $q_{0} := q_{r_0,k_0}^{\pi,\gamma}$ and $d_0 := d_0^{\pi,\gamma}$. 
 


Consider the quantity
\begin{equation}
\label{eqn::prooffullterm}
E_0[V_q^\pi(S)] - E_0[V_{q_0}^\pi(S)] 
+ E_0\!\left[d_{\rho,r}(A,S)\{r(A,S) + \gamma V_q^\pi(S') - q(A,S)\}\right]  
+ E_0\!\left[\rho(S)\Bigl\{\tfrac{\pi}{\pi_r}(A \mid S) - 1\Bigr\}\right].
\end{equation}

Adding and subtracting $r_0(A,S)$ inside the expectation, we obtain
\begin{align*} 
E_0\!\left[d_{\rho,r}(A,S)\{r(A,S) + \gamma V_q^\pi(S') - q(A,S)\}\right] 
&= E_0\!\left[d_{\rho,r}(A,S)\{r_0(A,S) + \gamma V_q^\pi(S') - q(A,S)\}\right] \\
&\quad + E_0\!\left[d_{\rho,r}(A,S)\{r(A,S) - r_0(A,S)\}\right].
\end{align*}
Hence, the quantity in \eqref{eqn::prooffullterm} can be written
\begin{equation}
\label{eqn::prooffullterm2}
\begin{aligned}
E_0[V_q^\pi(S)] - E_0[V_{q_0}^\pi(S)] 
+ E_0\!\left[d_{\rho,r}(A,S)\{r_0(A,S) + \gamma V_q^\pi(S') - q(A,S)\}\right] 
\\ + E_0\!\left[d_{\rho,r}(A,S)\{r(A,S) - r_0(A,S)\}\right] 
+ E_0\!\left[\rho(S)\Bigl\{\tfrac{\pi}{\pi_r}(A \mid S) - 1\Bigr\}\right].
\end{aligned}
\end{equation}
By Bellman’s equation,
\[
r_0(A,S) \;=\; q_0(A,S) \;-\; \gamma\,E_0\!\left[V_{q_0}^\pi(S') \mid A,S\right].
\]
Thus, by the law of iterated expectations, 
\begin{align*}
E_0\!\left[d_{\rho,r}^{\pi,\gamma}(A,S)\{r_0(A,S) + \gamma V_q^\pi(S') - q(A,S)\}\right] 
&= E_0\!\Big[d_{\rho,r}^{\pi,\gamma}(A,S)\{q_0(A,S) - \gamma V_{q_0}^\pi(S') + \gamma V_q^\pi(S') - q(A,S)\}\Big] \\
&= E_0\!\Big[d_{\rho,r}^{\pi,\gamma}(A,S)\{q_0(A,S) - q(A,S) + \gamma (V_q^\pi(S') - V_{q_0}^\pi(S'))\}\Big].
\end{align*}
Moreover, for any $q$,
\[
E_0[V_q^\pi(S)] \;=\; E_0\!\left[d_0(A,S)\{q(A,S) - \gamma V_q^\pi(S')\}\right].
\]
Therefore,
\begin{align*}
E_0[V_q^\pi(S)] - E_0[V_{q_0}^\pi(S)]
&= E_0\!\left[d_0(A,S)\{q(A,S) - q_0(A,S) - \gamma (V_q^\pi(S') - V_{q_0}^\pi(S'))\}\right].
\end{align*}
Combining these results, expression~\eqref{eqn::prooffullterm2} can be written as
\begin{equation}
\label{eqn::prooffullterm3}
\begin{aligned}
& E_0\!\left[d_0(A,S)\{q(A,S) - q_0(A,S) - \gamma (V_q^\pi(S') - V_{q_0}^\pi(S'))\}\right] \\
&\quad + E_0\!\left[d_{\rho,r}(A,S)\{q_0(A,S) - q(A,S) + \gamma (V_q^\pi(S') - V_{q_0}^\pi(S'))\}\right] \\
&\quad + E_0\!\left[d_{\rho,r}(A,S)\{r(A,S) - r_0(A,S)\}\right] 
   + E_0\!\left[\rho(S)\Bigl\{\tfrac{\pi}{\pi_r}(A \mid S) - 1\Bigr\}\right] \\
&= E_0\!\left[(d_0(A,S) - d_{\rho,r}(A,S))\{q(A,S) - q_0(A,S) - \gamma (V_q^\pi(S') - V_{q_0}^\pi(S'))\}\right] \\
&\quad + E_0\!\left[d_{\rho,r}(A,S)\{r(A,S) - r_0(A,S)\}\right] 
   + E_0\!\left[\rho(S)\Bigl\{\tfrac{\pi}{\pi_r}(A \mid S) - 1\Bigr\}\right].
\end{aligned}
\end{equation}

We turn to the term $E_0\!\left[d_{\rho,r}^{\pi,\gamma}(A,S)\{r(A,S) - r_0(A,S)\}\right]$. Since
\(
\max\{\|r\|_{L^\infty(\lambda)}, \|r_0\|_{L^\infty(\lambda)}\} < M,
\) 
it follows that
\[
\left|\exp\{r(A,S) - r_0(A,S)\} - 1 - \{r(A,S) - r_0(A,S)\}\right| 
   \;\lesssim\; \{r(A,S) - r_0(A,S)\}^2,
\]
with the constant depending only on $M$. Hence, recalling $\pi_r(a \mid s) = \exp\{r(a,s)\}$, we have
\begin{align*}
E_0\!\left[d_{\rho,r}^{\pi,\gamma}(A,S)\{r(A,S) - r_0(A,S)\}\right]
&= E_0\!\left[d_{\rho,r}^{\pi,\gamma}(A,S)\{\exp\{r(A,S) - r_0(A,S)\} - 1\}\right] \\
&\quad + O\!\left(E_0\!\left[d_{\rho,r}^{\pi,\gamma}(A,S)\{r(A,S) - r_0(A,S)\}^2\right]\right),
\end{align*}
where the first term on the right-hand side satisfies
\begin{align*}
E_0\!\left[d_{\rho,r}^{\pi,\gamma}(A,S)\{\exp\{r(A,S) - r_0(A,S)\} - 1\}\right] 
&= E_0\!\left[\rho(S)\Bigl\{\tfrac{\pi}{\pi_0}(A \mid S) - \tfrac{\pi}{\pi_r}(A \mid S)\Bigr\}\right] \\
&= E_0\!\left[\rho(S)\Bigl\{1 - \tfrac{\pi}{\pi_r}(A \mid S)\Bigr\}\right],
\end{align*}
where we used the law of iterated expectations and the fact that $E_0\!\left[\tfrac{\pi}{\pi_0}(A \mid S)\mid S\right] = 1$.
Hence,
\begin{align*}
    E_0\!\left[d_{\rho,r}^{\pi,\gamma}(A,S)\{r(A,S) - r_0(A,S)\}\right]
&= E_0\!\left[\rho(S)\Bigl\{1 - \tfrac{\pi}{\pi_r}(A \mid S)\Bigr\}\right]\\
&\quad + O\!\left(E_0\!\left[d_{\rho,r}^{\pi,\gamma}(A,S)\{r(A,S) - r_0(A,S)\}^2\right]\right).
\end{align*}
Equivalently,
\begin{equation}
\label{eqn::proofsecondterm}
\begin{aligned}
   E_0\!\left[d_{\rho,r}(A,S)\{r(A,S) - r_0(A,S)\}\right] 
   + E_0\!\left[\rho(S)\Bigl\{\tfrac{\pi}{\pi_r}(A \mid S) - 1\Bigr\}\right]\\
=O\!\left(E_0\!\left[d_{\rho,r}^{\pi,\gamma}(A,S)\{r(A,S) - r_0(A,S)\}^2\right]\right)
\end{aligned}
\end{equation}


Plugging \eqref{eqn::proofsecondterm} into \eqref{eqn::prooffullterm3}, we obtain
\begin{align*}
&\; E_0[V_q^\pi(S)] - E_0[V_{q_0}^\pi(S)] 
   + E_0\!\left[d_{\rho,r}^{\pi,\gamma}(A,S)\{r(A,S) + \gamma V_q^\pi(S') - q(A,S)\}\right]  
   + E_0\!\left[\rho(S)\Bigl\{\tfrac{\pi}{\pi_r}(A \mid S) - 1\Bigr\}\right] \\[6pt]
&= E_0\!\left[d_0(A,S)\{q(A,S) - q_0(A,S) - \gamma (V_q^\pi(S') - V_{q_0}^\pi(S'))\}\right] \\
&\quad + E_0\!\left[d_{\rho,r}^{\pi,\gamma}(A,S)\{q_0(A,S) - q(A,S) + \gamma (V_q^\pi(S') - V_{q_0}^\pi(S'))\}\right] \\
&\quad + O\!\left(E_0\!\left[d_{\rho,r}^{\pi,\gamma}(A,S)\{r(A,S) - r_0(A,S)\}^2\right]\right).
\end{align*}
We conclude that
\begin{align*}
&\; E_0[V_q^\pi(S)] - E_0[V_{q_0}^\pi(S)] 
   + E_0\!\left[d_{\rho,r}^{\pi,\gamma}(A,S)\{r(A,S) + \gamma V_q^\pi(S') - q(A,S)\}\right]  
   + E_0\!\left[\rho(S)\Bigl\{\tfrac{\pi}{\pi_r}(A \mid S) - 1\Bigr\}\right] \\[6pt]
&= E_0\!\left[\{d_0(A,S) - d_{\rho,r}^{\pi,\gamma}(A,S)\}\{q(A,S) - q_0(A,S) - \gamma (V_q^\pi(S') - V_{q_0}^\pi(S'))\}\right] \\
&\quad + O\!\left(E_0\!\left[d_{\rho,r}^{\pi,\gamma}(A,S)\{r(A,S) - r_0(A,S)\}^2\right]\right).
\end{align*}
Finally, since
\[
d_{\rho,r}^{\pi,\gamma}(a,s)
= \rho(s)\,\frac{\pi(a \mid s)}{\pi_r(a \mid s)}
= \rho(s)\,\frac{\pi(a \mid s)}{\pi_0(a \mid s)}\,\exp\{r_0(a,s)-r(a,s)\},
\]
the weight $d_{\rho,r}^{\pi,\gamma}$ is uniformly bounded by a constant depending only on
$\|\rho\|_{L^\infty(\mu)}$, $\|r\|_{L^\infty(\lambda)}$, and $\|r_0\|_{L^\infty(\lambda)}$.
Therefore,
\[
E_0\!\left[d_{\rho,r}^{\pi,\gamma}(A,S)\{r(A,S) - r_0(A,S)\}^2\right]
\lesssim \|r-r_0\|_{\mathcal{H},0}^2,
\]
which yields the stated remainder bound.

\end{proof}





\begin{lemma}[von Mises expansion for policy value with normalization]
\label{lemma::vonmisespolicyvaluenormalization}
Let $\rho \in L^\infty(\mu)$, $r, q^{\nu,\gamma}, q_{\nu}^{\pi,\gamma'} \in L^\infty(\lambda)$. Denote $r_{\nu,g} := g + (I - \nu)q^{\nu,\gamma}$ and $r_{0,\nu,g} := g + (I - \nu)q_0^{\nu,\gamma}$. Then
    \begin{align*} 
& E_0\!\left[V_{q_{\nu}^{\pi,\gamma'}}^\pi(S) - V_{q_{0,\nu}^{\pi,\gamma'}}^\pi(S)\right] 
+ E_0\!\left[d_{\rho,r}^{\pi,\gamma'}(A,S)\{r_{\nu,g}(A,S) + \gamma' V_{q_{\nu}^{\pi,\gamma'}}^\pi(S') - q_{\nu}^{\pi,\gamma'}(A,S)\}\right] \\
&\quad + E_0\!\left[d_{\rho,r}^{\pi,\gamma'}(A,S)\Bigl(1 - \tfrac{\nu}{\pi}(A \mid S)\Bigr)\{r(A,S)-g(S) + \gamma V_{q^{\nu,\gamma}}^{\nu}(S') - q^{\nu,\gamma}(A,S)\}\right] \\
&\quad + E_0\!\left[\rho(S)\Bigl(\tfrac{\pi}{\pi_r}(A\mid S) - \tfrac{\nu}{\pi_r}(A\mid S)\Bigr)\right] \\
&= E_0\!\Big[(d_{\rho,r}^{\pi,\gamma'}(A,S) - d_0^{\pi,\gamma'}(A,S))\{\mathcal{T}_{k_0, \pi, \gamma'}(q_{0,\nu}^{\pi,\gamma'} - q_{\nu}^{\pi,\gamma'})(A,S)\}\Big] \\
&\quad + O\!\left(\|r - r_0\|_{\mathcal{H},r_0}\,\|\mathcal{T}_{k_0,\nu,\gamma}(q^{\nu,\gamma} - q_0^{\nu,\gamma})\|_{\mathcal{H},r_0}\right) \\
&\quad + O\!\left(\|r - r_0\|_{\mathcal{H}, 0}^2\right).
\end{align*}
where the big-$O$ constant depends on $\|\rho\|_{L^\infty(\mu)}$, $\|r\|_{L^\infty(\lambda)}$, $\|r_0\|_{L^\infty(\lambda)}$, and $\|\tfrac{\nu}{\pi}\|_{L^\infty(\lambda)}$. 
\end{lemma}


\begin{proof}
Write
\[
d_0^{\pi,\gamma'}(a,s)
:= \rho_0^{\pi,\gamma'}(s)\,\frac{\pi(a \mid s)}{\pi_0(a \mid s)}.
\]

Applying the same decomposition as in the proof of
Lemma~\ref{lemma::vonmisespolicyvalue}, but with the reward term there replaced
by the surrogate normalized reward \(r_{\nu,g}\), yields
\begin{align} 
& E_0\!\left[V_{q_{\nu}^{\pi,\gamma'}}^\pi(S) - V_{q_{0,\nu}^{\pi,\gamma'}}^\pi(S)\right] 
+ E_0\!\left[d_{\rho,r}^{\pi,\gamma'}(A,S)\{r_{\nu,g}(A,S) + \gamma' V_{q_{\nu}^{\pi,\gamma'}}^\pi(S') - q_{\nu}^{\pi,\gamma'}(A,S)\}\right] \notag \\
&= E_0\!\Big[(d_{\rho,r}^{\pi,\gamma'}(A,S) - d_0^{\pi,\gamma'}(A,S))\{\mathcal{T}_{k_0, \pi, \gamma'}(q_{0,\nu}^{\pi,\gamma'} - q_{\nu}^{\pi,\gamma'})(A,S)\}\Big]  \notag \\
&\quad + E_0\!\left[d_{\rho,r}^{\pi,\gamma'}(A,S)\{r_{\nu,g}(A,S) - r_{0,\nu,g}(A,S)\}\right]. \label{eqn::nuproofstart}
\end{align}
 
 
Using that
\(\mathcal{T}_{k_0,\nu,\gamma}^{-1}\mathcal{T}_{k_0,\nu,\gamma} = I\), we have
\begin{align}
E_0\!\left[d_{\rho,r}^{\pi,\gamma'}(A,S)\{r_{\nu,g}(A,S) - r_{0,\nu,g}(A,S)\}\right] 
&= E_0\!\left[d_{\rho,r}^{\pi,\gamma'}(A,S)(I - \nu)(q^{\nu,\gamma} - q_0^{\nu,\gamma})(A,S)\right] \notag \\
&= E_0\!\left[d_{\rho,r}^{\pi,\gamma'}(A,S)(I - \nu)\mathcal{T}_{k_0,\nu,\gamma}^{-1}\mathcal{T}_{k_0,\nu,\gamma}(q^{\nu,\gamma} - q_0^{\nu,\gamma})(A,S)\right] \notag \\
&= E_0\!\left[d_{\rho,r}^{\pi,\gamma'}(A,S)(I - \nu)\mathcal{T}_{k_0,\nu,\gamma}(q^{\nu,\gamma} - q_0^{\nu,\gamma})(A,S)\right]. \label{eq:step1proof}
\end{align}
In the last equality, we argued as in the proof of Lemma~\ref{lemma::linearfuncidentnorm} that, for any $w$ and $f$,
\begin{equation*}
E_0\!\left[w(A,S)(I - \nu)\mathcal{T}_{k_0,\nu,\gamma}^{-1}f(A,S)\right] 
= E_0\!\left[w(A,S)(I - \nu)f(A,S)\right]. 
\end{equation*}
Continuing from \eqref{eq:step1proof} and adding and subtracting, we obtain
\begin{align}
& E_0\!\left[d_{\rho,r}^{\pi,\gamma'}(A,S)\{r_{\nu,g}(A,S) - r_{0,\nu,g}(A,S)\}\right] \notag \\
&= E_0\!\left[d_{\rho,r}^{\pi,\gamma'}(A,S)(1 - \tfrac{\nu}{\pi}(A \mid S))\mathcal{T}_{k_0,\nu,\gamma}(q^{\nu,\gamma} - q_0^{\nu,\gamma})(A,S)\right] \notag \\
&\quad + E_0\!\left[d_{\rho,r}^{\pi,\gamma'}(A,S)\left\{\tfrac{\nu(A \mid S)}{\pi(A \mid S)}\mathcal{T}_{k_0,\nu,\gamma}(q^{\nu,\gamma} - q_0^{\nu,\gamma})(A,S) - \bigl(\nu\mathcal{T}_{k_0,\nu,\gamma}(q^{\nu,\gamma} - q_0^{\nu,\gamma})\bigr)(S)\right\}\right]. \label{eq:step3proof}
\end{align}
Note that
\begin{equation*}
E_0\!\left[d_{\rho,r_0}^{\pi,\gamma'}(A,S)\left\{\tfrac{\nu(A \mid S)}{\pi(A \mid S)}\mathcal{T}_{k_0,\nu,\gamma}(q^{\nu,\gamma} - q_0^{\nu,\gamma})(A,S) - \bigl(\nu\mathcal{T}_{k_0,\nu,\gamma}(q^{\nu,\gamma} - q_0^{\nu,\gamma})\bigr)(S)\right\}\right] = 0,
\end{equation*}
by a change of measure. Hence,
\begin{align}
& E_0\!\left[d_{\rho,r}^{\pi,\gamma'}(A,S)(I - \nu)\mathcal{T}_{k_0,\nu,\gamma}(q^{\nu,\gamma} - q_0^{\nu,\gamma})(A,S)\right] \notag \\
&= E_0\!\left[d_{\rho,r}^{\pi,\gamma'}(A,S)(1 - \tfrac{\nu}{\pi}(A \mid S))\mathcal{T}_{k_0,\nu,\gamma}(q^{\nu,\gamma} - q_0^{\nu,\gamma})(A,S)\right] \notag \\
&\quad + E_0\!\left[(d_{\rho,r}^{\pi,\gamma'} - d_{\rho,r_0}^{\pi,\gamma'})(A,S)\left\{\tfrac{\nu(A \mid S)}{\pi(A \mid S)}\mathcal{T}_{k_0,\nu,\gamma}(q^{\nu,\gamma} - q_0^{\nu,\gamma})(A,S) - \bigl(\nu\mathcal{T}_{k_0,\nu,\gamma}(q^{\nu,\gamma} - q_0^{\nu,\gamma})\bigr)(S)\right\}\right]. \label{eq:step4proof}
\end{align}
Since $r_0 - g = \mathcal{T}_{k_0,\nu,\gamma}(q_0^{\nu,\gamma})$, by applying iterated expectations,
\begin{align}
E_0\!\left[d_{\rho,r}^{\pi,\gamma'}(A,S)(1 - \tfrac{\nu}{\pi}(A \mid S))\mathcal{T}_{k_0,\nu,\gamma}(q^{\nu,\gamma} - q_0^{\nu,\gamma})(A,S)\right] \notag
\\
= E_0\!\left[d_{\rho,r}^{\pi,\gamma'}(A,S)(1 - \tfrac{\nu}{\pi}(A \mid S))\{q^{\nu,\gamma}(A,S) - \gamma V_{q^{\nu,\gamma}}^{\nu}(S') - r_0(A,S) + g(S)\}\right].
\label{eq:step5proof}
\end{align}
Furthermore, by Cauchy--Schwarz,
\begin{align}
& E_0\!\left[(d_{\rho,r}^{\pi,\gamma'} - d_{\rho,r_0}^{\pi,\gamma'})(A,S)\left\{\tfrac{\nu(A \mid S)}{\pi(A \mid S)}\mathcal{T}_{k_0,\nu,\gamma}(q^{\nu,\gamma} - q_0^{\nu,\gamma})(A,S) - \bigl(\nu\mathcal{T}_{k_0,\nu,\gamma}(q^{\nu,\gamma} - q_0^{\nu,\gamma})\bigr)(S)\right\}\right] \notag \\
&= O\!\left( \|d_{\rho,r}^{\pi,\gamma'} - d_{\rho,r_0}^{\pi,\gamma'}\|_{\mathcal{H},r_0} 
              \,\Bigl\| \tfrac{\nu(A \mid S)}{\pi(A \mid S)}\mathcal{T}_{k_0,\nu,\gamma}(q^{\nu,\gamma} - q_0^{\nu,\gamma})(A,S) - \bigl(\nu\mathcal{T}_{k_0,\nu,\gamma}(q^{\nu,\gamma} - q_0^{\nu,\gamma})\bigr)(S) \Bigr\|_{\mathcal{H},r_0} \right) \notag \\
&= O\!\left( \|r - r_0\|_{\mathcal{H},r_0} 
              \,\|\mathcal{T}_{k_0,\nu,\gamma}(q^{\nu,\gamma} - q_0^{\nu,\gamma})\|_{\mathcal{H},r_0} \right),
\label{eq:step6proof}
\end{align}
where the big-$O$ constant depends on $\|\rho\|_{L^\infty(\mu)}$, $\|r\|_{L^\infty(\lambda)}$, $\|r_0\|_{L^\infty(\lambda)}$, and $\|\tfrac{\nu}{\pi}\|_{L^\infty(\lambda)}$. We also used that $\|\nu f\|_{\mathcal{H},r_0} \lesssim \|f\|_{\mathcal{H},r_0}$ by positivity of $\nu$ and $\pi_0$.
Combining \eqref{eq:step3proof}-\eqref{eq:step6proof}, we obtain
\begin{align}
& E_0\!\left[d_{\rho,r}^{\pi,\gamma'}(A,S)\{r_{\nu,g}(A,S) - r_{0,\nu,g}(A,S)\}\right] \notag \\
&= E_0\!\left[d_{\rho,r}^{\pi,\gamma'}(A,S)(1 - \tfrac{\nu}{\pi}(A \mid S))\{q^{\nu,\gamma}(A,S) - \gamma V_{q^{\nu,\gamma}}^{\nu}(S') - r_0(A,S) + g(S)\}\right] \notag \\
&\quad + O\!\left( \|r - r_0\|_{\mathcal{H},r_0}\,\|\mathcal{T}_{k_0,\nu,\gamma}(q^{\nu,\gamma} - q_0^{\nu,\gamma})\|_{\mathcal{H},r_0} \right).
\label{eq:finaleqn2}
\end{align}


Plugging \eqref{eq:finaleqn2} into \eqref{eqn::nuproofstart}, we obtain
\begin{align}
& E_0\!\left[V_{q_{\nu}^{\pi,\gamma'}}^\pi(S) - V_{q_{0,\nu}^{\pi,\gamma'}}^\pi(S)\right] \notag\\
&\quad + E_0\!\left[d_{\rho,r}^{\pi,\gamma'}(A,S)\{r_{\nu,g}(A,S) + \gamma' V_{q_{\nu}^{\pi,\gamma'}}^\pi(S') - q_{\nu}^{\pi,\gamma'}(A,S)\}\right]  \notag \\
&= E_0\!\Big[(d_{\rho,r}^{\pi,\gamma'}(A,S) - d_0^{\pi,\gamma'}(A,S))\{q_{0,\nu}^{\pi,\gamma'}(A,S) - \gamma' V_{q_{0,\nu}^{\pi,\gamma'}}^\pi(S') + \gamma' V_{q_{\nu}^{\pi,\gamma'}}^\pi(S') - q_{\nu}^{\pi,\gamma'}(A,S)\}\Big]    \notag \\
&\quad + E_0\!\left[d_{\rho,r}^{\pi,\gamma'}(A,S)\Bigl(1 - \tfrac{\nu}{\pi}(A \mid S)\Bigr)\{q^{\nu,\gamma}(A,S) - \gamma V_{q^{\nu,\gamma}}^{\nu}(S') - r_0(A,S) + g(S)\}\right] \notag \\
&\quad + O\!\left(\|r - r_0\|_{\mathcal{H},r_0}\,\|\mathcal{T}_{k_0,\nu,\gamma}(q^{\nu,\gamma} - q_0^{\nu,\gamma})\|_{\mathcal{H},r_0}\right) \notag \\
&= E_0\!\Big[(d_{\rho,r}^{\pi,\gamma'}(A,S) - d_0^{\pi,\gamma'}(A,S))\{q_{0,\nu}^{\pi,\gamma'}(A,S) - \gamma' V_{q_{0,\nu}^{\pi,\gamma'}}^\pi(S') + \gamma' V_{q_{\nu}^{\pi,\gamma'}}^\pi(S') - q_{\nu}^{\pi,\gamma'}(A,S)\}\Big]    \notag \\
&\quad + E_0\!\left[d_{\rho,r}^{\pi,\gamma'}(A,S)\Bigl(1 - \tfrac{\nu}{\pi}(A \mid S)\Bigr)\{q^{\nu,\gamma}(A,S) - \gamma V_{q^{\nu,\gamma}}^{\nu}(S') - r(A,S) + g(S)\}\right] \notag \\
&\quad + E_0\!\left[d_{\rho,r}^{\pi,\gamma'}(A,S)\Bigl(1 - \tfrac{\nu}{\pi}(A \mid S)\Bigr)\{r(A,S) - r_0(A,S)\}\right] \notag \\
&\quad + O\!\left(\|r - r_0\|_{\mathcal{H},r_0}\,\|\mathcal{T}_{k_0,\nu,\gamma}(q^{\nu,\gamma} - q_0^{\nu,\gamma})\|_{\mathcal{H},r_0}\right).
\label{eq:main_expansionnu}
\end{align}
where in the final equality we have added and subtracted \(r(A,S)-g(S)\).


Arguing as in the derivation of \eqref{eqn::proofsecondterm} in the proof of Lemma~\ref{lemma::vonmisespolicyvalue}, we obtain
\begin{align}
E_0\!\left[d_{\rho,r}^{\pi,\gamma'}(A,S)\Bigl(1 - \tfrac{\nu}{\pi}(A \mid S)\Bigr)\{r - r_0\}(A,S)\right]
&= E_0\!\left[\rho(S)\tfrac{\pi}{\pi_r}(A\mid S)\Bigl(1 - \tfrac{\nu}{\pi}(A \mid S)\Bigr)\{r - r_0\}(A,S)\right] \notag \\
&= E_0\!\Bigg[\rho(S)\Bigl(\tfrac{\pi}{\pi_r}(A\mid S) - \tfrac{\nu}{\pi_r}(A\mid S)\Bigr)\Bigl(\tfrac{\pi_r}{\pi_0} - 1\Bigr)(A,S)\Bigg] \notag \\
&\quad + O\!\left(\|r - r_0\|_{\mathcal{H}, 0}^2\right) \notag \\
&= E_0\!\left[\rho(S)\Bigl(\tfrac{\pi}{\pi_0}(A\mid S) - \tfrac{\nu}{\pi_0}(A\mid S)\Bigr)\right] \notag \\
&\quad - E_0\!\left[\rho(S)\Bigl(\tfrac{\pi}{\pi_r}(A\mid S) - \tfrac{\nu}{\pi_r}(A\mid S)\Bigr)\right] \notag \\
&\quad + O\!\left(\|r - r_0\|_{\mathcal{H}, 0}^2\right) \notag \\
&= -\,E_0\!\left[\rho(S)\Bigl(\tfrac{\pi}{\pi_r}(A\mid S) - \tfrac{\nu}{\pi_r}(A\mid S)\Bigr)\right] \notag \\
&\quad + O\!\left(\|r - r_0\|_{\mathcal{H}, 0}^2\right).
\label{eq:reward_diff_expansion}
\end{align}
where in the second equality we expanded $r - r_0$ around $r_0$ and collected higher-order terms. In the final step, we used the law of iterated expectations to show that 
\[
E_0\!\left[\rho(S)\Bigl(\tfrac{\pi}{\pi_0}(A\mid S) - \tfrac{\nu}{\pi_0}(A\mid S)\Bigr)\right] = 0.
\]


Plugging \eqref{eq:reward_diff_expansion} into \eqref{eq:main_expansionnu}, we conclude that
\begin{align*} 
& E_0\!\left[V_{q_{\nu}^{\pi,\gamma'}}^\pi(S) - V_{q_{0,\nu}^{\pi,\gamma'}}^\pi(S)\right] 
+ E_0\!\left[d_{\rho,r}^{\pi,\gamma'}(A,S)\{r_{\nu,g}(A,S) + \gamma' V_{q_{\nu}^{\pi,\gamma'}}^\pi(S') - q_{\nu}^{\pi,\gamma'}(A,S)\}\right]  \\
&= E_0\!\Big[(d_{\rho,r}^{\pi,\gamma'}(A,S) - d_0^{\pi,\gamma'}(A,S))\{q_{0,\nu}^{\pi,\gamma'}(A,S) - \gamma' V_{q_{0,\nu}^{\pi,\gamma'}}^\pi(S') + \gamma' V_{q_{\nu}^{\pi,\gamma'}}^\pi(S') - q_{\nu}^{\pi,\gamma'}(A,S)\}\Big] \\
&\quad + E_0\!\left[d_{\rho,r}^{\pi,\gamma'}(A,S)\Bigl(1 - \tfrac{\nu}{\pi}(A \mid S)\Bigr)\{q^{\nu,\gamma}(A,S) - \gamma V_{q^{\nu,\gamma}}^{\nu}(S') - r(A,S) + g(S)\}\right] \\
&\quad -\,E_0\!\left[\rho(S)\Bigl(\tfrac{\pi}{\pi_r}(A\mid S) - \tfrac{\nu}{\pi_r}(A\mid S)\Bigr)\right] \\
&\quad + O\!\left(\|r - r_0\|_{\mathcal{H},r_0}\,\|\mathcal{T}_{k_0,\nu,\gamma}(q^{\nu,\gamma} - q_0^{\nu,\gamma})\|_{\mathcal{H},r_0}\right) \\
&\quad + O\!\left(\|r - r_0\|_{\mathcal{H}, 0}^2\right).
\end{align*}
Rearranging terms,  
\begin{align*} 
& E_0\!\left[V_{q_{\nu}^{\pi,\gamma'}}^\pi(S) - V_{q_{0,\nu}^{\pi,\gamma'}}^\pi(S)\right] 
+ E_0\!\left[d_{\rho,r}^{\pi,\gamma'}(A,S)\{r_{\nu,g}(A,S) + \gamma' V_{q_{\nu}^{\pi,\gamma'}}^\pi(S') - q_{\nu}^{\pi,\gamma'}(A,S)\}\right] \\
&\quad + E_0\!\left[d_{\rho,r}^{\pi,\gamma'}(A,S)\Bigl(1 - \tfrac{\nu}{\pi}(A \mid S)\Bigr)\{r(A,S)-g(S) + \gamma V_{q^{\nu,\gamma}}^{\nu}(S') - q^{\nu,\gamma}(A,S)\}\right] \\
&\quad + E_0\!\left[\rho(S)\Bigl(\tfrac{\pi}{\pi_r}(A\mid S) - \tfrac{\nu}{\pi_r}(A\mid S)\Bigr)\right] \\
&= E_0\!\Big[(d_{\rho,r}^{\pi,\gamma'}(A,S) - d_0^{\pi,\gamma'}(A,S))\{q_{0,\nu}^{\pi,\gamma'}(A,S) - \gamma' V_{q_{0,\nu}^{\pi,\gamma'}}^\pi(S') + \gamma' V_{q_{\nu}^{\pi,\gamma'}}^\pi(S') - q_{\nu}^{\pi,\gamma'}(A,S)\}\Big] \\
&\quad + O\!\left(\|r - r_0\|_{\mathcal{H},r_0}\,\|\mathcal{T}_{k_0,\nu,\gamma}(q^{\nu,\gamma} - q_0^{\nu,\gamma})\|_{\mathcal{H},r_0}\right) \\
&\quad + O\!\left(\|r - r_0\|_{\mathcal{H}, 0}^2\right).
\end{align*}



 


 
 


 
 



\end{proof}




 



In the following lemma, for each $(r, v^\star)$, define
\[
\pi_{r,v^\star}^\star(a \mid s)
\;\propto\;
\mathbf{1}\{a \in \mathcal{A}^\star\}
\exp\!\left\{
    \tfrac{1}{\tau^\star}\bigl(r(a,s) + \gamma v^\star(a,s)\bigr)
\right\}.
\]
Further, define
\[
\Phi_{r,v^\star}^\star(s)
\;:=\;
\tau^\star \log\!\Biggl(
\sum_{a \in \mathcal{A}^\star}
\exp\!\left\{
    \tfrac{1}{\tau^\star}\bigl(r(a,s) + \gamma v^\star(a,s)\bigr)
\right\}
\Biggr),
\]
and let $\widetilde{d}_{r,v,\widetilde{\rho}^\star}^\star(a,s)
\;=\;
\widetilde{\rho}^\star(s)
\frac{\pi_{r,v^\star}^\star(a \mid s)}{\pi_r(a \mid s)}.$


\begin{lemma}[von Mises expansion for soft optimal value]
\label{lemma::vonmises_softstar}
Let $q^\star, r, v^\star \in L^\infty(\lambda)$ and let $\rho^\star, \widetilde{\rho}^\star \in L^\infty(\mu)$ be arbitrary functions. Then
    \begin{align*}
&\;E_0\!\Big[
    \sum_{a \in \mathcal{A}^\star} 
        \pi_{r,v^\star}^\star(a \mid S)\,q^\star(a,S) 
        - 
        \sum_{a \in \mathcal{A}^\star} 
        \pi_{r_0,k_0}^\star(a \mid S)\,q_{r_0,k_0}^\star(a,S)
\Big]
\\
&\quad+\;
E_0\!\Big[
    \rho^\star(S)\,
    \tfrac{\pi_{r,v^\star}^\star}{\pi_r}(A \mid S)\,
    \bigl\{ r(A,S)+\gamma\,V_{0,q^\star}^\star(S')-q^\star(A,S)\bigr\}
\Big]
\\
&\quad+\;
E_0\!\Big[
    \rho^\star(S)\,
    \Bigl\{1-\tfrac{\pi_{r,v^\star}^\star}{\pi_r}(A \mid S)\Bigr\}
\Big]
\\
&\quad+\;
\frac{\gamma}{\tau^\star}\,
E_0\!\Big[
       \widetilde{d}_{r,v,\widetilde{\rho}^\star}^\star(A,S)\,
       \Bigl(\Phi_{r,v^\star}^\star(S') - v^\star(A,S)\Bigr)
\Big]
\\
&\quad+\;
\frac{1}{\tau^\star}\,
E_0\!\Big[
        \widetilde{\rho}^\star(S)\,
        \Bigl(\tfrac{\pi_{r,v^\star}^\star}{\pi_r}-1\Bigr)(A \mid S)
\Big]
\\[6pt]
&=\;
O\!\Big(
     \|\rho^\star - \rho_{r_0,k_0}^\star\|_{\mathcal{H},r_0}\,
     \bigl\|\mathcal{T}_{r_0,k_0}^\star(q^\star - q_{r_0,k_0}^\star)\bigr\|_{\mathcal{H},r_0}
\\
&\hspace{2.2cm}
     +\;
     \bigl(\|r-r_0\|_{\mathcal{H},r_0} + \|v^\star - v_{r_0,k_0}^\star\|_{\mathcal{H},r_0}\bigr)\,
     \bigl\|\mathcal{T}_{r_0,k_0}^\star(q^\star - q_{r_0,k_0}^\star)\bigr\|_{\mathcal{H},r_0}
\\
&\hspace{2.2cm}
     +\;
     \|r-r_0\|_{\mathcal{H},r_0}^{2}
     +\;
     \|v^\star - v_{r_0,k_0}^\star\|_{\mathcal{H},r_0}^{2}
     +\;
     \|\widetilde{\rho}^\star - \widetilde{\rho}_{r_0,k_0}^\star\|_{\mathcal{H},r_0}\,
     \bigl(\|r-r_0\|_{\mathcal{H},r_0}+\|v^\star - v_{r_0,k_0}^\star\|_{\mathcal{H},r_0}\bigr)
\Big),
\end{align*}
where the implicit constant in the big-$O$ term depends only on 
\(\|q^\star\|_\infty\), 
\(\|r\|_\infty\), 
\(\|r_0\|_\infty\), 
\(\|v^\star\|_\infty\), 
\(\|v_{r_0,k_0}^\star\|_\infty\), 
\(\|\rho^\star\|_\infty\), 
\(\|\widetilde{\rho}^\star\|_\infty\), 
\(\gamma\), and \(\tau^\star\).

\end{lemma}


\begin{proof}[Proof of Lemma \ref{lemma::vonmises_softstar}]
    
 
In this proof, write $v_0 := v_{r_0,k_0}^\star$, so that
\(\pi_{r_0,v_0}^\star = \pi_{r_0,k_0}^\star\), and, for each $(r,v^\star)$,
\[
Q_{r,v^\star} := r + \gamma v^\star,
\qquad
\Phi_{r,v^\star}^\star(s)
:= \tau^\star \log\!\Biggl(
\sum_{a \in \mathcal{A}^\star}
\exp\!\left\{
    \tfrac{1}{\tau^\star}\bigl(r(a,s) + \gamma v^\star(a,s)\bigr)
\right\}
\Biggr).
\]
We have
\begin{align*}
E_0\!\left[\sum_{a \in \mathcal{A}^\star} \pi_{r,v^\star}^\star(a \mid S)\, q^\star(a,S)\right]
- E_0\!\left[\sum_{a \in \mathcal{A}^\star} \pi_{r_0,v_0}^\star(a \mid S)\, q_{r_0,k_0}^\star(a,S)\right]
&= E_0\!\left[
   \sum_{a \in \mathcal{A}^\star}
   \bigl(\pi_{r,v^\star}^\star - \pi_{r_0,v_0}^\star\bigr)(a \mid S)\,
   q_{r_0,k_0}^\star(a,S)\right] \\
&\quad + E_0\!\left[
   \sum_{a \in \mathcal{A}^\star} \pi_{r_0,v_0}^\star(a \mid S)\,
   \bigl(q^\star(a,S) - q_{r_0,k_0}^\star(a,S)\bigr)\right].
\end{align*}







\noindent \textbf{Bounding the second term.}
We begin by deriving a bound for the second term on the right-hand side of the above display. Note that
\[
E_0\!\left[
    \sum_{a \in \mathcal{A}^\star} 
        \pi_{r_0,v_0}^\star(a \mid S)\,
        \bigl(q^\star(a,S) - q_{r_0,k_0}^\star(a,S)\bigr)
\right]
= 
E_0\!\left[
    V_{0,q^\star}^\star(S) - V_{0,q_{r_0,k_0}^\star}^\star(S)
\right],
\]
where
\[
V_{0,q}^\star(s)
:= 
\sum_{a \in \mathcal{A}^\star} 
    \pi_{r_0,v_0}^\star(a \mid s)\, q(a,s).
\]
Applying the proof of Lemma~\ref{lemma::vonmisespolicyvalue} with $\pi := \pi_{r_0,v_0}^\star$ yields
\begin{align*}
& E_0\!\Big[
    \sum_{a \in \mathcal{A}^\star} 
        \pi_{r_0,v_0}^\star(a \mid S)\,
        \{ q^\star(a,S) - q_{r_0,k_0}^\star(a,S) \}
\Big]
\\[-2pt]
&\hspace{2.8cm}
+\;
E_0\!\Big[\rho^\star(S)\,
       \tfrac{\pi_{r_0,v_0}^\star}{\pi_r}(A \mid S)\,
        \bigl\{
            r_0(A,S)
            + \gamma\,V_{0,q^\star}^\star(S')
            - q^\star(A,S)
        \bigr\}
\Big]
\\[6pt]
&\quad=\;
E_0\!\Big[
    \pi_{r_0,v_0}^\star(A \mid S)
    \bigl\{
        \rho^\star(S)\,\pi_r(A \mid S)^{-1}
        - \rho_{r_0,k_0}^\star(S)\,\pi_0(A \mid S)^{-1}
    \bigr\}
\\[-2pt]
&\hspace{4.1cm}\times
	    \Bigl(
	        q^\star(A,S)
	        - q_{r_0,k_0}^\star(A,S)
	        - \gamma\{
	            V_{0,q^\star}^\star(S')
	            - V_{r_0,k_0}^\star(S')
	        \}
	    \Bigr)
	\Big]
\\[6pt]
&\quad=\;
O\!\Big(
    \bigl\|
        \rho^\star\,\pi_r^{-1}
        - \rho_{r_0,k_0}^\star\,\pi_0^{-1}
    \bigr\|_{\mathcal{H},r_0}\,
    \bigl\|
        \mathcal{T}_{r_0,k_0}^\star
        \bigl(q^\star - q_{r_0,k_0}^\star\bigr)
    \bigr\|_{\mathcal{H},r_0}
\Big).
\end{align*}
Next,
\begin{align*}
& E_0\!\Big[\rho^\star(S)\,
       \tfrac{\pi_{r_0,v_0}^\star}{\pi_r}(A \mid S)\,
        \bigl\{
            r_0(A,S)
            + \gamma\,V_{0,q^\star}^\star(S')
            - q^\star(A,S)
        \bigr\}
\Big]
\\
&\;=\;
E_0\!\Big[\rho^\star(S)\,
       \tfrac{\pi_{r,v^\star}^\star}{\pi_r}(A \mid S)\,
        \bigl\{
            r_0(A,S)
            + \gamma\,V_{0,q^\star}^\star(S')
            - q^\star(A,S)
        \bigr\}
\Big]
\\[-2pt]
&\hspace{1.2cm}
+\;
E_0\!\Big[\rho^\star(S)\,
       \tfrac{\pi_{r_0,v_0}^\star - \pi_{r,v^\star}^\star}{\pi_r}(A \mid S)\,
        \bigl\{
            r_0(A,S)
            + \gamma\,V_{0,q^\star}^\star(S')
            - q^\star(A,S)
        \bigr\}
\Big]
\\[6pt]
&\;=\;
E_0\!\Big[\rho^\star(S)\,
       \tfrac{\pi_{r,v^\star}^\star}{\pi_r}(A \mid S)\,
        \bigl\{
            r_0(A,S)
            + \gamma\,V_{0,q^\star}^\star(S')
            - q^\star(A,S)
        \bigr\}
\Big]
\\[-2pt]
&\hspace{1.2cm}
+\;
O\!\Big(
    \|Q_{r,v^\star} - Q_{r_0,v_{r_0,k_0}^\star}\|_{\mathcal{H},r_0}\,
    \bigl\|
        \mathcal{T}_{r_0,k_0}^\star(q^\star - q_{r_0,k_0}^\star)
    \bigr\|_{\mathcal{H},r_0}
\Big).
\end{align*}

Moreover, by adding and subtracting $r$ and using Lemma~\ref{lemma::vonmisespolicyvalue},
\begin{align*}
& E_0\!\Big[\rho^\star(S)\,
       \tfrac{\pi_{r,v^\star}^\star}{\pi_r}(A \mid S)\,
        \bigl\{
            r_0(A,S)
            + \gamma\,V_{0,q^\star}^\star(S')
            - q^\star(A,S)
        \bigr\}
\Big]
\\
&\;=\;
E_0\!\Big[\rho^\star(S)\,
       \tfrac{\pi_{r,v^\star}^\star}{\pi_r}(A \mid S)\,
        \bigl\{
            r(A,S)
            + \gamma\,V_{0,q^\star}^\star(S')
            - q^\star(A,S)
        \bigr\}
\Big]
\\[-2pt]
&\hspace{1.2cm}
+\;
E_0\!\Big[\rho^\star(S)\{r_0(A,S) - r(A,S)\}\Big]
+
O\!\bigl(\|r - r_0\|_{\mathcal{H},r_0}^2\bigr).
\end{align*}

Combining the above results, we obtain
\begin{align*}
& E_0\!\Big[
    \sum_{a \in \mathcal{A}^\star} 
        \pi_{r_0,v_0}^\star(a \mid S)\,
        \{ q^\star(a,S) - q_{r_0,k_0}^\star(a,S) \}
\Big]
\\[-2pt]
&\hspace{2.8cm}
+\;
E_0\!\Big[
    \rho^\star(S)\,
    \tfrac{\pi_{r,v^\star}^\star}{\pi_r}(A \mid S)\,
    \bigl\{
        r(A,S)
        + \gamma\,V_{0,q^\star}^\star(S')
        - q^\star(A,S)
    \bigr\}
\Big]
\\[-2pt]
&\hspace{2.8cm}
+\;
E_0\!\Big[
    \rho^\star(S)\,
    \Bigl\{
        1 - \tfrac{\pi_{r,v^\star}^\star}{\pi_r}(A \mid S)
    \Bigr\}
\Big]
\;+\;
O\!\bigl(\|r - r_0\|_{\mathcal{H},r_0}^2\bigr)
\\[6pt]
&\quad=\;
O\!\Big(
    \bigl\|
        \rho^\star\,\pi_r^{-1}
        - \rho_{r_0,k_0}^\star\,\pi_0^{-1}
    \bigr\|_{\mathcal{H},r_0}\,
    \bigl\|
        \mathcal{T}_{r_0,k_0}^\star
        \bigl(q^\star - q_{r_0,k_0}^\star\bigr)
    \bigr\|_{\mathcal{H},r_0}
\Big)
\;+\;
O\!\bigl(\|r - r_0\|_{\mathcal{H},r_0}^2\bigr)
\\[4pt]
&\qquad\qquad
+\;
O\!\Big(
    \|Q_{r,v^\star} - Q_{r_0,v_{r_0,k_0}^\star}\|_{\mathcal{H},r_0}\,
    \bigl\|
        \mathcal{T}_{r_0,k_0}^\star
        \bigl(q^\star - q_{r_0,k_0}^\star\bigr)
    \bigr\|_{\mathcal{H},r_0}
\Big)
\\[6pt]
&\quad=\;
O\!\Big(
    \Bigl\{
        \bigl\|\rho^\star - \rho_{r_0,k_0}^\star\bigr\|_{\mathcal{H},r_0}
        +
        \bigl\|r - r_0\bigr\|_{\mathcal{H},r_0}
        +
        \bigl\|v^\star - v_{r_0,k_0}^\star\bigr\|_{\mathcal{H},r_0}
    \Bigr\}
    \bigl\|
        \mathcal{T}_{r_0,k_0}^\star
        \bigl(q^\star - q_{r_0,k_0}^\star\bigr)
    \bigr\|_{\mathcal{H},r_0}
\Big)
\;+\;
O\!\bigl(\|r - r_0\|_{\mathcal{H},r_0}^2\bigr).
\end{align*}








\noindent \textbf{Bounding first term.}
\begin{align*}
&\sum_{a \in \mathcal{A}^\star} \bigl(\pi_{r,v^\star}^\star - \pi_{r_0,v_0}^\star\bigr)(a \mid S)\, q_{r_0,k_0}^\star(a,S)\\
&= \frac{1}{\tau^\star} \sum_{a \in \mathcal{A}^\star} q_{r_0,k_0}^\star(a,S)\,\pi_{r_0,v_0}^\star(a \mid S)
   \Bigl\{ \bigl(Q_{r,v^\star} - Q_{r_0,v_{r_0,k_0}^\star}\bigr)(a,S)
           - \sum_{\tilde a \in \mathcal{A}^\star} \pi_{r_0,v_0}^\star(\tilde a \mid S)\,
             \bigl(Q_{r,v^\star}- Q_{r_0,v_{r_0,k_0}^\star}\bigr)(\tilde a,S)
   \Bigr\}  \\
&\qquad +\, O\!\left(\|Q_{r,v^\star}- Q_{r_0,v_{r_0,k_0}^\star}(\cdot, S)\|_2^2\right) \\
&= \frac{1}{\tau^\star} \sum_{a \in \mathcal{A}^\star} \pi_{r_0,v_0}^\star(a \mid S)
   \Bigl\{ q_{r_0,k_0}^\star(a,S) - \sum_{\tilde a \in \mathcal{A}^\star} \pi_{r_0,v_0}^\star(\tilde a \mid S)\, q_{r_0,k_0}^\star(\tilde a,S) \Bigr\}
   \bigl(Q_{r,v^\star} - Q_{r_0,v_{r_0,k_0}^\star}\bigr)(a,S) \\
&\qquad +\, O\!\left(\|Q_{r,v^\star}(\cdot,S) - Q_{r_0,v_{r_0,k_0}^\star}(\cdot,S)\|^2\right).
\end{align*}
Taking expectations, we obtain
\begin{align*}
& E_0\!\left[\sum_{a \in \mathcal{A}^\star} \bigl(\pi_{r,v^\star}^\star - \pi_{r_0,v_0}^\star\bigr)(a \mid S)\, q_{r_0,k_0}^\star(a,S)\right]\\
&= \frac{1}{\tau^\star}\,
E_0\!\left[\sum_{a \in \mathcal{A}^\star} \pi_{r_0,v_0}^\star(a \mid S)
   \Bigl\{ q_{r_0,k_0}^\star(a,S) - \sum_{\tilde a \in \mathcal{A}^\star} \pi_{r_0,v_0}^\star(\tilde a \mid S)\, q_{r_0,k_0}^\star(\tilde a,S) \Bigr\}
   \bigl(Q_{r,v^\star} - Q_{r_0,v_{r_0,k_0}^\star}\bigr)(a,S)\right] \\
&\qquad +\, O\!\left(\|Q_{r,v^\star} - Q_{r_0,v_{r_0,k_0}^\star}\|^2_{\mathcal{H},r_0}\right).
\end{align*}





Denote
\[
F_{k_0}(q)(a,s)
:= q(a,s) - \gamma \int \tau^\star \log \!\sum_{a' \in \mathcal{A}^\star}
      \exp\!\left(\frac{q(a',s')}{\tau^\star}\right)
      k_0(s' \mid a,s)\,\mu(ds').
\]
Since \(Q_{r_0,v_{r_0,k_0}^\star}\) solves the soft Bellman equation, \(F_{k_0}(Q_{r_0,v_{r_0,k_0}^\star})=r_0\).
By Lemma \ref{lemma::bellmansmoothnessfull}, a first-order expansion around \(Q_{r_0,v_{r_0,k_0}^\star}\) yields
\begin{align*}
F_{k_0}(Q_{r,v^\star})
&= F_{k_0}(Q_{r_0,v_{r_0,k_0}^\star})
  + \mathcal{T}_{r_0,k_0}^\star\!\bigl(Q_{r,v^\star} - Q_{r_0,v_{r_0,k_0}^\star}\bigr)
  \;+\; \mathrm{Rem}(Q_{r,v^\star}),
\end{align*}
where the remainder satisfies
\[
\bigl\|\mathrm{Rem}(Q_{r,v^\star})\bigr\|_{L^2(\lambda)}
= O\!\bigl(\,\|Q_{r,v^\star} - Q_{r_0,v_{r_0,k_0}^\star}\|_{L^2(\lambda)}^{\,2}\bigr).
\]
Rearranging terms, we obtain
\begin{align*}
Q_{r,v^\star} - Q_{r_0,v_{r_0,k_0}^\star}
&= (\mathcal{T}_{r_0,k_0}^\star)^{-1}\!\bigl(F_{k_0}(Q_{r,v^\star}) - r_0\bigr)
  \;+\; \mathrm{Rem}_2(Q_{r,v^\star}),
\end{align*}
where
\[
\mathrm{Rem}_2(Q_{r,v^\star})
:= (\mathcal{T}_{r_0,k_0}^\star)^{-1}\mathrm{Rem}(Q_{r,v^\star}),
\qquad
\bigl\|\mathrm{Rem}_2(Q_{r,v^\star})\bigr\|_{L^2(\lambda)}
= O\!\bigl(\,\|Q_{r,v^\star} - Q_{r_0,v_{r_0,k_0}^\star}\|_{L^2(\lambda)}^{\,2}\bigr).
\]
By equivalence of the \(L^2(\lambda)\) and \(\|\cdot\|_{\mathcal{H},r_0}\) norms, it follows that
\[
\bigl\|\mathrm{Rem}_2(Q_{r,v^\star})\bigr\|_{\mathcal{H},r_0}
= O\!\bigl(\|Q_{r,v^\star} - Q_{r_0,v_{r_0,k_0}^\star}\|_{\mathcal{H},r_0}^{\,2}\bigr).
\]


We thus have
\begin{align*}
E_0\!\Bigg[\sum_{a \in \mathcal{A}^\star} \pi_{r_0,v_0}^\star(a \mid S)
   \Bigl\{ q_{r_0,k_0}^\star(a,S) - \sum_{\tilde a \in \mathcal{A}^\star} 
          \pi_{r_0,v_0}^\star(\tilde a \mid S)\, q_{r_0,k_0}^\star(\tilde a,S) \Bigr\}
   \bigl(Q_{r,v^\star} - Q_{r_0,v_{r_0,k_0}^\star}\bigr)(a,S)\Bigg] \\
=\;
E_0\!\Bigg[\sum_{a \in \mathcal{A}^\star} \pi_{r_0,v_0}^\star(a \mid S)
   \Bigl\{ q_{r_0,k_0}^\star(a,S) - \sum_{\tilde a \in \mathcal{A}^\star} 
          \pi_{r_0,v_0}^\star(\tilde a \mid S)\, q_{r_0,k_0}^\star(\tilde a,S) \Bigr\}
   (\mathcal{T}_{r_0,k_0}^\star)^{-1}\!\bigl(F_{k_0}(Q_{r,v^\star}) - r_0\bigr)(a,S)\Bigg] \\
\quad +\;
E_0\!\left[\sum_{a \in \mathcal{A}^\star} \pi_{r_0,v_0}^\star(a \mid S)
   \Bigl\{ q_{r_0,k_0}^\star(a,S) - \sum_{\tilde a \in \mathcal{A}^\star} 
          \pi_{r_0,v_0}^\star(\tilde a \mid S)\, q_{r_0,k_0}^\star(\tilde a,S) \Bigr\}
   \mathrm{Rem}_2(Q_{r,v^\star})(a,S)\right].
\end{align*}
By Cauchy--Schwarz and boundedness of \(Q_{r,v^\star}\), we have
\begin{align*}
&E_0\!\left[\sum_{a \in \mathcal{A}^\star} \pi_{r_0,v_0}^\star(a \mid S)
   \Bigl\{ q_{r_0,k_0}^\star(a,S) - \sum_{\tilde a \in \mathcal{A}^\star}
          \pi_{r_0,v_0}^\star(\tilde a \mid S)\, q_{r_0,k_0}^\star(\tilde a,S) \Bigr\}\right. \\
&\qquad\qquad \left.\times
   \mathrm{Rem}_2(Q_{r,v^\star})(a,S)\right]
&\;\le\; \bigl\|\mathrm{Rem}_2(Q_{r,v^\star})\bigr\|_{\mathcal{H},r_0} \\
&=\; O\!\left(\|Q_{r,v^\star} - Q_{r_0,v_{r_0,k_0}^\star}\|_{\mathcal{H},r_0}^{\,2}\right).
\end{align*}
Hence,
\begin{align*}
&E_0\!\left[\sum_{a \in \mathcal{A}^\star} \pi_{r_0,v_0}^\star(a \mid S)
   \Bigl\{ q_{r_0,k_0}^\star(a,S) - \sum_{\tilde a \in \mathcal{A}^\star}
          \pi_{r_0,v_0}^\star(\tilde a \mid S)\, q_{r_0,k_0}^\star(\tilde a,S) \Bigr\}\right. \\
&\qquad\qquad \left.\times
   \bigl(Q_{r,v^\star} - Q_{r_0,v_{r_0,k_0}^\star}\bigr)(a,S)\right] \\
&\;=\; E_0\!\left[\sum_{a \in \mathcal{A}^\star} \pi_{r_0,v_0}^\star(a \mid S)
   \Bigl\{ q_{r_0,k_0}^\star(a,S) - \sum_{\tilde a \in \mathcal{A}^\star}
          \pi_{r_0,v_0}^\star(\tilde a \mid S)\, q_{r_0,k_0}^\star(\tilde a,S) \Bigr\}\right. \\
&\qquad\qquad \left.\times
   (\mathcal{T}_{r_0,k_0}^\star)^{-1}\!\bigl(F_{k_0}(Q_{r,v^\star}) - r_0\bigr)(a,S)\right] \\
&\quad + O\!\left(\|Q_{r,v^\star} - Q_{r_0,v_{r_0,k_0}^\star}\|_{\mathcal{H},r_0}^{\,2}\right).
\end{align*}



We now use the identity
\[
\bigl((\mathcal{T}_{r_0,k_0}^\star)^{-1}f\bigr)(a,s)
    = \sum_{a' \in \mathcal{A}^\star} \int 
        \frac{\pi_{r_0,v_0}^\star(a' \mid s')}{\pi_0(a' \mid s')}
        \,\rho_{r_0,k_0}^\star(s' \mid a,s)\,
        f(a',s')\,\mu(ds').
\]
Substituting \(f = F_{k_0}(Q_{r,v^\star}) - r_0\) and defining
\[
C^\star(a,S) := q_{r_0,k_0}^\star(a,S) - \sum_{\tilde a \in \mathcal{A}^\star} 
          \pi_{r_0,v_0}^\star(\tilde a \mid S)\, q_{r_0,k_0}^\star(\tilde a,S),
\]
we obtain
\begin{align*}
&\;E_0\!\left[\sum_{a \in \mathcal{A}^\star} \pi_{r_0,v_0}^\star(a \mid S)\,
   C^\star(a,S)\;
   \bigl((\mathcal{T}_{r_0,k_0}^\star)^{-1}(F_{k_0}(Q_{r,v^\star})-r_0)\bigr)(a,S)\right] \\
&= E_0\!\left[\sum_{a \in \mathcal{A}^\star} \pi_{r_0,v_0}^\star(a \mid S)\, C^\star(a,S)\,
   \sum_{a' \in \mathcal{A}^\star} \int
        \frac{\pi_{r_0,v_0}^\star(a' \mid s')}{\pi_0(a' \mid s')}
        \,\rho_{r_0,k_0}^\star(s' \mid a,S)\right. \notag\\
&\qquad\qquad \left.\times
        \bigl(F_{k_0}(Q_{r,v^\star})-r_0\bigr)(a',s')\,
        \pi_0(a' \mid s')\rho_0(ds')\right].
\end{align*}
By Fubini--Tonelli, we may interchange the expectation, sum, and integral:
\begin{align*}
&\;= \sum_{a' \in \mathcal{A}^\star} \int
        \frac{\pi_{r_0,v_0}^\star(a' \mid s')}{\pi_0(a' \mid s')}
        \,\bigl(F_{k_0}(Q_{r,v^\star})-r_0\bigr)(a',s') \notag\\
&\qquad\qquad \times
        E_0\!\left[\sum_{a \in \mathcal{A}^\star} 
             \pi_{r_0,v_0}^\star(a \mid S)\, C^\star(a,S)\,
             \rho_{r_0,k_0}^\star(s' \mid a,S)\right]\,
        \pi_0(a' \mid s') \rho_0(ds').
\end{align*}
By the definition of \(\widetilde{\rho}_{r_0,k_0}^\star\) in \eqref{eqn::tilderho}, the inner expectation equals
\(\widetilde{\rho}_{r_0,k_0}^\star(s')\). Therefore,
\begin{align*}
&\;= \sum_{a' \in \mathcal{A}^\star} \int
        \frac{\pi_{r_0,v_0}^\star(a' \mid s')}{\pi_0(a' \mid s')}
        \,\widetilde{\rho}_{r_0,k_0}^\star(s')\,
        \bigl(F_{k_0}(Q_{r,v^\star})-r_0\bigr)(a',s') \,
        \pi_0(a' \mid s')\rho_0(ds') \\
&= E_0 \left[
        \frac{\pi_{r_0,v_0}^\star(A \mid S)}{\pi_0(A \mid S)}\,
        \widetilde{\rho}_{r_0,k_0}^\star(S)\,
        \bigl(F_{k_0}(Q_{r,v^\star})-r_0\bigr)(A,S)
      \right].
\end{align*}


We have now shown that
\begin{align*}
&E_0\!\left[\sum_{a \in \mathcal{A}^\star} \pi_{r_0,v_0}^\star(a \mid S)
   \Bigl\{ q_{r_0,k_0}^\star(a,S) - \sum_{\tilde a \in \mathcal{A}^\star} 
          \pi_{r_0,v_0}^\star(\tilde a \mid S)\, q_{r_0,k_0}^\star(\tilde a,S) \Bigr\}
   \bigl(Q_{r,v^\star} - Q_{r_0,v_{r_0,k_0}^\star}\bigr)(a,S)\right] \\
&\qquad= E_0 \left[
        \frac{\pi_{r_0,v_0}^\star(A \mid S)}{\pi_0(A \mid S)}\,
        \widetilde{\rho}_{r_0,k_0}^\star(S)\,
        \bigl(F_{k_0}(Q_{r,v^\star})-r_0\bigr)(A,S)
      \right] \\
&\quad +
O\!\left(\|Q_{r,v^\star} - Q_{r_0,v_{r_0,k_0}^\star}\|_{\mathcal{H},r_0}^{\,2}\right).
\end{align*}
Taking expectations then gives
\begin{align*}
&E_0\!\left[\sum_{a \in \mathcal{A}^\star}
\bigl(\pi_{r,v^\star}^\star - \pi_{r_0,v_0}^\star\bigr)(a \mid S)\,
q_{r_0,k_0}^\star(a,S)\right]
&= \frac{1}{\tau^\star}\,
	    E_0 \left[ 
	        \frac{\pi_{r_0,v_0}^\star(A \mid S)}{\pi_0(A \mid S)}\,
	        \widetilde{\rho}_{r_0,k_0}^\star(S)\,
	        \bigl(F_{k_0}(Q_{r,v^\star})-r_0\bigr)(A,S) \right] \\
	&\qquad +\;
	O\!\left(\|Q_{r,v^\star} - Q_{r_0,v_{r_0,k_0}^\star}\|_{\mathcal{H},r_0}^{\,2}\right).
\end{align*}

Denote $\widetilde{d}_{r,v,\widetilde{\rho}^\star}^\star(a,s) :=  \frac{\pi_{r,v^\star}^\star(a \mid s)}{\pi_r(a \mid s)}
        \,\widetilde{\rho}^\star(s)$. Next, adding and subtracting, we have
\begin{align*}
&\frac{1}{\tau^\star}\,
E_0 \!\left[
        \frac{\pi_{r_0,v_0}^\star(A \mid S)}{\pi_0(A \mid S)}\,
        \widetilde{\rho}_{r_0,k_0}^\star(S)\,
        \bigl(F_{k_0}(Q_{r,v^\star})-r_0\bigr)(A,S)
\right]
\\[6pt]
&=\frac{1}{\tau^\star}\,
E_0 \!\left[
       \bigl\{\widetilde{d}_{r_0,k_0}^\star(A,S) - \widetilde{d}_{r,v,\widetilde{\rho}^\star}^\star(A,S)\bigr\}\,
        \bigl(F_{k_0}(Q_{r,v^\star})-r_0\bigr)(A,S)
\right]
\\
&\quad + \frac{1}{\tau^\star}\,
E_0 \!\left[
      \widetilde{d}_{r,v,\widetilde{\rho}^\star}^\star(A,S)\,
        \bigl(F_{k_0}(Q_{r,v^\star})-r_0\bigr)(A,S)
\right].
\end{align*}
By Cauchy--Schwarz and the linearization of $F_{k_0}(Q_{r,v^\star})-r_0
= F_{k_0}(Q_{r,v^\star})-F_{k_0}(Q_{r_0,v_{r_0,k_0}^\star})$ we derived previously,
we obtain
\begin{align*}
    &E_0 \!\left[
       \bigl\{\widetilde{d}_{r_0,k_0}^\star(A,S) - \widetilde{d}_{r,v,\widetilde{\rho}^\star}^\star(A,S)\bigr\}\,
        \bigl(F_{k_0}(Q_{r,v^\star})-r_0\bigr)(A,S)
    \right] \\
    &\qquad \le\;
    O\!\left(
        \bigl\|\widetilde{d}_{r,v,\widetilde{\rho}^\star}^\star - \widetilde{d}_{r_0,k_0}^\star\bigr\|_{\mathcal{H},r_0}\,
        \bigl\|Q_{r,v^\star} - Q_{r_0,v_{r_0,k_0}^\star}\bigr\|_{\mathcal{H},r_0}
    \right)
    \;+\;
    O\!\left(
        \bigl\|Q_{r,v^\star} - Q_{r_0,v_{r_0,k_0}^\star}\bigr\|_{\mathcal{H},r_0}^2
    \right).
\end{align*}



Next, add and subtract \(r\) inside the expectation and apply the law of total expectation:
\begin{align*}
\frac{1}{\tau^\star}\,
E_0 \!\left[
       \widetilde{d}_{r,v,\widetilde{\rho}^\star}^\star(A,S)\,
       \bigl(F_{k_0}(Q_{r,v^\star})-r_0\bigr)(A,S)
\right]
&= \frac{\gamma}{\tau^\star}\,
E_0 \!\left[
       \widetilde{d}_{r,v,\widetilde{\rho}^\star}^\star(A,S)\,
       \bigl(v^\star(A,S) - \Phi_{r,v^\star}^\star(S')\bigr)
\right] \\
&\quad + \frac{1}{\tau^\star}\,
E_0 \!\left[
       \widetilde{d}_{r,v,\widetilde{\rho}^\star}^\star(A,S)\,
       \bigl(r - r_0\bigr)(A,S)
\right].
\end{align*}

As in the proof of Lemma~\ref{lemma::vonmisespolicyvalue},
\begin{align*}
\frac{1}{\tau^\star}\,
E_0 \!\left[
        \widetilde{d}_{r,v,\widetilde{\rho}^\star}^\star(A,S)\,
        \bigl(r - r_0\bigr)(A,S)
\right]
&=
\frac{1}{\tau^\star}\,
E_0 \!\left[
        \widetilde{\rho}^\star(S)\,
        \Bigl(1 - \tfrac{\pi_{r,v^\star}^\star}{\pi_r}\Bigr)(A \mid S)
\right]
\;+\;
O\!\left(\|r - r_0\|_{\mathcal{H},r_0}^{\,2}\right).
\end{align*}
Collecting terms, we obtain
\begin{align*}
	&E_0\!\left[
	        \sum_{a \in \mathcal{A}^\star}
	        \bigl(\pi_{r,v^\star}^\star - \pi_{r_0,v_0}^\star\bigr)(a \mid S)\,
	        q_{r_0,k_0}^\star(a,S)
	\right]
	\\[6pt]
	&=\;
	\frac{\gamma}{\tau^\star}\,
	E_0\!\left[
	       \widetilde{d}_{r,v,\widetilde{\rho}^\star}^\star(A,S)\,
	       \Bigl(v^\star(A,S) - \Phi_{r,v^\star}^\star(S')\Bigr)
	\right]
\\[4pt]
&\quad +\;
\frac{1}{\tau^\star}\,
E_0\!\left[
        \widetilde{\rho}^\star(S)\,
        \Bigl(1 - \tfrac{\pi_{r,v^\star}^\star}{\pi_r}\Bigr)(A \mid S)
\right]
\;+\;
O\!\left(\|r - r_0\|_{\mathcal{H},r_0}^{\,2}\right)
\;+\;
O\!\left(\|Q_{r,v^\star} - Q_{r_0,v_{r_0,k_0}^\star}\|_{\mathcal{H},r_0}^{\,2}\right)
\\[4pt]
&\quad +\;
O\!\left(
        \bigl\|\widetilde{d}_{r,v,\widetilde{\rho}^\star}^\star - \widetilde{d}_{r_0,k_0}^\star\bigr\|_{\mathcal{H},r_0}\,
        \bigl\|Q_{r,v^\star} - Q_{r_0,v_{r_0,k_0}^\star}\bigr\|_{\mathcal{H},r_0}
    \right).
\end{align*}
We also know
\[
\bigl\|\widetilde{d}_{r,v,\widetilde{\rho}^\star}^\star - \widetilde{d}_{r_0,k_0}^\star\bigr\|_{\mathcal{H},r_0}
= 
O\!\Bigl(
    \bigl\|\widetilde{\rho}^\star - \widetilde{\rho}_{r_0,k_0}^\star\bigr\|_{\mathcal{H},r_0}
    +  
    \bigl\|Q_{r,v^\star} - Q_{r_0,v_{r_0,k_0}^\star}\bigr\|_{\mathcal{H},r_0}
\Bigr),
\]
and
\[
\bigl\|Q_{r,v^\star} - Q_{r_0,v_{r_0,k_0}^\star}\bigr\|_{\mathcal{H},r_0}
=
O\!\bigl(\|r - r_0\|_{\mathcal{H},r_0} + \|v^\star - v_{r_0,k_0}^\star\|_{\mathcal{H},r_0}\bigr).
\]
Hence,
\begin{align*}
	&E_0\!\left[
	        \sum_{a \in \mathcal{A}^\star}
	        \bigl(\pi_{r,v^\star}^\star - \pi_{r_0,v_0}^\star\bigr)(a \mid S)\,
	        q_{r_0,k_0}^\star(a,S)
	\right]
	\\[6pt]
	&=\;
	\frac{\gamma}{\tau^\star}\,
	E_0\!\left[
	       \widetilde{d}_{r,v,\widetilde{\rho}^\star}^\star(A,S)\,
	       \Bigl(v^\star(A,S) - \Phi_{r,v^\star}^\star(S')\Bigr)
	\right]
\\[4pt]
&\quad +\;
\frac{1}{\tau^\star}\,
E_0\!\left[
        \widetilde{\rho}^\star(S)\,
        \Bigl(1 - \tfrac{\pi_{r,v^\star}^\star}{\pi_r}\Bigr)(A \mid S)
\right]
\;+\;
O\!\left(\|r - r_0\|_{\mathcal{H},r_0}^{\,2}\right)
\;+\;
O\!\left(\|v^\star - v_{r_0,k_0}^\star\|_{\mathcal{H},r_0}^{\,2}\right)
\\[4pt]
&\quad +\;
O\!\left(
       \bigl\|\widetilde{\rho}^\star - \widetilde{\rho}_{r_0,k_0}^\star\bigr\|_{\mathcal{H},r_0}\,
       \bigl(\|r - r_0\|_{\mathcal{H},r_0} + \|v^\star - v_{r_0,k_0}^\star\|_{\mathcal{H},r_0}\bigr)
    \right).
\end{align*}
Rearranging, we obtain
\begin{align*}
	&E_0\!\left[
	        \sum_{a \in \mathcal{A}^\star}
	        \bigl(\pi_{r,v^\star}^\star - \pi_{r_0,v_0}^\star\bigr)(a \mid S)\,
	        q_{r_0,k_0}^\star(a,S)
	\right]
	\\[6pt]
	&+\;
	\frac{\gamma}{\tau^\star}\,
	E_0\!\left[
	       \widetilde{d}_{r,v,\widetilde{\rho}^\star}^\star(A,S)\,
	       \Bigl(\Phi_{r,v^\star}^\star(S') - v^\star(A,S)\Bigr)
	\right]
\\[4pt]
&\quad +\;
\frac{1}{\tau^\star}\,
E_0\!\left[
        \widetilde{\rho}^\star(S)\,
        \Bigl(\tfrac{\pi_{r,v^\star}^\star}{\pi_r} - 1\Bigr)(A \mid S)
\right]
\\[4pt]
&\quad =\;
O\!\left(\|r - r_0\|_{\mathcal{H},r_0}^{\,2}\right)
\;+\;
O\!\left(\|v^\star - v_{r_0,k_0}^\star\|_{\mathcal{H},r_0}^{\,2}\right)
\\[4pt]
&\quad +\;
O\!\left(
       \bigl\|\widetilde{\rho}^\star - \widetilde{\rho}_{r_0,k_0}^\star\bigr\|_{\mathcal{H},r_0}\,
       \bigl(\|r - r_0\|_{\mathcal{H},r_0} + \|v^\star - v_{r_0,k_0}^\star\|_{\mathcal{H},r_0}\bigr)
    \right).
\end{align*}


\noindent \textbf{Combining both bounds.}  
Adding the two bounds together and simplifying the right-hand side, we obtain
\begin{align*}
&\;E_0\!\Big[
    \sum_{a \in \mathcal{A}^\star} 
        \pi_{r,v^\star}^\star(a \mid S)\,q^\star(a,S) 
        - 
        \sum_{a \in \mathcal{A}^\star} 
        \pi_{r_0,v_0}^\star(a \mid S)\,q_{r_0,k_0}^\star(a,S)
\Big]
\\
&\quad+\;
E_0\!\Big[
    \rho^\star(S)\,
    \tfrac{\pi_{r,v^\star}^\star}{\pi_r}(A \mid S)\,
    \bigl\{ r(A,S)+\gamma\,V_{0,q^\star}^\star(S')-q^\star(A,S)\bigr\}
\Big]
\\
&\quad+\;
E_0\!\Big[
    \rho^\star(S)\,
    \Bigl\{1-\tfrac{\pi_{r,v^\star}^\star}{\pi_r}(A \mid S)\Bigr\}
\Big]
\\
&\quad+\;
\frac{\gamma}{\tau^\star}\,
E_0\!\Big[
       \widetilde{d}_{r,v,\widetilde{\rho}^\star}^\star(A,S)\,
       \Bigl(\Phi_{r,v^\star}^\star(S') - v^\star(A,S)\Bigr)
\Big]
\\
&\quad+\;
\frac{1}{\tau^\star}\,
E_0\!\Big[
        \widetilde{\rho}^\star(S)\,
        \Bigl(\tfrac{\pi_{r,v^\star}^\star}{\pi_r}-1\Bigr)(A \mid S)
\Big]
\\[6pt]
&=\;
O\!\Big(
     \|\rho^\star - \rho_{r_0,k_0}^\star\|_{\mathcal{H},r_0}\,
     \bigl\|\mathcal{T}_{r_0,k_0}^\star(q^\star - q_{r_0,k_0}^\star)\bigr\|_{\mathcal{H},r_0}
\\
&\hspace{2.2cm}
     +\;
     \bigl(\|r-r_0\|_{\mathcal{H},r_0} + \|v^\star - v_{r_0,k_0}^\star\|_{\mathcal{H},r_0}\bigr)\,
     \bigl\|\mathcal{T}_{r_0,k_0}^\star(q^\star - q_{r_0,k_0}^\star)\bigr\|_{\mathcal{H},r_0}
\\
&\hspace{2.2cm}
     +\;
     \|r-r_0\|_{\mathcal{H},r_0}^{2}
     +\;
     \|v^\star - v_{r_0,k_0}^\star\|_{\mathcal{H},r_0}^{2}
     +\;
     \|\widetilde{\rho}^\star - \widetilde{\rho}_{r_0,k_0}^\star\|_{\mathcal{H},r_0}\,
     \bigl(\|r-r_0\|_{\mathcal{H},r_0}+\|v^\star - v_{r_0,k_0}^\star\|_{\mathcal{H},r_0}\bigr)
\Big).
\end{align*}

\end{proof}



\subsection{Proofs for Section \ref{sec::examplesestimators}}

 
\begin{proof}[Proof of Theorem \ref{theorem::ALex1}]
Note that $\psi_n = P_n \varphi_n$, where we define the estimated uncentered influence function by
\[
\varphi_n(s', a, s)
:= V_n^{\pi, \gamma}(s)  
\;+\; \rho_n^{\pi,\gamma}(s)\,\frac{\pi(a \mid s)}{\pi_n(a \mid s)}
    \big\{r_n(a, s) + \gamma V_n^{\pi, \gamma}(s') - q_n^{\pi,\gamma}(a, s)\big\}
\;+\; \rho_n^{\pi,\gamma}(s)\Big\{1 - \tfrac{\pi(a \mid s)}{\pi_n(a \mid s)}\Big\}.
\]
Let $\varphi_0$ denote the true uncentered influence function $\varphi_{r_0, k_0}$. Then
\begin{align*}
    \psi_n - \psi_0
    &= P_n \varphi_n - P_0 \varphi_0 \\
    &= P_n(\varphi_0 - P_0 \varphi_0) + P_n(\varphi_n - \varphi_0) \\
    &= P_n \chi_0 + P_n(\varphi_n - \varphi_0) \\
    &= P_n \chi_0 + (P_n - P_0)(\varphi_n - \varphi_0) + P_0(\varphi_n - \varphi_0) \\
    &= P_n \chi_0 + (P_n - P_0)(\varphi_n - \varphi_0) + \{P_0 \varphi_n - \psi_0\}.
\end{align*}
By Assumptions~\ref{assump:boundedness::policy} and~\ref{assump:consistency::policy},  
$\|\varphi_n - \varphi_0\|_{L^2(P_0)} = o_p(1)$. Moreover, by
Assumption~\ref{assump:splitting::policy} and Markov’s inequality,
\[
(P_n - P_0)(\varphi_n - \varphi_0)
= O_p\!\left(n^{-1/2}\,\|\varphi_n - \varphi_0\|_{L^2(P_0)}\right)
= o_p(n^{-1/2}).
\]
Thus,
\[
\psi_n - \psi_0
= P_n \chi_0 + o_p(n^{-1/2}) + \{P_0 \varphi_n - \psi_0\}.
\]
By Lemma~\ref{lemma::vonmisespolicyvalue}, recalling that  
$d_{\rho_n, r_n}^{\pi,\gamma} = \rho_n^{\pi,\gamma}\tfrac{\pi}{\pi_n}$,
\begin{align*}
   \{P_0 \varphi_n - \psi_0\}
   &= O\!\left(
       \big\langle
          d_{\rho_n, r_n}^{\pi,\gamma} - d_{\rho_0, r_0}^{\pi,\gamma},
          \mathcal{T}_{k_0,\pi,\gamma}(q_n^{\pi,\gamma} - q_{r_0, k_0}^{\pi,\gamma})
       \big\rangle_{\mathcal{H}, 0}
     \right)
     + O\bigl(\|r_n - r_0\|_{\mathcal{H}, 0}^2\bigr) \\
   &= O\!\left(
       \|d_{\rho_n, r_n}^{\pi,\gamma} - d_{\rho_0, r_0}^{\pi,\gamma}\|_{\mathcal{H}, 0}\,
       \|\mathcal{T}_{k_0,\pi,\gamma}(q_n^{\pi,\gamma} - q_{r_0, k_0}^{\pi,\gamma})\|_{\mathcal{H}, 0}
     \right)
     + O\bigl(\|r_n - r_0\|_{\mathcal{H}, 0}^2\bigr).
\end{align*}
By Assumption~\ref{assump:boundedness::policy},
\[
\|d_{\rho_n, r_n}^{\pi,\gamma} - d_{\rho_0, r_0}^{\pi,\gamma}\|_{\mathcal{H}, 0}
\lesssim
\|\rho_n^{\pi,\gamma} - \rho_0^{\pi,\gamma}\|_{\mathcal{H}, 0}
+ \|r_n - r_0\|_{\mathcal{H},0}.
\]
Hence,
\begin{align*}
   \{P_0 \varphi_n - \psi_0\}
   &= O\!\left(
        \{\|\rho_n^{\pi,\gamma} - \rho_0^{\pi,\gamma}\|_{L^2(\rho_0)}
          + \|r_n - r_0\|_{\mathcal{H}, 0}\}\,
        \|\mathcal{T}_{k_0,\pi,\gamma}(q_n^{\pi,\gamma} - q_0^{\pi,\gamma})\|_{\mathcal{H}, 0}
      \right)
      + O\bigl(\|r_n - r_0\|_{\mathcal{H}, 0}^2\bigr).
\end{align*}
The first term is $o_p(n^{-1/2})$ by Assumption~\ref{assump:rates::policy}, and the second is $o_p(n^{-1/2})$ by \ref{cond::rates}\ref{cond::reward}. Therefore,
\[
\psi_n - \psi_0
= P_n \chi_0 + o_p(n^{-1/2}),
\]
and the result follows.
\end{proof}



\begin{proof}[Proof of Theorem \ref{theorem::ALex1b}]
Note that $\psi_n = P_n \varphi_n$, where we define the estimated uncentered influence function by
\begin{align*}
\varphi_n(s',a,s)
&:= V_n^\star(s) \\
&\quad + \rho_n^\star(s)\,
        \frac{\pi_n^\star(a \mid s)}{\pi_n(a \mid s)}
        \bigl\{r_n(a,s) + \gamma\,V_n^\star(s') - q_n^\star(a,s)\bigr\} \\
&\quad + \frac{\gamma}{\tau^\star}\,
        \widetilde{\rho}_n^\star(s)\,
        \frac{\pi_n^\star(a \mid s)}{\pi_n(a \mid s)}
        \left[
           \tau^\star \log\!\left( 
                \sum_{\tilde a \in \mathcal{A}^\star} 
                \exp\!\left( 
                    \tfrac{1}{\tau^\star}\bigl( r_n(\tilde a,s') + \gamma\,v_n^\star(\tilde a,s') \bigr) 
                \right)
            \right)
            - v_n^\star(a,s)
        \right] \\
&\quad 
    + \left\{
            \tfrac{1}{\tau^\star}\,\widetilde{\rho}_n^\star(s)
            + \rho_n^\star(s)
      \right\}
      \left\{ 
            \tfrac{\pi_n^\star(a \mid s)}{\pi_n(a \mid s)} - 1 
      \right\}.
\end{align*}
Let $\varphi_0$ denote the true uncentered influence function corresponding to the nuisances 
$(r_0, v_{r_0,k_0}^\star, q_{r_0,k_0}^\star, \rho_{r_0,k_0}^\star, \widetilde{\rho}_{r_0,k_0}^\star)$, so that $P_0 \varphi_0 = \psi_0$ and $P_0\chi_0 = 0$, where $\chi_0$ is the influence function in Theorem~\ref{theorem::EIFsoftmax}. Then, arguing as in the proof of Theorem \ref{theorem::ALex1},
\begin{align*}
    \psi_n - \psi_0 = P_n \chi_0 + (P_n - P_0)(\varphi_n - \varphi_0) + \{P_0 \varphi_n - \psi_0\}.
\end{align*}

By Assumptions~\ref{assump:boundedness::softstar} and~\ref{assump:consistency::softstar}, 
$\|\varphi_n - \varphi_0\|_{L^2(P_0)} = o_p(1)$. Moreover, by
Assumption~\ref{assump:splitting::softstar} and Markov’s inequality,
\[
(P_n - P_0)(\varphi_n - \varphi_0)
= O_p\!\left(n^{-1/2}\,\|\varphi_n - \varphi_0\|_{L^2(P_0)}\right)
= o_p(n^{-1/2}).
\]
Thus
\[
\psi_n - \psi_0
= P_n \chi_0 + o_p(n^{-1/2}) + \{P_0 \varphi_n - \psi_0\}.
\]

It remains to control the remainder term $P_0 \varphi_n - \psi_0$.  
By construction of $\varphi_n$ and $\psi_n$, and using the notation of Lemma~\ref{lemma::vonmises_softstar}, we have
\[
P_0 \varphi_n - \psi_0
=
E_0\!\Big[
    \sum_{a \in \mathcal{A}^\star} 
        \pi_{r_n,v_n^\star}^\star(a \mid S)\,q_n^\star(a,S) 
        - 
        \sum_{a \in \mathcal{A}^\star} 
        \pi_{r_0,k_0}^\star(a \mid S)\,q_{r_0,k_0}^\star(a,S)
\Big]
\]
plus the three augmentation terms appearing in $\varphi_n$ with $(q^\star,r,v^\star,\rho^\star,\widetilde{\rho}^\star)$ replaced by $(q_n^\star,r_n,v_n^\star,\rho_n^\star,\widetilde{\rho}_n^\star)$.  
Applying Lemma~\ref{lemma::vonmises_softstar} with
\[
q^\star = q_n^\star,\quad
r = r_n,\quad
v^\star = v_n^\star,\quad
\rho^\star = \rho_n^\star,\quad
\widetilde{\rho}^\star = \widetilde{\rho}_n^\star,
\]
we obtain
\begin{align*}
\{P_0 \varphi_n - \psi_0\}
&= O\Big(
     \|\rho_n^\star - \rho_{r_0,k_0}^\star\|_{\mathcal{H},r_0}\,
     \bigl\|\mathcal{T}_{r_0,k_0}^\star(q_n^\star - q_{r_0,k_0}^\star)\bigr\|_{\mathcal{H},r_0} \\
&\qquad\qquad
     +\bigl(\|r_n-r_0\|_{\mathcal{H},r_0}
            +\|v_n^\star - v_{r_0,k_0}^\star\|_{\mathcal{H},r_0}\bigr)\,
      \bigl\|\mathcal{T}_{r_0,k_0}^\star(q_n^\star - q_{r_0,k_0}^\star)\bigr\|_{\mathcal{H},r_0} \\
&\qquad\qquad
     +\|r_n-r_0\|_{\mathcal{H},r_0}^{2}
     +\|v_n^\star - v_{r_0,k_0}^\star\|_{\mathcal{H},r_0}^{2} \\
&\qquad\qquad
     +\|\widetilde{\rho}_n^\star - \widetilde{\rho}_{r_0,k_0}^\star\|_{\mathcal{H},r_0}\,
      \bigl(\|r_n-r_0\|_{\mathcal{H},r_0}
            +\|v_n^\star - v_{r_0,k_0}^\star\|_{\mathcal{H},r_0}\bigr)
\Big).
\end{align*}

By Assumption~\ref{assump:rates::softstar}(c), the terms involving 
$\|\mathcal{T}_{r_0,k_0}^\star(q_n^\star - q_{r_0,k_0}^\star)\|_{\mathcal{H},r_0}$ are $o_p(n^{-1/2})$.  
By \ref{cond::rates}\ref{cond::reward} and 
Assumption~\ref{assump:rates::softstar}(a), we have
\[
\|r_n-r_0\|_{\mathcal{H},r_0}^{2} = o_p(n^{-1/2}),
\qquad
\|v_n^\star - v_{r_0,k_0}^\star\|_{\mathcal{H},r_0}^{2} = o_p(n^{-1/2}).
\]
Finally, by Assumption~\ref{assump:rates::softstar}(b),
\[
\|\widetilde{\rho}_n^\star - \widetilde{\rho}_{r_0,k_0}^\star\|_{\mathcal{H},r_0}\,
\bigl(\|r_n-r_0\|_{\mathcal{H},r_0}
      +\|v_n^\star - v_{r_0,k_0}^\star\|_{\mathcal{H},r_0}\bigr)
= o_p(n^{-1/2}).
\]
Hence $\{P_0 \varphi_n - \psi_0\} = o_p(n^{-1/2})$, and therefore
\[
\psi_n - \psi_0
= P_n \chi_0 + o_p(n^{-1/2}).
\]
Theorem~\ref{thm:asymptotic_linearity} then yields the claimed asymptotic linearity and efficiency with influence function $\chi_0$ from Theorem~\ref{theorem::EIFsoftmax}, and the result follows.
\end{proof}


\begin{proof}[Proof of Theorem \ref{theorem::ALex2}]
Note that $\psi_n = P_n \varphi_n$, where we define the estimated uncentered influence function by
\begin{align*}
\varphi_n(s',a,s)
&:= V_{n,\nu}^{\pi,\gamma'}(s) \\
&\quad + \rho_n^{\pi,\gamma'}(s)\,
        \frac{\pi(a \mid s)}{\pi_n(a \mid s)}
        \bigl\{r_{n,\nu,\gamma,g}(a,s) + \gamma' V_{n,\nu}^{\pi,\gamma'}(s') - q_{n,\nu}^{\pi,\gamma'}(a,s)\bigr\} \\
&\quad 
    + \rho_n^{\pi,\gamma'}(s)
        \left(
            \frac{\pi(a \mid s)}{\pi_n(a \mid s)} 
            - 
            \frac{\nu(a \mid s)}{\pi_n(a \mid s)}
        \right)
        \bigl\{r_n(a,s) - g(s) + \gamma\,V_n^{\nu,\gamma}(s') - q_n^{\nu,\gamma}(a,s)\bigr\} \\
&\quad 
    + \rho_n^{\pi,\gamma'}(s)
        \left(
            \frac{\pi(a \mid s)}{\pi_n(a \mid s)} 
            - 
            \frac{\nu(a \mid s)}{\pi_n(a \mid s)}
        \right).
\end{align*}
Let $\varphi_0$ denote the true uncentered influence function corresponding to the nuisances
$(r_0, q_0^{\nu,\gamma}, q_{0,\nu}^{\pi,\gamma'}, \rho_0^{\pi,\gamma'})$, so that
$P_0 \varphi_0 = \psi_0$ and $P_0 \chi_0 = 0$, where $\chi_0$ is the influence function in Theorem~\ref{theorem::EIFvaluenorm}. Write
\[
V_{0,\nu}^{\pi,\gamma'}(s) := V_{q_{0,\nu}^{\pi,\gamma'}}^\pi(s).
\]
Then
\[
\psi_n - \psi_0
= P_n \chi_0 + (P_n - P_0)(\varphi_n - \varphi_0) + \{P_0 \varphi_n - \psi_0\}.
\]
By Assumptions~\ref{assump:boundedness::policynorm} and~\ref{assump:consistency::policynorm},
$\|\varphi_n - \varphi_0\|_{L^2(P_0)} = o_p(1)$. Moreover, by
Assumption~\ref{assump:splitting::policynorm} and Markov’s inequality,
\[
(P_n - P_0)(\varphi_n - \varphi_0)
= O_p\!\left(n^{-1/2}\,\|\varphi_n - \varphi_0\|_{L^2(P_0)}\right)
= o_p(n^{-1/2}).
\]
Thus,
\[
\psi_n - \psi_0
= P_n \chi_0 + o_p(n^{-1/2}) + \{P_0 \varphi_n - \psi_0\}.
\]
It remains to control the remainder term $P_0 \varphi_n - \psi_0$.
By construction of $\varphi_n$ and $\psi_n$, and using the notation of
Lemma~\ref{lemma::vonmisespolicyvaluenormalization}, we may write
\[
P_0 \varphi_n - \psi_0
=
E_0\!\left[V_{n,\nu}^{\pi,\gamma'}(S) - V_{0,\nu}^{\pi,\gamma'}(S)\right]
\]
plus the three augmentation terms from Lemma~\ref{lemma::vonmisespolicyvaluenormalization}, with
\[
\rho = \rho_n^{\pi,\gamma'},\qquad
r = r_n,\qquad
q^{\nu,\gamma} = q_n^{\nu,\gamma},\qquad
q_{\nu}^{\pi,\gamma'} = q_{n,\nu}^{\pi,\gamma'}.
\]
Recalling that $d_{\rho_n^{\pi,\gamma'},r_n}^{\pi,\gamma'}(A,S)$ gathers
$\rho_n^{\pi,\gamma'}(S)$ and the relevant importance ratio (cf.\ Example~\ref{example::norm}), 
Lemma~\ref{lemma::vonmisespolicyvaluenormalization} yields
\begin{align*} 
\{P_0 \varphi_n - \psi_0\}
&= E_0\!\Big[(d_{\rho_n^{\pi,\gamma'},r_n}^{\pi,\gamma'}(A,S) - d_{\rho_0^{\pi,\gamma'},r_0}^{\pi,\gamma'}(A,S))\,
           \{\mathcal{T}_{k_0,\pi,\gamma'}(q_{0,\nu}^{\pi,\gamma'} - q_{n,\nu}^{\pi,\gamma'})(A,S)\}\Big] \\
&\quad + O\!\left(\|r_n - r_0\|_{\mathcal{H},r_0}\,
        \|\mathcal{T}_{k_0,\nu,\gamma}(q_n^{\nu,\gamma} - q_0^{\nu,\gamma})\|_{\mathcal{H},r_0}\right) \\
&\quad + O\!\left(\|r_n - r_0\|_{\mathcal{H},r_0}^2\right).
\end{align*}
By Cauchy–Schwarz and Assumption~\ref{assump:boundedness::policynorm},
\[
\|d_{\rho_n^{\pi,\gamma'},r_n}^{\pi,\gamma'} - d_{\rho_0^{\pi,\gamma'},r_0}^{\pi,\gamma'}\|_{\mathcal{H},r_0}
\lesssim
\|\rho_n^{\pi,\gamma'} - \rho_0^{\pi,\gamma'}\|_{L^2(\rho_0)}
+ \|r_n - r_0\|_{\mathcal{H},r_0},
\]
and therefore
\begin{align*}
\{P_0 \varphi_n - \psi_0\}
&= O\!\Big(
   \bigl\{\|\rho_n^{\pi,\gamma'} - \rho_0^{\pi,\gamma'}\|_{L^2(\rho_0)}
          + \|r_n - r_0\|_{\mathcal{H},r_0}\bigr\}
   \bigl\|\mathcal{T}_{k_0,\pi,\gamma'}(q_{n,\nu}^{\pi,\gamma'} - q_{0,\nu}^{\pi,\gamma'})\bigr\|_{\mathcal{H},r_0}
   \Big) \\
&\quad + O\!\left(\|r_n - r_0\|_{\mathcal{H},r_0}\,
        \|\mathcal{T}_{k_0,\nu,\gamma}(q_n^{\nu,\gamma} - q_0^{\nu,\gamma})\|_{\mathcal{H},r_0}\right)
        + O\!\left(\|r_n - r_0\|_{\mathcal{H},r_0}^2\right).
\end{align*}
By Assumption~\ref{assump:rates::policynorm}(a),
\[
\|\rho_n^{\pi,\gamma'} - \rho_0^{\pi,\gamma'}\|_{L^2(\rho_0)}\,
\bigl\|\mathcal{T}_{k_0,\pi,\gamma'}(q_{n,\nu}^{\pi,\gamma'} - q_{0,\nu}^{\pi,\gamma'})\bigr\|_{\mathcal{H},r_0}
= o_p(n^{-1/2}).
\]
By Assumption~\ref{assump:rates::policynorm}(b),
\[
\|r_n - r_0\|_{\mathcal{H},r_0}\,
\Bigl\{
   \|\mathcal{T}_{k_0,\nu,\gamma}(q_n^{\nu,\gamma} - q_0^{\nu,\gamma})\|_{\mathcal{H},r_0}
   + \|\mathcal{T}_{k_0,\pi,\gamma'}(q_{n,\nu}^{\pi,\gamma'} - q_{0,\nu}^{\pi,\gamma'})\|_{\mathcal{H},r_0}
\Bigr\}
= o_p(n^{-1/2}).
\]
Finally, by \ref{cond::rates}\ref{cond::reward},
\[
\|r_n - r_0\|_{\mathcal{H},r_0}^2 = o_p(n^{-1/2}).
\]
Combining these bounds yields
\[
\{P_0 \varphi_n - \psi_0\} = o_p(n^{-1/2}).
\]
Therefore,
\[
\psi_n - \psi_0
= P_n \chi_0 + o_p(n^{-1/2}),
\]
and Theorem~\ref{thm:asymptotic_linearity} implies the claimed asymptotic linearity and efficiency with influence function $\chi_0$ given in Theorem~\ref{theorem::EIFvaluenorm}. This completes the proof.
\end{proof}


 


 
 

\subsection{Structural coefficients under reward normalization}

\begin{proof}[Proof of Theorem \ref{theorem::EIFcoefnorm}]
For each coordinate \(j \in \{1,\dots,p\}\), define
\[
\tilde m_j(a,s,r,k) := x_j(a,s)\,r(a,s),
\qquad
b_j(P) := E_P[x_j(A,S)r_{P,\nu,\gamma,g}(A,S)].
\]
Since
\[
\partial_r \tilde m_j(a,s,r,k)[h] = x_j(a,s)h(a,s),
\qquad
\partial_k \tilde m_j(a,s,r,k)[\beta] = 0,
\]
the lower-level Riesz representers are
\[
\tilde \alpha_{0,j}(a,s) = x_j(a,s),
\qquad
\tilde \beta_{0,j} = 0.
\]
Applying Theorem~\ref{theorem::EIFnorm} coordinatewise therefore yields
\[
\alpha_{0,j}^x(a,s)
=
x_j(a,s)
- \frac{\nu(a\mid s)}{\pi_0(a\mid s)}
\sum_{\tilde a \in \mathcal A}\pi_0(\tilde a\mid s)x_j(\tilde a,s)
\]
and
\[
\frac{\beta_{0,j}^x}{k_0}(s'\mid a,s)
=
\alpha_{0,j}^x(a,s)
\bigl\{r_0(a,s)-g(s)+\gamma V_{r_0-g,k_0}^{\nu,\gamma}(s')
-q_{r_0-g,k_0}^{\nu,\gamma}(a,s)\bigr\}.
\]
Moreover,
\[
\sum_{\tilde a \in \mathcal A}\pi_0(\tilde a\mid s)\,\alpha_{0,j}^x(\tilde a,s)=0
\]
for every \(j\), so the generic EIF from Theorem~\ref{theorem::EIF} reduces to
\[
\begin{aligned}
\chi_{b_j,0}(s,a,s')
&=
\alpha_{0,j}^x(a,s) \\
&\quad
+ \alpha_{0,j}^x(a,s)\bigl\{r_0(a,s)-g(s)+\gamma V_{r_0-g,k_0}^{\nu,\gamma}(s')
-q_{r_0-g,k_0}^{\nu,\gamma}(a,s)\bigr\} \\
&\quad
+ x_j(a,s)\bar r_0(a,s)-b_{0,j}.
\end{aligned}
\]
Stacking the coordinates gives the stated vector influence function
\(\chi_{b,0}\) for \(b(P)\).

Since \(\theta(P)=G(P)^{-1}b(P)\) with \(G(P)=E_P[x(A,S)x(A,S)^\top]\), the
matrix delta method gives, for perturbations \((\Delta G,\Delta b)\),
\[
\partial \theta_{(G_0,b_0)}[\Delta G,\Delta b]
=
G_0^{-1}\Delta b - G_0^{-1}\Delta G\,\theta_0.
\]
The influence function of \(G(P)\) is \(x(A,S)x(A,S)^\top - G_0\). Hence
\[
\chi_{\theta,0}(s,a,s')
=
G_0^{-1}\Big[
\chi_{b,0}(s,a,s') - \{x(a,s)x(a,s)^\top - G_0\}\theta_0
\Big],
\]
which is the desired EIF.
\end{proof}


\begin{lemma}[von Mises expansion for structural coefficient moments]
\label{lemma::vonmisescoefnorm}
Let \(r,q \in L^\infty(\lambda)\), write
\[
q_0^{\nu,\gamma} := q_{r_0-g,k_0}^{\nu,\gamma},
\qquad
\bar r_q := g + q - V_q^\nu,
\]
and define
\[
\alpha_r^x(a,s)
:=
x(a,s)
- \frac{\nu(a\mid s)}{\pi_r(a\mid s)}
\sum_{\tilde a \in \mathcal A}\pi_r(\tilde a\mid s)x(\tilde a,s),
\qquad
\pi_r(a\mid s):=\exp\{r(a,s)\}.
\]
Then
\[
\begin{aligned}
\Big\|
&E_0\!\big[x(A,S)\bar r_q(A,S)\big] - b_0
+ E_0\!\big[\alpha_r^x(A,S)\big] \\
&\qquad
+ E_0\!\big[\alpha_r^x(A,S)\{r(A,S)-g(S)+\gamma V_q^\nu(S')-q(A,S)\}\big]
\Big\|_2
\end{aligned}
\]
\[
\lesssim
\|r-r_0\|_{\mathcal H,r_0}\,
\bigl\|\mathcal T_{k_0,\nu,\gamma}(q-q_0^{\nu,\gamma})\bigr\|_{\mathcal H,r_0}
+ \|r-r_0\|_{\mathcal H,0}^2,
\]
where the implicit constant depends only on \(\|x\|_{L^\infty(\lambda)}\),
\(\|r\|_{L^\infty(\lambda)}\), and \(\|r_0\|_{L^\infty(\lambda)}\).
\end{lemma}
\begin{proof}
Since
\[
\bar r_q - \bar r_0
=
(I-\nu)(q-q_0^{\nu,\gamma}),
\]
the same identity used in the proof of
Lemma~\ref{lemma::vonmisespolicyvaluenormalization} gives
\[
E_0\!\big[x(A,S)\{\bar r_q - \bar r_0\}(A,S)\big]
=
E_0\!\big[x(A,S)(I-\nu)\,
\mathcal T_{k_0,\nu,\gamma}(q-q_0^{\nu,\gamma})(A,S)\big].
\]
Define
\[
\delta_q
:=
\mathcal T_{k_0,\nu,\gamma}(q-q_0^{\nu,\gamma})
=
q - \gamma V_q^\nu - r_0 + g.
\]
By the definition of \(\alpha_0^x\),
\[
E_0\!\big[x(A,S)(I-\nu)\delta_q(A,S)\big]
=
E_0\!\big[\alpha_0^x(A,S)\delta_q(A,S)\big].
\]
Therefore,
\[
E_0\!\big[x(A,S)\{\bar r_q - \bar r_0\}(A,S)\big]
=
E_0\!\big[\alpha_r^x(A,S)\delta_q(A,S)\big]
+ E_0\!\big[(\alpha_0^x-\alpha_r^x)(A,S)\delta_q(A,S)\big].
\]
Since
\[
\delta_q
=
\{q - \gamma V_q^\nu - r + g\} + (r-r_0),
\]
we obtain
\[
\begin{aligned}
E_0\!\big[x(A,S)\{\bar r_q - \bar r_0\}(A,S)\big]
&=
E_0\!\big[\alpha_r^x(A,S)\{q(A,S)-\gamma V_q^\nu(S')-r(A,S)+g(S)\}\big] \\
&\quad
+ E_0\!\big[\alpha_r^x(A,S)\{r-r_0\}(A,S)\big] \\
&\quad
+ E_0\!\big[(\alpha_0^x-\alpha_r^x)(A,S)\delta_q(A,S)\big].
\end{aligned}
\]
By boundedness of \(x\) and the first-order expansion of \(\pi_r\) around
\(\pi_0\),
\[
\|\alpha_r^x - \alpha_0^x\|_{\mathcal H,r_0}
\lesssim
\|r-r_0\|_{\mathcal H,r_0},
\]
so Cauchy--Schwarz gives
\[
E_0\!\big[(\alpha_0^x-\alpha_r^x)(A,S)\delta_q(A,S)\big]
=
O\!\left(
\|r-r_0\|_{\mathcal H,r_0}\,
\bigl\|\mathcal T_{k_0,\nu,\gamma}(q-q_0^{\nu,\gamma})\bigr\|_{\mathcal H,r_0}
\right).
\]
It remains to control the reward term. By the second-order expansion of
\(\exp\{r-r_0\}\) around zero,
\[
\{r-r_0\}(A,S)
=
\frac{\pi_r}{\pi_0}(A\mid S)-1
+ O\!\left(\|r-r_0\|_{\mathcal H,0}^2\right),
\]
and therefore
\[
\begin{aligned}
E_0\!\big[\alpha_r^x(A,S)\{r-r_0\}(A,S)\big]
+ E_0\!\big[\alpha_r^x(A,S)\big]
&=
E_0\!\left[\alpha_r^x(A,S)\frac{\pi_r}{\pi_0}(A\mid S)\right]
+ O\!\left(\|r-r_0\|_{\mathcal H,0}^2\right) \\
&=
E_0\!\left[\sum_{a \in \mathcal A}\pi_r(a\mid S)\alpha_r^x(a,S)\right]
+ O\!\left(\|r-r_0\|_{\mathcal H,0}^2\right) \\
&=
O\!\left(\|r-r_0\|_{\mathcal H,0}^2\right),
\end{aligned}
\]
because \(\sum_{a \in \mathcal A}\pi_r(a\mid s)\alpha_r^x(a,s)=0\) by
construction. Combining the displays above yields the claim.
\end{proof}


\begin{proof}[Proof of Theorem \ref{theorem::ALex3}]
Define the estimated uncentered influence function for \(b_0\) by
\[
\begin{aligned}
\varphi_n^b(s',a,s)
:=
&x(a,s)\bar r_n(a,s)
+ \alpha_n^x(a,s) \\
&\quad
+ \alpha_n^x(a,s)\{r_n(a,s)-g(s)+\gamma V_n^{\nu,\gamma}(s')-q_n^{\nu,\gamma}(a,s)\},
\end{aligned}
\]
and let
\[
\varphi_0^b(s',a,s) := \chi_{b,0}(s,a,s') + b_0.
\]
Then \(b_n = P_n\varphi_n^b\), \(P_0\varphi_0^b = b_0\), and
\[
b_n - b_0
=
P_n\chi_{b,0}
+ (P_n-P_0)(\varphi_n^b-\varphi_0^b)
+ \{P_0\varphi_n^b-b_0\}.
\]
By Assumptions~\ref{assump:boundedness::coefnorm} and
\ref{assump:consistency::coefnorm},
\[
\|\varphi_n^b-\varphi_0^b\|_{L^2(P_0)} = o_p(1).
\]
Here \(\|f\|_{L^2(P_0)} := \{E_0[\|f(S',A,S)\|_2^2]\}^{1/2}\) for
vector-valued \(f\).
Since \(p\) is fixed, Assumption~\ref{assump:splitting::coefnorm} and
Markov's inequality imply
\[
\big\|(P_n-P_0)(\varphi_n^b-\varphi_0^b)\big\|_2
=
O_p\!\left(n^{-1/2}\,\|\varphi_n^b-\varphi_0^b\|_{L^2(P_0)}\right)
=
o_p(n^{-1/2}).
\]
Applying Lemma~\ref{lemma::vonmisescoefnorm} with \(r=r_n\) and \(q=q_n^{\nu,\gamma}\)
gives
\[
\|P_0\varphi_n^b-b_0\|_2
=
O_p\!\left(
\|r_n-r_0\|_{\mathcal H,r_0}\,
\bigl\|\mathcal T_{k_0,\nu,\gamma}(q_n^{\nu,\gamma}-q_0^{\nu,\gamma})\bigr\|_{\mathcal H,r_0}
+ \|r_n-r_0\|_{\mathcal H,0}^2
\right).
\]
By Assumption~\ref{assump:rates::coefnorm},
\[
\|P_0\varphi_n^b-b_0\|_2 = o_p(n^{-1/2}).
\]
Hence
\[
b_n-b_0 = P_n\chi_{b,0} + o_p(n^{-1/2}).
\]

Next, write
\[
G_n-G_0
=
\frac1n\sum_{i=1}^n \{x(A_i,S_i)x(A_i,S_i)^\top-G_0\}.
\]
Since each coordinate of \(x\) is bounded, \(G_n-G_0=O_p(n^{-1/2})\) entrywise.
Because \(\lambda_{\min}(G_0)>0\), it follows that \(G_n\) is invertible with
probability tending to one and
\[
G_n^{-1}
=
G_0^{-1} - G_0^{-1}(G_n-G_0)G_0^{-1} + o_p(n^{-1/2}).
\]
Therefore,
\[
\begin{aligned}
\theta_n-\theta_0
&=
G_n^{-1}b_n - G_0^{-1}b_0 \\
&=
G_0^{-1}(b_n-b_0) - G_0^{-1}(G_n-G_0)\theta_0 + o_p(n^{-1/2}) \\
&=
\frac1n\sum_{i=1}^n
G_0^{-1}\Big[
\chi_{b,0}(S_i,A_i,S_i')
- \{x(A_i,S_i)x(A_i,S_i)^\top-G_0\}\theta_0
\Big]
+ o_p(n^{-1/2}) \\
&=
\frac1n\sum_{i=1}^n \chi_{\theta,0}(S_i,A_i,S_i') + o_p(n^{-1/2}).
\end{aligned}
\]
The multivariate central limit theorem then gives
\[
\sqrt n(\theta_n-\theta_0)\rightsquigarrow N(0,\Sigma_0),
\qquad
\Sigma_0 = \operatorname{Var}_{P_0}\!\big(\chi_{\theta,0}(S,A,S')\big).
\]
This establishes asymptotic linearity and efficiency.
\end{proof}


\section{Smoothness of the Soft Bellman Equation and Its Solutions}

In this section, we establish several technical lemmas on the smoothness of the soft Bellman equation and its solutions. Some of these results are used in deriving the EIF and the von Mises expansion for Example~\ref{example::1b}. Other results are included for completeness. 


\subsection{Notation}

We study the functional smoothness of the soft Bellman equation and its fixed-point solution for a fixed temperature parameter $\tau \in (0,\infty)$. Let $\lambda := \# \otimes \mu$ denote the product measure on $\mathcal{A}\times\mathcal{S}$, where $\#$ is the uniform measure on the finite set $\mathcal{A}$ and $\mu$ is a finite dominating measure on $\mathcal{S}$. For any measurable $B \subseteq \mathcal{A}\times\mathcal{S}$,
\[
\lambda(B) = \tfrac{1}{|\mathcal{A}|}\sum_{a \in \mathcal{A}} \int 1_B(a,s)\,\mu(ds).
\]
For a linear operator $B$, define $\|B\|_{L^\infty \to L^2} := \sup_{\|f\|_{L^\infty(\lambda)} \le 1}\|Bf\|_{L^2(\lambda)}$ and more generally $\|B\|_{L^p \to L^q} := \sup_{\|f\|_{L^p(\lambda)} \le 1}\|Bf\|_{L^q(\lambda)}$.

For $r \in \mathcal{H}$ and $k \in \mathcal{K}$, let $v_{r,k}$ denote the unique solution in $L^\infty(\lambda)$ of the soft Bellman equation
\[
v_{r,k}(a,s) = \int \tau \log\!\left(\sum_{a'\in\mathcal{A}}\exp\!\left(\frac{r(a',s')+\gamma v_{r,k}(a',s')}{\tau}\right)\right)k(s'\mid a,s)\,\mu(ds').
\]
Define $F: L^\infty(\lambda)\times L^\infty(\lambda)\times\mathcal{K}\to L^\infty(\lambda)$ by
\[
F(r,v,k)(a,s) := v(a,s) - \int \tau \log\!\left(\sum_{a'\in\mathcal{A}}\exp\!\left(\frac{r(a',s')+\gamma v(a',s')}{\tau}\right)\right)k(s'\mid a,s)\,\mu(ds'),
\]
and for each $k \in \mathcal{K}$, let $F_k(r,v):=F(r,v,k)$.  

We write
\[
\operatorname{logsumexp}_{\tau}(q)(s) := \tau \log\!\left( \sum_{a' \in \mathcal{A}} \exp\!\left\{ \frac{q(a', s)}{\tau} \right\} \right),
\]
\[
\pi_{r,k}(a'\mid s):=\frac{\exp\!\left(\frac{r(a',s)+\gamma v_{r,k}(a',s)}{\tau}\right)}{\sum_{b\in\mathcal{A}}\exp\!\left(\frac{r(b,s)+\gamma v_{r,k}(b,s)}{\tau}\right)},
\qquad
q_{r,k}(a,s):=r(a,s)+\gamma v_{r,k}(a,s),
\]
and for general $q$ let
\[
\widetilde{\pi}_{q,\tau}(a'\mid s):=\frac{\exp\!\left(\frac{q(a',s)}{\tau}\right)}{\sum_{b\in\mathcal{A}}\exp\!\left(\frac{q(b,s)}{\tau}\right)}.
\]
For fixed $s'$, the conditional covariance under $\pi$ is
\[
\begin{aligned}
\operatorname{Cov}_{\pi(\cdot\mid s')}(g(\cdot,s'),h(\cdot,s'))
&:= \sum_{a\in\mathcal{A}}\pi(a\mid s')g(a,s')h(a,s') \\
&\quad - \Big(\sum_{a\in\mathcal{A}}\pi(a\mid s')g(a,s')\Big)
\Big(\sum_{a\in\mathcal{A}}\pi(a\mid s')h(a,s')\Big).
\end{aligned}
\]
Throughout this subsection, the temperature $\tau$ is fixed and all implicit
constants may depend on $\tau$.

 



\subsection{Technical lemmas}

\label{appendix::technicalsoftmax}
 
 

In this subsection, additive perturbations of the kernel are indexed by
elements of \(\addTK\), equipped with the norm \(\|\cdot\|_{\mathcal K}\),
and their additive Fr\'echet derivative is denoted by \(\adk\). The score
derivative \(\partial_k\) is reserved for the semiparametric statements above.


 



We begin with the following lemma, which ensures that the map $(r,k) \mapsto v_{r,k}$ is well defined.


 


\begin{lemma}[Existence and uniqueness of solution to soft Bellman equation]
    For each \(r \in L^\infty(\lambda)\) and $k \in \mathcal{K}$, there exists a unique solution \(v_{r,k} \in L^\infty(\lambda)\) to the soft Bellman equation $F_{k}(r, v) = 0$ over \(v \in L^\infty(\lambda)\).
\end{lemma}
\begin{proof}
For every \(r \in L^\infty(\lambda)\), the map \(v \mapsto v - F_{k}(r,v)\) is a contraction on \(L^\infty(\lambda)\). By Banach's fixed-point theorem, there exists a unique \(v_{r,k} \in L^\infty(\lambda)\) such that \(v_{r,k} = v_{r,k} - F_{k}(r,v_{r,k})\), and hence \(F_{k}(r,v_{r,k}) = 0\).
\end{proof}


The following lemma establishes Fréchet differentiability of $F_{k}$ for a fixed $k$.

 


\begin{lemma}[Fréchet differentiability of $F_k$]
\label{lemma::softbellmanissmooth}
Let $k \in \mathcal{K}$. The map $F_{k}$ is Fréchet differentiable in supremum norm, with partial derivatives at \((r,v) \in L^\infty(\lambda) \times L^\infty(\lambda)\), for directions $u,h \in L^\infty(\lambda)$, given by
\begin{align*}
\partial_v F_{k}(r,v)[h](a,s)
&= h(a,s) - \gamma \int \sum_{a' \in \mathcal{A}} \widetilde{\pi}_{\,r+\gamma v,\tau}(a' \mid s')\, h(a',s')\, k(s' \mid a,s) \mu(ds')\\
\partial_r F_{k}(r,v)[u](a,s)
&= - \int \sum_{a' \in \mathcal{A}} \widetilde{\pi}_{\,r+\gamma v,\tau}(a' \mid s')\, u(a',s')\, k(s' \mid a,s) \mu(ds').
\end{align*}
In particular, there exists a fixed constant \(C_\tau < \infty\), depending only on the fixed temperature $\tau$, such that, uniformly over $r,v,h,u$, 
\[
\|\,F_{k}(r+u,v+h) - F_{k}(r,v) - \partial_r F_{k}(r,v)[u] - \partial_v F_{k}(r,v)[h]\,\|_{L^\infty(\lambda)}
\;\le\; C_\tau \,\|u + \gamma h\|_{L^2(\lambda)}^{2}.
\]
\end{lemma}
\begin{proof}
It is straightforward to show that \(F_{k}\) is Gâteaux differentiable with partial derivatives given in the theorem statement. We now show that the map is Fréchet differentiable. 
 

By Taylor's theorem, we have, for each direction \(f \in L^\infty(\lambda)\), that
\begin{align*}
\big|\operatorname{logsumexp}_{\tau}(q+f)(s')-\operatorname{logsumexp}_{\tau}(q)(s')
-\sum_{a'\in\mathcal A}\widetilde{\pi}_{q,\tau}(a'\mid s')\,f(a',s')\big|
&\le \frac{1}{2\tau} \int_0^1 \sum_{a'\in\mathcal A}\widetilde{\pi}_{\,q+t f,\tau}(a'\mid s')\,|f(a',s')|^2\,dt\\
&\le \frac{1}{2\tau} \sum_{a'\in\mathcal A}|f(a',s')|^2.
\end{align*}
Substituting \(q:=r+\gamma v\) and \(f:=u+\gamma h\) in the above, we find that
\begin{align*}
&\big|F_{k}(r+u,v+h)(a,s)-F_{k}(r,v)(a,s)
-\partial_v F_{k}(r,v)[h](a,s)-\partial_r F_{k}(r,v)[u](a,s)\big|\\
&\qquad\le \frac{1}{2\tau}\int_0^1 \int \sum_{a'\in\mathcal A}
\widetilde{\pi}_{\,q+t f,\tau}(a'\mid s')\,|f(a',s')|^2 \, k(s'\mid a,s) \mu(ds')\,dt\\
&\qquad\le \frac{1}{2\tau}\int \sum_{a'\in\mathcal A} |u(a',s')+\gamma h(a',s')|^2 \, k(s'\mid a,s) \mu(ds').
\end{align*}
By definition of $k \in \mathcal{K}$, we have that $k(s'\mid a,s) \mu(ds') \lesssim \mu(ds')$. Hence,
\begin{align*}
&\big|F_{k}(r+u,v+h)(a,s)-F_{k}(r,v)(a,s)
-\partial_v F_{k}(r,v)[h](a,s)-\partial_r F_{k}(r,v)[u](a,s)\big|\\
&\qquad\lesssim  \int \sum_{a'\in\mathcal A} |u(a',s')+\gamma h(a',s')|^2 \, \mu(ds')\\
&\qquad\lesssim \| u + \gamma h\|_{L^2(\lambda)}^2.
\end{align*}
Taking the \(L^\infty\)-norm of both sides, we find
\begin{align*}
    \|F_{k}(r+u,v+h)-F_{k}(r,v)
-\partial_v F_{k}(r,v)[h] -\partial_r F_{k}(r,v)[u] \|_{L^\infty(\lambda)} &\lesssim \|u + \gamma h\|_{L^2(\lambda)}^2.
\end{align*} 

\end{proof}

 
 

\begin{lemma}[Fréchet differentiability of $F$]
\label{lemma::bellmansmoothnessfull}
The map $F: L^\infty(\lambda) \times L^\infty(\lambda) \times \mathcal{K} \to L^\infty(\lambda)$ is Fr\'echet differentiable in supremum norm, with partial derivatives at $(r,v,k)$ given by
\begin{align*}
\adk F(r,v,k)[k' - k](a,s)
&= - \int \operatorname{logsumexp}_{\tau}(r+\gamma v)(s')\, (k'-k)(s'\mid a,s)\, \mu(ds'),\\
\partial_v F(r,v,k)[h](a,s) &= \partial_v F_{k}(r,v)[h](a,s),\\
\partial_r F(r,v,k)[u](a,s) &= \partial_r F_{k}(r,v)[u](a,s),
\end{align*}
for $k' \in \mathcal{K}$ and $h,u \in L^\infty(\lambda)$. In particular, we have
\[
F(r + u, v+ h, k') - F(r , v, k)
- \partial_r F(r,v,k)[u] - \partial_v F(r,v,k)[h] - \adk F(r,v,k)[k' - k]
= R_2(u) + R_2(h) + R_{\infty},
\]
where 
\begin{align*}
\|R_2(u)\|_{L^2(\lambda)} &\lesssim \|u\|_{L^2(\lambda)} \, \|k - k'\|_{L^2(\mu \otimes \lambda)}, 
& \|R_2(u)\|_{L^\infty(\lambda)} &\lesssim \|u\|_{L^2(\lambda)} \, \|k - k'\|_{L^\infty(\mu \otimes \lambda)}, \\[4pt]
\|R_2(h)\|_{L^2(\lambda)} &\lesssim \|h\|_{L^2(\lambda)} \, \|k - k'\|_{L^2(\mu \otimes \lambda)}, 
& \|R_2(h)\|_{L^\infty(\lambda)} &\lesssim \|h\|_{L^2(\lambda)} \, \|k - k'\|_{L^\infty(\mu \otimes \lambda)}, \\[4pt]
\|R_{\infty}\|_{L^\infty(\lambda)} &\lesssim \|u + \gamma h\|_{L^2(\lambda)}^2. &&
\end{align*}

\end{lemma}
\begin{proof}
Adding and subtracting the term $ F(r , v, k') $, we have
\begin{align*}
    F(r + u, v+ h, k') - F(r , v, k) &=  F(r + u, v+ h, k') - F(r , v, k')  +  F(r , v, k')  -  F(r , v, k) .
\end{align*}
By Lemma \ref{lemma::softbellmanissmooth},
\[
F(r + u, v+ h, k') - F(r , v, k')
= \partial_v F_{k'}(r,v)[h] + \partial_r F_{k'}(r,v)[u]
+ R_{\infty}(\|u + \gamma h\|_{L^2(\lambda)}^2),
\]
where $\|R_{\infty}(\|u + \gamma h\|_{L^2(\lambda)}^2)\|_{L^\infty(\lambda)} \lesssim \|u + \gamma h\|_{L^2(\lambda)}^2$. By linearity of $k \mapsto F(r,v, k)$, we have the equality:
\[
F(r , v, k') - F(r , v, k) = \adk F(r,v, k)[k' - k].
\]
Hence,
\begin{align}
    F(r + u, v+ h, k') - F(r , v, k)
    &= \partial_v F_{k'}(r,v)[h] + \partial_r F_{k'}(r,v)[u] \nonumber \\
    &\quad + \adk F(r,v, k)[k' - k]
    + R_{\infty}(\|u + \gamma h\|_{L^2(\lambda)}^2) \nonumber \\
     &= \partial_v F_{k}(r,v)[h] + \partial_r F_{k}(r,v)[u]
     + \adk F(r,v, k)[k' - k]  \nonumber \\
     & \quad + \partial_v F_{k'}(r,v)[h]  - \partial_v F_{k}(r,v)[h] \nonumber \\
     & \quad + \partial_r F_{k'}(r,v)[u] - \partial_r F_{k}(r,v)[u] \nonumber \\
     & \quad + R_{\infty}(\|u + \gamma h\|_{L^2(\lambda)}^2). \label{eqnproof:addsub1}
\end{align}

 



We now consider the differences $\partial_v F_{k'}(r,v)[h]  - \partial_v F_{k}(r,v)[h] $ and $\partial_r F_{k'}(r,v)[u] - \partial_r F_{k}(r,v)[u]$. By the expressions of the partial derivatives in Lemma \ref{lemma::softbellmanissmooth}, we have that
\begin{align*}
  \left|  \partial_r F_{k'}(r,v)[u] - \partial_r F_{k}(r,v)[u]\right|(a,s) &\leq \left| \int  \sum_{a' \in \mathcal{A}} u(a', s') (k' - k)(s' \mid a, s) \mu(ds') \right|.
\end{align*}
For each $(a,s)$,
\begin{align*}
\left| \int  \sum_{a' \in \mathcal{A}} u(a', s') \,(k' - k)(s' \mid a, s)\, \mu(ds') \right|
&= \left| \int \sum_{a' \in \mathcal{A}} u(a',s')\,(\sqrt{k'}-\sqrt{k})(s'\mid a,s)\,(\sqrt{k'}+\sqrt{k})(s'\mid a,s)\,\mu(ds') \right| \\
&\le \left( \int \Big(\sum_{a' \in \mathcal{A}} u(a',s')\Big)^2 (\sqrt{k'}+\sqrt{k})^2 \, \mu(ds') \right)^{1/2} \\
&\quad \times \left( \int (\sqrt{k'}-\sqrt{k})^2(s' \mid a, s) \, \mu(ds') \right)^{1/2}\\
&\lesssim \|u\|_{L^2(\lambda)}  \times \left( \int (\sqrt{k'}-\sqrt{k})^2(s' \mid a, s) \, \mu(ds') \right)^{1/2},
\end{align*}
where we used that $\sqrt{k'} + \sqrt{k}$ is uniformly bounded by the definition of $\mathcal{K}$. Taking the $L^2(\lambda)$ norm over $(a,s)$, we have
\begin{align*}
  \left\|  \partial_r F_{k'}(r,v)[u] - \partial_r F_{k}(r,v)[u]\right\|_{L^2(\lambda)} &\lesssim  \|u\|_{L^2(\lambda)} \left\{  \int (\sqrt{k'}-\sqrt{k})^2(s' \mid a, s) \, \mu(ds')\, \lambda(da, ds) \right\}^{1/2}\\
  &\lesssim  \|u\|_{L^2(\lambda)}  \,\|k' - k\|_{L^2(\lambda \otimes \mu)},
\end{align*}
where we used that $(\sqrt{k'}-\sqrt{k})^2(s' \mid a, s) \;\lesssim\; (k'-k)^2(s' \mid a, s)$, due to the lower boundedness (positivity) of the elements of $\mathcal{K}$. 
Taking the supremum over $(a,s)$ instead, the same pointwise bound yields
\begin{align*}
  \left\|  \partial_r F_{k'}(r,v)[u] - \partial_r F_{k}(r,v)[u]\right\|_{L^\infty(\lambda)}
  &\lesssim  \|u\|_{L^2(\lambda)}  \,\|k' - k\|_{L^\infty(\lambda \otimes \mu)}.
\end{align*}


For the other partial derivative, an identical argument gives
\begin{align*}
  \left\|  \partial_v F_{k'}(r,v)[h] - \partial_v F_{k}(r,v)[h]\right\|_{L^2(\lambda)}  
  &\lesssim   \|h\|_{L^2(\lambda)} \, \|k' - k\|_{L^2(\lambda \otimes \mu)}.
\end{align*}
and
\begin{align*}
  \left\|  \partial_v F_{k'}(r,v)[h] - \partial_v F_{k}(r,v)[h]\right\|_{L^\infty(\lambda)}  
  &\lesssim   \|h\|_{L^2(\lambda)} \, \|k' - k\|_{L^\infty(\lambda \otimes \mu)}.
\end{align*}



 Returning to \eqref{eqnproof:addsub1}, we have that 
\begin{align*}
    F(r + u, v+ h, k') - F(r , v, k)    &= \partial_v F_{k}(r,v)[h] + \partial_r F_{k}(r,v)[u]   + \adk F(r,v, k)[k' - k] \\
     & \quad + R_2(u) + R_2(h) + R_{\infty}(\|u + \gamma h\|_{L^2(\lambda)}^2),
\end{align*}
where $\|R_2(u)\|_{L^2(\lambda)} \lesssim \|u\|_{L^2(\lambda)}\|k - k'\|_{L^2(\mu \otimes \lambda)} $, $\|R_2(u)\|_{L^\infty(\lambda)} \lesssim \|u\|_{L^2(\lambda)}\|k - k'\|_{L^\infty(\mu \otimes \lambda)} $, $\|R_2(h)\|_{L^2(\lambda)} \lesssim \|h\|_{L^2(\lambda)}\|k - k'\|_{L^2(\mu \otimes \lambda)} $, $\|R_2(h)\|_{L^\infty(\lambda)} \lesssim \|h\|_{L^2(\lambda)}\|k - k'\|_{L^\infty(\mu \otimes \lambda)} $, and $\|R_{\infty}(\|u + \gamma h\|_{L^2(\lambda)}^2)\|_{L^\infty(\lambda)} \lesssim \|u + \gamma h\|_{L^2(\lambda)}^2$. The result follows.

\end{proof}



 
\begin{lemma}[Sup-norm Fréchet differentiability of the soft value function]
\label{lemma::IFT}
Let $(r,k) \in \mathcal{H} \times \mathcal{K}$. Suppose $\gamma < 1$, so that $\mathcal{T}_{r,k}: L^\infty(\lambda) \to L^\infty(\lambda)$ is invertible. Then the map \((r,k) \mapsto v_{r,k}\) is Fréchet differentiable, with partial derivatives
\[
\partial_r v_{r,k} \;=\; \mathcal T_{r,k}^{-1}\, \mathcal P_{r,k},
\qquad
\adk v_{r,k}[w] \;=\; \mathcal T_{r,k}^{-1}\, \mathcal{P}_{w}\,\Phi(r + \gamma v),
\]
where $\Phi(r + \gamma v)(s') := \operatorname{logsumexp}_{\tau}(r + \gamma v)(s')$ and
\[
\mathcal{P}_{w}\,\Phi(r + \gamma v)(a,s)
:= \int \operatorname{logsumexp}_{\tau}(r + \gamma v)(s')\, w(s'\mid a,s)\,\mu(ds').
\]
\end{lemma}
\begin{proof}
Denote
$$\mathcal Q_{r,k}(w):= \int \operatorname{logsumexp}_{\tau}(r + \gamma v_{r,k})(s')\, w(s'\mid a,s)\,\mu(ds').$$

We apply the implicit function theorem to obtain joint differentiability of $v_{r,k}$ in $(r,k)$.  
Moreover, since $\gamma <  1$, the map $F$ is Fr\'echet differentiable as a map 
$L^\infty(\lambda)\times L^\infty(\lambda)\times \mathcal K \to L^\infty(\lambda)$, and its derivative depends jointly and continuously on $(r,v,k)$. By the invertibility assumption in the statement, the partial derivative
\[
\partial_v F_{k}(r,v_{r,k}): L^\infty(\lambda) \to L^\infty(\lambda)
\]
is a bounded linear isomorphism.  
The implicit function theorem for Banach spaces therefore implies that the solution map 
\((r,k) \mapsto v_{r,k}\) is Fr\'echet differentiable at \((r,k)\) as a map 
\(L^\infty(\lambda)\times \mathcal K \to L^\infty(\lambda)\), with derivatives
\begin{align*}
    \partial_r v_{r,k}[u]  
&= -\big(\partial_v F_{k}(r,v_{r,k})\big)^{-1}\!\big(\partial_r F_{k}(r,v_{r,k})[u]\big),\\
\adk v_{r,k}[w]  
&= -\big(\partial_v F_{k}(r,v_{r,k})\big)^{-1}\!\big(\adk F(r,v_{r,k},k)[w]\big).
\end{align*}
Since $\partial_v F_{k}(r,v)=\mathcal T_{r,k}$, 
$\partial_r F_{k}(r,v)=-\mathcal P_{r,k}$, and 
$\adk F(r,v,k)=-\mathcal Q_{r,k}$,
 we find that
 \[
\partial_r v_{r,k} \;=\; \mathcal T_{r,k}^{-1}\, \mathcal P_{r,k},
\qquad
\adk v_{r,k} \;=\; \mathcal T_{r,k}^{-1}\, \mathcal Q_{r,k}.
\]


 
 

 
\end{proof}

\begin{lemma}
\label{lemma::qderiv}
Under the conditions of Lemma \ref{lemma::IFT}, it holds that
\[
\partial_r q_{r,k}[u] = \mathcal{T}_{r,k}^{-1}u
\]
for every $u \in L^\infty(\lambda)$.
\end{lemma}




\begin{proof}
By Lemma \ref{lemma::IFT},
\begin{align*}
    \partial_r q_{r,k}[u] 
    &= u + \gamma \,\partial_r v_{r,k}(u) \\
    &= \left(I + \gamma \,\mathcal{T}_{r,k}^{-1}\mathcal{P}_{r,k}\right)(u) \\
    &= \mathcal{T}_{r,k}^{-1}\left(\mathcal{T}_{r,k} + \gamma \mathcal{P}_{r,k}\right)(u) \\
    &= \mathcal{T}_{r,k}^{-1}\left(I - \gamma \mathcal{P}_{r,k} + \gamma \mathcal{P}_{r,k}\right)(u) \\
    &= \mathcal{T}_{r,k}^{-1} u.
\end{align*}
\end{proof}


In the following lemma, we define the hard Bellman operator $\widetilde{\mathcal{T}}_k: L^\infty(\lambda) \to L^\infty(\lambda)$ by
\[
\widetilde{\mathcal{T}}_k(f)(a,s) 
:= f(a,s) - \gamma \int \max_{a' \in \mathcal{A}} f(a',s')\,k(s'\mid a,s)\,\mu(ds').
\] 


 

 
 


\begin{lemma}[Lipschitz continuity of value function map]
\label{lemma::lipschitzvalue}
Suppose that $\gamma < 1$. Then, for all $r,r' \in L^\infty(\lambda)$ and $k,k' \in \mathcal{K}$, we have  \begin{align*}
 \|v_{r',k}-v_{r,k}\|_{L^\infty(\lambda)} &\lesssim  (1-\gamma)^{-1}   \|r' - r\|_{L^2(\lambda)}\\
 \| \widetilde{\mathcal{T}}_{k} (v_{r,k'} - v_{r,k} ) \|_{L^2(\lambda)}  &\lesssim  ( 1 + \|r  \|_{L^\infty(\lambda)}) \cdot \|k' - k\|_{L^2(\mu \otimes \lambda)}.
\end{align*}
\end{lemma}
\begin{proof}  
By properties of the log-sum-exp map, we have
\begin{align*}
v_{r',k}(a,s) - v_{r,k}(a,s)
&= \int \Big[\operatorname{logsumexp}_{\tau}\big(q_{r',k}(\cdot,s')\big)
           - \operatorname{logsumexp}_{\tau}\big(q_{r,k}(\cdot,s')\big)\Big]\,
           k(s'\mid a,s)\,\mu(ds') \\[4pt]
&\le \int \max_{a'\in\mathcal A}\big\{r(a',s') - r'(a',s') + \gamma \, \{v_{r',k}(a', s') -v_{r,k}(a',s') \}\big\}\,
           k(s'\mid a,s)\,\mu(ds')\\
           &\le \int \max_{a'\in\mathcal A}\big\{r(a',s') - r'(a',s')\big\}\,
           k(s'\mid a,s)\,\mu(ds') +  \gamma \int \max_{a'\in\mathcal A}\big\{ v_{r',k}(a', s') -v_{r,k}(a',s') \big\}\\
           &\le \int \max_{a'\in\mathcal A}\big\{r(a',s') - r'(a',s')\big\}\,
           k(s'\mid a,s)\,\mu(ds') +  \gamma \|v_{r',k} - v_{r,k}\|_{L^\infty(\lambda)}.
\end{align*}
where $q_{r,k}(a', s') = r(a', s') + \gamma v_{r,k}(a',s')$. 
Hence, rearranging terms,
\begin{align*}
v_{r',k}(a,s) - v_{r,k}(a,s)-  \gamma \|v_{r',k} - v_{r,k}\|_{L^\infty(\lambda)}  &\leq \int \max_{a'\in\mathcal A}\big\{r(a',s') - r'(a',s')\big\}\,
           k(s'\mid a,s)\,\mu(ds') \\
& \lesssim \|r'-r\|_{L^2(\lambda)},
\end{align*}
where we used the boundedness of $k$ and the definition of $\lambda$. Hence, taking the supremum norm of both sides, we find
\[
 \|v_{r',k}-v_{r,k}\|_{L^\infty(\lambda)} \lesssim  (1-\gamma)^{-1}   \|r' - r\|_{L^2(\lambda)}.
\]



 

Next, we turn to $v_{r,k'} - v_{r,k}$. By definition of the fixed points and arguing as in the proof of Lemma \ref{lemma::bellmansmoothnessfull}, we have that
\begin{align*}
    v_{r,k'}(a,s) - v_{r,k}(a,s) 
    &= \int \operatorname{logsumexp}_{\tau}(r + \gamma v_{r,k'})(s') \,\big(k'-k\big)(s' \mid a,s)\,\mu(ds') \\
    &\quad + \int \Big(\operatorname{logsumexp}_{\tau}(r + \gamma v_{r,k'})(s') 
           - \operatorname{logsumexp}_{\tau}(r + \gamma v_{r,k})(s')\Big)\,k(s' \mid a,s)\,\mu(ds')\\
    & \lesssim \|r + \gamma v_{r, k'}\|_{L^\infty(\lambda)} \cdot \|(\sqrt{k'} - \sqrt{k})(\cdot \mid a, s)\|_{L^2(\mu)} \\
    & \quad + \gamma \int  \max_{a' \in \mathcal{A}}  \{v_{r,k'}(a', s') - v_{r,k}(a', s')\} \, k(s' \mid a,s)\,\mu(ds').
\end{align*}
 



Hence, using that $ \|r + \gamma v_{r, k'}\|_{L^\infty(\lambda)} \lesssim  1 + \|r  \|_{L^\infty(\lambda)} $,
\begin{align*}
    \widetilde{\mathcal{T}}_{k} (v_{r,k'} - v_{r,k} ) \lesssim  ( 1 + \|r  \|_{L^\infty(\lambda)}) \cdot \|(\sqrt{k'} - \sqrt{k})(\cdot \mid a, s)\|_{L^2(\mu)}.
\end{align*}
In fact, we can also show that $   \widetilde{\mathcal{T}}_{k} |v_{r,k'} - v_{r,k}| \lesssim  ( 1 + \|r  \|_{L^\infty(\lambda)}) \cdot \|(\sqrt{k'} - \sqrt{k})(\cdot \mid a, s)\|_{L^2(\mu)}.$
Thus,
\begin{align*}
   \| \widetilde{\mathcal{T}}_{k} (v_{r,k'} - v_{r,k} ) \|_{L^2(\lambda)}  \lesssim  ( 1 + \|r  \|_{L^\infty(\lambda)}) \cdot \|k' - k\|_{L^2(\mu \otimes \lambda)},
\end{align*}
where, as shown in the proof of Lemma~\ref{lemma::bellmansmoothnessfull}, the $L^2(\lambda)$-norm of $(a,s) \mapsto \|(\sqrt{k'} - \sqrt{k})(\cdot \mid a, s)\|_{L^2(\mu)}$ is bounded above by $\|k' - k\|_{L^2(\mu \otimes \lambda)}$.  

 
 

 

\end{proof}

 


 
\begin{lemma}[Quadratic remainder for partial Fréchet derivative of $r \mapsto v_{r,k}$]
 \label{lemma::valuefirstderivativeexpansion}
 Suppose $\gamma < 1$. For $r,r' \in \mathcal{H}$ and $k \in \mathcal{K}$, we have
\[
\big\|v_{r',k} - v_{r,k} - \partial_r v_{r,k}[\,r' - r\,]\big\|_{L^\infty(\lambda)}
\;\lesssim\; \|r' - r\|_{L^2(\lambda)}^2.
\]
\end{lemma}
\begin{proof}
Let $q_{r,k} := r + \gamma v_{r,k}$.  We apply the implicit function theorem to obtain differentiability of $v_{r,k}$ in $r$. By Lemma~\ref{lemma::bellmansmoothnessfull}, the map $F$ is Fr\'echet differentiable, satisfying the expansion:
\begin{align*}
0 
&= F(r', v_{r',k}, k) - F(r, v_{r,k}, k) \\
&= \partial_{r,v} F(r, v_{r,k}, k)\big[r' - r,\, v_{r',k} - v_{r,k}\big] + R(r' - r) \\
&= \partial_r F(r, v_{r,k}, k)[\,r' - r\,] 
  + \partial_v F(r, v_{r,k}, k)[\,v_{r',k} - v_{r,k}\,] + R(r' - r),
\end{align*}
where $\|R(r' - r)\|_{L^\infty(\lambda)} = O\!\big(\|q_{r',k} - q_{r,k}\|_{L^2(\lambda)}^2\big)$ by Lemma \ref{lemma::softbellmanissmooth}.
Moreover, differentiating both sides of $F(r, v_{r,k}, k) = 0$ with respect to $r$, we find, by the chain rule, that
\[
\partial_{r,v}F(r,v_{r,k},k)\big[r' - r,\,\partial_r v_{r,k}[\,r' - r\,]\big] 
= \partial_r F(r,v_{r,k},k)[\,r' - r\,] 
+ \partial_v F(r,v_{r,k},k)[\,\partial_r v_{r,k}[\,r' - r\,]\,] 
= 0.
\]
Hence,
\[
0 
= \partial_v F(r, v_{r,k}, k)\big[v_{r',k} - v_{r,k} - \partial_r v_{r,k}[\,r' - r\,]\big] 
   + R(r' - r).
\]
Equivalently,
\[
\mathcal{T}_{r,k}\,\big(v_{r',k} - v_{r,k} - \partial_r v_{r,k}[\,r' - r\,]\big) = -\,R(r' - r).
\]
Since $\gamma < 1$, the map $\mathcal{T}_{r,k}: L^\infty(\lambda) \rightarrow L^\infty(\lambda)$ is invertible by Banach's fixed point theorem. Hence,
\[
v_{r',k} - v_{r,k} - \partial_r v_{r,k}[\,r' - r\,]
= -\,\mathcal{T}_{r,k}^{-1} R(r' - r),
\]
so that
\[
\big\|v_{r',k} - v_{r,k} - \partial_r v_{r,k}[\,r' - r\,]\big\|_{L^\infty(\lambda)}
\;\le\; \|\mathcal{T}_{r,k}^{-1}\|_{L^\infty(\lambda)}\,\|R(r' - r)\|_{L^\infty(\lambda)}
\;\lesssim\; \|q_{r',k} - q_{r,k}\|_{L^2(\lambda)}^2.
\]
Finally, by Lemma \ref{lemma::lipschitzvalue}, $ \|q_{r',k} - q_{r,k}\|_{L^2(\lambda)}^2 \lesssim  \|r' - r\|_{L^2(\lambda)}^2$. Hence,
\[
\big\|v_{r',k} - v_{r,k} - \partial_r v_{r,k}[\,r' - r\,]\big\|_{L^\infty(\lambda)}
 \;\lesssim\; \|r' - r\|_{L^2(\lambda)}^2.
\]
\end{proof}

 
 
 
\begin{lemma}[Quadratic remainder for partial Fréchet derivative of $k \mapsto v_{r,k}$]
 \label{lemma::valuefirstderivativeexpansion::kernel}
 Suppose $\gamma < 1$. For $r \in \mathcal{H}$ and $k, k' \in \mathcal{K}$, we have
$$  \|\big(v_{r,k'} - v_{r,k} - \adk v_{r,k}[\,k' - k\,]\big)\|_{L^2(\lambda)}  \lesssim \|\mathcal{T}_{r,k}^{-1}\|_{L^2 \rightarrow L^2} \|v_{r,k'} - v_{r,k}\|_{L^2(\lambda)} \|k' - k\|_{L^2(\mu \otimes \lambda)}.$$
\end{lemma}
\begin{proof}

We apply the implicit function theorem to obtain differentiability of $v_{r,k}$ in $k$. By Lemma~\ref{lemma::bellmansmoothnessfull}, the map $F$ is Fr\'echet differentiable, satisfying the expansion:
\begin{align*}
0 
&= F(r, v_{r,k'}, k') - F(r, v_{r,k}, k) \\
&= \partial_v F(r, v_{r,k}, k)[\,v_{r,k'} - v_{r,k}\,] 
  + \adk F(r, v_{r,k}, k)[\,k' - k\,] + R_2,
\end{align*}
where $\|R_2\|_{L^2(\lambda)} = O\!\big(\|v_{r,k'} - v_{r,k}\|_{L^2(\lambda)} \|k' - k\|_{L^2(\mu \otimes \lambda)}\big)$.
Moreover, differentiating both sides of $F(r, v_{r,k}, k) = 0$ with respect to $k$, we find, by the chain rule, that
\[
\partial_v F(r,v_{r,k},k)[\,\adk v_{r,k}[\,k' - k\,]\,] + \adk F(r,v_{r,k},k)[\,k' - k\,]  
= 0.
\]
Hence,
\[
0 
= \partial_v F(r, v_{r,k}, k)\big[v_{r,k'} - v_{r,k} - \adk v_{r,k}[\,k' - k\,]\big] 
   + R_2.
\]
Equivalently,
\[
\mathcal{T}_{r,k}\,\big(v_{r,k'} - v_{r,k} - \adk v_{r,k}[\,k' - k\,]\big) = -\,R_2.
\]
Hence,
$$\| \mathcal{T}_{r,k}\,\big(v_{r,k'} - v_{r,k} - \adk v_{r,k}[\,k' - k\,]\big)\|_{L^2(\lambda)} \leq \|R_2\|_{L^2(\lambda)} \lesssim \|v_{r,k'} - v_{r,k}\|_{L^2(\lambda)} \|k' - k\|_{L^2(\mu \otimes \lambda)},$$
and, therefore,
$$  \|\big(v_{r,k'} - v_{r,k} - \adk v_{r,k}[\,k' - k\,]\big)\|_{L^2(\lambda)}  \lesssim \|\mathcal{T}_{r,k}^{-1}\|_{L^2 \rightarrow L^2} \|v_{r,k'} - v_{r,k}\|_{L^2(\lambda)} \|k' - k\|_{L^2(\mu \otimes \lambda)}.$$
%From Lemma \ref{lemma::lipschitzvalue}, we have that
%$$\|v_{r,k'} - v_{r,k}\|_{L^2(\lambda)}  \lesssim \|\widetilde{\mathcal{T}}_{k}^{-1}\|_{L^2 \rightarrow L^2} \|r\|_{L^\infty(\lambda)} \cdot \|k_1 - k\|_{L^2(\mu \otimes \lambda)}.$$
%Hence, 
%$$  \|\big(v_{r,k'} - v_{r,k} - \adk v_{r,k}[\,k' - k\,]\big)\|_{L^2(\lambda)}  \lesssim \|\mathcal{T}_{r,k}^{-1}\|_{L^2 \rightarrow L^2} \|v_{r,k'} - v_{r,k}\|_{L^2(\lambda)} \|k' - k\|_{L^2(\mu \otimes \lambda)}^2.$$


 
 



\end{proof}
 


 



\begin{lemma}[Lipschitz properties of softmax operator]
\label{lemma::softmaxLipschitz}
 For all $r, r' \in L^\infty(\lambda)$ and $k, k' \in \mathcal{K}$, we have that
$$\bigl\|\Pi_{r',k'}-\Pi_{r,k}\bigr\|_{L^2(\lambda) \rightarrow L^1(\mu)} \lesssim \|r'-r\|_{L^2(\lambda)}
       + \|v_{r,k'}-v_{r,k}\|_{L^2(\lambda)}.$$
\end{lemma}
\begin{proof}
By the $\tau^{-1}$-Lipschitz property of the temperature-$\tau$ softmax and by Cauchy--Schwarz, we have that
\begin{align*}
   \bigl|\Pi_{r',k'}(f)-\Pi_{r,k}(f)\bigr|(s)
   &= \Bigl|\sum_{a \in \mathcal{A}} f(a,s)\,\bigl(\pi_{r',k'}(a\mid s)-\pi_{r,k}(a\mid s)\bigr)\Bigr| \\
   &= \bigl|\langle f(\cdot,s),\,\pi_{r',k'}(\cdot\mid s)-\pi_{r,k}(\cdot\mid s)\rangle_{\ell_2(\mathcal A)}\bigr| \\
   &\le \|f(\cdot,s)\|_{\ell_2(\mathcal A)} \,\|\pi_{r',k'}(\cdot\mid s)-\pi_{r,k}(\cdot\mid s)\|_{\ell_2(\mathcal A)} \\
   &\lesssim \tau^{-1}\,\|f(\cdot,s)\|_{\ell_2(\mathcal A)} \,
   \Biggl(\sum_{a\in\mathcal A}
      \bigl(|r'(a,s)-r(a,s)| + |v_{r',k'}(a,s)-v_{r,k}(a,s)|\bigr)^2
   \Biggr)^{1/2},
\end{align*}
where $ \langle g,h\rangle_{\ell_2(\mathcal A)} := \sum_{a\in\mathcal A} g(a)h(a)$. 


Let $L^2(\mu;\ell_2(\mathcal A))$ denotes the space of measurable functions  $g:\mathcal{A}\times\mathcal{S}\to\mathbb{R}$ with finite mixed norm $\|g\|_{L^2(\mu;\ell_2(\mathcal A))}^2
   := \int_{\mathcal S} \sum_{a\in\mathcal A} |g(a,s)|^2 \, \mu(ds).$ Then, by Cauchy--Schwarz,
\begin{align*}
\bigl\|\Pi_{r',k'}(f)-\Pi_{r,k}(f)\bigr\|_{L^1(\mu)}
&= \int \bigl|\Pi_{r',k'}(f)-\Pi_{r,k}(f)\bigr|(s)\, \mu(ds) \\
&\le \int \|f(\cdot,s)\|_{\ell_2(\mathcal A)} \,
        \|\pi_{r',k'}(\cdot\mid s)-\pi_{r,k}(\cdot\mid s)\|_{\ell_2(\mathcal A)}\,\mu(ds) \\
&\le \|f\|_{L^2(\mu;\ell_2(\mathcal A))}\;
    \|\pi_{r',k'}-\pi_{r,k}\|_{L^2(\mu;\ell_2(\mathcal A))} \\
&\lesssim \|f\|_{L^2(\mu;\ell_2(\mathcal A))}\;
    \|(r'-r)+(v_{r',k'}-v_{r,k})\|_{L^2(\mu;\ell_2(\mathcal A))} \\
&\lesssim \|f\|_{L^2(\mu;\ell_2(\mathcal A))}\;
   \Big( \|r'-r\|_{L^2(\mu;\ell_2(\mathcal A))}
       + \|v_{r',k'}-v_{r,k}\|_{L^2(\mu;\ell_2(\mathcal A))} \Big).
\end{align*}
Finally, since $\mathcal{A}$ is finite, the $L^2(\mu;\ell_2(\mathcal A))$ and $L^2(\lambda)$ norms are equal up to a multiplicative factor depending on $\# |\mathcal{A}|$. Hence,
\begin{align*}
\bigl\|\Pi_{r',k'}(f)-\Pi_{r,k}(f)\bigr\|_{L^1(\mu)} &\lesssim \|f\|_{L^2(\lambda)}\;
   \Big( \|r'-r\|_{L^2(\lambda)}
       + \|v_{r',k'}-v_{r,k}\|_{L^2(\lambda)} \Big).
\end{align*}
Moreover, by Lemma~\ref{lemma::lipschitzvalue}, we have
\begin{align*}
\|v_{r',k'}-v_{r,k}\|_{L^2(\lambda)}
&\leq \|v_{r',k'}-v_{r,k'}\|_{L^2(\lambda)} 
   + \|v_{r,k'}-v_{r,k}\|_{L^2(\lambda)} \\
&\lesssim \|r'-r\|_{L^2(\lambda)} 
   + \|v_{r,k'}-v_{r,k}\|_{L^2(\lambda)}.
\end{align*}
Putting everything together, we conclude that
\begin{align*}
\bigl\|\Pi_{r',k'}(f)-\Pi_{r,k}(f)\bigr\|_{L^1(\mu)}
   \;\lesssim\; \|f\|_{L^2(\lambda)}
   \Big( \|r'-r\|_{L^2(\lambda)}
       + \|v_{r,k'}-v_{r,k}\|_{L^2(\lambda)} \Big),
\end{align*}
which implies, by definition, that $\bigl\|\Pi_{r',k'}-\Pi_{r,k}\bigr\|_{L^2(\lambda) \rightarrow L^1(\mu)} \lesssim \|r'-r\|_{L^2(\lambda)}
       + \|v_{r,k'}-v_{r,k}\|_{L^2(\lambda)}$.


\end{proof}


\begin{lemma}
\label{lemma::forwardopLipschitz}
There exists a constant $C<\infty$ (uniform on the parameter neighborhood) such that
\[
\|\mathcal{P}_{r',k'} - \mathcal{P}_{r,k}\|_{L^2(\lambda) \to L^2(\lambda)}
\;\le\; C\Big(\|k' - k\|_{\mathcal K} + \|r' - r\|_{L^2(\lambda)}
+ \|v_{r,k'}-v_{r,k}\|_{L^2(\lambda)}\Big).
\]
\end{lemma}

\begin{proof}
By definition, $\mathcal{P}_{r,k} = \mathcal{P}_k \,\Pi_{r,k}$. Hence
\begin{align*}
\|\mathcal{P}_{r',k'} - \mathcal{P}_{r,k}\|_{L^2(\lambda)\to L^2(\lambda)}
&= \|\mathcal{P}_{k'}\Pi_{r',k'} - \mathcal{P}_k \Pi_{r,k}\|_{L^2(\lambda)\to L^2(\lambda)} \\
&\le \|\mathcal{P}_{k'}-\mathcal{P}_k\|_{L^2(\lambda)\to L^2(\lambda)}\,
\|\Pi_{r',k'}\|_{L^2(\lambda)\to L^2(\lambda)} \\
&\qquad + \|\mathcal{P}_k\|_{L^2(\lambda)\to L^2(\lambda)}\,
\|\Pi_{r',k'}-\Pi_{r,k}\|_{L^2(\lambda)\to L^2(\lambda)},
\end{align*}
by the triangle inequality and submultiplicativity. By Lemma~\ref{lemma::bellmanfrechet},
\[
\|\mathcal{P}_{k'}-\mathcal{P}_k\|_{L^2(\lambda)\to L^2(\lambda)} \;\lesssim\; \|k' - k\|_{\mathcal K}.
\]
By Lemma~\ref{lemma::softmaxLipschitz},
\[
\|\Pi_{r',k'}-\Pi_{r,k}\|_{L^2(\lambda)\to L^2(\lambda)} \;\lesssim\; \|r'-r\|_{L^2(\lambda)} \;+\; \|v_{r,k'}-v_{r,k}\|_{L^2(\lambda)}.
\]
Using the uniform boundedness of $\|\Pi_{r',k'}\|_{L^2(\lambda)\to L^2(\lambda)}$ and $\|\mathcal{P}_k\|_{L^2(\lambda)\to L^2(\lambda)}$ on the neighborhood yields
\[
\|\mathcal{P}_{r',k'} - \mathcal{P}_{r,k}\|_{L^2(\lambda) \to L^2(\lambda)}
\;\lesssim\; \|k' - k\|_{\mathcal K} + \|r' - r\|_{L^2(\lambda)}
+ \|v_{r,k'}-v_{r,k}\|_{L^2(\lambda)}.
\]
\end{proof}




\begin{lemma}[Fréchet differentiability of linear operators]
\label{lemma::operatorderivatives}
Let $r, r' \in L^\infty(\lambda)$, $k, k' \in \mathcal{K}$, and $f \in L^\infty(\lambda)$.  
Suppose $\|\mathcal{T}_{r,k}^{-1}\|_{L^2 \rightarrow L^2} < \infty$.
Then the partial Gâteaux derivatives of the maps $(r,k) \mapsto \Pi_{r, k} f$ and $(r,k) \mapsto \mathcal{P}_{r, k} f$ in the direction $r'-r$ and $k' - k$ are
\begin{align*}
\partial_r \Pi_{r, k}[\,r'-r\,](f)(s') 
&= \frac{1}{\tau}\,\operatorname{Cov}_{\pi_{r, k}(\cdot \mid s')}
   \big(f(\cdot, s'),\, \partial_r q_{r,k}[\,r'-r\,](\cdot, s')\big), \\[0.8ex]
%
\partial_r \mathcal{P}_{r,k}[\,r'-r\,](f)(a,s)    
&= \frac{1}{\tau}\int \operatorname{Cov}_{\pi_{r, k}(\cdot\mid s')}\!\Big(
       f(\cdot,s'),\, \partial_r q_{r,k}[\,r'-r\,](\cdot,s')
     \Big)\,k(s'\mid a,s)\,\mu(ds'), \\[1.2ex]
%
\adk \Pi_{r, k}[\,k'-k\,](f)(s') 
&= \frac{1}{\tau}\,\operatorname{Cov}_{\pi_{r, k}(\cdot \mid s')}
   \big(f(\cdot, s'),\, \adk q_{r,k}[\,k'-k\,](\cdot, s')\big), \\[0.8ex]
%
\adk \mathcal{P}_{r,k}[\,k'-k\,](f)(a,s) 
&= \frac{1}{\tau}\int \operatorname{Cov}_{\pi_{r, k}(\cdot\mid s')}\!\Big(
       f(\cdot,s'),\, \adk q_{r,k}[\,k'-k\,](\cdot,s')
     \Big)\,k(s'\mid a,s)\,\mu(ds') \\
&\quad + \int \Pi_{r,k}(f)(s')\,(k'-k)(s' \mid a, s)\,\mu(ds').
\end{align*}
Moreover, the following bounds hold:
\begin{align*}
\|\Pi_{r', k} f - \Pi_{r, k} f 
   - \partial_r \Pi_{r, k}[\,r'-r\,](f)\|_{L^1(\mu)}
&\;\lesssim\; \|f\|_{L^\infty(\lambda)} \, \|r'-r\|_{L^2(\lambda)}^2, \\[1.0ex]
%
\big\| \mathcal{P}_{r', k} f - \mathcal{P}_{r, k} f 
      - \partial_r \mathcal{P}_{r, k}[\,r'-r\,](f) \big\|_{L^\infty(\lambda)}    
&\;\lesssim\; \|f\|_{L^\infty(\lambda)} \, \|r'-r\|_{L^2(\lambda)}^2, \\[1.2ex]
%
\|\Pi_{r, k'} f - \Pi_{r, k} f 
   - \adk \Pi_{r, k}[\,k'-k\,](f)\|_{L^1(\mu)} 
&\;\lesssim\; \|f\|_{L^\infty(\lambda)} \Big\{
      \|v_{r, k'} - v_{r,k}\|_{L^2(\lambda)}^2 \notag\\
&\qquad\qquad
      + \|\mathcal{T}_{r,k}^{-1}\|_{L^2(\lambda) \to L^2(\lambda)} \,
        \|v_{r,k'} - v_{r,k}\|_{L^2(\lambda)}\,
        \|k' - k\|_{L^2(\mu \otimes \lambda)}
      \Big\}, \\[1.2ex]
%
\big\|\mathcal{P}_{r,k'}(f) - \mathcal{P}_{r,k}(f)\notag\\
&\qquad - \adk \mathcal{P}_{r,k}[\,k'-k\,](f)\big\|_{L^2(\lambda)} 
\;\lesssim\; \|f\|_{L^\infty(\lambda)} \Big\{
      \|v_{r, k'} - v_{r,k}\|_{L^2(\lambda)}^2 \notag\\
&\qquad\qquad
      + \big(1 + \|\mathcal{T}_{r,k}^{-1}\|_{L^2(\lambda) \to L^2(\lambda)}\big)\,
        \|v_{r,k'} - v_{r,k}\|_{L^2(\lambda)}\,
        \|k' - k\|_{L^2(\mu \otimes \lambda)}
      \Big\}.
\end{align*}
\end{lemma}



\begin{proof}
\noindent \textbf{Differentiability of softmax operator.} Define the operator $\widetilde{\Pi}_{q,\tau}: L^2(\lambda) \rightarrow L^2(\mu)$ by
\[
(\widetilde{\Pi}_{q,\tau} f)(s) := \sum_{a' \in \mathcal{A}} \widetilde{\pi}_{q,\tau}(a' \mid s)\, f(a', s),
\qquad
\widetilde{\pi}_{q,\tau}(a' \mid s) := \frac{\exp\!\left(\frac{q(a',s)}{\tau}\right)}{\sum_{b \in \mathcal{A}} \exp\!\left(\frac{q(b,s)}{\tau}\right)}.
\]
For fixed $s' \in \mathcal{S}$, let
\[
\Phi(t) := (\widetilde{\Pi}_{q + t f_1,\tau} f)(s') 
= \sum_{a' \in \mathcal{A}} \widetilde{\pi}_{\,q + t f_1,\tau}(a' \mid s')\, f(a', s').
\]
The derivative of $\widetilde{\pi}_{\,q + t f_1,\tau}(a' \mid s')$ with respect to $t$ is
\[
\frac{d}{dt} \widetilde{\pi}_{\,q + t f_1,\tau}(a' \mid s') 
= \frac{1}{\tau}\,\widetilde{\pi}_{\,q + t f_1,\tau}(a' \mid s') \big[ f_1(a', s') - E_{\widetilde{\pi}_{\,q + t f_1,\tau}}[f_1(\cdot, s')] \big].
\]
Hence, $\Phi'(0) = \tfrac{1}{\tau}\operatorname{Cov}_{\widetilde{\pi}_{q,\tau}(\cdot \mid s')}(f(\cdot, s'), f_1(\cdot, s'))$.  
The second derivative $\Phi''(t)$ is uniformly bounded in absolute value by $2\,\tau^{-2}\,\|f(\cdot, s')\|_\infty\,\|f_1(\cdot, s')\|_{L^2(\widetilde{\pi}_{\,q + t f_1,\tau}(\cdot \mid s'))}^2.$
By Taylor’s theorem with integral remainder,
\begin{align*}
\left| \Phi(t) - \Phi(0) - t\,\Phi'(0) \right|
&\le t^2\,\tau^{-2}\,\|f(\cdot, s')\|_\infty 
   \sup_{u \in [0,1]} \|f_1(\cdot, s')\|_{L^2(\widetilde{\pi}_{\,q + u t f_1,\tau}(\cdot \mid s'))}^2 \\
&\le t^2\,\tau^{-2}\,\|f(\cdot, s')\|_\infty \sum_{a' \in \mathcal{A}} \{f_1(a', s')\}^2 .
\end{align*}
Setting $t=1$ and applying the previous bound pointwise in $s'$, we obtain
\begin{align*}
    \left| 
\sum_{a' \in \mathcal{A}} \widetilde{\pi}_{\,q + f_1,\tau}(a' \mid s') f(a', s')
- \sum_{a' \in \mathcal{A}} \widetilde{\pi}_{q,\tau}(a' \mid s') f(a', s')
- \frac{1}{\tau}\operatorname{Cov}_{\widetilde{\pi}_{q,\tau}(\cdot \mid s')}\big(f(\cdot, s'), f_1(\cdot, s')\big)
\right|\\
\le \tau^{-2}\max_{a' \in \mathcal{A}} |f(a', s')| \sum_{a' \in \mathcal{A}} \big\{ f_1(a', s') \big\}^2.
\end{align*}
Thus, letting 
\[
\partial_q \widetilde{\Pi}_{q,\tau}[f_1](f) := \big(s' \mapsto \tfrac{1}{\tau}\operatorname{Cov}_{\widetilde{\pi}_{q,\tau}(\cdot \mid s')}\big(f(\cdot, s'), f_1(\cdot, s')\big)\big),
\]
we obtain the pointwise bound
\begin{align*}
    \left|(\widetilde{\Pi}_{q + f_1,\tau} f)(s') - (\widetilde{\Pi}_{q,\tau} f)(s') - \partial_q \widetilde{\Pi}_{q,\tau}[f_1](f)(s') \right| 
&\lesssim  \max_{a' \in \mathcal{A}} |f(a', s')| \sum_{a' \in \mathcal{A}} \big\{ f_1(a', s') \big\}^2.
\end{align*}

\noindent \textbf{Partial differentiability of $\Pi_{r,k}$ in the reward.} We now turn to establishing partial differentiability of the map $(r,v) \mapsto \Pi_{r,k}$. We begin with the dependence on $r$, i.e., the map $r \mapsto \Pi_{r,k}$. By the chain rule, we have that $\partial_r \Pi_{r, k}[u] =  \partial_q \widetilde{\Pi}_{q_{r,k},\tau}[\partial_r q_{r,k}[u]]$ and, hence,
\begin{align*}
    \partial_r \Pi_{r, k}[u](f)(s') &=  \frac{1}{\tau}\operatorname{Cov}_{\pi_{r, k}(\cdot \mid s')}\big(f(\cdot, s'), \partial_r q_{r,k}[u](\cdot, s')\big) .
\end{align*} 
Next, we determine the remainder in the first-order Taylor expansion. Adding and subtracting terms and applying the triangle inequality, we have 
\begin{align*}
&\left| \Pi_{r + u, k} f - \Pi_{r, k} f - \partial_r \Pi_{r, k}[u](f) \right| \\
&= \left| \Pi_{r + u, k} f - \Pi_{r, k} f - \partial_q \widetilde{\Pi}_{q_{r,k},\tau}[\partial_r q_{r,k}[u]](f) \right| \\
&\leq \left| \widetilde{\Pi}_{q_{r + u, k},\tau} f - \widetilde{\Pi}_{q_{r,k},\tau} f - \partial_q \widetilde{\Pi}_{q_{r,k},\tau}[q_{r + u, k} - q_{r,k}] (f)\right| \\
&\quad + \left| \partial_q \widetilde{\Pi}_{q_{r,k},\tau}[q_{r + u, k} - q_{r,k} - \partial_r q_{r,k}[u]](f) \right|,
\end{align*}
where  $q_{r + u, k} - q_{r,k} - \partial_r q_{r,k}[u] 
= \gamma \big(v_{r + u, k} - v_{r,k} - \partial_r v_{r,k}[u]\big).$
Hence, applying our earlier bounds, we find that 
\begin{align*}
&\left| \Pi_{r + u, k} f - \Pi_{r, k} f - \partial_r \Pi_{r, k}[u](f) \right| \\
&\lesssim \|f\|_{L^\infty(\lambda)}  \sum_{a' \in \mathcal{A}} (q_{r + u, k} - q_{r,k})^2(a' , \cdot)
   + \left| \partial_q \widetilde{\Pi}_{q_{r,k},\tau}[q_{r + u, k} - q_{r,k} - \partial_r q_{r,k}[u]](f) \right| \\
&\lesssim \|f\|_{L^\infty(\lambda)} \,\sum_{a' \in \mathcal{A}} (q_{r + u, k} - q_{r,k})^2(a' , \cdot) 
   + \gamma\,\|f\|_{L^\infty(\lambda)} \cdot \Pi_{r,k}\left|v_{r + u, k} - v_{r,k} - \partial_r v_{r,k}[u] \right| ,
\end{align*}
where the final bound follows from the covariance properties in the definition of $\partial_q \widetilde{\Pi}_{q_{r,k},\tau}$. We have, by Lemma~\ref{lemma::lipschitzvalue}, 
\[
\|q_{r',k} - q_{r,k}\|_{L^2(\lambda)}  \leq   \|r' - r\|_{L^2(\lambda)}+ \gamma \|v_{r',k} - v_{r,k}\|_{L^\infty(\lambda)} \;\lesssim\; \|r' - r\|_{L^2(\lambda)}.
\]
Moreover, by Lemma~\ref{lemma::valuefirstderivativeexpansion},  $\big\|v_{r + u,k} - v_{r,k} - \partial_r v_{r,k}[\,u\,]\big\|_{L^\infty(\lambda)} \lesssim \|u\|_{L^2(\lambda)}^2.$ 
Hence, we conclude that 
\begin{align}
    \label{eqn::boundPi_reward}
    \| \Pi_{r + u, k} f - \Pi_{r, k} f - \partial_r \Pi_{r, k}[u](f) \|_{L^1(\mu)}
\;\lesssim\; \|f\|_{L^\infty(\lambda)} \,\|u\|_{L^2(\lambda)}^2
\end{align}
 

\noindent \textbf{Partial differentiability of $\Pi_{r,k}$ in the kernel.} Next we establish partial differentiability of the map $k \mapsto \Pi_{r,k}$. By the chain rule, we have that $\adk \Pi_{r, k}[k'-k] =  \partial_q \widetilde{\Pi}_{q_{r,k},\tau}[\adk q_{r,k}[k'-k]]$ and, hence,
\begin{align*}
    \adk \Pi_{r, k}[k'-k](f)(s') &=  \frac{1}{\tau}\operatorname{Cov}_{\pi_{r, k}(\cdot \mid s')}\big(f(\cdot, s'), \adk q_{r,k}[k'-k](\cdot, s')\big) .
\end{align*} 
Next, we determine the remainder in the first-order Taylor expansion. Adding and subtracting terms and applying the triangle inequality, we have 
\begin{align*}
&\left| \Pi_{r, k'} f - \Pi_{r, k} f - \adk \Pi_{r, k}[k'-k](f) \right| \\
&= \left| \Pi_{r , k'} f - \Pi_{r, k} f - \partial_q \widetilde{\Pi}_{q_{r,k},\tau}\big[\adk q_{r,k}[k'-k]\big](f) \right| \\
&\leq \left| \widetilde{\Pi}_{q_{r , k'},\tau} f - \widetilde{\Pi}_{q_{r,k},\tau} f 
        - \partial_q \widetilde{\Pi}_{q_{r,k},\tau}\big[q_{r, k'} - q_{r,k}\big](f)\right| \\
&\quad + \left| \partial_q \widetilde{\Pi}_{q_{r,k},\tau}\big[q_{r, k'} - q_{r,k} - \adk q_{r,k}[k'-k]\big](f) \right|.
\end{align*}
Since $q_{r,k}=r+\gamma v_{r,k}$ and $r$ is fixed here, we have
\[
q_{r, k'} - q_{r,k} - \adk q_{r,k}[k'-k]
= \gamma \big(v_{r, k'} - v_{r,k} - \adk v_{r,k}[k'-k]\big).
\]
Hence, we find that 
\begin{align*}
&\left| \Pi_{r, k'} f - \Pi_{r, k} f - \adk \Pi_{r, k}[k'-k](f) \right| \\
&\lesssim \|f\|_{L^\infty(\lambda)}  \sum_{a' \in \mathcal{A}} \big(q_{r, k'} - q_{r,k}\big)^2(a' , \cdot)
   + \left| \partial_q \widetilde{\Pi}_{q_{r,k},\tau}\big[q_{r, k'} - q_{r,k} - \adk q_{r,k}[k'-k]\big](f) \right| \\
&\lesssim \|f\|_{L^\infty(\lambda)} \,\sum_{a' \in \mathcal{A}} \big(v_{r, k'} - v_{r,k}\big)^2(a' , \cdot) 
   + \gamma\,\|f\|_{L^\infty(\lambda)}  \Pi_{r,k}\left|v_{r, k'} - v_{r,k} - \adk v_{r,k}[k'-k] \right| . 
\end{align*}
Lemma \ref{lemma::valuefirstderivativeexpansion::kernel} shows that $$  \|\big(v_{r,k'} - v_{r,k} - \adk v_{r,k}[\,k' - k\,]\big)\|_{L^2(\lambda)}  \lesssim \|\mathcal{T}_{r,k}^{-1}\|_{L^2 \rightarrow L^2} \|v_{r,k'} - v_{r,k}\|_{L^2(\lambda)} \|k' - k\|_{L^2(\mu \otimes \lambda)}.$$
Thus, applying Cauchy-Schwarz,
\begin{align}
\label{eqn::boundPi_kernel}
\| \Pi_{r, k'} f - \Pi_{r, k} f - \adk \Pi_{r, k}[k'-k](f) \|_{L^1(\mu)}
&\lesssim \|f\|_{L^\infty(\lambda)} \Big\{
    \|v_{r, k'} - v_{r,k}\|_{L^2(\lambda)}^2 \\
&\qquad\qquad
    + \|\mathcal{T}_{r,k}^{-1}\|_{L^2 \rightarrow L^2}
      \|v_{r,k'} - v_{r,k}\|_{L^2(\lambda)}
      \|k' - k\|_{L^2(\mu \otimes \lambda)}
  \Big\}
\end{align}
 

\noindent \textbf{Partial differentiability of $\mathcal{P}_{r,k}$ in the reward.} Building on the above, we turn to establishing partial differentiability of $(r, k) \mapsto \mathcal{P}_{r,k}$. We begin with the $r$ component. By definition, $(\mathcal{P}_{r,k}f)(a,s)=\int \Pi_{r,k}(f)(s')\,k(s'\mid a,s)\,\mu(ds')$.
By linearity and the chain rule,
\begin{align*}
\partial_r \mathcal{P}_{r,k}[\,u\,](f)(a,s)
&= \int \partial_r \Pi_{r,k}[\,u\,](f)(s')\,k(s'\mid a,s)\,\mu(ds') \\
&= \frac{1}{\tau}\int \operatorname{Cov}_{\pi_{r,k}(\cdot\mid s')}
   \!\Big(f(\cdot,s'),\,\partial_r q_{r,k}[\,u\,](\cdot,s')\Big)\,
   k(s'\mid a,s)\,\mu(ds').
\end{align*}
Therefore,
\begin{align*}
&\left| \mathcal{P}_{r+u,k}f - \mathcal{P}_{r,k}f - \partial_r \mathcal{P}_{r,k}[\,u\,](f) \right|(a,s) \\
&\qquad= \left|\int \Big(\Pi_{r+u,k}f - \Pi_{r,k}f - \partial_r \Pi_{r,k}[\,u\,](f)\Big)(s')\,k(s'\mid a,s)\,\mu(ds')\right| \\
&\qquad\le \|k(\cdot\mid a,s)\|_{L^\infty(\mu)}\,
          \left\|\Pi_{r+u,k}f - \Pi_{r,k}f - \partial_r \Pi_{r,k}[\,u\,](f)\right\|_{L^1(\mu)} \\
&\qquad\lesssim \left\|\Pi_{r+u,k}f - \Pi_{r,k}f - \partial_r \Pi_{r,k}[\,u\,](f)\right\|_{L^1(\mu)}
\;\lesssim\; \|f\|_{L^\infty(\lambda)}\,\|u\|_{L^2(\lambda)}^2,
\end{align*}
where the last step uses \eqref{eqn::boundPi_reward} and the uniform bound
$\sup_{k \in \mathcal{K}} \sup_{(a,s)}\|k(\cdot\mid a,s)\|_{L^\infty(\mu)} < \infty$.

\noindent In particular, taking the supremum over $(a,s)$ yields
\[
\big\|\mathcal{P}_{r+u,k}f - \mathcal{P}_{r,k}f - \partial_r \mathcal{P}_{r,k}[\,u\,](f)\big\|_{L^\infty(\lambda)}
\;\lesssim\; \|f\|_{L^\infty(\lambda)}\,\|u\|_{L^2(\lambda)}^2.
\]


\noindent \textbf{Partial differentiability of $\mathcal{P}_{r,k}$ in the kernel.}
By the chain rule,
\begin{align*}
\adk \mathcal{P}_{r,k}[\,k'-k\,](f)(a,s)
&= \int \adk \Pi_{r,k}[\,k'-k\,](f)(s')\,k(s'\mid a,s)\,\mu(ds') \\
&\quad + \int \Pi_{r,k}(f)(s')\,(k'-k)(s'\mid a,s)\,\mu(ds') \\
&= \frac{1}{\tau}\int \operatorname{Cov}_{\pi_{r,k}(\cdot\mid s')}\!\Big(f(\cdot,s'),\,\adk q_{r,k}[\,k'-k\,](\cdot,s')\Big)\,k(s'\mid a,s)\,\mu(ds') \\
&\quad + \int \Pi_{r,k}(f)(s')\,(k'-k)(s'\mid a,s)\,\mu(ds').
\end{align*}
Hence,
\begin{align*}
&\Big|\mathcal{P}_{r,k'}(f)(a,s)-\mathcal{P}_{r,k}(f)(a,s)-\adk \mathcal{P}_{r,k}[\,k'-k\,](f)(a,s)\Big| \\
&\qquad= \Bigg|\int \Big(\Pi_{r,k'}(f)-\Pi_{r,k}(f)-\adk \Pi_{r,k}[\,k'-k\,](f)\Big)(s')\,k(s'\mid a,s)\,\mu(ds') \\
&\qquad\qquad\quad + \int \big(\Pi_{r,k'}(f)-\Pi_{r,k}(f)\big)(s')\,(k'-k)(s'\mid a,s)\,\mu(ds')\Bigg|.
\end{align*}
By \eqref{eqn::boundPi_kernel}, we have
\begin{align*}
&\Bigg|\int \Big(\Pi_{r,k'}(f)-\Pi_{r,k}(f)-\adk \Pi_{r,k}[\,k'-k\,](f)\Big)(s')\,k(s'\mid a,s)\,\mu(ds')\Bigg| \\
&\qquad\lesssim \|f\|_{L^\infty(\lambda)}\Big\{ \|v_{r,k'}-v_{r,k}\|_{L^2(\lambda)}^2
+ \|\mathcal{T}_{r,k}^{-1}\|_{L^2(\lambda)\to L^2(\lambda)}\,\|v_{r,k'}-v_{r,k}\|_{L^2(\lambda)}\,\|k'-k\|_{L^2(\mu\otimes\lambda)}\Big\}.
\end{align*}
For the second integral, using that the temperature-$\tau$ softmax map is $\tau^{-1}$-Lipschitz from $(\ell_\infty,\ell_1)$,
\[
\big|\Pi_{r,k'}(f)(s')-\Pi_{r,k}(f)(s')\big|
\le \|f\|_{L^\infty(\lambda)}\,
\big\|\pi_{r,k'}(\cdot\mid s')-\pi_{r,k}(\cdot\mid s')\big\|_1
\le \frac{\gamma}{\tau}\,\|f\|_{L^\infty(\lambda)}\,
\sum_{a' \in \mathcal{A}}\big|v_{r,k'}(a',s')-v_{r,k}(a',s')\big|.
\]
Arguing as in the proof of Lemma \ref{lemma::bellmansmoothnessfull}, we find that
\begin{align*}
\Big\|(a,s)\mapsto \int &\big(\Pi_{r,k'}(f)-\Pi_{r,k}(f)\big)(s')\,(k'-k)(s'\mid a,s)\,\mu(ds')\Big\|_{L^2(\lambda)} \\
&\lesssim \|f\|_{L^\infty(\lambda)}\,
\|v_{r,k'}-v_{r,k}\|_{L^2(\lambda)}\,
\|k'-k\|_{L^2(\mu\otimes\lambda)}.
\end{align*}
Combining the two bounds,
\begin{align*}
\Big\|\mathcal{P}_{r,k'}(f)-\mathcal{P}_{r,k}(f)-\adk \mathcal{P}_{r,k}[\,k'-k\,](f)\Big\|_{L^2(\lambda)}
\;\lesssim\; &\|f\|_{L^\infty(\lambda)}\Big\{ 
\|v_{r,k'}-v_{r,k}\|_{L^2(\lambda)}^2 \\
&+ \big(1+\|\mathcal{T}_{r,k}^{-1}\|_{L^2(\lambda)\to L^2(\lambda)}\big)\,
\|v_{r,k'}-v_{r,k}\|_{L^2(\lambda)}\,
\|k'-k\|_{L^2(\mu\otimes\lambda)}\Big\}.
\end{align*}



\end{proof}






 
\begin{lemma}
\label{lemma::valuestarderivativereward}
Let $r \in L^2(\lambda)$ and $h \in L^2(\lambda)$, and fix $k$.  
Suppose the linear operators $\mathcal{T}_{r,k}:L^2(\lambda) \to L^2(\lambda)$ 
and $\mathcal{T}_{r+h,k}:L^2(\lambda) \to L^2(\lambda)$ are invertible. Then
\[
\big\|\Pi_{r+h,k}\,\mathcal{T}_{r+h,k}^{-1}(r+h) 
 - \Pi_{r,k}\,\mathcal{T}_{r,k}^{-1}(r) 
 - \partial_r\!\big[\Pi_{r,k}\,\mathcal{T}_{r,k}^{-1}(r)\big][h]\big\|_{L^1(\lambda)} 
= O(\|h\|_{L^2(\lambda)}^2),
\]
with the big-$O$ constant depending on $\|r\|_{L^\infty(\lambda)}$ 
and $\|\mathcal{T}_{r,k}^{-1}\|_{L^2(\lambda)\to L^2(\lambda)}$.
\end{lemma}
\begin{proof}
By the chain rule,
\begin{align*}
\partial_r\!\big[\Pi_{r,k}\,\mathcal{T}_{r,k}^{-1}(r)\big][h]
&= \big(\partial_r \Pi_{r,k}[h]\big)\,\mathcal{T}_{r,k}^{-1}(r)
\;+\;\Pi_{r,k}\,\big(\partial_r \mathcal{T}_{r,k}^{-1}[h]\big)(r)
\;+\;\Pi_{r,k}\,\mathcal{T}_{r,k}^{-1}(h),
\end{align*}
where $\partial_r \mathcal{T}_{r,k}^{-1}[h] = -\,\mathcal{T}_{r,k}^{-1}\big(\partial_r \mathcal{T}_{r,k}[h]\big)\mathcal{T}_{r,k}^{-1}$.
By adding and subtracting, we obtain
\begin{align*}
&\Pi_{r+h,k}\,\mathcal{T}_{r+h,k}^{-1}(r+h) - \Pi_{r,k}\,\mathcal{T}_{r,k}^{-1}(r) - \partial_r\!\big[\Pi_{r,k}\,\mathcal{T}_{r,k}^{-1}(r)\big][h] \\
&= \big(\Pi_{r+h,k}-\Pi_{r,k}-\partial_r \Pi_{r,k}[h]\big)\,\mathcal{T}_{r,k}^{-1}(r) \\
&\quad+\;\Pi_{r,k}\,\big(\mathcal{T}_{r+h,k}^{-1}-\mathcal{T}_{r,k}^{-1}-\partial_r \mathcal{T}_{r,k}^{-1}[h]\big)(r) \\
&\quad+\;\big(\Pi_{r+h,k}-\Pi_{r,k}\big)\big(\mathcal{T}_{r+h,k}^{-1}-\mathcal{T}_{r,k}^{-1}\big)(r) \\
&\quad+\;\big(\Pi_{r+h,k}\,\mathcal{T}_{r+h,k}^{-1}-\Pi_{r,k}\,\mathcal{T}_{r,k}^{-1}\big)(h).
\end{align*}
Taking norms and applying Lemma \ref{lemma::operatorderivatives}, we have
\begin{align*}
& \big\|\Pi_{r+h,k}\,\mathcal{T}_{r+h,k}^{-1}(r+h) - \Pi_{r,k}\,\mathcal{T}_{r,k}^{-1}(r) - \partial_r\!\big[\Pi_{r,k}\,\mathcal{T}_{r,k}^{-1}(r)\big][h]\big\|_{L^1(\lambda)} \\
&\lesssim \big\|\Pi_{r+h,k}-\Pi_{r,k}-\partial_r \Pi_{r,k}[h]\big\|_{L^\infty(\lambda) \to L^1(\lambda)} \,\big\|\mathcal{T}_{r,k}^{-1}(r)\big\|_{L^\infty(\lambda)} \\
&\quad+ \big\|\mathcal{T}_{r+h,k}^{-1}-\mathcal{T}_{r,k}^{-1}-\partial_r \mathcal{T}_{r,k}^{-1}[h]\big\|_{L^\infty(\lambda) \to L^2(\lambda)} \,\|r\|_{L^\infty(\lambda)} \\
&\quad+\;\big\|\Pi_{r+h,k}-\Pi_{r,k}\big\|_{L^2(\lambda) \to L^1(\mu)}\,\big\|\big(\mathcal{T}_{r+h,k}^{-1}-\mathcal{T}_{r,k}^{-1}\big)(r)\big\|_{L^2(\lambda)} \\
&\quad+\;\big\|\big(\Pi_{r+h,k}\,\mathcal{T}_{r+h,k}^{-1}-\Pi_{r,k}\,\mathcal{T}_{r,k}^{-1}\big)(h)\big\|_{L^1(\lambda)}.
\end{align*}

We will show that each term is bounded above by $\|h\|_{L^2(\lambda)}^2$, up to constants depending on $\gamma$, $r$, $\|\mathcal{T}_{r+h,k}^{-1}\|_{L^2(\lambda)\to L^2(\lambda)}$, and $\|\mathcal{T}_{r,k}^{-1}\|_{L^2(\lambda)\to L^2(\lambda)}$. If so, the claim follows since
\[
\big\|\Pi_{r+h,k}\,\mathcal{T}_{r+h,k}^{-1}(r+h) - \Pi_{r,k}\,\mathcal{T}_{r,k}^{-1}(r) - \partial_r\!\big[\Pi_{r,k}\,\mathcal{T}_{r,k}^{-1}(r)\big][h]\big\|_{L^1(\lambda)} 
\;\lesssim\; \|h\|_{L^2(\lambda)}^2.
\]



We study each term in turn.

\medskip
\noindent \textbf{Term 1.}
By Lemma \ref{lemma::operatorderivatives},
\[
\big\|\Pi_{r+h,k}-\Pi_{r,k}-\partial_r \Pi_{r,k}[h]\big\|_{L^\infty(\lambda) \to L^1(\lambda)} \lesssim \|h\|_{L^2(\lambda)}^2.
\]
Hence,
\[
\big\|\Pi_{r+h,k}-\Pi_{r,k}-\partial_r \Pi_{r,k}[h]\big\|_{L^\infty(\lambda) \to L^1(\lambda)} \,\big\|\mathcal{T}_{r,k}^{-1}(r)\big\|_{L^\infty(\lambda)}
\lesssim \big\|\mathcal{T}_{r,k}^{-1}(r)\big\|_{L^\infty(\lambda)} \,\|h\|_{L^2(\lambda)}^2.
\]

\medskip
\noindent \textbf{Term 2.}
Using the resolvent identity, the inverse–operator remainder admits the second–order expansion
\[
\begin{aligned}
&\mathcal{T}_{r+h,k}^{-1} - \mathcal{T}_{r,k}^{-1}
+ \mathcal{T}_{r,k}^{-1}\big(\partial_r \mathcal{T}_{r,k}[h]\big)\mathcal{T}_{r,k}^{-1} \\
&\qquad=- \mathcal{T}_{r,k}^{-1}\Big( \mathcal{T}_{r+h,k} - \mathcal{T}_{r,k} - \partial_r \mathcal{T}_{r,k}[h] \Big)\mathcal{T}_{r,k}^{-1} \\
&\qquad\quad +\; \mathcal{T}_{r,k}^{-1}\big(\mathcal{T}_{r+h,k} - \mathcal{T}_{r,k}\big)\mathcal{T}_{r,k}^{-1}\big(\mathcal{T}_{r+h,k} - \mathcal{T}_{r,k}\big)\mathcal{T}_{r+h,k}^{-1}.
\end{aligned}
\]
Hence,
\[
\begin{aligned}
& \big\|\mathcal{T}_{r+h,k}^{-1} - \mathcal{T}_{r,k}^{-1}
+ \mathcal{T}_{r,k}^{-1}(\partial_r \mathcal{T}_{r,k}[h])\mathcal{T}_{r,k}^{-1}\big\|_{L^\infty \to L^\infty} \\
&\;\;\lesssim \|\mathcal{T}_{r,k}^{-1}\|_{L^\infty \to L^\infty}\,
   \big\|\,\mathcal{T}_{r+h,k} - \mathcal{T}_{r,k} - \partial_r \mathcal{T}_{r,k}[h]\,\big\|_{L^\infty \to L^\infty}
   \;+\; \|\mathcal{T}_{r,k}^{-1}\|_{L^\infty \to L^\infty}^2\,
   \|\mathcal{T}_{r+h,k}^{-1}\|_{L^\infty \to L^\infty}\,\|h\|_{L^2(\lambda)}^2 .
\end{aligned}
\]
The remainder bound in Lemma \ref{lemma::operatorderivatives} for $\mathcal{P}_{r,k}$ implies
\[
\big\|\,\mathcal{T}_{r+h,k} - \mathcal{T}_{r,k} - \partial_r \mathcal{T}_{r,k}[h]\,\big\|_{L^\infty \to L^\infty} \lesssim \|h\|_{L^2(\lambda)}^2.
\]
Consequently,
\[
\big\|\mathcal{T}_{r+h,k}^{-1} - \mathcal{T}_{r,k}^{-1}
+ \mathcal{T}_{r,k}^{-1}(\partial_r \mathcal{T}_{r,k}[h])\mathcal{T}_{r,k}^{-1}\big\|_{L^\infty \to L^\infty} \lesssim \|h\|_{L^2(\lambda)}^2,
\]
and thus
\[
\big\|\mathcal{T}_{r+h,k}^{-1}-\mathcal{T}_{r,k}^{-1}-\partial_r \mathcal{T}_{r,k}^{-1}[h]\big\|_{L^\infty(\lambda) \to L^2(\lambda)} \,\|r\|_{L^\infty(\lambda)} \lesssim \|h\|_{L^2(\lambda)}^2 \,\|r\|_{L^\infty(\lambda)}.
\]

\medskip
\noindent \textbf{Term 3.}
We have
\[
\|\mathcal{T}_{r+h,k}^{-1} - \mathcal{T}_{r,k}^{-1}\|_{L^\infty \to L^\infty}
\;\le\; \|\mathcal{T}_{r,k}^{-1}\|_{L^\infty \to L^\infty}\,
        \|\mathcal{T}_{r+h,k} - \mathcal{T}_{r,k}\|_{L^\infty \to L^\infty}\,
        \|\mathcal{T}_{r+h,k}^{-1}\|_{L^\infty \to L^\infty}.
\]
Furthermore, by Lemma~\ref{lemma::operatorderivatives},
\[
\big\| \mathcal{P}_{r + h, k}  - \mathcal{P}_{r, k} \big\|_{L^\infty \to L^\infty}   
     \;\lesssim\; \|\partial_r \mathcal{P}_{r, k}[h]\|_{L^\infty \to L^\infty} 
     + \|h\|_{L^2(\lambda)}^2,
\]
and $\|\partial_r \mathcal{P}_{r, k}[h]\|_{L^\infty \to L^\infty} \lesssim \|h\|_{L^2(\lambda)}$. Hence,
\[
\|\mathcal{P}_{r+h,k} - \mathcal{P}_{r,k}\|_{L^\infty \to L^\infty}
\;\lesssim\; \|h\|_{L^2(\lambda)}, 
\qquad
\|\mathcal{T}_{r+h,k} - \mathcal{T}_{r,k}\|_{L^\infty \to L^\infty}
\;\lesssim\; \|h\|_{L^2(\lambda)}.
\]
Therefore,
\[
\|\mathcal{T}_{r+h,k}^{-1} - \mathcal{T}_{r,k}^{-1}\|_{L^\infty \to L^\infty} 
\;\lesssim \|h\|_{L^2(\lambda)},
\]
and
\[
\big\|\big(\mathcal{T}_{r+h,k}^{-1}-\mathcal{T}_{r,k}^{-1}\big)(r)\big\|_{L^2(\lambda)} 
\lesssim \|h\|_{L^2(\lambda)} \,\|r\|_{L^\infty(\lambda)}.
\]

Finally, by Lemma \ref{lemma::softmaxLipschitz},
\[
\big\|\Pi_{r+h,k}-\Pi_{r,k}\big\|_{L^2(\lambda) \to L^1(\mu)} \lesssim \|h\|_{L^2(\lambda)}.
\]
Putting everything together,
\[
\big\|\Pi_{r+h,k}-\Pi_{r,k}\big\|_{L^2(\lambda) \to L^1(\mu)} 
\,\big\|\big(\mathcal{T}_{r+h,k}^{-1}-\mathcal{T}_{r,k}^{-1}\big)(r)\big\|_{L^2(\lambda)} 
\;\lesssim \|h\|_{L^2(\lambda)}^2 \,\|r\|_{L^\infty(\lambda)}.
\]


\medskip
\noindent \textbf{Term 4.}
\begin{align*}
\big\|\big(\Pi_{r+h,k}\,\mathcal{T}_{r+h,k}^{-1}-\Pi_{r,k}\,\mathcal{T}_{r,k}^{-1}\big)(h)\big\|_{L^1(\lambda)}
&= \big\|\big(\Pi_{r+h,k}-\Pi_{r,k}\big)\,\mathcal{T}_{r+h,k}^{-1}(h)
\;+\;\Pi_{r,k}\,\big(\mathcal{T}_{r+h,k}^{-1}-\mathcal{T}_{r,k}^{-1}\big)(h)\big\|_{L^1(\lambda)} \\
&\le \big\|\Pi_{r+h,k}-\Pi_{r,k}\big\|_{L^2(\lambda)\to L^1(\mu)}\;
        \big\|\mathcal{T}_{r+h,k}^{-1}\big\|_{L^2(\lambda)\to L^2(\lambda)}\;\|h\|_{L^2(\lambda)} \\
&\quad\;+\;\big\|\Pi_{r,k}\big\|_{L^2(\lambda)\to L^1(\mu)}\;
        \big\|\mathcal{T}_{r+h,k}^{-1}-\mathcal{T}_{r,k}^{-1}\big\|_{L^2(\lambda)\to L^2(\lambda)}\;\|h\|_{L^2(\lambda)}.
\end{align*}
By Lemma~\ref{lemma::softmaxLipschitz},
$\big\|\Pi_{r+h,k}-\Pi_{r,k}\big\|_{L^2(\lambda)\to L^1(\mu)}\lesssim \|h\|_{L^2}$. Moreover, arguing as in Term~3 and applying Lemma~\ref{lemma::forwardopLipschitz}, we find
$\big\|\mathcal{T}_{r+h,k}^{-1}-\mathcal{T}_{r,k}^{-1}\big\|_{L^2(\lambda)\to L^2(\lambda)}\lesssim \|h\|_{L^2}$,
using that $\|\mathcal{T}_{r+h,k}^{-1}\|_{L^2(\lambda)\to L^2(\lambda)}$ and $\|\Pi_{r,k}\|_{L^2(\lambda)\to L^1(\mu)}$ are uniformly bounded under the standing assumptions.
Therefore,
\[
\big\|\big(\Pi_{r+h,k}\,\mathcal{T}_{r+h,k}^{-1}
- \Pi_{r,k}\,\mathcal{T}_{r,k}^{-1}\big)(h)\big\|_{L^1(\lambda)}
\;\lesssim\; \|h\|_{L^2(\lambda)}^2.
\]
The implicit constant depends on $\|\mathcal{T}_{r+h,k}^{-1}\|_{L^2(\lambda)\to L^2(\lambda)}$. 
However, we showed that
\begin{align*}
\|\mathcal{T}_{r+h,k}^{-1}\|_{L^2(\lambda)\to L^2(\lambda)} 
&\;\lesssim\; \|\mathcal{T}_{r,k}^{-1}\|_{L^2(\lambda)\to L^2(\lambda)} 
+ \|\mathcal{T}_{r+h,k}^{-1} - \mathcal{T}_{r,k}^{-1}\|_{L^2(\lambda)\to L^2(\lambda)} \\
&\;\lesssim\; \|\mathcal{T}_{r,k}^{-1}\|_{L^2(\lambda)\to L^2(\lambda)} 
+ O\!\big(\|h\|_{L^2(\lambda)}^{2}\big).
\end{align*}
Thus,
\[
\big\|\Pi_{r+h,k}\,\mathcal{T}_{r+h,k}^{-1}(r) 
- \Pi_{r,k}\,\mathcal{T}_{r,k}^{-1}(r) 
- \partial_r\!\big[\Pi_{r,k}\,\mathcal{T}_{r,k}^{-1}(r)\big][h]\big\|_{L^1(\lambda)} 
= O\!\left(\|h\|_{L^2(\lambda)}^{2}\right),
\]
where the big-$O$ constant depends on 
$\|\mathcal{T}_{r,k}^{-1}\|_{L^2(\lambda)\to L^2(\lambda)}$ 
but not directly on 
$\|\mathcal{T}_{r+h,k}^{-1}\|_{L^2(\lambda)\to L^2(\lambda)}$.

\end{proof}


 


\begin{lemma}
\label{lemma::valuestarderivativekernel}
Let $r \in L^2(\lambda)$, $k \in \mathcal{K}$, and $g \in \addTK$. 
Suppose the linear operators $\mathcal{T}_{r,k}:L^2(\lambda) \to L^2(\lambda)$ 
and $\mathcal{T}_{r,k+g}:L^2(\lambda) \to L^2(\lambda)$ are invertible. Then
\[
\big\|\Pi_{r,k+g}\,\mathcal{T}_{r,k+g}^{-1}(r) 
- \Pi_{r,k}\,\mathcal{T}_{r,k}^{-1}(r) 
- \adk\!\big[\Pi_{r,k}\,\mathcal{T}_{r,k}^{-1}(r)\big][g]\big\|_{L^1(\lambda)} 
= O\!\left(\|v_{r,k+g} - v_{r,k}\|_{L^2(\lambda)}^{2} + \|g\|_{\mathcal K}^{2}\right)
\]
with the big-$O$ constant depending on $\gamma$, $\|r\|_{L^\infty(\lambda)}$, 
and $\|\mathcal{T}_{r,k}^{-1}\|_{L^2(\lambda)\to L^2(\lambda)}$.
\end{lemma}


\begin{proof}
By the chain rule,
\begin{align*}
\adk\!\big[\Pi_{r,k}\,\mathcal{T}_{r,k}^{-1}(r)\big][g]
&= \big(\adk \Pi_{r,k}[g]\big)\,\mathcal{T}_{r,k}^{-1}(r)
\;+\;\Pi_{r,k}\,\big(\adk \mathcal{T}_{r,k}^{-1}[g]\big)(r),
\end{align*}
where $\adk \mathcal{T}_{r,k}^{-1}[g] = -\,\mathcal{T}_{r,k}^{-1}\big(\adk \mathcal{T}_{r,k}[g]\big)\mathcal{T}_{r,k}^{-1}$.
By adding and subtracting, we obtain
\begin{align*}
&\Pi_{r,k+g}\,\mathcal{T}_{r,k+g}^{-1}(r) - \Pi_{r,k}\,\mathcal{T}_{r,k}^{-1}(r) - \adk\!\big[\Pi_{r,k}\,\mathcal{T}_{r,k}^{-1}(r)\big][g] \\
&= \big(\Pi_{r,k+g}-\Pi_{r,k}-\adk \Pi_{r,k}[g]\big)\,\mathcal{T}_{r,k}^{-1}(r) \\
&\quad+\;\Pi_{r,k}\,\big(\mathcal{T}_{r,k+g}^{-1}-\mathcal{T}_{r,k}^{-1}-\adk \mathcal{T}_{r,k}^{-1}[g]\big)(r) \\
&\quad+\;\big(\Pi_{r,k+g}-\Pi_{r,k}\big)\big(\mathcal{T}_{r,k+g}^{-1}-\mathcal{T}_{r,k}^{-1}\big)(r).
\end{align*}
Taking norms and applying the $k$–direction analogues of Lemma \ref{lemma::operatorderivatives}, we have
\begin{align*}
& \big\|\Pi_{r,k+g}\,\mathcal{T}_{r,k+g}^{-1}(r) - \Pi_{r,k}\,\mathcal{T}_{r,k}^{-1}(r) - \adk\!\big[\Pi_{r,k}\,\mathcal{T}_{r,k}^{-1}(r)\big][g]\big\|_{L^1(\lambda)} \\
&\lesssim \big\|\Pi_{r,k+g}-\Pi_{r,k}-\adk \Pi_{r,k}[g]\big\|_{L^\infty(\lambda) \to L^1(\lambda)} \,\big\|\mathcal{T}_{r,k}^{-1}(r)\big\|_{L^\infty(\lambda)} \\
&\quad+ \big\|\mathcal{T}_{r,k+g}^{-1}-\mathcal{T}_{r,k}^{-1}-\adk \mathcal{T}_{r,k}^{-1}[g]\big\|_{L^\infty(\lambda) \to L^2(\lambda)} \,\|r\|_{L^\infty(\lambda)} \\
&\quad+\;\big\|\Pi_{r,k+g}-\Pi_{r,k}\big\|_{L^2(\lambda) \to L^1(\mu)}\,\big\|\big(\mathcal{T}_{r,k+g}^{-1}-\mathcal{T}_{r,k}^{-1}\big)(r)\big\|_{L^2(\lambda)}.
\end{align*}



\noindent \textbf{Term 1.}
By Lemma \ref{lemma::operatorderivatives}, we have
\begin{align*}
    \|\Pi_{r, k+g} f - \Pi_{r, k} f 
   - \adk \Pi_{r, k}[\,g\,](f)\|_{L^1(\mu)} 
&\;\lesssim\; \|f\|_{L^\infty(\lambda)} \Big\{ 
      \|v_{r, k+g} - v_{r,k}\|_{L^2(\lambda)}^2 \\
&\hspace{6em}
      + \|\mathcal{T}_{r,k}^{-1}\|_{L^2(\lambda) \to L^2(\lambda)} \,
        \|v_{r,k+g} - v_{r,k}\|_{L^2(\lambda)}\,
        \|g\|_{L^2(\mu \otimes \lambda)} \Big\}.
\end{align*}
Hence,
\begin{align*}
 \big\|\Pi_{r,k+g}-\Pi_{r,k}-\adk \Pi_{r,k}[g]\big\|_{L^\infty(\lambda) \to L^1(\lambda)} \,
 \big\|\mathcal{T}_{r,k}^{-1}(r)\big\|_{L^\infty(\lambda)}     
&\lesssim  \|v_{r, k+g} - v_{r,k}\|_{L^2(\lambda)}^2 \\
&\quad+  \|v_{r,k+g} - v_{r,k}\|_{L^2(\lambda)}\,\|g\|_{\mathcal K},
\end{align*}
where the implicit constant depends on $\|\mathcal{T}_{r,k}^{-1}\|_{L^2 \to L^2}$ and $\|r\|_\infty$.






\noindent \textbf{Term 2.}
Using the resolvent identity, the inverse–operator remainder admits the second–order expansion
\[
\begin{aligned}
&\mathcal{T}_{r,k+g}^{-1} - \mathcal{T}_{r,k}^{-1}
+ \mathcal{T}_{r,k}^{-1}\big(\adk \mathcal{T}_{r,k}[g]\big)\mathcal{T}_{r,k}^{-1} \\
&\qquad=- \mathcal{T}_{r,k}^{-1}\Big( \mathcal{T}_{r,k+g} - \mathcal{T}_{r,k} - \adk \mathcal{T}_{r,k}[g] \Big)\mathcal{T}_{r,k}^{-1} \\
&\qquad\quad +\; \mathcal{T}_{r,k}^{-1}\big(\mathcal{T}_{r,k+g} - \mathcal{T}_{r,k}\big)\mathcal{T}_{r,k}^{-1}\big(\mathcal{T}_{r,k+g} - \mathcal{T}_{r,k}\big)\mathcal{T}_{r,k+g}^{-1}.
\end{aligned}
\]
Hence,
\[
\begin{aligned}
& \big\|\mathcal{T}_{r,k+g}^{-1} - \mathcal{T}_{r,k}^{-1}
+ \mathcal{T}_{r,k}^{-1}(\adk \mathcal{T}_{r,k}[g])\mathcal{T}_{r,k}^{-1}\big\|_{L^\infty(\lambda) \to L^2(\lambda)} \\
&\;\;\lesssim \|\mathcal{T}_{r,k}^{-1}\|_{L^2(\lambda) \to L^2(\lambda)}\,
   \big\|\,\mathcal{T}_{r,k+g} - \mathcal{T}_{r,k} - \adk \mathcal{T}_{r,k}[g]\,\big\|_{L^\infty(\lambda) \to L^2(\lambda)} \,\|\mathcal{T}_{r,k}^{-1}\|_{L^2(\lambda) \to L^2(\lambda)} \\
&\hspace{6em}
   +\; \|\mathcal{T}_{r,k}^{-1}\|_{L^2(\lambda) \to L^2(\lambda)}^2\,
   \|\mathcal{T}_{r,k+g}^{-1}\|_{L^2(\lambda) \to L^2(\lambda)}\,
   \big\|\mathcal{T}_{r,k+g} - \mathcal{T}_{r,k}\big\|_{L^\infty(\lambda) \to L^2(\lambda)}^{\,2}.
\end{aligned}
\]
The remainder bound in Lemma \ref{lemma::operatorderivatives} for $\mathcal{P}_{r,k}$ implies
\[
\big\|\,\mathcal{T}_{r,k+g} - \mathcal{T}_{r,k} - \adk \mathcal{T}_{r,k}[g]\,\big\|_{L^\infty(\lambda) \to L^2(\lambda)} 
\;\lesssim\;  \|v_{r, k+g} - v_{r,k}\|_{L^2(\lambda)}^2
+  \|v_{r,k+g} - v_{r,k}\|_{L^2(\lambda)}\,\|g\|_{\mathcal{K}},
\]
and similarly, by Lemma \ref{lemma::forwardopLipschitz},
\[
\big\|\,\mathcal{T}_{r,k+g} - \mathcal{T}_{r,k}\,\big\|_{L^\infty(\lambda) \to L^2(\lambda)} 
\;\lesssim\;  \|v_{r, k+g} - v_{r,k}\|_{L^2(\lambda)}
+  \|g\|_{\mathcal{K}}.
\]
Consequently,
\[
\begin{aligned}
&\big\|\mathcal{T}_{r,k+g}^{-1} - \mathcal{T}_{r,k}^{-1}
+ \mathcal{T}_{r,k}^{-1}(\adk \mathcal{T}_{r,k}[g])\mathcal{T}_{r,k}^{-1}\big\|_{L^\infty(\lambda) \to L^2(\lambda)} \\
&\lesssim\; \|\mathcal{T}_{r,k}^{-1}\|_{L^2(\lambda) \to L^2(\lambda)}^2 \notag\\
&\qquad \times \Big( \|v_{r, k+g} - v_{r,k}\|_{L^2(\lambda)}^2
+  \|v_{r,k+g} - v_{r,k}\|_{L^2(\lambda)}\,\|g\|_{\mathcal{K}} \Big) \\
&\quad +\; \|\mathcal{T}_{r,k}^{-1}\|_{L^2(\lambda) \to L^2(\lambda)}^2
   \,\|\mathcal{T}_{r,k+g}^{-1}\|_{L^2(\lambda) \to L^2(\lambda)} \notag\\
&\qquad \times
   \Big(
      \|v_{r, k+g} - v_{r,k}\|_{L^2(\lambda)}
      + \|g\|_{\mathcal{K}}
   \Big)^{\!2}.
\end{aligned}
\]
Therefore,
\[
\begin{aligned}
\big\|\mathcal{T}_{r,k+g}^{-1}-\mathcal{T}_{r,k}^{-1}
&\qquad - \adk \mathcal{T}_{r,k}^{-1}[g]\big\|_{L^\infty(\lambda) \to L^2(\lambda)} \\
&\qquad \times \|r\|_{L^\infty(\lambda)}
\lesssim\;& \|v_{r, k+g} - v_{r,k}\|_{L^2(\lambda)}^2 \\
&+  \|v_{r,k+g} - v_{r,k}\|_{L^2(\lambda)}\,\|g\|_{\mathcal{K}}.
\end{aligned}
\]

 
\noindent \textbf{Term 3.}
We have
\[
\|\mathcal{T}_{r,k+g}^{-1} - \mathcal{T}_{r,k}^{-1}\|_{L^2(\lambda) \to L^2(\lambda)}
\;\le\; \|\mathcal{T}_{r,k}^{-1}\|_{L^2(\lambda) \to L^2(\lambda)}\,
        \|\mathcal{T}_{r,k+g} - \mathcal{T}_{r,k}\|_{L^2(\lambda) \to L^2(\lambda)}\,
        \|\mathcal{T}_{r,k+g}^{-1}\|_{L^2(\lambda) \to L^2(\lambda)}.
\]
Furthermore, by Lemma~\ref{lemma::operatorderivatives},
\[
\|\mathcal{P}_{r,k+g} - \mathcal{P}_{r,k}\|_{L^2(\lambda) \to L^2(\lambda)} 
\;\lesssim\; \|v_{r,k+g} - v_{r,k}\|_{L^2(\lambda)} 
+ \|g\|_{\mathcal K},
\]
which implies
\[
\|\mathcal{T}_{r,k+g} - \mathcal{T}_{r,k}\|_{L^2(\lambda) \to L^2(\lambda)}
\;\lesssim\; \|v_{r,k+g} - v_{r,k}\|_{L^2(\lambda)} + \|g\|_{\mathcal K}.
\]
Therefore,
\[
\|\mathcal{T}_{r,k+g}^{-1} - \mathcal{T}_{r,k}^{-1}\|_{L^2(\lambda) \to L^2(\lambda)} 
\;\lesssim\; \big(\|v_{r,k+g} - v_{r,k}\|_{L^2(\lambda)} + \|g\|_{\mathcal K}\big).
\]
Consequently,
\[
\big\|\big(\mathcal{T}_{r,k+g}^{-1}-\mathcal{T}_{r,k}^{-1}\big)(r)\big\|_{L^2(\lambda)}
\;\lesssim\; \big(\|v_{r,k+g} - v_{r,k}\|_{L^2(\lambda)} + \|g\|_{\mathcal K}\big)\,
\|r\|_{L^\infty(\lambda)}.
\]

Finally, by Lemma \ref{lemma::softmaxLipschitz},
\[
\big\|\Pi_{r,k+g}-\Pi_{r,k}\big\|_{L^2(\lambda) \to L^1(\mu)} 
\;\lesssim\; \|v_{r,k+g} - v_{r,k}\|_{L^2(\lambda)} + \|g\|_{\mathcal K}.
\]
Putting everything together,
\[
\begin{aligned}
\big\|\Pi_{r,k+g}-\Pi_{r,k}\big\|_{L^2(\lambda) \to L^1(\mu)}
&\times \big\|\big(\mathcal{T}_{r,k+g}^{-1}-\mathcal{T}_{r,k}^{-1}\big)(r)\big\|_{L^2(\lambda)} \\
&\lesssim\; \big(\|v_{r,k+g} - v_{r,k}\|_{L^2(\lambda)} + \|g\|_{\mathcal K}\big)^2\,
\|r\|_{L^\infty(\lambda)}.
\end{aligned}
\]


\noindent \textbf{Final bound.}  
Combining the above estimates, we obtain
\[
\begin{aligned}
\big\|\Pi_{r,k+g}\,\mathcal{T}_{r,k+g}^{-1}(r) 
- \Pi_{r,k}\,\mathcal{T}_{r,k}^{-1}(r) 
- \adk\!\big[\Pi_{r,k}\,\mathcal{T}_{r,k}^{-1}(r)\big][g]\big\|_{L^1(\lambda)}
&\;\lesssim\; \big(\|v_{r,k+g} - v_{r,k}\|_{L^2(\lambda)} + \|g\|_{\mathcal K}\big)^{2}\,
\|r\|_{L^\infty(\lambda)}.
\end{aligned}
\]
Equivalently (absorbing $\|r\|_{L^\infty(\lambda)}$ into the implicit constant),
\[
\begin{aligned}
\big\|\Pi_{r,k+g}\,\mathcal{T}_{r,k+g}^{-1}(r) 
- \Pi_{r,k}\,\mathcal{T}_{r,k}^{-1}(r) 
- \adk\!\big[\Pi_{r,k}\,\mathcal{T}_{r,k}^{-1}(r)\big][g]\big\|_{L^1(\lambda)}
&\;\lesssim\; \|v_{r,k+g} - v_{r,k}\|_{L^2(\lambda)}^{2} + \|g\|_{\mathcal K}^{2}.
\end{aligned}
\]
The implicit constant depends on $\|\mathcal{T}_{r,k+g}^{-1}\|_{L^2(\lambda)\to L^2(\lambda)}$. 
However, we showed that
\begin{align*}
\|\mathcal{T}_{r,k+g}^{-1}\|_{L^2(\lambda)\to L^2(\lambda)} 
&\;\lesssim\; \|\mathcal{T}_{r,k}^{-1}\|_{L^2(\lambda)\to L^2(\lambda)} 
+ \|\mathcal{T}_{r,k+ g}^{-1} - \mathcal{T}_{r,k}^{-1}\|_{L^2(\lambda)\to L^2(\lambda)} \\
&\;\lesssim\; \|\mathcal{T}_{r,k}^{-1}\|_{L^2(\lambda)\to L^2(\lambda)} 
+ O\!\big(\|v_{r,k+g} - v_{r,k}\|_{L^2(\lambda)}^{2} + \|g\|_{\mathcal K}^{2}\big).
\end{align*}
Thus,
\[
\begin{aligned}
\big\|\Pi_{r,k+g}\,\mathcal{T}_{r,k+g}^{-1}(r) 
- \Pi_{r,k}\,\mathcal{T}_{r,k}^{-1}(r) 
- \adk\!\big[\Pi_{r,k}\,\mathcal{T}_{r,k}^{-1}(r)\big][g]\big\|_{L^1(\lambda)}
&= O\!\left(\|v_{r,k+g} - v_{r,k}\|_{L^2(\lambda)}^{2} + \|g\|_{\mathcal K}^{2}\right),
\end{aligned}
\]
where the big-$O$ constant depends on 
$\|\mathcal{T}_{r,k}^{-1}\|_{L^2(\lambda)\to L^2(\lambda)}$ 
but not directly on 
$\|\mathcal{T}_{r,k+g}^{-1}\|_{L^2(\lambda)\to L^2(\lambda)}$.

 

\end{proof}



\begin{lemma}
\label{lemma::jointderivative}
Let $r \in L^2(\lambda)$, $k \in \mathcal{K}$, 
$h \in L^2(\lambda)$, and $g \in \addTK$. 
Suppose the linear operator
$\mathcal{T}_{r,k}:L^2(\lambda) \to L^2(\lambda)$ is invertible and $\|h\|_{L^2(\lambda)}+\|v_{r,k+g}-v_{r,k}\|_{L^2(\lambda)}+\|g\|_{\mathcal K} \;\to\; 0.$
Then
\begin{align*}
&\big\|\Pi_{r+h,k+g}\,\mathcal{T}_{r+h,k+g}^{-1}(r+h) 
- \Pi_{r,k}\,\mathcal{T}_{r,k}^{-1}(r) 
- \partial_r\!\big[\Pi_{r,k}\,\mathcal{T}_{r,k}^{-1}(r)\big][h]
- \adk\!\big[\Pi_{r,k}\,\mathcal{T}_{r,k}^{-1}(r)\big][g]\big\|_{L^1(\lambda)} \\
&\;\;=\; O\!\left(\|h\|_{L^2(\lambda)}^{2} 
+ \|v_{r,k+g}-v_{r,k}\|_{L^2(\lambda)}^{2} 
+ \|g\|_{\mathcal K}^{2}\right),
\end{align*}
with the big-$O$ constant depending on $\gamma$, $\|r\|_{L^\infty(\lambda)}$, 
and $\|\mathcal{T}_{r,k}^{-1}\|_{L^2(\lambda)\to L^2(\lambda)}$; moreover, the inverses appearing above exist for all sufficiently small perturbations.
\end{lemma}
\begin{proof}
\noindent \textbf{Invertibility of the perturbed operators.}
We first show that $\mathcal{T}_{r+h,k}:L^2(\lambda)\to L^2(\lambda)$ and 
$\mathcal{T}_{r+h,k+g}:L^2(\lambda)\to L^2(\lambda)$ are invertible whenever 
$O\!\left(\|h\|_{L^2(\lambda)}^{2}+\|v_{r,k+g}-v_{r,k}\|_{L^2(\lambda)}^{2}+\|g\|_{\mathcal K}^{2}\right)$ 
is sufficiently small.  

Since $\mathcal{T}_{r,k}$ is invertible, it suffices to control the perturbations
\[
\Delta_1:=\mathcal{T}_{r+h,k}-\mathcal{T}_{r,k}, 
\qquad 
\Delta_2:=\mathcal{T}_{r+h,k+g}-\mathcal{T}_{r,k}.
\]
By Lemma \ref{lemma::forwardopLipschitz},
\begin{align*}
\|\Delta_1\|_{L^2(\lambda)\to L^2(\lambda)} &= O\!\big(\|h\|_{L^2(\lambda)}\big),\\
\|\Delta_2\|_{L^2(\lambda)\to L^2(\lambda)} &= O\!\Big(\|h\|_{L^2(\lambda)}+\|v_{r,k+g}-v_{r,k}\|_{L^2(\lambda)}+\|g\|_{\mathcal K}\Big).
\end{align*}
If $\|\mathcal{T}_{r,k}^{-1}\|_{L^2(\lambda)\to L^2(\lambda)}\,\|\Delta_i\|_{L^2(\lambda)\to L^2(\lambda)}<1$ for $i=1,2$,
then by the Neumann–series inversion lemma, both $\mathcal{T}_{r+h,k}$ and 
$\mathcal{T}_{r+h,k+g}$ are invertible, with
\begin{align*}
    \|\mathcal{T}_{r+h,k}^{-1}\|_{L^2(\lambda)\to L^2(\lambda)}
\le \frac{\|\mathcal{T}_{r,k}^{-1}\|_{L^2(\lambda)\to L^2(\lambda)}}{1-\|\mathcal{T}_{r,k}^{-1}\|_{L^2(\lambda)\to L^2(\lambda)}\,\|\Delta_1\|_{L^2(\lambda)\to L^2(\lambda)}},\\
\|\mathcal{T}_{r+h,k+g}^{-1}\|_{L^2(\lambda)\to L^2(\lambda)}
\le \frac{\|\mathcal{T}_{r,k}^{-1}\|_{L^2(\lambda)\to L^2(\lambda)}}{1-\|\mathcal{T}_{r,k}^{-1}\|_{L^2(\lambda)\to L^2(\lambda)}\,\|\Delta_2\|_{L^2(\lambda)\to L^2(\lambda)}}.
\end{align*}
As $\|h\|_{L^2(\lambda)}+\|v_{r,k+g}-v_{r,k}\|_{L^2(\lambda)}+\|g\|_{\mathcal K} \rightarrow 0$, it holds that $\|\mathcal{T}_{r,k}^{-1}\|_{L^2(\lambda)\to L^2(\lambda)}\,\|\Delta_i\|_{L^2(\lambda)\to L^2(\lambda)}<1$ for $i=1,2$. 
Hence,  $\mathcal{T}_{r+h,k}:L^2(\lambda)\to L^2(\lambda)$ and 
$\mathcal{T}_{r+h,k+g}:L^2(\lambda)\to L^2(\lambda)$ become invertible.
Without loss of generality, we proceed assuming these inverses exist.
 
 

\noindent \textbf{Proof of bound.} Lemma \ref{lemma::valuestarderivativereward} establishes that
\[
\big\|\Pi_{r+h,k}\,\mathcal{T}_{r+h,k}^{-1}(r+h) 
 - \Pi_{r,k}\,\mathcal{T}_{r,k}^{-1}(r) 
 - \partial_r\!\big[\Pi_{r,k}\,\mathcal{T}_{r,k}^{-1}(r)\big][h]\big\|_{L^1(\lambda)} 
= O(\|h\|_{L^2(\lambda)}^2).
\]
Lemma \ref{lemma::valuestarderivativekernel} establishes that
\[
\big\|\Pi_{r,k+g}\,\mathcal{T}_{r,k+g}^{-1}(r) 
- \Pi_{r,k}\,\mathcal{T}_{r,k}^{-1}(r) 
- \adk\!\big[\Pi_{r,k}\,\mathcal{T}_{r,k}^{-1}(r)\big][g]\big\|_{L^1(\lambda)} 
= O\!\left(\|v_{r,k+g} - v_{r,k}\|_{L^2(\lambda)}^{2} + \|g\|_{\mathcal K}^{2}\right).
\]
Therefore, adding and subtracting,
\begin{align*}
&\big\|\Pi_{r+h,k+g}\,\mathcal{T}_{r+h,k+g}^{-1}(r+h) 
- \Pi_{r,k}\,\mathcal{T}_{r,k}^{-1}(r) 
- \partial_r\!\big[\Pi_{r,k}\,\mathcal{T}_{r,k}^{-1}(r)\big][h]
- \adk\!\big[\Pi_{r,k}\,\mathcal{T}_{r,k}^{-1}(r)\big][g]\big\|_{L^1(\lambda)} \\
&\;\;\le 
\big\|\Pi_{r+h,k+g}\,\mathcal{T}_{r+h,k+g}^{-1}(r+h) 
      - \Pi_{r+h,k}\,\mathcal{T}_{r+h,k}^{-1}(r+h)
      - \adk\!\big[\Pi_{r+h,k}\,\mathcal{T}_{r+h,k}^{-1}(r+h)\big][g]\big\|_{L^1(\lambda)} \\
&\qquad + 
\big\|\Pi_{r+h,k}\,\mathcal{T}_{r+h,k}^{-1}(r+h) 
      - \Pi_{r,k}\,\mathcal{T}_{r,k}^{-1}(r) 
      - \partial_r\!\big[\Pi_{r,k}\,\mathcal{T}_{r,k}^{-1}(r)\big][h]\big\|_{L^1(\lambda)} \\
&\qquad + 
\big\|\big(\adk\!\big[\Pi_{r+h,k}\,\mathcal{T}_{r+h,k}^{-1}(r+h)\big]
      - \adk\!\big[\Pi_{r,k}\,\mathcal{T}_{r,k}^{-1}(r)\big]\big)[g]\big\|_{L^1(\lambda)}.
\end{align*}
Applying the lemmas, we obtain
\begin{align*}
&\big\|\Pi_{r+h,k+g}\,\mathcal{T}_{r+h,k+g}^{-1}(r+h) 
- \Pi_{r,k}\,\mathcal{T}_{r,k}^{-1}(r) 
- \partial_r\!\big[\Pi_{r,k}\,\mathcal{T}_{r,k}^{-1}(r)\big][h]
- \adk\!\big[\Pi_{r,k}\,\mathcal{T}_{r,k}^{-1}(r)\big][g]\big\|_{L^1(\lambda)} \\
&\;\;\le 
O\!\left(\|v_{r,k+g} - v_{r,k}\|_{L^2(\lambda)}^{2} + \|g\|_{\mathcal K}^{2}\right) 
+ O\!\left(\|h\|_{L^2(\lambda)}^2\right) \\
&\qquad + 
\big\|\big(\adk\!\big[\Pi_{r+h,k}\,\mathcal{T}_{r+h,k}^{-1}(r+h)\big]
      - \adk\!\big[\Pi_{r,k}\,\mathcal{T}_{r,k}^{-1}(r)\big]\big)[g]\big\|_{L^1(\lambda)}.
\end{align*}
We claim that 
\[
\big\|\big(\adk\![\Pi_{r+h,k}\mathcal{T}_{r+h,k}^{-1}(r{+}h)]
      - \adk\![\Pi_{r,k}\mathcal{T}_{r,k}^{-1}(r)]\big)[g]\big\|_{L^1(\lambda)}
= O\!\left(\|h\|_{L^2(\lambda)}^2 + \|v_{r,k+g}-v_{r,k}\|_{L^2(\lambda)}^2 + \|g\|_{\mathcal K}^2\right).
\]
Assuming this claim, it follows that
\begin{align*}
&\big\|\Pi_{r+h,k+g}\,\mathcal{T}_{r+h,k+g}^{-1}(r+h) 
- \Pi_{r,k}\,\mathcal{T}_{r,k}^{-1}(r) 
- \partial_r\!\big[\Pi_{r,k}\,\mathcal{T}_{r,k}^{-1}(r)\big][h]
- \adk\!\big[\Pi_{r,k}\,\mathcal{T}_{r,k}^{-1}(r)\big][g]\big\|_{L^1(\lambda)} \\
&\;\;=\; O\!\left(\|h\|_{L^2(\lambda)}^2 + \|v_{r,k+g}-v_{r,k}\|_{L^2(\lambda)}^2 + \|g\|_{\mathcal K}^2\right),
\end{align*}
as desired.


\noindent \textbf{Proof of claim.}
We bound
\[
\big\|\big(\adk\!\big[\Pi_{r+h,k}\,\mathcal{T}_{r+h,k}^{-1}(r+h)\big]
      - \adk\!\big[\Pi_{r,k}\,\mathcal{T}_{r,k}^{-1}(r)\big]\big)[g]\big\|_{L^1(\lambda)}.
\]
By the product rule,
\[
\adk\!\big[\Pi_{r,k}\,\mathcal{T}_{r,k}^{-1}(r)\big][g]
= \big(\adk \Pi_{r,k}[g]\big)\,\mathcal{T}_{r,k}^{-1}(r)
\;+\;\Pi_{r,k}\,\big(\adk \mathcal{T}_{r,k}^{-1}[g]\big)(r),
\]
and similarly at $(r{+}h,k)$. Hence, adding and subtracting intermediate terms,
\begin{align*}
&\big(\adk\![\Pi_{r+h,k}\mathcal{T}_{r+h,k}^{-1}(r{+}h)]
      - \adk\![\Pi_{r,k}\mathcal{T}_{r,k}^{-1}(r)]\big)[g] \\
&= \underbrace{\big(\adk \Pi_{r+h,k}[g]-\adk \Pi_{r,k}[g]\big)\,\mathcal{T}_{r+h,k}^{-1}(r{+}h)}_{A_1}
 \;+\; \underbrace{\adk \Pi_{r,k}[g]\;\big(\mathcal{T}_{r+h,k}^{-1}(r{+}h)-\mathcal{T}_{r,k}^{-1}(r)\big)}_{A_2} \\
&\quad+\;\underbrace{(\Pi_{r+h,k}-\Pi_{r,k})\,\big(\adk \mathcal{T}_{r+h,k}^{-1}[g]\big)(r{+}h)}_{A_3}
 \;+\; \underbrace{\Pi_{r,k}\,\Big(\big(\adk \mathcal{T}_{r+h,k}^{-1}[g]\big)-\big(\adk \mathcal{T}_{r,k}^{-1}[g]\big)\Big)(r{+}h)}_{A_4} \\
&\quad+\;\underbrace{\Pi_{r,k}\,\big(\adk \mathcal{T}_{r,k}^{-1}[g]\big)(h)}_{A_5}.
\end{align*}

\medskip
\noindent \textbf{Bounds for $A_1$--$A_5$.}
Using Lemma~\ref{lemma::operatorderivatives}, Lemma~\ref{lemma::forwardopLipschitz}, and the bounds derived in the proofs of Lemmas~\ref{lemma::valuestarderivativereward} and~\ref{lemma::valuestarderivativekernel}, we obtain:
\begin{align*}
\|A_1\|_{L^1(\lambda)}
&\;\lesssim\; \|\adk \Pi_{r+h,k}[g]-\adk \Pi_{r,k}[g]\|_{L^2(\lambda)\to L^1(\mu)}\,
               \|\mathcal{T}_{r+h,k}^{-1}(r{+}h)\|_{L^2(\lambda)} \\
&= O\!\left(\|h\|_{L^2(\lambda)}\Big(\|v_{r,k+g}-v_{r,k}\|_{L^2(\lambda)}+\|g\|_{\mathcal K}\Big)\right), \\[0.8ex]
%
\|A_2\|_{L^1(\lambda)}
&\;\lesssim\; \|\adk \Pi_{r,k}[g]\|_{L^2(\lambda)\to L^1(\mu)}\,
               \|\mathcal{T}_{r+h,k}^{-1}(r{+}h)-\mathcal{T}_{r,k}^{-1}(r)\|_{L^2(\lambda)} \\
&= O\!\left(\|h\|_{L^2(\lambda)}\Big(\|v_{r,k+g}-v_{r,k}\|_{L^2(\lambda)}+\|g\|_{\mathcal K}\Big)\right), \\[0.8ex]
%
\|A_3\|_{L^1(\lambda)}
&\;\lesssim\; \|\Pi_{r+h,k}-\Pi_{r,k}\|_{L^2(\lambda)\to L^1(\mu)}\,
               \|\adk \mathcal{T}_{r+h,k}^{-1}[g]\|_{L^2(\lambda)\to L^2(\lambda)}\,\|r{+}h\|_{L^\infty(\lambda)} \\
&= O\!\left(\|h\|_{L^2(\lambda)}\Big(\|v_{r,k+g}-v_{r,k}\|_{L^2(\lambda)}+\|g\|_{\mathcal K}\Big)\right), \\[0.8ex]
%
\|A_4\|_{L^1(\lambda)}
&\;\lesssim\; \|\Pi_{r,k}\|_{L^2(\lambda)\to L^1(\mu)}\,
               \|\adk \mathcal{T}_{r+h,k}^{-1}[g]-\adk \mathcal{T}_{r,k}^{-1}[g]\|_{L^2(\lambda)\to L^2(\lambda)}\,\|r{+}h\|_{L^\infty(\lambda)} \\
&= O\!\left(\|h\|_{L^2(\lambda)}\Big(\|v_{r,k+g}-v_{r,k}\|_{L^2(\lambda)}+\|g\|_{\mathcal K}\Big)\right), \\[0.8ex]
%
\|A_5\|_{L^1(\lambda)}
&\;\lesssim\; \|\Pi_{r,k}\|_{L^2(\lambda)\to L^1(\mu)}\,\|\adk \mathcal{T}_{r,k}^{-1}[g]\|_{L^2(\lambda)\to L^2(\lambda)}\,\|h\|_{L^2(\lambda)} \\
&= O\!\left(\|h\|_{L^2(\lambda)}\Big(\|v_{r,k+g}-v_{r,k}\|_{L^2(\lambda)}+\|g\|_{\mathcal K}\Big)\right).
\end{align*}

Combining the bounds and using that all operator norms involved are uniformly bounded on the neighborhood, we conclude
\[
\big\|\big(\adk\![\Pi_{r+h,k}\mathcal{T}_{r+h,k}^{-1}(r{+}h)]
      - \adk\![\Pi_{r,k}\mathcal{T}_{r,k}^{-1}(r)]\big)[g]\big\|_{L^1(\lambda)}
= O\!\left(\|h\|_{L^2(\lambda)}^2 + \|v_{r,k+g}-v_{r,k}\|_{L^2(\lambda)}^2 + \|g\|_{\mathcal K}^2\right).
\]

 


\end{proof}
 

\newpage 

\newpage 

 
\begin{lemma}[Second-order partial Fréchet differentiability of $r \mapsto v_{r,k}$]
\label{lemma::valuederivsecond}
Suppose $\gamma < 1$. 
Then the map $r \mapsto \partial_r v_{r,k} = \mathcal{T}_{r,k}^{-1}\mathcal{P}_{r,k}$
is Fréchet differentiable as a map from $L^\infty(\lambda)$ to $L^\infty(\lambda)$, with derivative
\[
\partial_r^2 v_{r,k}[h_1,h_2] := \partial_r\!\big[\mathcal{T}_{r,k}^{-1}\mathcal{P}_{r,k}\big][h_1](h_2)
= \frac{1}{\tau}\,\mathcal{T}_{r,k}^{-1}\,\mathcal{P}_{r,k}\!\left[
\operatorname{Cov}_{\pi_{r,k}}\!\Big(\mathcal{T}_{r,k}^{-1}h_1,\, \mathcal{T}_{r,k}^{-1}h_2
\Big)\right],
\]
where $\partial_r q_{r,k}[h](\cdot,s) := h(\cdot,s) + \gamma\,\partial_r v_{r,k}[h](\cdot,s)$. 
In particular, we have the quadratic expansion
\begin{align*}
 \|\partial_{r} v_{r+h,k} - \partial_{r} v_{r,k} - \partial_r^2 v_{r,k}[\cdot, h]\|_{L^\infty \to L^\infty} \;\lesssim\; \|h\|_{L^2(\lambda)}^2 .
\end{align*}
\end{lemma}
\begin{proof}
By the chain rule, we have that \( r \mapsto \mathcal{T}_{r,k} \) is Gateaux differentiable with derivative operator
\[
\partial_r \mathcal{T}_{r,k}[h] \;=\; -\,\gamma\,\partial_r \mathcal{P}_{r,k}[h],
\qquad\text{i.e.}\quad
\big(\partial_r \mathcal{T}_{r,k}[h]\big)(f)
= -\,\gamma\,\big(\partial_r \mathcal{P}_{r,k}[h]\big)(f).
\]
Applying the inverse–map differentiation formula, we define the derivative of the inverse $\mathcal{T}_{r,k}^{-1}$ as
\[
\partial_r \mathcal{T}_{r,k}^{-1}[h]
= -\,\mathcal{T}_{r,k}^{-1}\,\big(\partial_r \mathcal{T}_{r,k}[h]\big)\,\mathcal{T}_{r,k}^{-1}.
\]
Applying the product rule and differentiability of
\(r\mapsto \mathcal{T}_{r,k}\) and \(r\mapsto \mathcal{P}_{r,k}\), we obtain that
\[
\partial_r\!\big[\mathcal{T}_{r,k}^{-1}\mathcal{P}_{r,k}\big][h]
= - \mathcal{T}_{r,k}^{-1}\,\big(\partial_r \mathcal{T}_{r,k}[h]\big)\,\mathcal{T}_{r,k}^{-1}\,\mathcal{P}_{r,k}
\;+\;
\mathcal{T}_{r,k}^{-1}\,\big(\partial_r \mathcal{P}_{r,k}[h]\big).
\]
Later, we will show that $r \mapsto \partial_r v_{r,k}$ is Fr\'echet differentiable, in an appropriate sense, with the second derivative given above. Before doing so, we simplify the expression for this second derivative as follows:
\begin{align*}
\partial_r^2 v_{r,k}[h_1,h_2]
&= -\,\mathcal{T}_{r,k}^{-1}\,(\partial_r \mathcal{T}_{r,k}[h_2])\,\mathcal{T}_{r,k}^{-1}\,\mathcal{P}_{r,k}[h_1]
\;+\;
\mathcal{T}_{r,k}^{-1}\,(\partial_r \mathcal{P}_{r,k}[h_2])[h_1] \\
&= \gamma\,\mathcal{T}_{r,k}^{-1}\,(\partial_r \mathcal{P}_{r,k}[h_2])\,\mathcal{T}_{r,k}^{-1}\,\mathcal{P}_{r,k}[h_1]
\;+\;
\mathcal{T}_{r,k}^{-1}\,(\partial_r \mathcal{P}_{r,k}[h_2])[h_1] \\
&= \mathcal{T}_{r,k}^{-1}\,(\partial_r \mathcal{P}_{r,k}[h_2])\,
\big(I+\gamma\,\mathcal{T}_{r,k}^{-1}\mathcal{P}_{r,k}\big)[h_1]\\
&= \mathcal{T}_{r,k}^{-1}\,(\partial_r \mathcal{P}_{r,k}[h_2])\,
\big(I+\gamma\,\partial_r v_{r,k}\big)[h_1]\\
&= \frac{1}{\tau}\,\mathcal{T}_{r,k}^{-1}\!\Big[(a,s)\mapsto 
\int \operatorname{Cov}_{\pi_{r,k}(\cdot\mid s')}\!\Big(
\partial_r q_{r,k}[h_1](\cdot,s'),\; \partial_r q_{r,k}[h_2](\cdot,s')
\Big)\,k(s'\mid a,s) \mu(ds')\Big],
\end{align*}
$\partial_r q_{r,k}[h](\cdot,s):=h(\cdot,s)+\gamma\,\partial_r v_{r,k}[h](\cdot,s).$
We can write:
\[
\partial_r^2 v_{r,k}[h_1,h_2]
= \frac{1}{\tau}\,\mathcal{T}_{r,k}^{-1}\,\mathcal{P}_{r,k}\!\left[(a',s')\mapsto 
\operatorname{Cov}_{\pi_{r,k}(\cdot\mid s')}\!\Big(
\partial_r q_{r,k}[h_1](\cdot,s'),\; \partial_r q_{r,k}[h_2](\cdot,s')
\Big)\right].
\]
Finally, by Lemma \ref{lemma::qderiv}, \(\partial_r q_{r,k}[h] = \mathcal{T}_{r,k}^{-1}h\). 
Hence,
\[
\partial_r^2 v_{r,k}[h_1,h_2]
= \frac{1}{\tau}\,\mathcal{T}_{r,k}^{-1}\,\mathcal{P}_{r,k}\!\left[
\operatorname{Cov}_{\pi_{r,k}}\!\Big(\mathcal{T}_{r,k}^{-1}h_1,\, \mathcal{T}_{r,k}^{-1}h_2
\Big)\right].
\]

 
Now, we establish a quadratic bound for the remainder of the first-order Taylor expansion of $r \mapsto \mathcal{T}_{r,k}^{-1}\mathcal{P}_{r,k}$; namely,
$\mathcal{T}_{r+h,k}^{-1}\mathcal{P}_{r+h,k}
 - \mathcal{T}_{r,k}^{-1}\mathcal{P}_{r,k}
 - \partial_r\!\big[\mathcal{T}_{r,k}^{-1}\mathcal{P}_{r,k}\big][h]$. Adding and subtracting 
$\mathcal{T}_{r+h,k}^{-1}\,\mathcal{P}_{r,k}$ 
and 
$\mathcal{T}_{r,k}^{-1}\,\mathcal{P}_{r+h,k}$, 
we obtain
\begin{align*}
&\mathcal{T}_{r+ h,k}^{-1}\mathcal{P}_{r+ h,k}
 - \mathcal{T}_{r,k}^{-1}\mathcal{P}_{r,k}
 - \partial_r\!\big[\mathcal{T}_{r,k}^{-1}\mathcal{P}_{r,k}\big][h] \\
&= \Big(\mathcal{T}_{r+h,k}^{-1} - \mathcal{T}_{r,k}^{-1}\Big)\mathcal{P}_{r,k}
   + \mathcal{T}_{r,k}^{-1}\Big(\mathcal{P}_{r+h,k} - \mathcal{P}_{r,k}\Big) \\
&\quad + \Big(\mathcal{T}_{r+h,k}^{-1} - \mathcal{T}_{r,k}^{-1}\Big)
       \Big(\mathcal{P}_{r+h,k} - \mathcal{P}_{r,k}\Big) \\
&\quad + \mathcal{T}_{r,k}^{-1}\mathcal{P}_{r,k}
     - \mathcal{T}_{r,k}^{-1}\mathcal{P}_{r,k}
     - \partial_r\!\big[\mathcal{T}_{r,k}^{-1}\mathcal{P}_{r,k}\big][h].
\end{align*}


Thus,
\begin{align*}
&\mathcal{T}_{r+ h,k}^{-1}\mathcal{P}_{r+ h,k}
 - \mathcal{T}_{r,k}^{-1}\mathcal{P}_{r,k}
 - \partial_r\!\big[\mathcal{T}_{r,k}^{-1}\mathcal{P}_{r,k}\big][h] \\
&= \underbrace{\Big[\mathcal{T}_{r+h,k}^{-1} - \mathcal{T}_{r,k}^{-1}
         + \mathcal{T}_{r,k}^{-1}\,\big(\partial_r \mathcal{T}_{r,k}[h]\big)\,
           \mathcal{T}_{r,k}^{-1}\Big]\mathcal{P}_{r,k}}_{\text{Inverse-operator remainder}} \\
&\quad + \underbrace{\mathcal{T}_{r,k}^{-1}\Big[\mathcal{P}_{r+h,k} - \mathcal{P}_{r,k}
           - \partial_r \mathcal{P}_{r,k}[h]\Big]}_{\text{Direct $\mathcal{P}$ remainder}} \\
&\quad + \underbrace{\Big(\mathcal{T}_{r+h,k}^{-1} - \mathcal{T}_{r,k}^{-1}\Big)
       \Big(\mathcal{P}_{r+h,k} - \mathcal{P}_{r,k}\Big)}_{\text{Mixed interaction term}}.
\end{align*}


We analyze each term, beginning with the direct $\mathcal{P}$ remainder and the mixed interaction term, and then returning to the inverse-operator remainder. We will use that the Bellman operators $\mathcal{T}_{r+h,k}$ and $\mathcal{T}_{r,k}$ admit bounded inverses on $L^\infty(\lambda)$ since $\gamma < 1$.

\noindent \textbf{For the direct $\mathcal{P}$ remainder,} by Lemma~\ref{lemma::operatorderivatives} and the definition of operator norms, we obtain
\begin{align*}
\big\|\mathcal{T}_{r,k}^{-1}\big[\mathcal{P}_{r+h,k} - \mathcal{P}_{r,k}
           - \partial_r \mathcal{P}_{r,k}[h]\big]\big\|_{L^\infty \to L^\infty}
&\leq \|\mathcal{T}_{r,k}^{-1}\|_{L^\infty \to L^\infty}\;
   \|\mathcal{P}_{r+h,k} - \mathcal{P}_{r,k}
           - \partial_r \mathcal{P}_{r,k}[h]\|_{L^\infty \to L^\infty} \\
&\lesssim 
   \|h\|_{L^2(\lambda)}^2,
\end{align*}
where we used that $ \|\mathcal{T}_{r,k}^{-1}\|_{L^\infty \to L^\infty} < \infty$ since $\gamma < 1$.


\noindent \textbf{For the mixed interaction term,} we note that
\begin{align*}
   \big\|\big(\mathcal{T}_{r+h,k}^{-1} - \mathcal{T}_{r,k}^{-1}\big)
       \big(\mathcal{P}_{r+h,k} - \mathcal{P}_{r,k}\big)\big\|_{L^\infty \to L^\infty} 
   &\leq  \|\mathcal{T}_{r+h,k}^{-1} - \mathcal{T}_{r,k}^{-1} \|_{L^\infty \to L^\infty}\;
       \|\mathcal{P}_{r+h,k} - \mathcal{P}_{r,k}\|_{L^\infty \to L^\infty}.
\end{align*}
Moreover,
\[
\|\mathcal{T}_{r+h,k}^{-1} - \mathcal{T}_{r,k}^{-1}\|_{L^\infty \to L^\infty}
\;\le\; \|\mathcal{T}_{r,k}^{-1}\|_{L^\infty \to L^\infty}\,
        \|\mathcal{T}_{r+h,k} - \mathcal{T}_{r,k}\|_{L^\infty \to L^\infty}\,
        \|\mathcal{T}_{r+h,k}^{-1}\|_{L^\infty \to L^\infty}.
\]
 Furthermore, by Lemma~\ref{lemma::operatorderivatives}, we have
\begin{align*}
    \big\| \mathcal{P}_{r + h, k}  - \mathcal{P}_{r, k}  
     \big\|_{L^\infty \to L^\infty}   
     \;\lesssim\; \|\partial_r \mathcal{P}_{r, k}[h]\|_{L^\infty \to L^\infty} 
     + \|h\|_{L^2(\lambda)}^2.
\end{align*}
Moreover, $\|\partial_r \mathcal{P}_{r, k}[h]\|_{L^\infty \to L^\infty} \lesssim \|h\|_{L^2(\lambda)}$. Applying this bound, we obtain
\[
\|\mathcal{P}_{r+h,k} - \mathcal{P}_{r,k}\|_{L^\infty \to L^\infty}
\;\lesssim\; \|h\|_{L^2(\lambda)}, 
\qquad
\|\mathcal{T}_{r+h,k} - \mathcal{T}_{r,k}\|_{L^\infty \to L^\infty}
\;\lesssim\; \|h\|_{L^2(\lambda)}.
\]
Hence,
\begin{align*}
   \big\|\big(\mathcal{T}_{r+h,k}^{-1} - \mathcal{T}_{r,k}^{-1}\big)
       \big(\mathcal{P}_{r+h,k} - \mathcal{P}_{r,k}\big)\big\|_{L^\infty \to L^\infty} 
   &\;\lesssim\; \big(1 + \|\mathcal{T}_{r,k}^{-1}\|_{L^\infty \to L^\infty}\big)\,
      \|h\|_{L^2(\lambda)}^2\\
       &\;\lesssim \|h\|_{L^2(\lambda)}^2.
\end{align*}



\noindent \textbf{We now turn to the inverse-operator remainder.} Using the resolvent identity, the inverse–operator remainder can itself be expressed using the second–order expansion:
\[
\begin{aligned}
&\mathcal{T}_{r+h,k}^{-1} - \mathcal{T}_{r,k}^{-1}
  + \mathcal{T}_{r,k}^{-1}\big(\partial_r \mathcal{T}_{r,k}[h]\big)\mathcal{T}_{r,k}^{-1} \\
&\qquad=-  \mathcal{T}_{r,k}^{-1}\Big( \mathcal{T}_{r+h,k} - \mathcal{T}_{r,k}  - \partial_r \mathcal{T}_{r,k}[h] \Big)\mathcal{T}_{r,k}^{-1} \\
&\qquad\quad +\; \mathcal{T}_{r,k}^{-1}\big(\mathcal{T}_{r+h,k} - \mathcal{T}_{r,k}\big)\mathcal{T}_{r,k}^{-1}\big(\mathcal{T}_{r+h,k} - \mathcal{T}_{r,k}\big)\mathcal{T}_{r+h,k}^{-1}.
\end{aligned}
\]
Hence, taking operator norms, we have that
\[
\begin{aligned}
& \big\|\mathcal{T}_{r+h,k}^{-1} - \mathcal{T}_{r,k}^{-1}
+ \mathcal{T}_{r,k}^{-1}(\partial_r \mathcal{T}_{r,k}[h])\mathcal{T}_{r,k}^{-1}\big\|_{L^\infty \to L^\infty} \\
&\;\;\le  \|\mathcal{T}_{r,k}^{-1}\|_{L^\infty \to L^\infty}\,
   \big\|\,\mathcal{T}_{r+h,k} - \mathcal{T}_{r,k} - \partial_r \mathcal{T}_{r,k}[h]\,\big\|_{L^\infty \to L^\infty}  \\
&\qquad + \|\mathcal{T}_{r,k}^{-1}\|_{L^\infty \to L^\infty}\,
   \|\mathcal{T}_{r+h,k} - \mathcal{T}_{r,k}\|_{L^\infty \to L^\infty}\,
   \|\mathcal{T}_{r,k}^{-1}\|_{L^\infty \to L^\infty}\,
   \|\mathcal{T}_{r+h,k} - \mathcal{T}_{r,k}\|_{L^\infty \to L^\infty}\,
   \|\mathcal{T}_{r+h,k}^{-1}\|_{L^\infty \to L^\infty} \\
&\;\;\lesssim \|\mathcal{T}_{r,k}^{-1}\|_{L^\infty \to L^\infty}\,
   \big\|\,\mathcal{T}_{r+h,k} - \mathcal{T}_{r,k} - \partial_r \mathcal{T}_{r,k}[h]\,\big\|_{L^\infty \to L^\infty}
   \;+\; \|\mathcal{T}_{r,k}^{-1}\|_{L^\infty \to L^\infty}^2\,
   \|\mathcal{T}_{r+h,k}^{-1}\|_{L^\infty \to L^\infty}\,\|h\|_{L^2(\lambda)}^2 .
\end{aligned}
\]
The remainder bound in Lemma \ref{lemma::operatorderivatives} for $\mathcal{P}_{r, k}$ implies that $\big\|\,\mathcal{T}_{r+h,k} - \mathcal{T}_{r,k} - \partial_r \mathcal{T}_{r,k}[h]\,\big\|_{L^\infty \to L^\infty}   
 \lesssim \|h\|_{L^2(\lambda)}^2.$
We conclude that
\[
\begin{aligned}
& \big\|\mathcal{T}_{r+h,k}^{-1} - \mathcal{T}_{r,k}^{-1}
+ \mathcal{T}_{r,k}^{-1}(\partial_r \mathcal{T}_{r,k}[h])\mathcal{T}_{r,k}^{-1}\big\|_{L^\infty \to L^\infty} \lesssim   \|h\|_{L^2(\lambda)}^2.
\end{aligned}
\]


\noindent \textbf{Combining the bounds for all three terms} and using that $ \|\mathcal{T}_{r,k}^{-1}\|_{L^2(\lambda) \to L^2(\lambda)} < \infty$, we conclude that
 \begin{align*}
 \|\partial_{r} v_{r + h,k} - \partial_{r} v_{r,k} - \partial_r^2 v_{r, k}[\cdot, h]  \|_{L^\infty \rightarrow L^\infty}  &= \|\mathcal{T}_{r+ h,k}^{-1}\mathcal{P}_{r+ h,k}
 - \mathcal{T}_{r,k}^{-1}\mathcal{P}_{r,k}
 - \partial_r\!\big[\mathcal{T}_{r,k}^{-1}\mathcal{P}_{r,k}\big][h]\|_{L^\infty \rightarrow L^\infty}\\
 & \lesssim \|h\|_{L^2(\lambda)}^2 .
 \end{align*}

 



\end{proof}





\begin{lemma}[Cross Fréchet differentiability of $(r,k) \mapsto v_{r,k}$]
\label{lemma::valuederivsecond::kernel}
Let $k, k' \in \mathcal{K}$. Suppose $\mathcal{T}_{r,k}$ and $\mathcal{T}_{r,k'}$ are both invertible as maps from $L^2(\lambda) \to L^2(\lambda)$ with $\|\mathcal{T}_{r,k}^{-1}\|_{L^2 \rightarrow L^2} + \|\mathcal{T}_{r,k'}^{-1}\|_{L^2 \rightarrow L^2} < \infty$. Then, the map  $k \mapsto \partial_r v_{r,k}[h]$ is Fréchet differentiable with derivative
\[
\begin{aligned}
\adk\partial_r v_{r,k}\big[\,h,\,k'-k\,\big]
&= \mathcal{T}_{r,k}^{-1}\Bigg[
\frac{1}{\tau}\,\mathcal{P}_{r,k}\!\left[(a',s')\mapsto
\operatorname{Cov}_{\pi_{r,k}(\cdot\mid s')}\!\Big(
\partial_r q_{r,k}[h](\cdot,s'),\; \adk q_{r,k}[\,k'-k\,](\cdot,s')
\Big)\right] \\
&\qquad\qquad
+ \mathcal{P}_{k'-k}\,\Pi_{r,k}\,\partial_r q_{r,k}[h]
\Bigg].
\end{aligned}
\]
Here $\partial_r q_{r,k}[h](\cdot,s') := \mathcal{T}_{r,k}^{-1}h(\cdot,s')$ and $\adk q_{r,k}[k'-k](\cdot,s') := \gamma\,\adk v_{r,k}[k'-k](\cdot,s')$.
In particular, we have the expansion:
\begin{align*}
 \|\partial_r v_{r,k'} - \partial_r v_{r,k} -  \adk \partial_r v_{r,k}[\cdot, k' - k]  \|_{L^\infty \rightarrow L^2}
 &\lesssim  \|k'-k\|_{L^2(\mu\otimes\lambda)}^2
      + \|v_{r,k'}-v_{r,k}\|_{L^2(\lambda)}^2 .
\end{align*}
 
\end{lemma}
\begin{proof}
By the chain rule, \(k \mapsto \mathcal{T}_{r,k}\) is Gâteaux differentiable with
\[
\adk \mathcal{T}_{r,k}[k'-k] \;=\; -\,\gamma\,\adk \mathcal{P}_{r,k}[k'-k],
\qquad\text{i.e.}\quad
\big(\adk \mathcal{T}_{r,k}[k'-k]\big)(f)
= -\,\gamma\,\big(\adk \mathcal{P}_{r,k}[k'-k]\big)(f).
\]
Applying the inverse–map differentiation formula,
\[
\adk \mathcal{T}_{r,k}^{-1}[k'-k]
= -\,\mathcal{T}_{r,k}^{-1}\,\big(\adk \mathcal{T}_{r,k}[k'-k]\big)\,\mathcal{T}_{r,k}^{-1}.
\]
By the product rule and differentiability of \(k\mapsto \mathcal{T}_{r,k}\) and \(k\mapsto \mathcal{P}_{r,k}\),
\[
\adk\!\big[\mathcal{T}_{r,k}^{-1}\mathcal{P}_{r,k}\big][k'-k]
= -\,\mathcal{T}_{r,k}^{-1}\,\big(\adk \mathcal{T}_{r,k}[k'-k]\big)\,\mathcal{T}_{r,k}^{-1}\,\mathcal{P}_{r,k}
\;+\;
\mathcal{T}_{r,k}^{-1}\,\big(\adk \mathcal{P}_{r,k}[k'-k]\big).
\]
Hence, for directions $h\in L^\infty(\lambda)$ and
$k' - k$, the mixed partial satisfies
\begin{align*}
\adk\partial_r v_{r,k}[\,h,\,k' - k\,]
&= -\,\mathcal{T}_{r,k}^{-1}\,\big(\adk \mathcal{T}_{r,k}[k' - k]\big)\,
     \mathcal{T}_{r,k}^{-1}\,\mathcal{P}_{r,k}[h]
 \;+\; \mathcal{T}_{r,k}^{-1}\,\big(\adk \mathcal{P}_{r,k}[k' - k]\big)[h] \\
&= \gamma\,\mathcal{T}_{r,k}^{-1}\,\big(\adk \mathcal{P}_{r,k}[k' - k]\big)\,
      \mathcal{T}_{r,k}^{-1}\,\mathcal{P}_{r,k}[h]
 \;+\; \mathcal{T}_{r,k}^{-1}\,\big(\adk \mathcal{P}_{r,k}[k' - k]\big)[h] \\
&= \mathcal{T}_{r,k}^{-1}\,\big(\adk \mathcal{P}_{r,k}[k' - k]\big)\,
   \big(I+\gamma\,\mathcal{T}_{r,k}^{-1}\mathcal{P}_{r,k}\big)[h] \\
&= \mathcal{T}_{r,k}^{-1}\,\big(\adk \mathcal{P}_{r,k}[k' - k]\big)\,
   \big(I+\gamma\,\partial_r v_{r,k}\big)[h].
\end{align*}
Here we used the identity $\adk \mathcal{T}_{r,k}[k' - k] = -\,\gamma\,\adk \mathcal{P}_{r,k}[k' - k]$ and
$\big(I+\gamma\,\mathcal{T}_{r,k}^{-1}\mathcal{P}_{r,k}\big)=\big(I+\gamma\,\partial_r v_{r,k}\big)$.
Plugging in the definitions of the operators from Lemma \ref{lemma::operatorderivatives}, we conclude that
\[
\begin{aligned}
\adk\partial_r v_{r,k}\big[\,h,\,k'-k\,\big]
&= \mathcal{T}_{r,k}^{-1}\Bigg[
\frac{1}{\tau}\,\mathcal{P}_{r,k}\!\left[(a',s')\mapsto
\operatorname{Cov}_{\pi_{r,k}(\cdot\mid s')}\!\Big(
\partial_r q_{r,k}[h](\cdot,s'),\; \adk q_{r,k}[\,k'-k\,](\cdot,s')
\Big)\right] \\
&\qquad\qquad
+ \mathcal{P}_{k'-k}\,\Pi_{r,k}\,\partial_r q_{r,k}[h]
\Bigg].
\end{aligned}
\]
where $\adk q_{r,k}[k' - k](\cdot,s') := \gamma\,\adk v_{r,k}[k' - k](\cdot,s')$ and
$\pi_{r,k}(\cdot\mid s')$ denotes the softmax distribution induced by $q_{r,k}(\cdot,s')$.






\noindent \textbf{Establish second-order bound for Taylor series remainder.} Arguing exactly as in the proof of Lemma \ref{lemma::valuederivsecond}, but differentiating in \(k\), we have
\begin{align*}
&\mathcal{T}_{r,k'}^{-1}\mathcal{P}_{r,k'}
 - \mathcal{T}_{r,k}^{-1}\mathcal{P}_{r,k}
 - \adk\!\big[\mathcal{T}_{r,k}^{-1}\mathcal{P}_{r,k}\big][\,k' - k\,] \\
&= \underbrace{\Big[\mathcal{T}_{r,k'}^{-1} - \mathcal{T}_{r,k}^{-1}
         + \mathcal{T}_{r,k}^{-1}\,\big(\adk \mathcal{T}_{r,k}[\,k' - k\,]\big)\,
           \mathcal{T}_{r,k}^{-1}\Big]\mathcal{P}_{r,k}}_{\text{Inverse-operator remainder}} \\
&\quad + \underbrace{\mathcal{T}_{r,k}^{-1}\Big[\mathcal{P}_{r,k'} - \mathcal{P}_{r,k}
           - \adk \mathcal{P}_{r,k}[\,k' - k\,]\Big]}_{\text{Direct $\mathcal{P}$ remainder}} \\
&\quad + \underbrace{\Big(\mathcal{T}_{r,k'}^{-1} - \mathcal{T}_{r,k}^{-1}\Big)
       \Big(\mathcal{P}_{r,k'} - \mathcal{P}_{r,k}\Big)}_{\text{Mixed interaction term}}.
\end{align*}
We can write the first term as
\[
\begin{aligned}
&\mathcal{T}_{r,k'}^{-1} - \mathcal{T}_{r,k}^{-1}
  + \mathcal{T}_{r,k}^{-1}\big(\adk \mathcal{T}_{r,k}[\,k' - k\,]\big)\mathcal{T}_{r,k}^{-1} \\
&\qquad= - \mathcal{T}_{r,k}^{-1}\Big(\mathcal{T}_{r,k'} - \mathcal{T}_{r,k}  - \adk \mathcal{T}_{r,k}[\,k' - k\,]\big) \mathcal{T}_{r,k}^{-1} \\
&\qquad\quad +\; \mathcal{T}_{r,k}^{-1}\big(\mathcal{T}_{r,k'} - \mathcal{T}_{r,k}\big)\mathcal{T}_{r,k}^{-1}\big(\mathcal{T}_{r,k'} - \mathcal{T}_{r,k}\big)\mathcal{T}_{r,k'}^{-1}.
\end{aligned}
\]
Combining the previous two displays, we find that
\begin{align*}
&\mathcal{T}_{r,k'}^{-1}\mathcal{P}_{r,k'}
 - \mathcal{T}_{r,k}^{-1}\mathcal{P}_{r,k}
 - \adk\!\big[\mathcal{T}_{r,k}^{-1}\mathcal{P}_{r,k}\big][\,k' - k\,] \\
&=
\underbrace{- \mathcal{T}_{r,k}^{-1}\Big(\mathcal{T}_{r,k'} - \mathcal{T}_{r,k}  - \adk \mathcal{T}_{r,k}[\,k' - k\,]\big) \mathcal{T}_{r,k}^{-1} }_{\text{Inverse linear remainder}} \\
&\qquad+\;
\underbrace{\mathcal{T}_{r,k}^{-1}\big(\mathcal{T}_{r,k'} - \mathcal{T}_{r,k}\big)\mathcal{T}_{r,k}^{-1}\big(\mathcal{T}_{r,k'} - \mathcal{T}_{r,k}\big)\mathcal{T}_{r,k'}^{-1}\,\mathcal{P}_{r,k}}_{\text{Inverse quadratic product}} \\
&\qquad+\;
\underbrace{\mathcal{T}_{r,k}^{-1}\Big[\mathcal{P}_{r,k'} - \mathcal{P}_{r,k}
           - \adk \mathcal{P}_{r,k}[\,k' - k\,]\Big]}_{\text{Direct $\mathcal{P}$ remainder}} \\
&\qquad+\;
\underbrace{\Big(\mathcal{T}_{r,k'}^{-1} - \mathcal{T}_{r,k}^{-1}\Big)
       \Big(\mathcal{P}_{r,k'} - \mathcal{P}_{r,k}\Big)}_{\text{Mixed interaction term}}.
\end{align*}




\noindent \textbf{For the first term}, we have
\begin{align*}
  &\big\| \mathcal{T}_{r,k}^{-1}\Big(\mathcal{T}_{r,k'} - \mathcal{T}_{r,k} 
      - \adk \mathcal{T}_{r,k}[\,k' - k\,]\Big)\mathcal{T}_{r,k}^{-1}
   \big\|_{L^\infty \rightarrow L^2}  \\
  &\leq \|\mathcal{T}_{r,k}^{-1}\|_{L^2 \rightarrow L^2}\;
       \big\|\mathcal{T}_{r,k'} - \mathcal{T}_{r,k} 
          - \adk \mathcal{T}_{r,k}[\,k' - k\,]\big\|_{L^\infty \rightarrow L^2}\;
       \|\mathcal{T}_{r,k}^{-1}\|_{L^\infty \rightarrow L^\infty} \\
  &\lesssim \big\|\mathcal{T}_{r,k'} - \mathcal{T}_{r,k} 
          - \adk \mathcal{T}_{r,k}[\,k' - k\,]\big\|_{L^\infty \rightarrow L^2}.
\end{align*}
Here we used that $\|\mathcal{T}_{r,k}^{-1}\|_{L^2 \rightarrow L^2} < \infty$ and absorbed its dependence into the notation $\lesssim$. By Lemma~\ref{lemma::operatorderivatives}, we thus obtain
\begin{align*}
  &\big\| \mathcal{T}_{r,k}^{-1}\Big(\mathcal{T}_{r,k'} - \mathcal{T}_{r,k} 
      - \adk \mathcal{T}_{r,k}[\,k' - k\,]\Big)\mathcal{T}_{r,k}^{-1}
   \big\|_{L^\infty \rightarrow L^2} \\
  &\lesssim \|v_{r,k'} - v_{r,k}\|_{L^2(\lambda)}^2
          + \|v_{r,k'} - v_{r,k}\|_{L^2(\lambda)}\,
            \|k' - k\|_{L^2(\mu \otimes \lambda)}.
\end{align*}

\noindent \textbf{For the second term.} At the end of this proof, we establish the claim that $\|\mathcal{P}_{r,k'}-\mathcal{P}_{r,k}\|_{L^\infty\to L^2}  \lesssim \|\mathcal{P}_{r,k'}-\mathcal{P}_{r,k}\|_{L^2(\lambda)\to L^2(\lambda)} \lesssim \|k'-k\|_{L^2(\mu\otimes\lambda)}
    +  \|v_{r,k'}-v_{r,k}\|_{L^2(\lambda)}$ and, hence, also $\|\mathcal{T}_{r,k'}-\mathcal{T}_{r,k}\|_{L^2(\lambda)\to L^2(\lambda)} \lesssim \|k'-k\|_{L^2(\mu\otimes\lambda)}
    +  \|v_{r,k'}-v_{r,k}\|_{L^2(\lambda)}$. Using submultiplicativity of operator norms,
\begin{align*}
&\big\|\mathcal{T}_{r,k}^{-1}\big(\mathcal{T}_{r,k'}-\mathcal{T}_{r,k}\big)
  \mathcal{T}_{r,k}^{-1}\big(\mathcal{T}_{r,k'}-\mathcal{T}_{r,k}\big)
  \mathcal{T}_{r,k'}^{-1}\,\mathcal{P}_{r,k}\big\|_{L^2 \rightarrow L^2} \\
&\qquad\le 
\|\mathcal{T}_{r,k}^{-1}\|_{L^2(\lambda)\to L^2(\lambda)}^2\,
\|\mathcal{T}_{r,k'}^{-1}\|_{L^2(\lambda)\to L^2(\lambda)}\,
\|\mathcal{P}_{r,k}\|_{L^2(\lambda)\to L^2(\lambda)}\,
\|\mathcal{T}_{r,k'}-\mathcal{T}_{r,k}\|_{L^2(\lambda)\to L^2(\lambda)}^2 \\
&\qquad\lesssim 
\Big(\|k'-k\|_{L^2(\mu\otimes\lambda)}
      + \|v_{r,k'}-v_{r,k}\|_{L^2(\lambda)}\Big)^2\\
      &\qquad\lesssim 
\|k'-k\|_{L^2(\mu\otimes\lambda)}^2
      + \|v_{r,k'}-v_{r,k}\|_{L^2(\lambda)}^2
\end{align*}
where we absorbed $\|\mathcal{T}_{r,k}^{-1}\|_{L^2(\lambda)\to L^2(\lambda)}$, $\|\mathcal{T}_{r,k'}^{-1}\|_{L^2(\lambda)\to L^2(\lambda)}$, and $\|\mathcal{P}_{r,k}\|_{L^2(\lambda)\to L^2(\lambda)}$ into the $\lesssim$ notation.




\noindent \textbf{For the third term}, by an identical argument to the first term, we have that
\begin{align*}
  &\big\| \mathcal{T}_{r,k}^{-1}\Big(\mathcal{P}_{r,k'} - \mathcal{P}_{r,k} 
      - \adk \mathcal{P}_{r,k}[\,k' - k\,]\Big) 
   \big\|_{L^\infty \rightarrow L^2} \\
  &\lesssim \|v_{r,k'} - v_{r,k}\|_{L^2(\lambda)}^2
          + \|v_{r,k'} - v_{r,k}\|_{L^2(\lambda)}\,
            \|k' - k\|_{L^2(\mu \otimes \lambda)}.
\end{align*}



\noindent \textbf{For the fourth term}, we note that
\begin{align*}
   \big\|\big(\mathcal{T}_{r,k'}^{-1} - \mathcal{T}_{r,k}^{-1}\big)
       \big(\mathcal{P}_{r,k'} - \mathcal{P}_{r,k}\big)\big\|_{L^\infty \to L^2} 
   &\lesssim  \|\mathcal{T}_{r,k'}^{-1} - \mathcal{T}_{r,k}^{-1} \|_{L^2(\lambda) \to L^2(\lambda)}\;
       \|\mathcal{P}_{r,k'} - \mathcal{P}_{r,k}\|_{L^\infty \to L^2}.
\end{align*}
Moreover, by the resolvent identity, we have
\[
\begin{aligned}
\|\mathcal{T}_{r,k'}^{-1} - \mathcal{T}_{r,k}^{-1}\|_{L^2(\lambda) \to L^2(\lambda)}
&= \|\mathcal{T}_{r,k}^{-1}\big(\mathcal{T}_{r,k} - \mathcal{T}_{r,k'}\big)\mathcal{T}_{r,k'}^{-1}\|_{L^2(\lambda) \to L^2(\lambda)} \\
&\le \|\mathcal{T}_{r,k}^{-1}\|_{L^2(\lambda) \to L^2(\lambda)}\,
\|\mathcal{T}_{r,k'} - \mathcal{T}_{r,k}\|_{L^2(\lambda) \to L^2(\lambda)}\,
\|\mathcal{T}_{r,k'}^{-1}\|_{L^2(\lambda) \to L^2(\lambda)}.
\end{aligned}
\]
Hence,
\begin{align*}
   \big\|\big(\mathcal{T}_{r,k'}^{-1} - \mathcal{T}_{r,k}^{-1}\big)
       \big(\mathcal{P}_{r,k'} - \mathcal{P}_{r,k}\big)\big\|_{L^\infty \to L^2} 
   &\lesssim  \|\mathcal{T}_{r,k}^{-1}\|_{L^2(\lambda) \to L^2(\lambda)}
       \|\mathcal{T}_{r,k'}^{-1}\|_{L^2(\lambda) \to L^2(\lambda)} \\
   &\qquad \times \|\mathcal{T}_{r,k'} - \mathcal{T}_{r,k} \|_{L^2(\lambda) \to L^2(\lambda)}\;
       \|\mathcal{P}_{r,k'} - \mathcal{P}_{r,k}\|_{L^\infty \to L^2}.
\end{align*}
At the end of this proof, we establish the claim that
\[
\begin{aligned}
\|\mathcal{P}_{r,k'}-\mathcal{P}_{r,k}\|_{L^\infty\to L^2}
&\lesssim \|\mathcal{P}_{r,k'}-\mathcal{P}_{r,k}\|_{L^2(\lambda)\to L^2(\lambda)} \\
&\lesssim \|k'-k\|_{L^2(\mu\otimes\lambda)}
    +  \|v_{r,k'}-v_{r,k}\|_{L^2(\lambda)},
\end{aligned}
\]
and, hence, also
\[
\|\mathcal{T}_{r,k'}-\mathcal{T}_{r,k}\|_{L^2(\lambda)\to L^2(\lambda)}
\lesssim \|k'-k\|_{L^2(\mu\otimes\lambda)}
    +  \|v_{r,k'}-v_{r,k}\|_{L^2(\lambda)}.
\]
Assuming this claim holds, we conclude that
\begin{align*}
   \big\|\big(\mathcal{T}_{r,k'}^{-1} - \mathcal{T}_{r,k}^{-1}\big)
&\qquad \times \big(\mathcal{P}_{r,k'} - \mathcal{P}_{r,k}\big)\big\|_{L^\infty \to L^2} \\
   &\lesssim  \|\mathcal{T}_{r,k}^{-1}\|_{L^2(\lambda) \to L^2(\lambda)}
       \|\mathcal{T}_{r,k'}^{-1}\|_{L^2(\lambda) \to L^2(\lambda)} \\
   &\qquad \times \left\{ \|k'-k\|_{L^2(\mu\otimes\lambda)}
    +  \|v_{r,k'}-v_{r,k}\|_{L^2(\lambda)} \right\}^2 \\
     &\lesssim    \|k'-k\|_{L^2(\mu\otimes\lambda)}^2
    +  \|v_{r,k'}-v_{r,k}\|_{L^2(\lambda)}^2.
\end{align*}


\noindent \textbf{Combining bounds for all four terms}, we conclude that
\begin{align*}
 \|\partial_r v_{r,k'} - \partial_r v_{r,k}
&\qquad -  \adk \partial_r v_{r,k}[\cdot, k' - k]  \|_{L^\infty \rightarrow L^2} \\
&= \left \|\mathcal{T}_{r,k'}^{-1}\mathcal{P}_{r,k'}
 - \mathcal{T}_{r,k}^{-1}\mathcal{P}_{r,k}
 - \adk\!\big[\mathcal{T}_{r,k}^{-1}\mathcal{P}_{r,k}\big][\,k' - k\,] \right\|_{L^\infty \rightarrow L^2} \\
 & \lesssim  \|k'-k\|_{L^2(\mu\otimes\lambda)}^2
      + \|v_{r,k'}-v_{r,k}\|_{L^2(\lambda)}^2,
\end{align*}
as desired.
   


\noindent \textbf{Establishing claimed bound for $\|\mathcal{P}_{r,k'}-\mathcal{P}_{r,k}\|_{L^2(\lambda)\to L^2(\lambda)}$.} We can write
\[
\begin{aligned}
\big((\mathcal{P}_{r,k'}-\mathcal{P}_{r,k})f\big)(a,s)
&= \int \Big\langle \pi_{r,k'}(\cdot\mid s'),\, f(\cdot,s')\Big\rangle\,
   \underbrace{\big(k'-k\big)(s'\mid a,s)}_{:=\,\Delta k}\,\mu(ds') \\
&\quad + \int \underbrace{\Big\langle \pi_{r,k'}(\cdot\mid s')-\pi_{r,k}(\cdot\mid s'),\, f(\cdot,s')\Big\rangle}_{:=\,\langle\Delta\pi,\,f\rangle(s')}\,
   k(s'\mid a,s)\,\mu(ds').
\end{aligned}
\]
By Cauchy–Schwarz and Jensen’s inequality, the first term satisfies
\begin{align*}
    \Big\|\int \langle \pi_{r,k'}(\cdot\mid s'), f(\cdot,s')\rangle\,\Delta k(s'\mid a,s)\,\mu(ds')\Big\|_{L^2(\lambda)}
&\le \|\Delta k\|_{L^2(\mu\otimes\lambda)}\,
    \Big\|\Big(\sum_{a' \in \mathcal{A}} \pi_{r,k'}(a'\mid s')\,|f(a',s')|^2\Big)^{1/2}\Big\|_{L^2(\mu)} \\
&\le \sqrt{|\mathcal A|}\,\|\Delta k\|_{L^2(\mu\otimes\lambda)}\,\|f\|_{L^2(\lambda)}.
\end{align*}
Similarly, viewing the second term as a Hilbert–Schmidt operator with kernel $K^*((a,s),(a',s')) = k(s'\mid a,s)\,\Delta\pi(a'\mid s'),$
we obtain
\[
\Big\|\int \langle \Delta\pi(\cdot\mid s'), f(\cdot,s')\rangle\,k(s'\mid a,s)\,\mu(ds')\Big\|_{L^2(\lambda)}
\;\le\; \|k\|_{L^2(\mu\otimes\lambda)}\,\|\Delta\pi(\cdot\mid\cdot)\|_{L^2(\lambda)}\,\|f\|_{L^2(\lambda)}.
\]
Hence, using the Lipschitz property of the softmax map,
\[
\|\pi_{r,k'}-\pi_{r,k}\|_{L^2(\lambda)}
\;\lesssim\; \|v_{r,k'}-v_{r,k}\|_{L^2(\lambda)}.
\]
we conclude that
\begin{align*}
\|\mathcal{P}_{r,k'}-\mathcal{P}_{r,k}\|_{L^2(\lambda)\to L^2(\lambda)}
&\le \|k'-k\|_{L^2(\mu\otimes\lambda)}
    + \|k\|_{L^2(\mu\otimes\lambda)}\,\|\pi_{r,k'}-\pi_{r,k}\|_{L^2(\lambda)} \\
&\lesssim \|k'-k\|_{L^2(\mu\otimes\lambda)}
    +  \|v_{r,k'}-v_{r,k}\|_{L^2(\lambda)}.
\end{align*}


 



 \end{proof}

 
 

{
\singlespacing
\bibliographystyle{abbrvnat}
\bibliography{ref}
}

\end{document}
